% !TEX program = xelatex
% 数学认知空间 MCS：学习者状态外置的完整理论专稿
% 由 edition/MCS-完整理论专稿.md 经 validation/build_latex.py 生成。
% 编译：xelatex（连续运行两次以生成目录与页码）。
\documentclass[UTF8,a4paper,zihao=-4,fontset=windows,openany]{ctexbook}

\usepackage{amsmath}
\usepackage{amssymb}
\usepackage{array}
\usepackage{longtable}
\usepackage{booktabs}
\usepackage{enumitem}
\usepackage{geometry}
\usepackage{graphicx}
\usepackage{xcolor}
\usepackage{fancyhdr}
\usepackage[hidelinks,unicode]{hyperref}

\geometry{a4paper,left=2.6cm,right=2.6cm,top=2.7cm,bottom=2.5cm,headsep=0.5cm,headheight=15pt}
\linespread{1.32}

\setmainfont{Times New Roman}
\setsansfont{Arial}
\setmonofont{Consolas}[Scale=0.88]

\setcounter{secnumdepth}{-2}
\setcounter{tocdepth}{1}
\setlength{\parindent}{2em}
\tolerance=2000
\emergencystretch=2em
\renewcommand{\arraystretch}{1.18}

\newcommand{\mcspow}[1]{\ensuremath{^{\mathrm{#1}}}}
\newcommand{\mcssub}[1]{\ensuremath{_{\mathrm{#1}}}}
\newcommand{\mcsnum}[1]{\textcircled{\scriptsize #1}}
\newcommand{\mcsorigin}[1]{\par\noindent{\small\color{black!55}#1}\par\addvspace{0.8em}}
\newcommand{\mcssep}{\par\addvspace{0.6em}\noindent\rule{\linewidth}{0.2pt}\par\addvspace{0.6em}}

\newsavebox{\mcsmathbox}
\newcommand{\mcsdisplay}[1]{%
  \par\addvspace{0.7em}%
  \sbox{\mcsmathbox}{\ensuremath{\displaystyle #1}}%
  \ifdim\wd\mcsmathbox>\linewidth
    \noindent\resizebox{\linewidth}{!}{\usebox{\mcsmathbox}}%
  \else
    \noindent\hspace*{\fill}\usebox{\mcsmathbox}\hspace*{\fill}%
  \fi
  \par\addvspace{0.7em}%
}

\pdfstringdefDisableCommands{%
  \def\ensuremath#1{#1}%
  \def\text#1{#1}%
  \def\operatorname#1{#1}%
  \def\mathrm#1{#1}%
}

\providecommand{\llbracket}{\mathopen{[\![}}
\providecommand{\rrbracket}{\mathclose{]\!]}}
\newcommand{\mcsbrk}{\penalty 200\relax}

\pagestyle{fancy}
\fancyhf{}
\fancyhead[L]{\small\itshape 数学认知空间 MCS}
\fancyhead[R]{\small\thepage}
\renewcommand{\headrulewidth}{0.25pt}
\fancypagestyle{plain}{\fancyhf{}\fancyhead[R]{\small\thepage}\renewcommand{\headrulewidth}{0pt}}

\title{\Huge 数学认知空间 MCS\\[0.7em]\Large 新时代数学教育的基础设施\\[0.45em]\large 可认证的数学知识本体与个性化学习路线}
\author{刘沛恒}
\date{}

\hypersetup{%
  pdftitle={数学认知空间 MCS：新时代数学教育的基础设施},
  pdfauthor={刘沛恒},
  pdfsubject={可认证的数学知识本体与个性化学习路线},
}

\begin{document}
\maketitle
\thispagestyle{empty}
整合阅读版。共享本体 M、外部认知模型 E 与派生结果 D 分离。本文件由分章正文生成，数学论证、引用元定理、有限核验与经验效果分别标注。
分章正文为修订源；本版不另外修改其数学内容。阅读目录见 主入口。索引中的行号指向分章源文件。
研究地位：本稿是一套附条件数学结果的知识组织与认知模型接入规范。“完整”表示章节覆盖，不表示工程实现或学习机制研究已经完成；接口本身不能确定教学效果，具体反例与检验条件见第 16.7 节。
\tableofcontents
\clearpage
\clearpage\chapter{D 跨章接口契约}
\mcsorigin{原章节：D-接口契约.md}
本附录规定全文各部分怎样交换信息、分别承担哪些检查义务。附录里混有两类陈述，须分开读：一类是记号与对象的定义（证据状态 \textbf{DEF}）；另一类是以“须 / 不得 / 保留”写出的接口义务，即对后文章节与实现的验收条款。义务条款本身不是定义，也不自带证据，只有在后文给出定义、证明或实现检查时才逐条兑现；本附录不能凭这些约定宣称整套理论已获证明或软件已经合规。正因如此，本附录在第 04 章给出坐标定义、在后文兑现各项义务之前，不能独立完成核验。下面的小例子用于解释约束的作用，不证明实际教学效果。

\textbf{为什么需要这份契约。} 从数学内部看，“某人相信一个命题”不能混进“该命题已有证明”的推理；从学习应用看，同一份数学内容又应允许不同的人从不同入口接近。设计因此从知识条目及连线，发展到带证据的公共行动，再把个人状态与派生路线分开。分开以后，仍须说明接口传递的是材料、证明、使用条件还是个人判断。

本节依次处理六件事：信息归属、表达与证明的基础、行动和路线的形成、支持与关系的含义、个人视图与计划的产生、证据责任。这里的“公共”表示共享版本与规则，不表示所有登记内容都已证实或对每个人都有效。

\section{D.1 三层与记号}
直观上，先保存共同讨论的知识档案，再另行描述一个学习者，最后计算适合此次任务的视图和候选计划。三层是数据与判断的分工，不是三条已经验证的心理规律。

\begin{itemize}
\item \textbf{本体} \(M\) 是给定版本的公共知识结构：签名 \(\Sigma\)、背景理论 \(T\)、节点集 \(N_M\)（每个节点具有有限描述，节点集本身不必有限）、不可变负载 \(\mathrm{Payload}\)、表征 \(\mathrm{Rep}_M\)、认证证据 \(\mathrm{Cert}_M\)、支持族 \(\mathrm{Supp}_M\)、公开行动契约族 \(\Lambda_M\)、通用关系描述 \(R^{0}_M\)、聚合 \(\mathrm{Agg}\) 与通用模板 \(\mathrm{Templ}\)、环境版本边界 \(\partial M\)（共十二个坐标，见定义 4.1）。推导出的硬关系也只依赖这些公开材料及声明的结构背景。
\item \textbf{外部模型} \(E\) 指定状态空间 \(X_E\)、观察/转移/评价接口及其版本；\(s\in X_E\) 是某学习者的状态。目标 \(g\)、预算 \(b\)、偏好/策略 \(\rho\) 是外部规划输入。个体标识、作答记录、掌握估计不属于 \(M\)。
\item \textbf{派生结果} \(D(M,E,s,g,b,\rho)\) 包括适配视图、评价、路线与呈现，不自动写回 \(M\)。随机接口返回分布，或另显式给随机种子。
\item 一般误区模式、题目、方法与能力维度是共享对象；谁持有误区、谁答过题、谁掌握何种能力是外部实例。描述认知模型的节点不等于该模型运行时的状态。
\end{itemize}
“学习者无关”是一条卫生约定，不是泄漏检测结论：它指固定公共版本后，替换学习者输入不暗中改写该版本，且任何对公共材料的修改都必须走显式版本化修订路径，不作为个人状态更新的副作用。它挡不住“把学习者数据误登记进 \(M\)”这类内容错误，只规定 \(M\) 的变更路径必须可审计；它也不否认公共材料的选取曾受到研究判断、教学经验或认知研究的影响。是否真的做到无关，须由实现检查 \(M\) 的更新入口是否只接受显式修订来核验，本约定本身不提供该保证。

公共性、正确性与个体适用性分别判断。比如“用面积解释积分”可以登记为共享表征，其数学适用条件由相应证据负责，而它是否适合当前学习者由 \(E\) 评价。\(M\) 也可保存未决主张、候选关联和已反驳命题；登记只说明这些对象可被引用，不把它们统一认证为真。

\section{D.2 形式基础}
选择形式语言，是为了把想讨论的对象写清，并使推理可以检查。例如，本稿不只记录两节点之间有没有关系，还要在对象语言内部量化关系：要写出“若 \(r\) 是 \(s\) 的子关系，凡由 \(r\) 许可的步骤也由 \(s\) 许可”这类跨关系的保持断言，就需要以关系为取值的变元和对关系的量词；只给每对节点存一条布尔边，写不出这种断言。这是采用关系类型和更高类型的一个具体动机。该动机成立的前提是后文确实依赖这类断言（由第 02、05–08 章的使用兑现），本处不单独证明其必要性；仅仅保存节点和连线，并不足以单独说明必须采用高阶语言。

元理论取 ZFC，用来定义语法、记录、函数和模型等对象。自动生成部分采用有效可数签名。每个具体语义模型为集合大小结构；所有这类模型的总体按元层的类处理，不假定它是另一个集合。对象语言采用经典外延简单类型论，含真值类型 \(o\)、基本类型和箭头类型。三种判断分别回答：

\begingroup\small\setlength{\tabcolsep}{4pt}
\begin{longtable}{>{\raggedright\arraybackslash}p{\dimexpr 0.1700\linewidth-2\tabcolsep\relax}>{\raggedright\arraybackslash}p{\dimexpr 0.4594\linewidth-2\tabcolsep\relax}>{\raggedright\arraybackslash}p{\dimexpr 0.3706\linewidth-2\tabcolsep\relax}}
\toprule
\textbf{判断} & \textbf{所问的问题} & \textbf{检查位置} \\
\midrule
\endfirsthead
\toprule
\textbf{判断} & \textbf{所问的问题} & \textbf{检查位置} \\
\midrule
\endhead
\bottomrule
\endfoot
\bottomrule
\endlastfoot
\(\Gamma\vdash_T \varphi\) & 按指定演算，在背景 \(T\) 与假设 \(\Gamma\) 下能否推出 \(\varphi\)？ & 第 02 章固定演算与证书条件 \\
\(T\cup\Gamma\vDash_H \varphi\) & 在指定 Henkin 模型类及相应赋值下，前提成立时结论是否成立？ & 第 03 章的 Henkin 语义 \\
\(T\cup\Gamma\vDash_\mathrm{std} \varphi\) & 换成全函数标准模型类后，上述后承是否成立？ & 第 03 章的标准语义 \\
\end{longtable}
\endgroup

后两者改变的是模型类，不能互换；前者是推导判断，与语义后承的连接要靠可靠性和相应完备性证明。采用这些基础是明确的工作选择，不是基础越强便越能解释学习的结论。

对象表达、语法代码、节点引用分别分型。节点可生成不意味着命题为真；节点同一不定义为逻辑等价；已证定理不统一压成一个节点。

本稿不假设 \(V=L\)。可构造宇宙、绝对性和凝结结论只在声明的范围内作背景比较，不能仅凭形式相似就把它们移植成节点生成、编号或认知效果的证明。

\section{D.3 先有行动契约，后有路线与先修}
要判断一条路线是否合法，应先知道有哪些合法动作；不能先用“已有路线”定义动作，再用动作反过来证明路线存在。这里先规定公共资源的使用接口，再讨论怎样把行动接起来。

\(\Lambda_M\) 的每个公开行动契约 \(a\) 具有有限输入 \(I_a\subseteq N_M\)、有限输出 \(O_a\subseteq N_M\)、实现/结构见证 \(w_a\) 和作用类型（定义引入、认证推导、表征转换、任务/方法呈现等）。每类见证有独立的合法性规则；方法呈现只认证其呈现接口，不认证学习成功。契约不以“已存在一条规划路线”作为形成条件。

\textbf{输出含义与衔接义务（接口义务，非定义）。} \(I_a\)、\(O_a\) 给出节点端点；每项输入要求的使用条件，以及每项输出实际提供的资源和证据，须能由行动类型、\(w_a\) 与公开节点/证据引用逐项确定，不能只由一个行动总标签猜测。需要时明确理论、开放假设、表征类型、认证状态与版本。背景 \(B\) 中的资源也须有相应来源与声明条件；允许作为假设使用的命题不因此成为无条件定理。

例如，“展示链式法则的陈述”“提供它在指定背景下的证明”“判断某人已会使用它”分别是呈现、认证和外部评价。后续行动若要求认证结论，就必须查到匹配的证据及假设，不能只检查该节点曾被展示。个人掌握不由公共输出自动产生。这里规定的是契约的见证检查义务；2026-09-29 的独立行动适配器已在四个局部证书上落实理论、开放假设、类型、版本与结论核对，覆盖范围见案例表，不等于全契约库已经认证。

先按契约和有限事件见证定义结构性路线族 \(\mathrm{Rte}_M(B,G)\)，其中 \(B\) 是公开声明假定可用的背景、\(G\) 是结构目标。每个事件的输入须指定来自 \(B\) 的具体来源，或某个严格更早事件的对应输出，并检查该来源满足上述使用条件；目标来源同样检查。这一层的资源可重复引用、使用后不消耗；数量限制、时间占用与个人遗忘另列约束，不能从集合端点直接推出。

硬先修相对于\textbf{固定版本的契约库、背景、目标与声明的非空路线族}定义。它记录指定范围内路线的共同输入义务，不自动表示学习者在起点必须额外具备这些资源。例如 \(\emptyset\to\{m\}\to\{g\}\) 的路线途中使用 \(m\)，但 \(m\) 由路线本身提供；研究起点必要条件时须另看背景来源。个人状态筛选产生的路线子族及其必经点属于派生结果，不回写为本体普遍关系。结构可行不蕴含真实学习有效。共同输入义务只能在声明路线族内报告；未经该族枚举完整性证明时，它是可被新登记路线推翻的搜索结论，对外呈现必须附带这一范围与稳定性限定（判据见 D.4）。

\section{D.4 概念与关系}
“写到了什么”与“非用它不可”是不同问题。在经典逻辑中，\(P\) 与 \(P\land(Q\lor\neg Q)\) 等价；后者出现 \(Q\)，并不因此使 \(Q\) 成为前者的必要前提。这个例子说明出现记录不能直接充当必要性证明。

\(C_\mathrm{all}\) 是全等价表达概念出现的并集，仅在明确的语言、等价关系及出现函数下讨论其饱和。实际分析采用带来源和用途的支持族，分别记录一份表达出现了什么、一份证明使用了什么、一条路线要求了什么；是否所有允许路线都绕不开某项内容，须另外按声明的路线范围判断。支持记录不自动具有全局最小性或必要性。

精确支持接口返回 \(\mathrm{Unknown}\) 或 \(\mathrm{Known}(A)\)，其中 \(A\) 是指定用途下的支持集合；\(\mathrm{Known}(\emptyset)\) 是已经查明空支持的特例。未知不按空集补值。部分支持界由独立接口保存，不能与精确的 \(\mathrm{Known}(A)\) 混为一类。

公共路线的可达与共同先修均相对于固定版本的行动库及声明的路线族。若目前只登记了经 \(a\) 到达目标的路线，只能报告该范围内的必经条件；新增合法的绕行路线可能使它不再必经。尤其不能把“本次搜索找到的全部路线”未经完整性证明就称为“所有允许路线”。先修类结论的稳定性判据因此只有两条：声明范围（属于哪个路线族）与枚举完整性（该族是否已被完整列出）。两条未同时满足时，任何“先修 / 必经”结论都是可修订的搜索结论，对外呈现与下游引用必须携带这一限定；只有证明路线族已完整列出后，才能改称该范围内的必要条件。

硬演绎携带联合前提与证明；硬泛化携带同载体蕴含或指定类型的结构遗忘/解释见证。特化是逆关系。一般理论解释不自动等同于对象类包含。

\(R^{0}_M\) 保存通用的应用/类比/对偶等结构描述和候选认知关联；\(\mathrm{Eval}_E(s,r,\mathrm{ctx})\) 才产生个体适用性、成本或效果判断。通用描述存在不表示个体软关系评价成立。

\section{D.5 局部化和规划}
局部化回答“此次显示哪些内容”，规划回答“在声明的输入、目标和约束下怎样安排事件”。显示时隐藏了某个前提，不表示执行时可以不满足它。

\(\mathrm{Loc}_\lambda(M)=(V,\partial V,\tau,\mathrm{Pres})\)；输出可以是带来源映射的聚合视图，\(\tau\) 在聚合情形是到节点有限子集的溯源对应，不强称节点单值映射。若引用闭包无限，边界保留符号化环境引用，不承诺有限展开。

认知适配器的规范签名是 \(\mathrm{Adapt}(E,s,g,b)\)，参数中\textbf{不含 \(M\)}；\(E\) 自行声明所绑定的本体版本与只读访问范围，\(\mathrm{Adapt}\) 只产生显式参数 \(\lambda\)，再由 \(\mathrm{Loc}_\lambda(M)\) 读取公共结构。需要直接读取 \(M\) 内容的适配器不得沿用 \(\mathrm{Adapt}\) 这一记号，必须改用显式签名 \(\mathrm{Adapt}_M(M,E,s,g,b)\)，同样声明版本与只读范围；两种签名指同一适配器职责的两种入口，不增加局部化算子，且分开命名、分开检查，\(\mathrm{Adapt}\) 不因增加参数而悄悄变成 \(\mathrm{Adapt}_M\)。任一签名都不得暗中查询范围之外的本体，也不得改写 \(M\)。状态、目标或参数含节点引用时，还须检查它们与当前 \(M\) 的版本相容；无法解析的旧引用须报告，不能当作相同节点或空依赖。所有 36 个算子分别给数学定义、有效域、损失、保持声明及计算边界；“定义存在”不等于可判。

个人目标 \(g\) 须明确其结构目标 \(G\) 及能力或评价要求；外部模型从 \(s\) 提出的可用背景 \(B\) 也须保留确认依据或未决标记。把 \(B\) 用作结构路线的入口，是接受这些声明条件，不是公共本体已经证明该学习者具备它们。

\textbf{事件与次序（接口义务，含定义约定）。} 同一行动可以多次执行，同一节点可以反复学习；每次执行使用不同的事件标识。先选择 AND/OR 方案及输入来源，再加入策略次序；\textbf{检查合并后的事件图无环，通过后才取自反传递闭包形成偏序}，否则拒绝该候选。不能只检查来源图本身无环。

\textbf{删去条件会怎样（PROOF：有限反例）。} 若来源要求 \(e_{1}\to e_{2}\)，策略要求 \(e_{2}\to e_{1}\)，二者分别无环，合并却成环。自反传递闭包会同时包含 \(e_{1}\preceq e_{2}\) 和 \(e_{2}\preceq e_{1}\)，而 \(e_{1}\ne e_{2}\)，违反反对称性，因而不能称为偏序。

对固定版本 \(M\)，\(\mathrm{Plan}_{M,E,s,g,b,\rho}:\mathrm{View}_M \rightharpoonup \mathrm{FinPos}\) 仍表示到有限事件偏序同构类型的部分映射。完整结果包 \(\Pi\) 保留具体事件集 \(P\)、偏序 \(\preceq\)，以及定义在 \(P\) 上的行动、输入与目标来源、关注节点、排程、外部约束、见证和版本；\(\mathrm{Plan}\) 的值是该包忘掉附加数据后的投影 \([(P,\preceq)]\)。事件改名时，全部附加数据须沿同一双射同步迁移。两条行动完全不同的二事件链可以有同一偏序类型，不能因此把它们的结果包当作同一个计划。

有限搜索须限制候选行动、事件数/时域及需要枚举的标注和策略约束；各候选域能在有限时间完整列出，所有使用条件、联合无环性、预算与外部判据均有可判检查。接口另行区分已找到、附条件、界内无解与未知；只有完整检查该界内所有候选、均失败且无待决项时，才能报告界内无解。Pareto 筛选产生展示子族，不声称保留每个可行候选；有限搜索结论也不证明现实学习成功。

\section{D.6 证据与交付}
证据负责支撑一项具体断言，标签只帮助读者找到相应检查方式。数学定义、公理/建模约定、完整证明、引用定理、经验结果、未知项分开标注，不组成单一可信度全序；不同证据回答不同问题，并非可以互相替代。本稿所说的“可检查”指断言—证据可追溯、机读清单可核对、核验可复现，不是“数学证明已被机器验证”；正文出现“可检查”时按此含义读，两者不混用。

每项关键断言须能查到：所说的结论、使用的假设、适用范围、证据或引用位置、核验方式与实际状态，以及相关版本。\(\mathrm{PROOF}\) 指向可核对步骤的论证；\(\mathrm{REF}\) 说明引用结果的假设怎样匹配本稿；\(\mathrm{FINITE}\) 说明实际检查的输入域和边界；经验判断另给观测或实验范围。未知和尚未完成的检查保留原状态，不以登记或标注本身代替认证。

机读入口为断言—证据清单。\(\mathrm{check}_\mathrm{status}\) 独立使用 \(\mathrm{not}_\mathrm{run} / \mathrm{passed} / \mathrm{failed} / \mathrm{unsupported} / \mathrm{resource}_\mathrm{exhausted}\)；另列正文有证明、具体证书被接受、检查器已形式验证三个不同字段。清单检查发现缺失、悬空和过期证据时失败，不自动判断自然语言论证或文献适用性。

有限模型检查只核对所列实例或有限枚举域，不替代一般证明。引用页码准确不自动证明引用适用，脚本通过不自动证明脚本覆盖了所声称的全部性质。本稿不据此宣称已经机械形式化整部专稿或验证教学效果。

每章用 \texttt{定义 x.y}、\texttt{命题 x.y} 等可索引标题；跨章引用优先给章节与明确记号，避免引用尚未确定的定理编号。旧稿的完成标记不作证据；实际正文核对及修订后的双页码由附录 B 记录。

\textbf{本节的正向保证（可依赖项）。} 撇开否定句，本契约正面承诺五件事：固定版本下，学习者输入的替换不改变 \(M\) 的版本（版本不变式）；每个行动契约的输入使用条件与输出资源可由见证逐项核对（可审计）；路线、先修与可达结论都携带声明范围，未经完整性证明不升格为必要条件（范围性）；计划的结果包保留事件、来源、见证与版本，事件改名时附加数据同步迁移（可复现）；每项关键断言有可追踪的证据入口、状态与版本（可追溯）。超出这五条的结论不在本契约保证之内。

\textbf{用直白语言串起来。} 先保存公共数学档案，另行描述学习者；公共档案给出合法动作，动作及其来源连接成结构路线，路线在声明范围内决定相对先修；再根据个人输入选视图、评价候选并安排事件，最后逐项说明结果的依据和限度。每一步得到的东西都不同，不能把提供材料、完成证明和已经学会当作同一次输出。

\textbf{思想与方法。} 本节的思想是“跨越边界时交代保持义务”：从知识到行动、从行动到路线、从公共路线到个人计划，每次都说明传递了什么，以及证据允许推出什么。具体方法是 D.3 的输出与来源逐项核对、D.4 的出现/支持/必要性分辨、D.5 的合并图反例检验和 D.6 的断言—证据追溯。更高层的满足保持、语义保守与路线输送仍须在对应章节证明，不能由接口名称直接取得。

\textbf{回看问题。} 一个行动说自己“产出了某节点”时，究竟增加了材料、证明、使用条件，还是对学习者能力的判断？哪一条见证支持这个解释？如果暂时答不清，就保留该处未决，不继续把它当作后续结论的已知前提。

\clearpage\chapter{01 元理论、表达范围与可检查性}
\mcsorigin{原章节：01-元理论与可构造宇宙.md}
本章说明 MCS 用什么数学语言描述知识，以及每项证明与计算承诺依赖什么条件。元理论采用 ZFC；有限描述、有限前提生成和有界搜索各有不同的能力范围。文件名保留“可构造宇宙”以便追溯，相关比较集中在 1.4 的历史辨析，不进入 MCS 的推导链。

\textbf{从一条结论开始（ILLUSTRATION）。} 我们希望保存群的左消去律：在群中，\(a\cdot b=a\cdot c\) 蕴涵 \(b=c\)。先用语言写出结论，注明群公理；整数加法群、实数加法群等结构解释这些符号；公共本体保存结论与证明记录。学习者甲已经会使用逆元，乙还不会，因此两人的练习安排可以不同，而同一版本中的证明记录保持不变。这个例子解释分工，不声称相应安排已验证有效。

这里有两条动机：数学上要分清“写出了什么、在哪些结构中成立、凭什么证明”；学习应用上要让个人路线变化，同时使公共结论可追溯。设计由一条知识记录出发，补上语言、背景与证据，再接入外部状态和派生结果。本章先用 1.1 区分几种保证，1.2 固定各类对象的角色，1.3 说明描述与生成的范围，1.5 汇总实际可检查的事项；1.4 可在主线之后阅读。

\section{1.1 基础约定、两种完备性与判定性}
\textbf{约定 1.1（元理论，DEF）。} 全稿在 ZFC 中进行。自然数、有限字符串、有限树、函数、关系、商集及集合大小模型均为元理论中的集合对象。ZFC 是工作前提，不是本稿已经证明一致或唯一正确的体系。它使我们可以共同描述公式、证明、模型与档案；元层假设不自动成为对象理论的公理。

相对一致性主张须分别给出证明所用的元理论、前提理论 \(S\) 与目标理论 \(T\)，例如 \(\mathrm{ZFC}\vdash(\mathrm{Con}(S)\Rightarrow\mathrm{Con}(T))\)，一致性按指定演算及公理呈现编码。“存在传递集合模型”等更强假设不能改写成单纯一致性；尚未证明的蕴含保留为条件性主张或开放问题。实际使用的附加前提按第 15 章逐条登记。

\textbf{定义 1.2（两种完备性与判定性，DEF）。} 对指定演算 \(\vdash\)、模型类 \(\mathcal K\)、理论 \(\mathcal T\) 和前提集 \(\Gamma\)，区分以下三个问题：

\begin{itemize}
\item \textbf{逻辑完备性：所有语义后果是否都有证明？} 强形式为：若 \(\mathcal T\cup\Gamma\models_{\mathcal K}\varphi\)，则 \(\Gamma\vdash_{\mathcal T}\varphi\)，其中允许任意前提集。\(\mathcal T=\Gamma=\varnothing\) 的情形称弱完备性，即“有效式皆可证”。模型类必须明确；第 03 章定理 3.9 与第 06 章的联合前提演绎使用相应强形式。
\item \textbf{理论的句子完备性：理论能否决定每个句子的正反？} 对每个句子 \(\varphi\)，都有 \(\mathcal T\vdash\varphi\) 或 \(\mathcal T\vdash\neg\varphi\)。一般另要求理论一致，排除爆炸造成的平凡性。
\item \textbf{判定性：指定问题是否有总能结束的正确算法？} 存在对每个输入都终止并正确判定该性质的算法。它是计算性质，不是第三种完备性；可枚举证明不等于可判定可证性。
\end{itemize}
\textbf{可靠性（DEF）} 则要求：若 \(\Gamma\vdash_{\mathcal T}\varphi\)，则 \(\mathcal T\cup\Gamma\models_{\mathcal K}\varphi\)。它与逻辑完备性分别给出“可证”与“语义后承”之间的两个方向，不与句子完备性或判定性构成共同的对偶。对群的例子，可靠性保证按群公理正确证明的左消去律在群模型中成立；它不保证任何关于群的问题都能判定。指定演算的可靠性与 Henkin 完备性见第 03 章定理 3.9（PROOF+REF）；一般概念的外部核对入口见本章末的 Schimmerling 教材。

本稿不要求背景理论句子完备。凡称“可判定”，都须说明判定的是类型、证书、定理资格还是有界候选，不说“整个框架可判定”。“完整框架”另指节点、关系、视图、认知接口、路径和宏观层级获得定义及验证义务；它是覆盖范围声明，不是逻辑完备性，也不表示已经穷尽数学或认知机制。各项状态与验收映射见第 15 章。

\textbf{范畴性的边界（REF+PROOF）。} 含无限模型的一阶理论不能在不限制模型基数的情况下，把全部模型确定为同一个同构类型；Löwenheim–Skolem 定理给出不同大小的模型。但它可以在每个无限基数上分别范畴。“在基数 \(\kappa\) 上范畴”指该大小的模型恰有一个同构类型。例如，纯等号语言中的公理“至少有 \(n\) 个元素”（\(n\) 遍历正整数）恰好规定无限集合；同基数集合之间任意双射都是同构，因而在每个无限基数上分别范畴。反例的核对只需上述双射论证，文献对照见章末 Gallinaro 例 2.2.3；不限制基数时的结论采用 Weiss–D’Mello 定理 6：对一个无限模型，取大于其大小且不小于语言大小的基数，即可获得该大小的初等扩张。这段用于防止混淆，不是 MCS 的推导规则。

\section{1.2 语言、模型、本体与外部状态}
先分清两条关系：语言与数学模型通过解释连接；公共本体、个人状态与派生结果按数据职责分工。元理论同时描述这些对象，不把它们排成一条统一的抽象阶梯。

\textbf{定义 1.3（四类角色，DEF）。}

\begin{enumerate}
\item \textbf{对象语言} \(\mathcal L_\Sigma\) 规定怎样表达。有效可数签名 \(\Sigma\) 声明基本类型（至少含真值 \(o\) 与个体 \(\iota\)）和非逻辑常元类型。项由变量、常元、应用与 \(\lambda\) 抽象生成；公式为 \(o\) 型项，句子为闭公式。形成、绑定与替换见第 02 章定义 2.1–2.9。认证证明为有限语法对象，其规则定义在 \(\lambda\) 提升后的多排序一阶语言 \(\mathcal F_\Sigma\) 上（定义 2.10–2.13）。
\item \textbf{语义结构} \(\mathfrak A\) 解释这些表达。每个类型 \(\sigma\) 有非空域 \(D_\sigma\)，其中 \(D_o=\{0,1\}\)；常元有指定解释，应用与抽象有相应求值规则。\(\mathfrak A,v\models\varphi\) 表示 \(\llbracket\varphi\rrbracket_v=1\)，满足关系在元理论中定义。本稿只对集合大小结构讨论满足。
\item \textbf{公共本体} \(\mathcal M\) 保存给定版本的节点、表征、证据、行动契约、通用关系与模板，并有签名与背景公理槽（D.1）。它是一份数学内容档案，不等于其中某个领域理论的模型。左消去律的节点不是某个群元素，节点引用与群元素的解释须分型。
\item \textbf{外部输入} \(\mathcal E,s,g,b,\rho\) 分别给出认知模型、状态、目标、预算与偏好，参与产生 \(D(\mathcal M,\mathcal E,s,g,b,\rho)\)。典型流程是适配器生成 \(\lambda\)，局部化给出 \(\operatorname{Loc}_\lambda(\mathcal M)=(V,\partial V,\tau,Pres)\)，规划在其定义域内返回 \(\operatorname{Plan}_{\mathcal M,\mathcal E,s,g,b,\rho}(V)\in\mathrm{FinPos}\)。该值仅为事件偏序的同构类型；具体事件、行动、来源、排程与外部约束保存在完整结果包中。无合法结果或检查未决时不能强行填入一个偏序。适配器读取本体的方式、版本相容及结果状态遵守 D.5。
\end{enumerate}
视图与路线是派生结果的分量，不是上式额外的外部输入，也不自动写回 \(\mathcal M\)。是否把派生结果另编号为一层不影响理论；重要的是其来源、依赖与更新权限。\(\Xi\vdash_\Sigma t:\sigma\) 表示类型判断，\(\Gamma\vdash_T\varphi\) 表示推导，元层另检查有限证书；这些判断不能仅凭共用符号而混同。字体只辅助辨认，不替代类型与角色声明。

本稿中的“模型”须按用途读取：

\begin{itemize}
\item \(\mathfrak A\) 是解释对象语言的结构。箭头域为真实函数集并满足抽象闭合等条件时采用 Henkin 语义；进一步取全函数集时采用标准语义。标准模型构成 Henkin 模型的子类，\(\models_H\) 与 \(\models_{\mathrm{std}}\) 不能互换（第 03 章）。
\item 某结构要成为 MCS 结构理论的模型，还须满足单独列明的结构公理。它对节点和行动类型的解释，不等于群论模型对群元素的解释。具体软件声称实现某片段时，须另履行附录 E 的义务。
\item \textbf{外部认知模型} \(\mathcal E\) 是状态空间与接口规格，不是上述意义的语义模型。它给出 \(X_{\mathcal E}\)、条目与结果域、外部行动域，以及观察、转移、评价接口和版本（定义 7.8）。它允许学习、遗忘与修正，评价可含未知和随机性；个体记录与掌握估计保存在该侧。接口只读公共本体，软判断不进入数学演绎闭包，相关性不产生公共硬先修（第 07 章及 D.1）。
\end{itemize}
\textbf{定义 1.4（大小约定，DEF）。} 固定的签名、理论、节点集、证据集和一份本体版本均为集合。“所有集合大小模型”可以作为元层的类讨论，不作为集合类型的一个元素；有关范畴采用局部小范畴或明确限定的集合大小片段。模型类 \(\mathcal K\) 为真类时，相应语义声明须以 ZFC 中可表达的逐模型断言或明确的模式给出。不假定有“所有数学对象组成的集合”，本稿不需要 Grothendieck 宇宙或不可达基数。

\(\mathcal M\) 在元理论中是集合记录；其可内部化部分另以 Node、Code、CertCode 等类型编码。内部编码只解释被编码部分，不自动得到覆盖自身语言的全局真理谓词。

\textbf{命题 1.5（外置的基本非干扰，PROOF）。} 固定本体版本 \(\mathcal M\)，令 \(\bar a=\mathsf{pub}(\mathcal M)\) 为定义 4.1 的十二个公开坐标。固定一个集合论公式 \(\theta(\bar x,y)\)，展开辅助定义并显式列出参数后，假设


\mcsdisplay{\operatorname{FV}(\theta)\subseteq\{x_1,\ldots,x_k,y\},
\qquad \exists!y\,\theta(\bar a,y).}
\noindent 其中 \(\operatorname{FV}\) 为自由变元集，约束变量不计入；\(\theta\) 及其解释、公开参数均不随外部输入更换。令 \(h(\mathcal M)\) 为上述唯一输出。则替换外部输入 \((\mathcal E,s)\) 不改变 \(h(\mathcal M)\)。唯一性只要求在当前 \(\bar a\) 成立，不承诺 \(h\) 在所有本体版本上都有定义。

\textbf{证明概括。} 两次调用代入相同公开参数，得到同一个定义条件；条件的唯一解也相同。

\textbf{证明。} 固定当前公式 \(\theta\)。展开所有辅助定义后，外部输入不是它的自由参数。两次调用均使用同一式子 \(\theta(\bar a,y)\)，因而给出同一组候选 \(y\)；唯一性保证两次输出相同。这里对每个固定公式分别陈述此论证，不在 ZFC 内部引入对全体元理论公式求真的函数。\(\square\)

\textbf{条件与检查。} 若允许额外自由参数，例如 \(\theta\equiv(y=s)\)，更换状态便可改变输出。若删去唯一性，例如只要求 \(y\in\{0,1\}\)，可选结果的集合仍不变，但某个依赖个人状态的选择器可能返回不同元素。相反，\((y=x_1)\land\forall s(s=s)\) 并不依赖外部状态：这里的 \(s\) 是约束变量，换名不改变意义。禁止某个约束变量的拼写不是非干扰条件。

实际核对要展开辅助定义，追踪公开参数以外的数据和接口依赖。只把量词范围限定在节点集内仍不够，例如条件 \(s\in n\) 依旧读取自由参数 \(s\)。\(\operatorname{Eval}_{\mathcal E}\)、\(\operatorname{Obs}_{\mathcal E}\)、\(\operatorname{Trans}_{\mathcal E}\) 的调用若依赖外部输入，便不能直接套用本命题；若认为其结果实际上不变，需要另给消除该依赖的论证。

本命题给出定义层的基本不变性，不证明软件已实现隔离，不比较不同集合论模型，也不保证不同公共版本的输出相同。V09 的哈希检查只核验指定实现夹具。个体成本、适配视图与学习效果可以随状态变化，属于派生侧；群例子中变化的是路线，而非固定档案中的证明记录。

\section{1.3 有效描述与有限生成}
先分开三个问题：有多少个描述，能否有效列出描述，以及规则能产生哪些结果。可数性回答第一个问题，编码与枚举条件回答第二个，闭包条件回答第三个。

\textbf{约定 1.6（有效签名，DEF）。} 自动生成部分采用有效呈现的签名 \(\Sigma\)：符号有一一自然数编码，合法符号、基本类型和常元声明可判，并据此检查有限语法。仅递归可枚举的材料须另附在指定枚举过程中的有限输出见证，引用检查才可据此终止。非有效理论仍可语义研究，但不附带算法承诺。

\textbf{可弱化的范围。} 正文默认采用符号与声明可判的强形式。若整体退到符号和声明仅可枚举、每次引用都携带见证，则对已给见证的形成检查仍可判；对任意裸代码的良构判断不再作同样承诺。\(\alpha\) 等价与无捕获替换对已给有限项仍可计算。证书检查还须满足公理、开放假设等引用有可判依据或有限见证的条件（命题 2.14），不能仅凭签名有效便断言所有检查终止。

\textbf{命题 1.7（有限描述的计数边界，PROOF+REF）。} 设 \(\Sigma\) 为带指定单射编码的可数签名；描述、公式及证明树中的全部原子字段也来自可编码材料，不夹带任意实数或集合参数。令 \(\mathrm{Term}_\Sigma,\mathrm{Fml}_\Sigma,\mathrm{Prf}_\Sigma,\mathrm{Desc}_\Sigma\) 分别为有限项、有限公式、有限证明树与有限节点描述的集合。则其带类别标签的不交并存在单射


\mcsdisplay{e:\ \mathrm{Term}_\Sigma\sqcup\mathrm{Fml}_\Sigma\sqcup\mathrm{Prf}_\Sigma\sqcup\mathrm{Desc}_\Sigma
\hookrightarrow\mathbb N^{<\omega},
\qquad \mathbb N^{<\omega}=\bigcup_{k<\omega}\mathbb N^k.}
\noindent 因此四类对象至多可数，不存在从 \(\mathrm{Desc}_\Sigma\) 到 \(\mathbb R\) 的满射。任何以描述集为定义域的集合函数，其值域仍为集合，也不能覆盖全体集合。这里不要求编码可计算；算法承诺另由约定 1.6 及具体枚举条件提供。

\textbf{证明概括。} 把每个有限对象编码成有限自然数序列，再将这些序列分成可依次列出的有限层。

\textbf{证明。} 给原子字段、对象类别、树的分支及边界分别编码，递归序列化有限形成树，得到可唯一恢复原对象的编码。对序列 \((a_1,\ldots,a_k)\) 按 \(k+\sum_i a_i\) 分层，层内按长度与字典序排列；每层有限，每个序列恰在一层出现，因此 \(\mathbb N^{<\omega}\) 可数。此步已有显式枚举，不需另用选择原则；也不能先耗尽长度一的序列再处理长度二。结合 Cantor 的实数不可数性，得到不能满射命名全部实数。集合函数的像为集合，而全体集合不构成集合，得到最后一项边界。\(\square\)

编码与有限序列的计数论证已在正文给出（PROOF）；实数不可数性采用经典结果（REF），沿用附录 B 的 R07-b 定位及其采用范围。结论限于固定可数编码材料：若为每个实数加入独立常元，符号集就不可数，计数前提失效。另一方面，一个“实数”概念节点和量词即可谈论不可数域，无须逐个命名其成员。群公理的有限表达同样可以解释在不可数的实数加法群中。

\textbf{命题 1.8（有限前提生成在 \(\omega\) 达到闭包，PROOF）。} 固定集合 \(U\)、种子 \(S\subseteq U\) 和规则集


\mcsdisplay{R\subseteq[U]^{<\omega}\times U,
\qquad
\mathcal F(X)=\{u\in U:\exists A\subseteq X,\ (A,u)\in R\}.}
\noindent \([U]^{<\omega}\) 表示 \(U\) 的有限子集族；\((A,u)\) 表示已有全部前提 \(A\) 时可以加入 \(u\)。允许零前提规则，也允许无限多个规则。固定规则不因当前资源增加而失效，因此该定义给出单调算子。令


\mcsdisplay{X_0=S,\quad X_{n+1}=X_n\cup\mathcal F(X_n),\quad
X_\omega=\bigcup_{n<\omega}X_n.}
\noindent 则 \(X_\omega\) 是包含 \(S\) 的最小 \(\mathcal F\)-闭集合，其中闭合指 \(\mathcal F(X)\subseteq X\)。

\textbf{证明概括。} 有限前提总能在某个共同的有限阶段凑齐，所以并集已经闭合；任何包含种子的闭集又必须包含每一步生成结果。

\textbf{证明。} 若 \(u\in\mathcal F(X_\omega)\)，取有限 \(A\subseteq X_\omega\) 使 \((A,u)\in R\)。各前提出现阶段的最大值给出 \(A\subseteq X_n\)；\(A=\varnothing\) 时取 \(n=0\)。于是 \(u\in X_{n+1}\subseteq X_\omega\)。若 \(Y\supseteq S\) 且 \(\mathcal F(Y)\subseteq Y\)，由单调性归纳得到每个 \(X_n\subseteq Y\)，故 \(X_\omega\subseteq Y\)。\(\square\)

\textbf{有限前提为何必要（PROOF：反例）。} 在 \(U=\mathbb N\cup\{\ast\}\) 中，从空集出发，始终允许加入 \(0\)，有 \(n\) 就加入 \(n+1\)，另规定“所有自然数都已出现才加入 \(\ast\)”。该算子单调，但最后一条有无限前提，超出上述 \(R\) 的形式。有限阶段得到 \(X_n=\{0,\ldots,n-1\}\)，并集是 \(\mathbb N\)，此时才可生成 \(\ast\)，故 \(X_\omega\) 未闭合。单调性不能替代有限前提。

若推广到任意算子而放弃单调性，最小性也会失败：在 \(\{a,b\}\) 上令 \(\mathcal F(\varnothing)=\{a\}\)、\(\mathcal F(\{a\})=\{b\}\)，其余取空集。从空种子迭代得到 \(\{a,b\}\)，但 \(\{b\}\) 已是包含种子的闭集。这个反例超出了固定正向规则形式；该形式本身已经蕴含单调性。

\textbf{数学闭包与算法。} \(X_\omega\) 的构造不保证任何一个有限阶段已收齐全部结果，也不保证每个阶段都能在有限时间内算完。只有给定有效编码、种子与规则的有效枚举、可判的有限形成见证等条件，才可由公平枚举半判定成员资格；非成员输入仍可能一直等待。取不可递归枚举的 \(S\subseteq\mathbb N\) 和空规则集，就有 \(X_\omega=S\)，说明计数和闭包存在都不自动带来有效枚举。每个已生成对象有有限形成深度，不等于该深度可对任意输入判定。自动枚举的具体条件见第 04 章。

\section{1.4 历史辨析：可构造宇宙的类比限度}
这一节可跳读。旧稿将逐阶段形成的知识结构与 Gödel 的可构造宇宙 \(L\) 类比；两者有阶段累积的外观相似，但不足以转移节点编号、计算或学习难度的结论。以下保留识别这项类比所需的背景，完整文献定位集中在附录 B。

\textbf{定义 1.9（可构造层级，DEF：引用背景）。} \(\operatorname{Def}(A)\) 是在 \((A,\in)\) 中由一阶公式和来自 \(A\) 的有限参数定义的子集全体。令


\mcsdisplay{L_0=\varnothing,\qquad L_{\alpha+1}=\operatorname{Def}(L_\alpha),\qquad
L_\lambda=\bigcup_{\alpha<\lambda}L_\alpha\quad(\lambda\text{ 为非零极限序数}),
\qquad L=\bigcup_{\alpha\in\mathrm{Ord}}L_\alpha.}
\noindent 每个阶段为集合，末项为类记号。阶段的集合大小满足关系在元理论中定义；初始空阶段约定 \(\operatorname{Def}(\varnothing)=\{\varnothing\}\)。这不是有限前提规则的 \(\omega\) 步闭包，不能由命题 1.8 替代。

\textbf{引用事实 1.10（内模型背景，REF）。} \(L\) 是包含全部序数的传递真类，并为 ZFC 的内模型；“满足 ZFC”按每条公理的相对化分别解释。这里不把它当作一个集合大小模型，也不用它证明 MCS 的后续结论。所引背景为 Manin IV.1–3；其集合运算构造与本节定义不预设逐阶段相同。

\textbf{反例 1.11（可定义不等于可计算，REF）。} 标准自然数结构中可定义“程序 \(e\) 在输入 \(e\) 上停机”，但不能统一判定全部实例（可计算性来源见附录 B 的 R02-d）。因此，从 \(L\) 的良序可定义，不能直接得到关于其表示的比较或成员资格算法；有效输入表示和检查程序必须另给。这一计算障碍与命题 1.7 的计数边界不同：程序本来就可有效编号，不可判定性仍然存在。

\textbf{边界声明 1.12（本稿的采用范围，NOT-CLAIMED）。} 本稿不假设 \(V=L\)。公共内容由签名和背景理论明确声明；预算和事件界由外部输入给出，不由 GCH 估算；语法编码与替换由第 02 章的规则支撑，不借凝结引理证明；个人状态变更下的非干扰由显式依赖条件支撑，不借集合论绝对性保证。未来若比较不同基础，须另核对参数、模型关系、公式复杂度与量词域。

\section{1.5 可检查性与真理层次}
给出一份证明后逐步核对，和寻找是否存在证明，是不同任务。核对一个计划，与断言整个候选范围内都无解，也需要不同的保证。

\begingroup\small\setlength{\tabcolsep}{4pt}
\begin{longtable}{>{\raggedright\arraybackslash}p{\dimexpr 0.2427\linewidth-2\tabcolsep\relax}>{\raggedright\arraybackslash}p{\dimexpr 0.1824\linewidth-2\tabcolsep\relax}>{\raggedright\arraybackslash}p{\dimexpr 0.2555\linewidth-2\tabcolsep\relax}>{\raggedright\arraybackslash}p{\dimexpr 0.3194\linewidth-2\tabcolsep\relax}}
\toprule
\textbf{判断} & \textbf{归属} & \textbf{承诺} & \textbf{限制} \\
\midrule
\endfirsthead
\toprule
\textbf{判断} & \textbf{归属} & \textbf{承诺} & \textbf{限制} \\
\midrule
\endhead
\bottomrule
\endfoot
\bottomrule
\endlastfoot
类型、良构、无捕获替换 & 元层与对象语言 & 有限递归检查与计算 & 有效签名、显式类型 \\
给定有限证明正确 & 元层（证书检查） & 可判证书检查 & 规则有效；实际开放假设有限显式；公理或更大假设族的引用资格可判或附有限见证 \\
是否存在证明 & 元层（证明搜索） & 半判定 & 有效背景理论与证明枚举；无证明时一般不保证终止 \\
定义保守 & 元层（理论扩张） & 对指定定义模式证明 & 不承诺任意自然语言定义的通用自动判定 \\
全部标准高阶有效式 & 语义层 & 不承诺有效完备枚举 & 第 03 章给出相应语义限度 \\
真实认知状态与教学收益 & \(\mathcal E\) 侧 & 不属本体判定 & 另需外部模型与证据；允许未知及不确定性 \\
有界规划的有限穷举 & \(D\) 侧 & 所声明候选域内终止 & 所有参与搜索的行动、事件、参数、排程、标注和见证候选均可有效有限枚举；各候选检查均可判 \\
有限结构实例核验 & 元层脚本 & 所列实例或有限域的核验（FINITE） & 不替代一般证明，范围见核验报告 \\
机器内核认证 & 实现层 & 不承诺（NOT-CLAIMED） & 需另实现或接入可信证明内核 \\
\end{longtable}
\endgroup

各行的正式依据分别见命题 2.9、2.14，定理 3.9 及证明枚举，命题 3.13 与定义 3.14，推论 3.10–3.11，第 07 章，定理 12.13 及 D.5，第 15 章和第 16 章。本表是条件与检查入口，不表示所有程序已经实现。有限规划只有完整核验所声明候选域且没有未决项，才能报告界内无解；有限域外仍不作断言。

\textbf{为何还要限制参数域（PROOF：归约，采用停机不可判定性）。} 只设一个行动、一次事件，但允许无上界的参数 \(n\in\mathbb N\)，并要求“程序 \(e\) 在输入 \(e\) 上于 \(n\) 步内停机”。每个候选只需模拟有限步，检查可判；存在合法候选却等价于该程序最终停机。若仅凭行动数、事件数有限和逐候选检查可判，就能对所有程序编码 \(e\) 判定有无方案，便可判定停机问题，与反例 1.11 的引用结果矛盾。有限穷举因此还须控制全部候选数据域；特殊无限域若有其他完备判定程序，应另证明其适用性。

\textbf{命题 1.13（真理与证明，PROOF+REF）。} “\(p\) 是 \(\mathcal T\) 中 \(\varphi\) 的证明”是关于有限代码的关系，其有效检查须满足上表条件；“\(\varphi\) 在结构 \(\mathfrak A\) 与赋值 \(v\) 下真”则是满足关系 \(\mathfrak A,v\models\varphi\)。可靠性给出


\mcsdisplay{\Gamma\vdash_{\mathcal T}\varphi
\quad\Longrightarrow\quad
\mathcal T\cup\Gamma\models_{\mathcal K}\varphi.}
\noindent Henkin 语义下，反方向在定理 3.9 的条件下成立；它使用所有允许的 Henkin 模型，不是某一个预期模型。在有效背景下，证明枚举能在有证明时最终找到证明，无证明时一般不保证停止。标准语义下，推论 3.10–3.11 进一步限制有效演算的覆盖能力：不存在覆盖全部标准有效式的有效可靠完备演算。这不表示标准语义的后承较弱；由于标准模型是 Henkin 模型的子类，Henkin 后承必蕴涵标准后承。

\textbf{论证与引用范围。} 对 \(\mathcal T\) 的模型及满足 \(\Gamma\) 的赋值，可靠性把证明连接到真；反向的完备性与有效性边界按上述定理分别使用。Tarski 不可定义性不否认元理论为集合大小结构定义满足；Gödel 第二不完全性也须带有效性、足够强度、一致性和适当可证性谓词等假设，不能概括为“任何体系不能证明自己一致”。这些一般限度作为引用结果采用，定位与采用范围见附录 B 的 R02-b；本章没有给出其完整证明。冯琦第 13.3 节关于定义扩充的论述只支持相应条件核对，不替代这些元定理。

\textbf{思想与方法。} 本章的思想是“每项保证都对应一个明确问题”：固定讨论对象，再交代前提，最后说明证据与算法覆盖到哪里。具体方法是 1.1 的判断分类、1.2 的自由参数与接口依赖审计、1.3 的有限前提反例检验，以及 1.5 对候选域和终止条件的逐项核对。更高层的语义与保守性结果必须在这些边界内使用。

\textbf{回看问题。} 当一条结论说“可以生成、可以证明、可以检查”时，输入是什么，输出是什么，是否保证结束？换一个学习者、公共版本或模型类后，哪项结论仍有依据，哪项需要重验？定义、证明、引用、有限核验和未知状态分别记录，不把它们排成单一可信度全序。

\textbf{参考与核验入口。} 本章新增核对仅使用下列明确章节或例子；沿用文献的版次、页码与已核范围见附录 B，不把定位信息当作本轮已重读全文的记录。

\begin{itemize}
\item Ernest Schimmerling，在线逻辑教材，第 3 章：可靠性、理论完备性与逻辑完备性；本章使用其相应定义与定理陈述。\href{https://www.math.cmu.edu/\textasciitilde{}eschimme/300Book/300-Page-3.html}{章节}
\item William Weiss、Cherie D’Mello，Fundamentals of Model Theory，定理 6，印刷第 20 页／PDF 第 22 页：向上 Löwenheim–Skolem 定理；本章核对其假设与结论。\href{https://people.math.sc.edu/mcnulty/modeltheory/WeissDMello.pdf}{讲义}
\item Francesco Paolo Gallinaro，Around Exponential-Algebraic Closedness，例 2.2.3：纯等号无限集理论在每个无限基数上分别范畴。\href{https://etheses.whiterose.ac.uk/id/eprint/31077/1/Gallinaro\_FP\_Mathematics\_PhD\_2022.pdf}{学位论文}
\item Manin IV.1–3 的可构造宇宙背景，以及冯琦的计数、可计算性与定义扩充相关定位，统一见附录 B。其中定义扩充的局部核对不替代 Tarski 或 Gödel 元定理的证明。
\end{itemize}
\clearpage\chapter{02 对象语言与可检查的经典外延 HOL 演算}
\mcsorigin{原章节：02-对象语言：简单类型论.md}
本章选定一个演算，不在“等式内核”“自然推演”和“高阶语义”之间省略接口。对象语言是带 \(\lambda\) 的经典外延简单类型论；有限认证证书采用第 2.6 节明确规定的多排序自然推演。这是对高阶语言的证明编码，不把关系量化删除为只量化数学个体。

\section{2.1 类型、签名与语法层次}
本节固定对象语言的静态部分（类型、签名、字母表），并划出语法层次：哪些判断属于对象语言、哪些属于元层检查（定义 2.2）。

\textbf{定义 2.1（简单类型、签名与字母表）。} 给定有效可数的基本类型集 \(B\)，含 \(o\)（真值）；本稿默认另含 \(\iota\)（个体），以便讨论个体与关系的一般情形，具体签名可按需省略 \(\iota\)。类型是对箭头构造闭合的\textbf{最小}集合


\mcsdisplay{\mathrm{Ty}=B\cup\{\sigma\to\tau:\sigma,\tau\in\mathrm{Ty}\}.}
\noindent 箭头右结合；最小性支撑后文对类型结构的归纳。本稿的类型构造子只有箭头 \(\to\)，类型层不再引入别的构造，故\textbf{签名} \(\Sigma=(B,(c:\sigma_c)_{c\in C})\) 只由 \(B\) 与一个非逻辑常元集 \(C\) 及其唯一声明类型 \(\sigma_c\in\mathrm{Ty}\) 组成。语言的\textbf{字母表}由三部分并成：逻辑符号 \(\{\top,\bot,\neg,\wedge,\vee,\Rightarrow,=_\sigma,\forall_\sigma,\exists_\sigma:\sigma\in\mathrm{Ty}\}\)、按类型给出的变量、以及常元 \(C\)；三部分都按约定 1.6 有效呈现（逻辑符号由 \(\mathrm{Ty}\) 生成，仍可数；本节取强层，若改用见证层则相应判断降为“有见证时可判”）。类型不是项，语言也没有类型层的量化：定义 2.4 的 \(\forall_\sigma,\exists_\sigma\) 是对每个类型 \(\sigma\) 各给一个常元，属于\textbf{模式}（schema）。给定 \(\Sigma\) 即有对象语言 \(\mathcal L_\Sigma\)（定义 1.3）。

每个类型有可数个变量，变量名可按类型有效枚举，且\textbf{变量名与常元符号不重叠}；常元允许任意有限简单类型，不人为限制阶数。非逻辑常元与定义 2.4 的逻辑常元相对，后者是语言内置的真值、连接词、等式与量词。

节点、表达式代码、证书、事件等按一条判据决定是否内部化：\textbf{凡要在对象语言内被量化、被证书检查或作为节点负载出现的成分，设为互异基本类型}（结果即第 04 章与附录 E 的清单，例如 \(\mathsf{Node},\mathsf{Code},\mathsf{Cert},\mathsf{Evt}\)）；纯版本与簿记信息可留在元层。内部编码只解释被编码的部分，不增添覆盖自身语言的全局真理谓词（呼应 1.2 的本体双重身份）。类型名本身不是该类型的元素。若需编码依赖类型、子类型或多态，另给翻译；不暗中把它们加入简单类型核。

\textbf{定义 2.2（三种判断与一项元层检查）。} 围绕对象语言区分三种判断：类型判断 \(\Xi\vdash_\Sigma t:\sigma\)、公式推导 \(\Gamma\vdash_T\varphi\)、满足 \(\mathfrak A,v\models\varphi\)。前两者是关于有限语法及推导的判断，满足是元理论中定义的语义关系；这些判断本身都不是对象语言中的 HOL 公式。此外，类型判断与公式推导的有限证书可在元层检查；这不承诺给任意模型中的满足判断配备可判证书。类型语境 \(\Xi\) 与公式假设集 \(\Gamma\) 不混用；\(\Gamma\) 是开放假设，背景公理 \(\mathcal T\) 另行给出（用法见第 03 章定理 3.9 的 \(\mathcal T\cup\Gamma\)）。

\section{2.2 项、函数、谓词和公式}
\textbf{定义 2.3（项形成）。} 固定签名 \(\Sigma\)，以下省略下标，记 \(\Xi\vdash t:\sigma\)（需要强调签名时写作 \(\Xi\vdash_\Sigma t:\sigma\)）；类型语境 \(\Xi\) 是\textbf{有限}声明集合，其成员资格可判（命题 2.9 的可判性以此为前件）。


\mcsdisplay{\frac{x:\sigma\in\Xi}{\Xi\vdash x:\sigma}\quad
\frac{c:\sigma\in\Sigma}{\Xi\vdash c:\sigma}\quad
\frac{\Xi\vdash f:\sigma\to\tau\quad\Xi\vdash u:\sigma}
{\Xi\vdash fu:\tau}\quad
\frac{\Xi,x:\sigma\vdash t:\tau}
{\Xi\vdash\lambda x:\sigma.t:\sigma\to\tau}.}
\noindent 四条规则依次称为 \textbf{Var（变量）、Const（常元）、App（应用）、Abs（抽象）}。项集是上述四条规则生成的\textbf{最小}集合；由此可对项的形成树归纳，命题 2.8 的证明即用此原理。公式就是 \(o\) 型项，闭公式称句子；\(k\) 元谓词是类型 \(\sigma_1\to\cdots\to\sigma_k\to o\) 的项，可以是常元、变量或 \(\lambda\) 表达式。函数与谓词不另设语法范畴，只是值域是否为 \(o\) 之别。\textbf{解析约定}：应用左结合且比 \(\lambda\) 结合更紧，\(\lambda\) 的体向右延伸至不能延伸；由此良构串的解析树唯一（命题 2.9 的解析步骤以此为前件）。

\textbf{定义 2.4（原始逻辑常元）。}


\mcsdisplay{\top,\bot:o;\quad \neg:o\to o;\quad
\land,\lor,\Rightarrow:o\to o\to o;}

\mcsdisplay{=_\sigma:\sigma\to\sigma\to o;\qquad
\forall_\sigma,\exists_\sigma:(\sigma\to o)\to o
\quad(\sigma\in\mathrm{Ty}).}
\noindent \(p\leftrightarrow q\) 缩写为 \((p\Rightarrow q)\land(q\Rightarrow p)\)；\(\forall x:\sigma.\varphi\) 缩写为\(\forall_\sigma(\lambda x:\sigma.\varphi)\)，存在量词同理。本稿的官方语法以定义 2.1 的字母表为准：这些缩写须可按上述规则消去，认证与检查在消去后的形式上核对。逻辑常元的经典含义由第 2.5 节的公理（证明侧）与第 03 章的语义解释（定义 3.3）共同确定，不能仅凭 \(\beta\eta\) 等式宣布经典逻辑成立。

\textbf{例 2.5（关系和关系之间的量化）。} 令\(\mathsf{Rel}=\mathsf{Node}\to\mathsf{Node}\to o\)，则


\mcsdisplay{\forall R:\mathsf{Rel}.\;
\bigl[\forall x,y,z:\mathsf{Node}.\,
Rxy\land Ryz\Rightarrow Rxz\bigr]
\Rightarrow \operatorname{Transitive}(R)}
\noindent 良构，其中 \(\operatorname{Transitive}:\mathsf{Rel}\to o\) 可定义为方括号内的 \(\lambda\) 抽象；按该定义展开后此式是重言式——例子主张的是\textbf{良构性}，不是非平凡内容。对照：若把关系作用于谓词，如 \(R\,P\)（\(R:\mathsf{Rel}\)、\(P:\mathsf{Rel}\to o\)），第一个参数要求 \(\mathsf{Node}\) 而 \(P\) 不是该类型，故\textbf{不良构}。更高一层可量化 \(Q:\mathsf{Rel}\to\mathsf{Rel}\to o\)，表达关系间的包含、复合或证据保持。

另须记住：\textbf{裸谓词的真值不携带证明对象}——证据与证书另有类型（定义 2.10–2.13 与第 04 章的 \(\mathsf{Cert}\)）。

\section{2.3 绑定、替换与有限语法}
绑定名是这套语法里第一个真正的设计问题：名字本身是任意的（\(\lambda x.x\) 与 \(\lambda y.y\) 应当视为同一个模式），而替换又可能\textbf{捕获}自由变量——把 \(u\) 代入 \(\lambda y.t\) 时，若 \(y\) 出现在 \(u\) 中，原意就被改变。本节的三个决定都是为“可检查”服务的：\(\alpha\) 等价用生成式定义固定下来，替换写成确定性的捕获避免运算，比较落到规范绑定代码上。项与证明都是有限串（形成树有限），与 1.3 的有限描述一致。

\textbf{定义 2.6（自由变量与 \(\alpha\) 等价）。} 自由变量集递归定义为


\mcsdisplay{FV(x)=\{x\},\ FV(c)=\varnothing,\quad
FV(tu)=FV(t)\cup FV(u),\quad
FV(\lambda x.t)=FV(t)\setminus\{x\}.}
\noindent 记 \(Vars(t)\) 为 t 中出现的全部变量名（含自由名与绑定名）；其递归与 FV 相同，唯抽象情形为 \(Vars(\lambda x.t)=Vars(t)\cup\{x\}\)。\textbf{\(\alpha\) 等价}是由绑定改名 \(\lambda x.t\sim\lambda y.t\{y/x\}\) \textbf{生成的最小同余}，要求 y 与 x 同类型且 \(y\notin Vars(t)\cup\{x\}\)。其中 \(t\{y/x\}\) 只改 t 中自由出现的 x：变量 x 改成 y，其他变量与常元不变，应用逐支递归，遇到内层 \(\lambda x.u\) 停止进入，其余抽象递归进入。这些出现恰由外层的 \(\lambda x\) 约束；由于新名不曾在 t 中出现，改名不会被内层绑定捕获，定义不依赖 2.7。若只排除自由名并作直接改名，\(\lambda x.\lambda y.x\) 就可能误成 \(\lambda y.\lambda y.y\)，把“返回第一个参数”改成“返回第二个参数”。

为代码与身份检查，实现采用 de Bruijn 索引或等价的规范绑定表示，保留自由变量的名称与类型；证明按命名式书写，实现按规范代码比较——两者相差 \(\alpha\) 等价。

\textbf{定义 2.7（无捕获替换）。} 替换 \(t[u/x]\) 递归定义：

\begin{enumerate}
\item \(x[u/x]=u\)；其他变量与常元不变；
\item \((t_1t_2)[u/x]=(t_1[u/x])(t_2[u/x])\)；
\item \((\lambda x.t)[u/x]=\lambda x.t\)；
\item 若 \(y\ne x\)：当 \(y\notin FV(u)\) 时 \((\lambda y.t)[u/x]=\lambda y.(t[u/x])\)；否则先取与 y 同类型的新变量 \(y'\notin Vars(t)\cup Vars(u)\cup\{x,y\}\)，再令 \((\lambda y.t)[u/x]=\lambda y'.(t\{y'/y\}[u/x])\)。这里先按定义 2.6 改名，再对改名后的体递归替换。
\end{enumerate}
第 4 条的新名取\textbf{代码最小可用者}；替换因此是确定性运算，而不是“任选一个足够新的名字”。尤其必须避开待替换变量 x：在 \((\lambda y.y)[y/x]\) 中，若允许新名恰为 x，就会错误得到 \(\lambda x.y\)，使原来封闭的恒等项出现自由变量 y。上述规则得到 \(\lambda y'.y'\)，与原项 \(\alpha\) 等价。

\textbf{命题 2.8（替换保持类型）。} 若\(\Xi,x:\sigma\vdash t:\tau\) 且 \(\Xi\vdash u:\sigma\)，则\(\Xi\vdash t[u/x]:\tau\)，且替换尊重 \(\alpha\) 等价。

\textbf{证明。} 对 \(t\) 的形成树归纳，四条规则分别对应定义 2.3 的 Var、Const、App、Abs。变量分 \(t=x\) 与其他变量，分别用给定类型判断和原判断；常元不变。应用的两个归纳结论经 App 合成。抽象绑定 \(x\) 时替换不进入；绑定其他变量时，不捕获分支直接用 Abs，捕获分支先取新名 \(y'\)，扩展类型语境后对体应用归纳假设，再用 Abs。两种合法新名所得结果仅相差绑定名，故替换尊重 \(\alpha\) 等价。\(\square\)

\textbf{命题 2.9（语法有效性）。} 固定有效签名 \(\Sigma\)（约定 1.6），记 \(\mathsf{Str}=\mathbb N^{<\omega}\) 为按符号代码写成的有限串集合，\(\mathsf{Ctx}_{\mathrm{fin}}\) 为有限类型语境集合。则下述判定问题与运算都有一个对每个输入都终止的算法，且该算法一致地对 \(\Sigma\) 的代码可用：

\begin{enumerate}
\item \textbf{良构判定}：给定 \(s\in\mathsf{Str}\) 与 \(\Xi\in\mathsf{Ctx}_{\mathrm{fin}}\)，判定是否存在 \(t,\sigma\)，使 \(s\) 按 2.2 的解析约定解析为 \(t\) 且 \(\Xi\vdash t:\sigma\)；
\item \textbf{类型合成}：在 1 判定为良构时，输出该 \(\sigma\)；
\item \textbf{类型唯一}：若 \(\Xi\vdash t:\sigma\) 且 \(\Xi\vdash t:\sigma'\)，则 \(\sigma=\sigma'\)——这正是 2 的合成结果唯一的意思；
\item \textbf{\(\alpha\) 判定与替换}：给定 \(t_1,t_2\) 判定是否 \(t_1\equiv_\alpha t_2\)；给定 \(t,u,x\) 计算 \(t[u/x]\)。
\end{enumerate}
\textbf{证明。} 由 2.2 的解析约定，串的候选解析树唯一，解析是有限语法运算。类型合成对形成树自底向上：Var、Const 查有效签名的表（约定 1.6 保证声明可判）；App 检查函数域与实参类型相等；Abs 在语境中增加一项声明。子项类型按归纳唯一确定，故合成唯一（第 3 条）。\(\alpha\) 判定在规范绑定代码上退化为代码比较。替换按定义 2.7 递归；改名不增加体的树大小，所以递归的主要项严格小于原抽象。禁用变量名集合有限，而同类型变量无限且可有效枚举，按代码查找首个可用新名只需有限步，故替换终止。所有签名查询统一调用所给有效呈现的判定程序，因此算法相对于该呈现一致可用。\(\square\)

这里不需要 \(\beta\eta\) 正规化来决定节点身份（节点身份见第 04 章定义 4.10 的两种相等）。若将 \(\beta\eta\) 正规化作为表达检索工具，必须分别说明\(“\eta\) 缩约正规形”与\(“\eta\) 长形”；二者不能同时含混称为同一归约的正规形。

\section{2.4 \(\lambda\) 提升的多排序认证语言}
\textbf{定义 2.10（认证翻译）。} 本节的翻译把 \(\lambda\) 抽象提升为一阶函数符号，故称 \(\lambda\) 提升；目的是让认证演算落在多排序一阶逻辑里。建立多排序一阶语言 \(\mathcal F_\Sigma\)：每个简单类型 \(\sigma\) 给一个非空排序；每个常元保留其排序；每对类型加入


\mcsdisplay{\operatorname{app}_{\sigma,\tau}:
(\sigma\to\tau)\times\sigma\longrightarrow\tau .}
\noindent 函数类型排序的元素在这里是待解释对象，不是“任意元理论函数”的默认全体：排序只要求非空，箭头域的具体范围由模型给出，即定义 3.1 的 \(D_{\sigma\to\tau}\subseteq D_\tau^{D_\sigma}\) 在语法侧的对应。

对每个 \(\alpha\) 规范抽象 \(a=\lambda x:\sigma.t:\sigma\to\tau\)，\textbf{按变量代码递增次序}列出其自由变量\(\bar y=(y_1:\rho_1,\ldots,y_m:\rho_m)\)，加入一阶函数符号\(\ell_a:\rho_1\times\cdots\times\rho_m\to(\sigma\to\tau)\)。\(m=0\) 时它是常元；不同的 \(\alpha\) 规范抽象取不同的符号（\(a\mapsto\ell_a\) 单射），故检查器可由抽象代码唯一重算该符号。所有此类符号及其公理（定义 2.11 第 6 条，必要时连同第 5 条函数外延）是有效可数模式——这是命题 1.7 的计数边界加约定 1.6 的有效编码的直接推论。

项翻译递归定义为


\mcsdisplay{x^\#=x,\quad c^\#=c,\quad
(fu)^\#=\operatorname{app}(f^\#,u^\#),\quad
a^\#=\ell_a(\bar y).}
\noindent 对 \(o\) 型项 \(p\)，公式翻译为 \(\widehat p:=(p^\#=\top)\)。翻译可计算，且因 \(a\) 取 \(\alpha\) 规范形式而对 \(\alpha\) 等价良定义。

两套逻辑符号须分开读：对象语言的形式符号（含 \(\top,\neg,\wedge,\vee,\Rightarrow,\forall_\sigma,\exists_\sigma,=_\sigma\)）翻译后是 \(\mathcal F_\Sigma\) 各排序上的\textbf{项}；\(\mathcal F_\Sigma\) 自身的连接词、量词与等式（\(\wedge,\vee,\Rightarrow,\neg,\forall,\exists,\leftrightarrow,=\)）只在元层用于写公理与推导。第 2.6 节的推导规则沿用这一区分：认证逻辑的假常元记作 \(\bot_{\mathrm{FOL}}\)，以便与对象语言的 \(\bot\) 区分。

\section{2.5 认证背景公理 \(H_\Sigma\)}
本节把对象逻辑的经典解释与 \(\lambda\) 闭合写成成文条款。没有这些条款，\(\mathcal F_\Sigma\) 只是中性的等式理论；第 03 章的可靠性与完备性以它们为前提。

\textbf{定义 2.11（Boolean、外延与抽象公理）。} \(H_\Sigma\) 由下列一阶句子模式组成，所有未束缚变量均全称闭包。

\begin{enumerate}
\item \textbf{两值}：
\mcsdisplay{\top\ne\bot,\qquad \forall p:o.\;(p=\top\lor p=\bot).}
\item \textbf{逻辑常元的真值}：记 \(B(u)\equiv(u=\top)\)，则
\mcsdisplay{B(\operatorname{app}(\neg,p))\leftrightarrow\neg B(p),}
\mcsdisplay{B(\operatorname{app}(\operatorname{app}(\land,p),q))
\leftrightarrow(B(p)\wedge B(q)).}
对 \(\lor,\Rightarrow\) 分别以一阶析取、蕴含给出同形公理。由两值公理，\(o\) 排序恰有两个元素，故这些双条件唯一确定 \(\neg,\wedge,\vee,\Rightarrow\) 的真值表。
\item \textbf{等式}：
\mcsdisplay{B(\operatorname{app}(\operatorname{app}(=_\sigma,x),y))
\leftrightarrow x=y.}
左边 \(\operatorname{app}(\operatorname{app}(=_\sigma,x),y)\) 中的 \(=_\sigma\) 是对象语言的常元（翻译后是 \(\mathcal F_\Sigma\) 的项），右边的 \(=\) 是 \(\mathcal F_\Sigma\) 自身的等式符号；两者字形相同、身份不同。
\item \textbf{任意类型量词}：
\mcsdisplay{B(\operatorname{app}(\forall_\sigma,P))
\leftrightarrow\forall x:\sigma.\,B(\operatorname{app}(P,x)),}
存在量词将右侧换为 \(\exists x:\sigma\)。
\item \textbf{函数外延}：
\mcsdisplay{\forall f,g:\sigma\to\tau.\,
[\forall x:\sigma.\operatorname{app}(f,x)=\operatorname{app}(g,x)]
\Rightarrow f=g.}
\item \textbf{抽象解释}：对定义 2.10 的每个 \(a=\lambda x:\sigma.t:\sigma\to\tau\)，取 \(\bar y\) 为 \(a\) 的全部自由变量并按变量代码递增排序，即 \(\{y_1,\ldots,y_m\}=FV(t)\setminus\{x\}\)，
\mcsdisplay{\forall\bar y\,\forall x:\sigma.\,
\operatorname{app}(\ell_a(\bar y),x)=t^\# .}
右侧在同一自由变量上下文下读出，即 \(FV(t^\#)\subseteq\{x,\bar y\}\)。嵌套抽象按有限语法树逐层翻译，不产生自指真理定义。
\end{enumerate}
这些模式同时规定命题外延性、函数外延性和 \(\lambda\) 闭合。它们没有选择公理或无穷公理。本节的原始等式只有第 6 条的 \(\beta\)；\(\eta\) 是导出结果（证明见命题 2.15：两边对任意 \(x\) 应用，由抽象公理与函数外延得到相等）。\(H_\Sigma\) 的可满足性由第 03 章的模型构造保证：命题 3.7 把每个 Henkin 模型变为满足 \(H_\Sigma\) 的认证模型，例 3.17 给出非空实例；本稿不在此另证一致性。

\section{2.6 有限认证证明的全部规则}
\textbf{定义 2.12（多排序经典自然推演 ND）。} 判断为 \(A_1,\ldots,A_n\vdash_{\mathrm{ND}}B\)，语境是有限假设集合。规则如下；所有替换须同排序且避免捕获。

\begingroup\small\setlength{\tabcolsep}{4pt}
\begin{longtable}{>{\raggedright\arraybackslash}p{\dimexpr 0.1211\linewidth-2\tabcolsep\relax}>{\raggedright\arraybackslash}p{\dimexpr 0.8789\linewidth-2\tabcolsep\relax}}
\toprule
\textbf{规则} & \textbf{允许的推导} \\
\midrule
\endfirsthead
\toprule
\textbf{规则} & \textbf{允许的推导} \\
\midrule
\endhead
\bottomrule
\endfoot
\bottomrule
\endlastfoot
假设、弱化 & \(A\in\Gamma\Rightarrow\Gamma\vdash A\)；可增加未使用假设 \\
合取 & 从 \(A,B\) 引入 \(A\wedge B\)；从合取消去任一支 \\
析取 & 从任一支引入；从 \(A\vee B\) 及分别假设 \(A,B\) 均得 \(C\) 推出 \(C\)，解除这两个分支假设 \\
蕴含 & 从 \(\Gamma,A\vdash B\) 得 \(\Gamma\vdash A\to B\)；从 \(A\to B,A\) 得 \(B\) \\
假、否定 & \(\bot_{\mathrm{FOL}}\) 推出任意 \(A\)；\(\neg A\) 是 \(A\to\bot_{\mathrm{FOL}}\) 的缩写 \\
经典反证 & 从 \(\Gamma,\neg A\vdash\bot_{\mathrm{FOL}}\) 得 \(\Gamma\vdash A\) \\
全称 & 从 \(\Gamma\vdash A(x)\) 得 \(\Gamma\vdash\forall x A(x)\)，要求 \(x\notin FV(\Gamma)\)；可代入任意同排序项 \\
存在 & 从 \(A(t)\) 得 \(\exists x A(x)\)；从存在式及 \(\Gamma,A(y)\vdash C\) 得 \(C\)，要求新变元 \(y\notin FV(\Gamma,C,\exists xA)\) \\
等式 & \(t=t\)；从 \(t=u\) 和 \(F(t)\) 得 \(F(u)\)，替换指定的同排序自由出现并避免捕获 \\
\end{longtable}
\endgroup

表中每行同时给出引入与消去（\(\forall\) 的代入与 \(\exists\) 的引入并入对应行）。\textbf{证明集是上述规则生成的有限树的最小集合}——证书即这种树，故可对证明树归纳；命题 2.14 的自底向上核对与定理 3.9 的逐规则保真都以此为前件。一阶真常式为 \(\bot_{\mathrm{FOL}}\to\bot_{\mathrm{FOL}}\)。对称、传递、函数同余、双条件均为派生规则。无限理论只允许在一棵证明树中使用有限多条公理。

\textbf{定义 2.13（本稿的 HOL 可证性）。}


\mcsdisplay{\Gamma\vdash_T\varphi
\quad\Longleftrightarrow\quad
H_\Sigma\cup\{\widehat\theta:\theta\in T\}
\cup\{\widehat\gamma:\gamma\in\Gamma\}
\vdash_{\mathrm{ND}}\widehat\varphi.}
\noindent \(T\) 是闭句子的集合；开放公理模式须先明确全称闭包。\(\Gamma\) 可以是无限的开放公式集合，但每份证书只使用其中有限多式，记这个有限子集为 \(\Delta\subseteq\Gamma\) 并保留它们作为未解除假设：\textbf{证书内只在 \(\Delta\) 上工作}，理论层的 \(\Gamma\) 不要求有限。\(\Xi\) 为这些式及结论的自由变量指定类型；不能把局部开放假设自动全称化。证书包含有限 ND 树、实际使用的开放假设 \(\Delta\) 及其成员见证、所用 \(H_\Sigma\) 模式实例和所用理论公理的标识/枚举见证。该定义固定了全文的 \(\vdash\)；正文可用 HOL 的可读派生推理书写，但认证意义始终是这一有限证书。

\textbf{有效回译（定义 2.13 的续段）。} 为使第 05 章能扫描认证证明的概念，固定有效回译 \(r\)。一阶项的变量、原常元原样返回，\(r(\operatorname{app}(s,t))=r(s)r(t)\)，而 \(r(\ell_a(\bar s))=a[r(\bar s)/\bar y]\)，其中 \(\bar y\) 是定义 2.10 保存的自由变量列表，采用同时无捕获替换。一阶等式、连接词和量词分别回译为对应的 HOL 等式、连接词和带类型量词，\(\bot_{\mathrm{FOL}}\) 回译为对象常元 \(\bot\)。这是对有限语法树的递归，保持排序/类型；它把证书表达转换为可供 \(C_{\mathrm{occ}}\) 扫描的 HOL 项，而不把辅助符号 \(\ell_a\) 当成新的领域概念。回译后的语法载体依赖所选编码，不宣称等价证明具有同一出现集。

\textbf{命题 2.14（认证检查终止）。} 若签名有效，实际开放假设 \(\Delta\) 有限且显式，理论公理引用带有可判见证，则给定证书是否证明 \(\Delta\vdash_T\varphi\) 可判。要进一步认证 \(\Gamma\vdash_T\varphi\)，还须使 \(\Delta\subseteq\Gamma\) 的核验可判，例如 \(\Gamma\) 也是有限显式集合，或每项引用携带可判成员见证；任意无限假设集不自动获得这一算法承诺。

\textbf{证明。} 按树自底向上核对结论、排序、替换、假设解除与特征变量条件；这些均为有限语法运算。开放假设与显式 \(\Delta\) 作有限比较；若另声明 \(\Gamma\)，则核验所附成员见证。抽象公理匹配按抽象代码重算其自由变量、排序和体翻译。枚举见证只重放有限步。所有递归均严格沿有限树或给定步数进行。\(\square\)

仅判断“某式恰属一个任意给定理论或假设集”可能不可判，不能从本命题删除见证假设。定理 3.9 对任意 \(\Gamma\) 的数学完备性与这里受有效呈现限制的检查算法是不同结论。

\section{2.7 可读的 HOL 规则及两个附加理论}
\textbf{引理（翻译—替换交换；命题 2.15 的预备）。} 记 \(\#\) 为定义 2.10 的翻译。设 \(\Xi,x:\sigma\vdash t:\tau\) 且 \(\Xi\vdash u:\sigma\)，即 u 与被替换变量 x 同类型，则等式


\mcsdisplay{(t[u/x])^{\#}=t^{\#}[u^{\#}/x]}
\noindent 可在一阶等式演算中由 \(H_\Sigma\) 的第 5、6 条证明；两侧均为 \(\tau\) 排序项，替换按定义 2.7 避免捕获。这是可证相等，因而在所有 \(H_\Sigma\) 模型中相等，一般不是字面相同。

\textbf{证明。} 对 \(t\) 的形成树归纳。变量、常元按定义成立；应用由两个子项的归纳假设及一阶函数同余得到。抽象 \(t=\lambda y.t'\)：若 \(y=x\)，x 不在该抽象的自由参数表中，故两侧是同一项。否则先按定义 2.6 合法改名，使 \(y\notin FV(u)\)，并取同类型的新变量 z，避开当前所有自由参数及 x、y。记 \(a=\lambda y.t'\)、\(b=\lambda y.(t'[u/x])\)，对应自由参数表为 \(\bar v,\bar w\)。待证相等的两项为


\mcsdisplay{L=\ell_b(\bar w),\qquad
R=\ell_a(\bar v)[u^{\#}/x]
=\ell_a(\bar v[u^{\#}/x]).}
\noindent 右侧只替换函数符号的实参，不改写符号名 a 中记录的原抽象代码；a 与 b 也未必不同。由抽象公理及其同排序实例化，


\mcsdisplay{\operatorname{app}(L,z)=(t'[u/x])^{\#}[z/y],\qquad
\operatorname{app}(R,z)=t'^{\#}[u^{\#}/x][z/y].}
\noindent 归纳假设在 \([z/y]\) 下的实例使右端相等；再对新变量 z 全称引入并用函数外延，得到 \(L=R\)。这完成抽象情形。\(\square\)

\textbf{命题 2.15（常用规则的认证）。} 上述演算可推导经典连接词与任意类型量词的自然推演、等式替换，以及


\mcsdisplay{(\lambda x.t)u=t[u/x],\qquad
\lambda x.fx=f\quad(x\notin FV(f)),}

\mcsdisplay{(\forall x.fx=gx)\Rightarrow f=g,\qquad
(p\leftrightarrow q)\Rightarrow p=_o q.}
\noindent 本命题即可读推理的合法性依据：正文引用这些派生规则时，认证意义仍展开为定义 2.13 的有限证书。

\textbf{证明。} 逻辑与量词规则经 \(H_\Sigma\) 的真值双条件转为 ND 对应规则。\(\beta\) 等式由抽象公理代入 \(u^\#\)，再用上面的交换引理把 \((t[u/x])^{\#}\) 换成 \(t^{\#}[u^{\#}/x]\)。\(\eta\) 两边对任意 \(x\) 的应用由抽象公理相等，再用函数外延。函数外延就是模式 5 的对象等式重述。命题外延从 \(p,q\) 均处于 \(\{\top,\bot\}\) 且真值相同得到。一阶等式替换和等式常元的解释给出各类型替换规则。\(\square\)

\textbf{定义 2.16（显式可选公理）。}

\begin{itemize}
\item \(\mathrm{Choice}_\sigma\)：增添 \(\varepsilon_\sigma:(\sigma\to o)\to\sigma\)，并加入
\mcsdisplay{\forall P:\sigma\to o.\,(\exists x:\sigma.Px)\Rightarrow P(\varepsilon_\sigma P).}
\item \(\mathrm{Inf}_\iota\)：
\mcsdisplay{\exists f:\iota\to\iota.\,
[\forall x,y.\,fx=fy\Rightarrow x=y]\land
[\exists z.\,\forall x.\,fx\ne z].}
\end{itemize}
使用这些公理须写入背景理论版本。ZFC 元理论有选择不意味着每个 Henkin 函数域已经包含相应选择算子；追加 Choice 会限制模型类（Choice 本身是逐类型给出的\textbf{模式族}）。Inf 排除有限个体域；基础演算本身不排除它们。加入可选公理集 \(S\) 后，定理 3.9 的可靠性与完备性仍适用于 \(H_\Sigma\cup S\cup\widehat T\)（它们是演算的性质），但 \(H_\Sigma\cup S\) 的\textbf{可满足性}不再由命题 3.7 自动保证，须另给模型或相容性论据。

\textbf{例 2.17（\(\beta\) 与谓词节点）。} 若 \(P:\iota\to o\)，则 \(\lambda x.Px\) 与 \(P\) \(\eta\) 等价，但可分别作为“谓词抽象表达”与“已命名性质”的内容。数学等式或逻辑等价不直接决定第 04 章的节点身份（见定义 4.10 的两种相等）。

\section{2.8 有限实现入口与规则子集}
2026-09-29 的检查器实现有限有类型项、de Bruijn 绑定、无捕获替换、按完整抽象结构命名的 \(\lambda\) 提升，以及定义2.11各类背景模式的实例重建。ND 规则支持假设、蕴含、合取、等式、全称与存在的列明子集，具体白名单、资源界及可信基础见设计记录。未实现规则返回“未支持”，不由本章的完整数学演算定义自动补齐。

开放假设、公理依赖和定理引用由证书重算；接受记录对应具体文件哈希及检查器版本。量词变量条件有合法/非法证书与受控缺陷测试，见核验报告。这些是 FINITE 记录，不是检查器一般可靠性或完备性的机器元证明。

\clearpage\chapter{03 Henkin 语义、标准语义与认证演算}
\mcsorigin{原章节：03-语义：Henkin与标准.md}
本章给第 02 章指定演算的语义，证明语言翻译与模型对应，再使用经典多排序一阶完备性。没有把标准高阶语义的全部有效式宣布为可有效证明。

\section{3.1 类型框架及解释}
\textbf{定义 3.1（函数框架）。} 对每个简单类型 \(\sigma\) 给非空集合 \(D_\sigma\)，其中 \(D_o=\{0,1\}\)；非空性取与定义 2.10 相同的约定。这里排除空域是为了匹配本稿的 ND：该演算可由等式自反和存在引入证明 \(\exists x:\sigma.x=x\)，而这一句在空的 \(\sigma\) 域上为假。允许空域的逻辑可以另行建立，但须相应调整规则与语义，不能直接沿用本稿的可靠性结论。并要求


\mcsdisplay{D_{\sigma\to\tau}\subseteq D_\tau^{D_\sigma}.}
\noindent 箭头域元素是真实集合函数；应用按函数求值。\textbf{非逻辑常元}解释 \(I(c)\in D_\sigma\)（逻辑常元由定义 3.3 固定）；对类型语境 \(\Xi\) 的\textbf{合法赋值} v 指给 \(\Xi\) 中每个声明指定其类型域中的元素。变量、常元和应用的解释分别为 v(x)、I(c)、\(\llbracket f\rrbracket_v(\llbracket u\rrbracket_v)\)。\(D_o\) 中的 0、1 是元理论对象（如 \(\varnothing\) 与 \(\{\varnothing\}\)），与对象语言的符号 \(\top,\bot\) 不同——后者由定义 3.3 解释为 1、0。

\textbf{定义 3.2（抽象闭合）。} 对每个良类型项 \(\Xi,x:\sigma\vdash t:\tau\) 和对 \(\Xi\) 的合法赋值 v，要求


\mcsdisplay{a\longmapsto\llbracket t\rrbracket_{v[x:=a]}
\quad\text{属于 }D_{\sigma\to\tau}.}
\noindent 它就是 \(\llbracket\lambda x.t\rrbracket_v\)。此条件对全部合法赋值而非仅闭项成立。只取任意函数子集不够；例如若排除恒等函数，则 \(\lambda x.x\) 无解释。

\textbf{定义 3.3（逻辑解释与满足）。} \(\top,\bot\) 解释为 1,0；连接词为经典真值函数；\(=_\sigma\) 为 \(D_\sigma\) 上相等的特征函数；\(\forall_\sigma(P)=1\) 当且仅当每个 \(a\in D_\sigma\) 均有 P(a)=1，存在量词同理。这些逻辑函数须属于其对应类型域。定义 \(\mathfrak A,v\models\varphi\) 当且仅当 \(\llbracket\varphi\rrbracket_v=1\)。

满足关系位于元理论。对象语言的 o 型项不等于一个涵盖自身语法的真理谓词。

\textbf{定义 3.4（两类模型）。} 满足以上条件者为 Henkin 模型；另要求每个箭头域都是全函数集者为标准模型。标准模型是 Henkin 模型的子类。语义后承分别记为 \(\vDash_H\) 与 \(\vDash_\mathrm{std}\)。若讨论本体档案的模型，还须满足其单独列明的结构公理（定义 4.1 与附录 E）；不能把任意函数框架都叫作 MCS 知识结构。

对闭句子理论 \(T\)、可能开放的假设集 \(\Gamma\) 及公式 \(\varphi\)，记 \(\mathfrak A\models T\) 当且仅当对每个 \(\theta\in T\) 有 \(\mathfrak A\models\theta\)（闭句子的满足与赋值无关）；定义 1.2 的模型类 \(K\) 在本章取 Henkin 类或标准类之一。\(T\cup\Gamma\models_H\varphi\) 的含义是：对每个 \(\mathfrak A\models T\) 和每个给相关自由变量赋值的 \(v\)，若 \(\mathfrak A,v\models\gamma\) 对所有 \(\gamma\in\Gamma\) 成立，则 \(\mathfrak A,v\models\varphi\)。标准语义使用同一赋值约定，只限制模型类。开放假设在当前赋值下满足，与其全称闭包在模型中成立不同；例如 \(P(x)\vDash_H\) P(x)，一般却没有 \(P(x)\vDash_H\forall x.P(x)\)。

\section{3.2 项解释的基本引理}
由定义 3.1–3.2 的递归，解释 \(\llbracket t\rrbracket_v\) 对每个良类型项 \(t\) 与每个合法赋值有定义（\(\llbracket\cdot\rrbracket\) 沿 \(t\) 的形成树递归）。

\textbf{命题 3.5（赋值局部性及语义替换）。} 设 \(v,v'\) 是类型语境 \(\Xi\) 的合法赋值，\(t\) 满足 \(\Xi\vdash t:\tau\)。

\begin{enumerate}
\item \textbf{局部性}：若 \(v\) 与 \(v'\) 在 \(FV(t)\) 上取值相同，则 \(\llbracket t\rrbracket_v=\llbracket t\rrbracket_{v'}\)。
\item \textbf{语义替换}：若 \(\Xi,x:\sigma\vdash t:\tau\) 且 \(\Xi\vdash u:\sigma\)，则
\end{enumerate}

\mcsdisplay{\llbracket t[u/x]\rrbracket_v
=\llbracket t\rrbracket_{v[x:=\llbracket u\rrbracket_v]}.}
\noindent \textbf{证明。} 对 \(t\) 的形成树作同时归纳（两条一起证）。变量情形：若 \(t=x\)，第 2 条两侧都是 \(\llbracket u\rrbracket_v\)；其他变量与常元在两条中都直接成立。应用情形由两个子项的归纳式与函数应用合成。抽象情形 \(t=\lambda y.t':\rho\to\tau'\)：先取不出现在 \(u\) 的自由变量中的同类型新绑定名（\(\alpha\) 改名不改变解释）。要比较的两个值都落在 \(D_{\rho\to\tau'}\) 中；对每个 \(a\in D_\rho\)，由第 2 条的归纳假设两边的体在同一扩展赋值下取值相同，而第 1 条的归纳假设保证更新 \(y\) 不改动 \(\llbracket u\rrbracket_v\)。由定义 3.1，\(D_{\rho\to\tau'}\) 的元素就是函数，逐点相等即同一元素——这一步只用定义 3.1，不用 \(H_\Sigma\) 的函数外延公理。局部性在抽象情形同样逐点应用归纳假设，其余情形直接。\(\square\)

本引理在命题 3.7、定理 3.9 与第 01 章命题 1.5 的证明中被使用。

\section{3.3 与多排序认证模型的对应}
本节把第 02 章的翻译与本章的语义对接：命题 3.7 把 Henkin 模型送进认证模型，命题 3.8 反向把认证模型表示成 Henkin 模型。两者合起来是定理 3.9 的技术前提——只有两个模型类对应，经典的一阶可靠性与完备性才搬得过来。

\textbf{定义 3.6（认证语言）。} 使用第 02 章的 \(\mathcal F_\Sigma\)、app、抽象函数符号 \(\ell_a\)、公理 \(H_\Sigma\) 与翻译 \(t\mapsto t^\#\)、\(\varphi\mapsto\widehat\varphi\)。每个简单类型对应一个一阶\textbf{非空排序}；\(\ell_a\) 是函数符号（本身不可被量词约束），其值属于排序 \(\sigma\to\tau\)，因而可被该排序的一阶变元量化——新增符号只是有限语法的命名，不改变原高阶语言的量词类型。

\textbf{命题 3.7（Henkin 到认证模型）。} 每个 Henkin 模型 \(\mathfrak A\) 唯一确定一个满足 \(H_\Sigma\) 的认证模型 \(\mathfrak A^{\#}\)（唯一性由 \(H_\Sigma\) 第 5、6 条给出），且


\mcsdisplay{\llbracket t^\#\rrbracket_{\mathfrak A^\#,v}
=\llbracket t\rrbracket_{\mathfrak A,v},\qquad
\mathfrak A^\#,v\models\widehat\varphi\iff \mathfrak A,v\models\varphi.}
\noindent \textbf{证明。} 保留各类型域、常元及真实应用。把 \(\ell_a(\bar{y})\) 解释为抽象 a 在参数 \(\bar{y}\) 赋值下的函数；局部性使定义不依赖其他变量，抽象闭合使其取值在正确域中。真值公理、量词公理与外延性按模型定义成立；抽象方程由解释定义成立。由定义 2.3 的最小生成性对项的形成树归纳，得到翻译相等，o 型项等于 \(\top\) 当且仅当其值为 1。\(\ell_a\) 的值由逐点抽象方程及外延性唯一决定。\(\blacksquare\)

\textbf{命题 3.8（认证模型的外延化）。} 每个 \(F_\Sigma\) 模型 \(\mathfrak B\models H_\Sigma\) 都与一个 Henkin 模型的认证扩张按各排序同构，因而保留所有翻译公式的真值。

\textbf{证明。} 按类型树复杂度构造双射 \(j_\sigma\)。对基本类型取原域的拷贝；Boolean 排序由两值公理恰有两个元素，选取使 \(\top\mapsto1\)、\(\bot\mapsto0\) 且与真值公理相符的对应。已构造 \(\sigma,\tau\) 后，对 \(f\in\mathfrak B_{\sigma\to\tau}\) 定义函数


\mcsdisplay{F_f(j_\sigma(a))=j_\tau(app^{\mathfrak B}(f,a)).}
\noindent \(j_\sigma\) 为双射，故良定义。令 \(D_{\sigma\to\tau}=\{F_f:f\in\mathfrak B_{\sigma\to\tau}\}\)，\(j_{\sigma\to\tau}(f)=F_f\)。满射由定义；若 \(F_f=F_g\)，则所有 app(f,a)=app(g,a)，由 \(H_\Sigma\) 外延公理得到 f=g，故单射。将常元按 j 输送。由定义 2.3 的最小生成性对原项的形成树归纳，抽象公理保证每个赋值下所需函数由相应 \(\ell_a\) 的值提供；因此抽象闭合成立。\(H_\Sigma\) 的真值、等式和量词公理给出定义 3.3 的逻辑解释。递归构造同时保持应用和 \(\ell_a\)，故是认证结构的排序同构，并保持翻译公式。\(\blacksquare\)

这个对应解释了为何 Henkin 语义可借多排序一阶模型论研究。它不表示箭头域已包含全部元理论函数，因而不能把相应结果自动提升到标准语义。

\section{3.4 可靠性、完备性及其限度}
\textbf{定理 3.9（指定演算的可靠性与 Henkin 完备性）。}


\mcsdisplay{\Gamma\vdash_T\varphi
\quad\Longleftrightarrow\quad
T\cup\Gamma\models_H\varphi.}
\noindent 其中 \(\models_H\) 见 3.1 末段的语义后承定义。此处使用第 02 章的认证定义，并以经典多排序一阶自然推演的可靠性与完备性为引用元定理；这正是定义 1.2 的强形式完备性在模型类 \(\mathcal K\) 取 Henkin 类时的兑现。

\textbf{证明。} 正向：任取满足 \(T\) 的 Henkin 模型 \(\mathfrak A\) 及使 \(\Gamma\) 成立的赋值 \(v\)，命题 3.7 给出满足 \(H_\Sigma\cup \hat{T}\) 的认证模型，且 \(\hat{\Gamma}\) 在对应赋值下成立。有限 ND 证明逐条保持满足：假设与弱化直接；连接词规则由真值表；全称引入使用本征变量不自由于假设，故可任意改变其赋值；存在消去用一个见证扩展赋值并以新变量条件保证结论不依赖见证名；等式替换由同一对象可替代，经典反证由两值逻辑成立（以上逐条对应定义 2.12 的规则表，其中本征变量与见证变元两条 side condition 是可靠性论证的关键）。故 \(\hat{\varphi}\) 在该赋值下为真，再由\textbf{命题 3.7 的对译等价}得到 \(\varphi\)。

反向：若 \(T\cup\Gamma\vDash_H\varphi\)，则由命题 3.8，任何满足 \(H_\Sigma\cup \hat{T}\) 的认证模型及使 \(\hat{\Gamma}\) 成立的赋值都使 \(\hat{\varphi}\) 成立。对一阶理论 \(H_\Sigma\cup\widehat T\cup\widehat\Gamma\) 应用多排序一阶完备性，得到一份有限 ND 证明，恰为第 02 章定义的 \(\Gamma\vdash_T\varphi\)。\(\blacksquare\)

多排序完备性可由通常 Henkin 见证构造按排序进行得到；本地 Prestel–Delzell §1.5 与 EFT 完备性章节提供一阶基准，Henkin（1950）给出类型论一般模型思想（印刷 pp.84–85 / 所核扫描 PDF pp.5–6；原文链接与采用边界见附录 B 的 S01）。这里的转移已明确翻译和模型条件，不能仅以作者姓名代替匹配论证。

\textbf{推论 3.10（语义的单向关系）。}


\mcsdisplay{T\models_H\varphi\ \Longrightarrow\ T\models_{\rm std}\varphi.}
\noindent 反向一般不成立。标准模型构成较小模型类，可能满足一般模型不满足的句子；例 3.16 给出一个具体见证。标准高阶有效性不存在覆盖全部有效式的有效可靠完备演算；此为引用的不可有效公理化结果，见 EFT IX §1、X §5 及 Prestel–Delzell 附录 B。不是本稿自证的不完全性定理。

\textbf{推论 3.11（Henkin 紧致性）。} 若 T 的每个有限子集都有 Henkin 模型，则 T 有 Henkin 模型。

\textbf{证明。} \(H_\Sigma\cup \hat{T}\) 的任意有限子集只涉及 T 的有限部分，依命题 3.7 有模型；一阶紧致性给出其全体模型，再由命题 3.8 外延化。\(\blacksquare\)

注意：\textbf{推论 3.11 对标准语义不成立}——标准语义不满足同样的结论。经编码使用初等子结构、超积或 Löwenheim–Skolem 时，只能对对应一般模型类及编码语言断言；全函数域条件、标准实数域或特定外部概率空间不因此保持。

\section{3.5 定义扩张的精确条件}
\textbf{定义 3.12（项定义）。} 两种情形：

\begin{enumerate}
\item \textbf{常元定义}：取旧语言闭项 \(t:\sigma\)，添加新常元 \(c:\sigma\) 及公理 \(c=t\)；
\item \textbf{函数定义}：添加新常元 \(c:\rho_1\to\cdots\to\rho_m\to\tau\) 及公理 \(c=\lambda\bar y.u\)，要求右侧只含旧符号与列明的参数 \(\bar y\)。
\end{enumerate}
\textbf{命题 3.13（项定义保守）。} 每个旧 Henkin 模型在原类型域上唯一扩张为该定义模型；对旧语言的语义及本稿可证性都保守。

\textbf{证明。} 令 c 解释为 \(\llbracket t\rrbracket\)，它已在旧域中。含 c 的任何新抽象 \(s\) 先作\textbf{常元替换} \(s[c:=t]\) 得到旧项；对 \(s\) 作与命题 3.5 同一形式的归纳（把变量情形换成常元情形，用 \(I(c)=\llbracket t\rrbracket\)）得 \(\llbracket s\rrbracket=\llbracket s[c:=t]\rrbracket\)，故该抽象的解释已在旧域中，抽象闭合未被破坏。c 的解释由定义唯一决定，旧公式解释未变，所以每个旧模型都可扩张（模型扩张保守性）。可证性保守分三步：设 \(T'\) 为扩张理论、\(\varphi\) 为旧语言句子，由定理 3.9 的\textbf{可靠性}方向 \(T'\vdash\varphi\Rightarrow T'\models\varphi\)；由每个旧模型可扩张，\(T'\models\varphi\Rightarrow T\models\varphi\)；再由定理 3.9 的\textbf{完备性}方向，\(T\models\varphi\Rightarrow T\vdash\varphi\)。\(\blacksquare\)

谓词定义 \(P(\bar{y})\leftrightarrow\varphi(\bar{y})\) 可写成 \(P=\lambda\bar{y}.\varphi\)，适用本命题。

\textbf{定义 3.14（规格定义的额外义务）。} 设旧语言关系 R(x,y) 除 \(x:\sigma\)、\(y:\tau\) 外无其他自由变量；额外参数须先列明并闭合到拟定义函数的参数类型。若想由 R 引入新常元 \(f:\sigma\to\tau\)，除了逐点存在唯一，还须保证其实现落在当前函数域。稳妥认证前提是两条：\textbf{（域内实现）}


\mcsdisplay{T\vdash \exists f:\sigma\to\tau.\ \forall x\,R(x,fx)}
\noindent \textbf{（逐点唯一）}


\mcsdisplay{T\vdash\forall x\forall y\forall z.\,((R(x,y)\land R(x,z))\Rightarrow y=z).}
\noindent 二者共同给出当前函数域中唯一的实现。

\textbf{充分性（定义 3.14 的证明纲要）。} 可在每个旧模型中取该已有函数 \(F\) 解释新常元 f。新语言抽象中的 f 全部改为同类型的新自由变量 z，得到旧语言项，再在 z:=F 的赋值下解释；定义 3.2 对全部赋值的闭合条件保证结果仍在原域。因此这项实现条件足以保证唯一模型扩张与旧语言保守性。这里不必假设 \(F\) 有旧语言闭项名称。

另一充分方案是在带相应 Choice 的背景中给出旧语言项 \(f=\lambda x.\varepsilon(\lambda y.R(x,y))\)，再用逐点存在性证明规格，转回命题 3.13。逐点存在唯一、函数域内存在实现、存在可消去闭项名称是三个不同条件，不能混为一谈。

\textbf{边界 3.15。} 在全函数标准模型中，逐点唯一关系的函数图自然属于全函数域。任意 Henkin 子函数域则不保证这一点。不能把一阶教材中“增加一个函数符号”的保守扩张无条件移植为“向已经量化的 HOL 函数类型增加一个元素”。本稿自动认证的默认定义扩张采用旧语言项定义；更广的规格定义必须附相应实现及闭合证书。

\textbf{例 3.16（边界 3.15 的具体反模型）。} 本例如下，既说明定义 3.14 的实现条件不可省，又见证推论 3.10 的反向失败。取 \(D_o=\{0,1\}\)、\(D_\iota=\{a,b\}\) 且 \(a\ne b\)。仅用于构造模型，为 \(D_o\) 配置辅助等价关系 \(\sim_o=D_o\times D_o\)，为 \(D_\iota\) 配置辅助相等关系；其他基本类型可取单点集。递归令 \(D_{\sigma\to\tau}\) 为所有保持辅助等价的函数，即 \(x\sim_\sigma y\) 蕴含 \(F(x)\sim_\tau F(y)\)，并令 \(F\sim_{\sigma\to\tau}G\) 当且仅当对每个 \(x\in D_\sigma\) 都有 \(F(x)\sim_\tau G(x)\)。所有函数域非空，因为包含常值函数。辅助等价不是对象语言的等式；对象等式始终解释为真实相等。

这些域构成 Henkin 模型：最终值为 o 的谓词类型具有全辅助等价关系，故经典逻辑、真实相等与量词常元均保持辅助等价，属于所需域。按项结构归纳，辅助等价的赋值产生辅助等价的结果；抽象情形因此得到保持辅助等价的函数，且两组赋值所得函数逐点辅助等价。这同时证明所有赋值下的 \(\lambda\) 闭合。应用是真实求值，所以函数外延性也成立。

但是 \(D_{o\to\iota}\) 只有常值函数。给 a,b 同名个体常元，旧语言关系


\mcsdisplay{R(p,y)\equiv(p\land y=a)\lor(\neg p\land y=b)}
\noindent 满足 \(\forall p:o.\exists!y:\iota.R(p,y)\)，却不满足 \(\exists f:o\to\iota.\forall p.R(p,\mathrm{fp})\)。唯一的元理论实现必须分别把 1、0 映为 a、b，因而不在该箭头域中。这证明“逐点唯一”不能替代定义 3.14 的实现条件。等价地，句子


\mcsdisplay{\forall a,b:\iota.\;\exists f:o\to\iota.\,
(f\top=a\land f\bot=b)}
\noindent 在每个标准模型中成立，在上述 Henkin 模型中失败，具体见证了推论 3.10 的反向不成立。

\section{3.6 模型见证与计算边界}
\textbf{例 3.17（基础核的标准模型；存在性见证）。} 记\textbf{基础核}为以 \(H_\Sigma\) 为全部公理、不含额外理论公理的最小理论。令每个非 Boolean 基本类型是一个非空单点集，\(D_o=\{0,1\}\)，各箭头类型取全函数域，逻辑常元按定义解释。所有项及抽象有解释，因此基础核\textbf{可满足}——这是 §2.5 与命题 3.7 所依赖的存在性见证。加入 Inf 后此单点个体模型不再适用；加入其他领域公理必须另查模型存在性。

一个有限登记网络可以引用此模型，也可以引用无限数学结构；\textbf{“注册有限”不等于“语义有限”}——这与 1.3“描述的大小与语义的大小必须分开”是同一条政策。无限有效理论可枚举有限证书，却不保证任意真伪、等价、最优证明或全部硬关系可判；\textbf{“有模型”不等于“可判”}（见命题 2.14 末注：数学完备性与受有效呈现限制的检查算法是不同结论）。

第 07 章\textbf{命题 7.17}（外部扩张的语义保守性）以\textbf{每个核心模型可扩张}为条件。MCS 全体理论并不因存在一个玩具实例就已证明任何给定数学背景的一致性。

\section{3.7 实现验证的语义边界}
独立有限参考求值器比较35个有界项在175个域/赋值实例中的对象解释与\(\lambda\)提升解释，并检查四案例片段的42个赋值及列明的背景模式实例。此范围只支持 FINITE；定义3.2对全部项和全部合法赋值的抽象闭合，仍由本章一般论证承担。实现保真及模型对应的机器证明义务见义务表。

当前内核的定义扩张仅接受新符号的旧语言闭项定义；规格定义返回未支持，不由逐点存在唯一省略域内实现。例3.16的方向保持为“标准语义成立、给定Henkin模型失败”，不改写成Henkin后承不能推出标准后承。

\clearpage\chapter{04 节点生成、知识本体与稳定身份}
\mcsorigin{原章节：04-节点生成与身份.md}
本章从有类型有限语法生成候选，再定义合法知识本体。节点的表达结构与证据可以严格检查；“对谁是最小认知单元”移至第 07 章，不进入生成规则或身份判定。

\section{4.1 整个本体的结构签名}
\textbf{定义 4.1（本体版本）。} 一份集合大小的本体版本记为


\mcsdisplay{M=(\Sigma,T,N_M,Payload,Rep,Cert,Supp,\Lambda_M,R_M^0,Agg,Templ,\partial M).}
\noindent \(\Sigma,T\) 是公开语言与背景模块；\(N_M\) 是节点引用域；Payload 是不可变负载解释；Rep 给表达及公开表征；Cert 是已登记证书与认证状态；Supp 是第 05 章的支持接口；\(\Lambda_M\) 是第 06 章的公开行动契约；\(R^{0}_M\) 是通用关系描述；Agg、Templ 是层级聚合和模板；\(\partial M\) 是对环境版本的引用边界。\textbf{十二个坐标都是公开分量，不含具体学习者}；\(\partial M\) 声明环境版本引用，也属公开数据。硬演绎、泛化与契约先修由这些数据按第 06 章派生，不依赖个人状态。

这些分量\textbf{按 2.1 的内部化判据}在第 02 章高阶语言中分型：Node、Code、CertCode、Action、Aggregate、Template 与角色标号域 Role 为互异基本类型（核心签名的完整基本类型清单见附录 E，本段只列本节用到的名字）；Payload 为到代码的函数，\(\mathrm{HasRole}:\mathrm{Node}\to\mathrm{Role}\to o\)，\(\mathrm{Input}:\mathrm{Action}\to\mathrm{Node}\to o\)，\(\mathrm{Produces}:\mathrm{Action}\to\mathrm{Node}\to o\) 等为谓词。有限性和检查器的标准含义先在 ZFC 元层给出；内部编码只解释其被编码部分，不添覆盖自身语言的全局真理谓词。

\textbf{定义 4.2（合法本体条件）。} 结构 M 合法当且仅当满足以下条件；这是条件清单的合取，不是一组无需检验的标签。

\begingroup\small\setlength{\tabcolsep}{4pt}
\begin{longtable}{>{\raggedright\arraybackslash}p{\dimexpr 0.1509\linewidth-2\tabcolsep\relax}>{\raggedright\arraybackslash}p{\dimexpr 0.8491\linewidth-2\tabcolsep\relax}}
\toprule
\textbf{条件} & \textbf{精确要求} \\
\midrule
\endfirsthead
\toprule
\textbf{条件} & \textbf{精确要求} \\
\midrule
\endhead
\bottomrule
\endfoot
\bottomrule
\endlastfoot
身份功能性 & 每个引用在一个版本中恰有一份不可变负载；同名新内容必须用新版本引用 \\
负载良构 & 每个节点属于定义 4.4 的一个构造类型并带所需形成见证 \\
引用封闭 & 每个引用要么指向 \(N_M\) 的条目，要么指向 \(\partial M\) 中声明的环境版本，不能悬空伪装成空依赖 \\
认证可靠 & 被标作检查通过的证书确实满足固定检查器 \(\mathrm{Check}_T\)（见命题 2.14）；未编码的数学论证使用不同的证据标签 \\
行动独立 & \(\Lambda\) 的形成不调用路径规划或个体掌握估计，见第 06 章 \\
支持有型 & 表达/证明/路线概念支持与输入资源分开；Unknown 与精确空值分开 \\
层级与溯源 & 聚合保持成员及其层级类型，派生视图保留到环境的来源 \\
外部隔离 & Payload、Cert、公共契约和硬关系定义不含个人标识、观察或学习状态字段 \\
\end{longtable}
\endgroup

表中各条的判据位置：身份功能性与负载良构见定义 4.3–4.4；引用封闭见定义 4.1 的 \(\partial M\)；认证可靠见命题 2.14；行动独立见第 06 章定义 6.5–6.6；支持有型见第 05 章定义 5.9–5.12；层级与溯源见第 09 章与 11.1；外部隔离见第 07 章与命题 1.5 的检查义务。

这些公理不要求每份有限档案存储所有语法候选或全部真后果。生成完整性与存储完整性不同；稀疏登记实例通过边界引用说明未展开部分。若声称某实例是生成闭包，须另证该闭包性质。

\section{4.2 节点、角色和形成证据}
\textbf{定义 4.3（不可变节点记录）。} 节点为


\mcsdisplay{n=(id,k,d),\qquad d\in Data_k.}
\noindent 其中 \textbf{\(\mathrm{Data}_k\) 是定义 4.4 中 k 型构造的有限描述负载所成的域}，每份负载描述有限，域本身可以无限；id 是版本内的唯一标识符，k 是构造类型，d 是有类型负载，映射 id \(\mapsto\) (k,d) 在一个版本内\textbf{函数式}（这正是定义 4.2 的身份功能性）。同一有限描述可以有不同来源记录；固定 id 不允许同时指向不同 k,d。若有形式表达，背景记录为 \((T,\Xi,\Gamma,\tau)\)：\(\Xi\) 为类型环境，\(\Gamma\) 为公式假设，\(\tau\) 为表达类型。

\textbf{定义 4.4（节点构造与形成条件）。}

\begingroup\small\setlength{\tabcolsep}{4pt}
\begin{longtable}{>{\raggedright\arraybackslash}p{\dimexpr 0.2776\linewidth-2\tabcolsep\relax}>{\raggedright\arraybackslash}p{\dimexpr 0.3368\linewidth-2\tabcolsep\relax}>{\raggedright\arraybackslash}p{\dimexpr 0.3856\linewidth-2\tabcolsep\relax}}
\toprule
\textbf{构造类型} & \textbf{有限负载} & \textbf{WF 条件（以及哪些留给另行认证）} \\
\midrule
\endfirsthead
\toprule
\textbf{构造类型} & \textbf{有限负载} & \textbf{WF 条件（以及哪些留给另行认证）} \\
\midrule
\endhead
\bottomrule
\endfoot
\bottomrule
\endlastfoot
Symbol & 声明、类型、签名版本 & 声明合法；不要求它已有特定真实解释 \\
Term & \((T,\Xi,t,\tau)\) & \(\Xi\vdash t:\tau\)；不要求 t 处于某逻辑正规形 \\
Concept & 对象类型、参数、谓词 P & P 类型正确；参数经明确 \(\lambda\) 闭合形成概念表示 \\
Definition & 新符号、旧项、扩张记录 & 默认使用第 03 章的旧语言项定义；更广规格另附实现/闭合证书 \\
Claim & \((T,\Xi,\Gamma,\varphi)\) & \(\varphi\) 为 o 型；合法不等于可证 \\
Proof & 目标 Claim、有限证书代码 & 代码语法合法；是否 Check=1 是另一个认证字段 \\
Example & 概念引用、对象规格、满足断言 & 类型及引用正确；满足性另证，不能假称任意模型都可判 \\
Counterexample & 精确目标、对象规格、失败见证 & 引用和量词类型正确；确实反驳另需证书 \\
Problem & 输入输出、目标规格、约束 & 任务接口类型正确；是否已解决另有公共来源记录 \\
Theory & 语言、演算、公理呈现/模块引用 & 有限呈现代码或明确的环境理论引用；公理一致性不由 WF 保证 \\
Construction & 输入输出、步骤规格、验证目标 & 有限步骤及引用类型正确；正确性及终止性分别认证 \\
Method & 范围、公共接口、启发体、案例引用 & 记录格式与引用合法；不要求启发体成为决定算法 \\
Representation & 对象、表达介质、对应说明 & 介质/对象匹配；等价或保真另行认证 \\
Miscon\mcsbrk{}ceptio\mcsbrk{}nPatte\mcsbrk{}rn & 错误规则、任务、反例、范围 & 公共描述，不包含某个人是否持有它 \\
\end{longtable}
\endgroup

公理、定理、性质是 Claim 的角色：公理角色需要相应理论的成员见证；定理角色需要认证证明或清楚标记的引用证据；性质角色说明用途。全局方法和局部方法是 Method 的范围标签。上述公理、定理、性质角色是 Claim 的标签，\textbf{不改变构造类型 k}；其他构造也可拥有各自的合法角色，例如第 05 章的概念与符号锚定。角色可以重叠，但不把不同构造类型混成一个无型节点。

\textbf{命题 4.5（良构与真理分离）。} 在有效呈现及有限形成见证条件下 WF 可判，但\textbf{仅由 WF(n) 不能推出}节点命题为真或可证（WF 再加上相应的认证字段可能给出更多，见定义 4.4 各行的区分）。

\textbf{证明。} 表格中的 WF 只执行有限解析、类型、引用与证书格式检查；理论成员需要可判的引用见证。Claim 构造允许任一良类型 \(\varphi\)，因此也允许 \(\varphi\) 与 \(\neg\varphi\)；它不可能仅凭语法就同时认证二者为数学真理。\(\blacksquare\)

对任意无有效呈现的理论仍可定义语义节点，但不能沿用本命题的算法承诺。

\section{4.3 有限生成与自动枚举}
\textbf{定义 4.6（候选宇宙）。} 固定有效可数符号表、有限字符串编码、规则表与公开理论呈现集合。种子 S 包含符号声明、有限描述原子及背景模块代码。自动生成所用的种子、规则索引与理论引用须有有效枚举，成员调用携带有限形成或枚举见证；仅有可数性不足以保证可枚举。F 是单调的有限元构造规则族（\(X\subseteq Y\) \(\Longrightarrow\) \(F(X)\subseteq F(Y)\)），实例通过可判形成见证核验。置


\mcsdisplay{U_0=S,\qquad U_{n+1}=U_n\cup F(U_n),\qquad
U_\Sigma=\bigcup_{n<\omega}U_n.}
\noindent 每份节点描述是有限树或有限带引用记录；引用目标采用已有合法记录或明确环境代码。相互引用的公共导航边可以后登记，不把未证的循环定义当作形式内容。

\textbf{定理 4.7（最小生成）。} 设 F 单调，则 \(U_\Sigma\) 是包含 S 的最小 F-闭集合；可数种子与可数有效规则产生可数 \(U_\Sigma\)。

\textbf{证明。} 每次构造仅用有限多个输入，其阶段编号有最大值，输出位于下一阶段，故 \(U_\Sigma\) 闭合。设 \(Y\supseteq S\) 且 \(F(Y)\subseteq Y\)：\(U_0=S\subseteq Y\)；若 \(U_n\subseteq Y\)，由单调性 \(F(U_n)\subseteq F(Y)\subseteq Y\)，故 \(U_{n+1}\subseteq Y\)；归纳得各 \(U_n\subseteq Y\)，故 \(U_\Sigma\subseteq Y\)。有限元规则作用于可数域的有限元组所得结果可数，可数并仍可数。\(\blacksquare\)

这里使用 \(\omega\) 阶段就足够，不将 \(L_\alpha\) 的超穷层级移植到有限语法。

\textbf{定义 4.8（公平候选生成）。} 按自然数代码递增枚举全部有限描述，或先编码到固定有限字母表，再按编码长度及层内字典序枚举；逐一检查合法的形成树，输出通过者。不能只按可数无限符号表中的符号个数逐层耗尽，因为长度为 1 的一层就可能无限。对公理枚举见证和证明搜索采用交错枚举。每个有限合法候选最终被访问。避免重复可用规范代码集合，但不把逻辑等价当作去重判断。

\textbf{命题 4.9（枚举覆盖与边界）。} 每个有有限合法描述及有限认证见证的候选最终输出；不保证在有限时间输出全部候选，也不保证识别全部标准语义真理。

\textbf{证明。} 其全部代码、规则索引和见证长度有有限最大界，枚举到足够预算时包含该对象。不存在的证据可能被永久搜索，因此整体决定性不随之成立。\(\blacksquare\)

自动生成的是基础表达、节点及证书候选；有效检查进一步产生认证层。人工选择研究重点及添加公开解释，不负责穷举基础语法。

\section{4.4 稳定身份与版本}
\textbf{定义 4.10（两种相等）。} 节点记录相等按 id,k,d；表达可另按 \(\alpha\)、\(\beta\eta\)、可证等价或模型外延相等比较，各自带类型和背景。后四者都不自动推出节点相等。

规范表达锚定可以对 \(\alpha\) 代码采用一个确定的生成引用；带不同来源、角色或材料的节点仍独立存在。这样允许机械识别绑定名差异，又保留不同内容记录。

独立的来源记录不必全部成为第 05 章 \(\kappa\) 的同码别名。该查询须有限返回，因此在固定登记版本中，同一规范表达只登记有限个概念锚点；更多来源经独立的溯源关联保存，不据此合并节点身份。

\textbf{命题 4.11（学习者替换保持身份）。} 若 M 不变，任何外部模型 \(E_{1},E_{2}\) 及状态 \(s_{1},s_{2}\) 对节点身份、WF 和已登记证明的 Check 结果给出相同值。

\textbf{证明。} 三个判断均只读取固定的 id,k,d、签名和证书，不读取外部状态。函数性给出相等。\(\blacksquare\)

新增证明应建立新 Proof 记录并追加认证关联，不改旧证明内容；改写定义正文则建立新版本并给兼容/翻译说明。学习者答对题只更新外部记录，不能借用“版本更新”绕过此隔离。

\section{4.5 方法节点的严格外壳}
\textbf{定义 4.12（方法接口）。}


\mcsdisplay{meth=(name,scope,In,Out,Pre,Post,Use,Demo,Fail,body).}
\noindent In/Out 为任务与候选结果的类型；Pre/Post 是有明确语义的条件，允许标注尚不能完全形式化的范围；Use 是结构性概念锚定；Demo/Fail 是公共案例及失败范围；body 是自然语言启发体。WF 检查字段格式与引用，不谎称可以判定一切 Pre/Post。

全局方法包括直觉建立与失效、内省、从特例切入、反例构造、逆向分析、不变量、应用替代练习、可视化。局部方法如“用夹逼估计处理特定极限”，可实例化全局策略。一个方法若后来取得严格算法和正确性证明，可另建 Construction 与 Proof 节点；启发内容与形式实现通过显式关系连接，不偷换。

一般方法呈现可作为 \(\Lambda\) 中的公开行动，但其证书仅保证任务接口匹配及内容合法。学习成功率、难度、掌握变化均由第 07 章外部模型评价。

\section{4.6 节点与认知原子的接口}
本体节点是可独立引用并具有明确输入输出接口的知识描述，\textbf{不是预设对所有人的最小心理单位}。本体保存分解候选及重构契约；认知原子性


\mcsdisplay{Atom_E(s,n;g,\mathcal R,\mathcal D)}
\noindent 由第 07 章定义在外部。给定学习目标和表征，某节点可整体呈现，也可分解成多个目标；这是外部组块及派生视图的差异。

有限可判的分解域可供穷举最小性检验；一般“没有想到更小分解”不是极小性证明。方法节点也可参与这类分解，但不能强制其启发体成为封闭形式公式。

\section{4.7 从本体到具体网络}
知识网络是合法本体 M 的关系呈现；学习网络是 \((M,E,s,g,b,\rho)\) 的派生结果；可视化网络是带溯源和保持声明的显示投影。画图中的一个点可以是节点、聚合或事件，必须标注其类型。不同图的点数、布局和显示标签不决定本体同构。

概念函数、第 06 章关系、第 08 章结构保持映射、第 11 章视图、第 12 章规划都以本章接口为输入。外部模型的知识空间可以成为一个 Theory/Model 描述节点，但其某位学习者的实际状态始终不是 M 的一项运行数据。

\clearpage\chapter{05 概念函数与支持族}
\mcsorigin{原章节：05-概念函数与支持族.md}
本章区分三个问题：表达\textbf{写到了什么}；定义、证明或路线\textbf{使用了什么}；完成指定用途\textbf{至少需要什么}。分别对应出现集、带见证支持族和极小充分支持。全部函数只读取本体材料，不读取学习者状态。

\section{5.1 概念域、命名与出现}
\textbf{定义 5.1（概念登记）。} 固定有效签名 \(\Sigma\)、理论 T 与节点生成规则。令 \(K_\Sigma\subseteq N_M\) 是赋予概念、符号或谓词/构造接口角色的节点集合。每个可形式命名概念 c 带闭的有类型表示 \(t_c:\tau_c\)；有参数的概念先以 \(\lambda\) 抽象闭合。谓词的表示类型为 \(\alpha_{1}\to\ldots\to\alpha_k\to o\)，故谓词本身可以成为概念节点。闭项仍可含签名常元。

登记函数 \(\kappa\) 把有限有类型语法的 \(\alpha\) 规范代码映到有限个概念引用；默认使用生成的表达节点作锚定，别名显式登记。要求查询可计算、返回有限集合，并满足 \(c\in\kappa(\mathrm{code}_\alpha(t_c))\)。这要求固定登记版本中每个规范表示的概念纤维 \(\{c:\mathrm{code}_\alpha(t_c)=q\}\) 有限；有限纤维本身还不保证可计算，查询终止是另一个明确条件。若同一表达有无限多个来源记录，只把有限个规范概念锚点放入该查询，其他来源通过溯源关系保存。这是\textbf{表示约定}，不是自动识别全部语义概念的算法。相同表示可以锚定不同来源记录，不因此改变节点身份。

\textbf{定义 5.2（语法概念出现）。} 给定 \(\Xi\vdash e:\tau\)，遍历有限语法树。对位置 p 的子项 \(u_p\)，按其位置环境中的变量声明次序，把自由变量 \(\lambda\) 闭合为 \(\mathrm{close}(u_p)\)。定义


\mcsdisplay{C_{\mathrm{occ}}(e)=\bigcup_{p\in\operatorname{Pos}(e)}
\kappa(\operatorname{code}_{\alpha}(\operatorname{close}(u_p))).}
\noindent \(\mathrm{code}_\alpha\) 只消除约束变量改名歧义，\textbf{不}对逻辑等价、\(\beta\eta\) 等价或高阶代入实例取商。类型环境与闭合次序固定。概念角色由公开登记决定，不由个人掌握程度决定。

\textbf{命题 5.3（有限性与相对可计算性）。} 在定义 5.1 的条件下，\(C_\mathrm{occ}(e)\) 是可计算的有限集合，并在一致的 \(\alpha\) 改名下不变。

\textbf{证明。} 有限语法树仅有有限多个位置。各位置的自由变量集及闭合项由结构递归计算，绑定索引给出 \(\alpha\) 代码；每次 \(\kappa\) 查询终止且输出有限。有限次取并仍终止且有限。\(\alpha\) 改名保持绑定索引，因而保持查询结果。\(\blacksquare\)

这不是高阶模式匹配或语义等价判定定理。把“经任意定义展开后出现”混入此查询时，终止性须重新证明。

\section{5.2 原始全等价概念集及其退化}
\textbf{定义 5.4（相对等价）。} 对固定 T、类型环境 \(\Xi\) 和类型 \(\tau\) 的项，记


\mcsdisplay{e\equiv_{T,\Xi,\Gamma_0}e'\quad\Longleftrightarrow\quad
\Xi;T;\Gamma_0\vdash e=_\tau e'.}
\noindent 其中 \(\Xi\) 只负责项的类型检查，\(\Gamma_{0}\) 是固定公式假设族；后文省写 \(\Gamma_{0}\) 时沿用同一假设族。证明的非逻辑假设也须固定，不能为了得到等价临时加入新假设。命题类型的等式对应外延 HOL 中的逻辑等价。标准语义等价可以另行研究，但须另设记号。

\textbf{定义 5.5（全等价形式概念并集）。} 对节点 n 的公开有限表达列表 \(\mathrm{Rep}_{0}(n)\)，定义


\mcsdisplay{C_{\mathrm{all}}(n)=
\bigcup_{e_0\in Rep_0(n)}
\ \bigcup_{e\equiv_{T,\Xi}e_0}C_{\mathrm{occ}}(e).}
\noindent 每个分量保持自身的背景、类型和登记域。跨理论比较先提供明确翻译，不直接跨背景判等。无形式表达时第一分量为空；这不表示方法节点全部支持已知为空，见 §5.5。

\textbf{定理 5.6（全等价概念集饱和）。} 假设 \(\mathrm{Rep}_{0}(n)\ne\emptyset\)，等价允许普通 \(\beta\) 变换，出现函数扫描归约前的完整语法树。则当前签名中每个有闭表示的登记概念都属于 \(C_\mathrm{all}(n)\)。若 \(K_\Sigma\) 全部按定义 5.1 登记，则


\mcsdisplay{C_{\mathrm{all}}(n)=K_\Sigma.}
\noindent \textbf{证明。} 任取 \(\Xi\vdash e_{0}:\tau\) 和概念表示 \(t_c:\beta\)。选不自由出现于 \(e_{0}\) 的新变元 \(z:\beta\)，置


\mcsdisplay{e_c=(\lambda z^\beta.e_0)t_c.}
\noindent \(e_c\) 仍为 \(\tau\) 型，\(\beta\) 规则给出 \(e_c=e_{0}\)。\(t_c\) 是其闭子项，因此 \(c\in C_\mathrm{occ}(e_c)\subseteq C_\mathrm{all}(n)\)。任取 c 得到 \(K_\Sigma\subseteq C_\mathrm{all}(n)\)；反向包含由出现函数值域成立。\(\blacksquare\)

命题版本为 \(\varphi\land(t_c=t_c)\)。\textbf{逻辑有效式完全可能含非逻辑概念}，例如 \(P(a)\lor\neg P(a)\)，不能用“恒真式不含新概念”排除此退化。

此结论相对于指定语言、登记和等价规则，不覆盖不可命名的语义对象。若先删除无用项再扫描，或限制允许变换，得到的是另一函数。

\textbf{推论 5.7（不能恢复依赖）。} 定理条件下，两个毫不相关而均有形式表示的节点也有相同 \(C_\mathrm{all}\)。其包含关系不足以推出演绎、先修、泛化或认知相似。

\(C_\mathrm{all}:N_M\to\mathcal{P}(K_\Sigma)\) 不是 \(\mathcal{P}(K_\Sigma)\) 上的自映射，\(C_\mathrm{all}(C_\mathrm{all}(n))\) 通常类型错误。本稿不称它为闭包算子。

\section{5.3 表达变换与带见证支持}
\textbf{定义 5.8（允许变换系统）。} \(\Theta\) 是带有限见证的表达变换规则族，一步记录为 (e,e',w)。每类规则明确其保持对象：类型、定义等价、可证明等价、特定模型中的外延，或仅启发性表征。只有有相应等价证书者进入形式等价分量。


\mcsdisplay{\mathcal S^{\mathrm{exp}}_\Theta(n)=
\{(\pi,e,C_{\mathrm{occ}}(e)):
e_0\in Rep_0(n),\ \pi:e_0\leadsto_\Theta e\}.}
\noindent 保留变换链 \(\pi\)，而非仅取全部支持之并。若 \(\Theta\) 含任意无关垫加，仍可能饱和；增加索引本身不消除问题。实际局部研究选用有界展开、特定表征库或显式任务相关规则。

\textbf{定义 5.9（四类结构支持）。} 索引 i 包含种类、背景、用途、来源与见证。

\begingroup\small\setlength{\tabcolsep}{4pt}
\begin{longtable}{>{\raggedright\arraybackslash}p{\dimexpr 0.0945\linewidth-2\tabcolsep\relax}>{\raggedright\arraybackslash}p{\dimexpr 0.4482\linewidth-2\tabcolsep\relax}>{\raggedright\arraybackslash}p{\dimexpr 0.4573\linewidth-2\tabcolsep\relax}}
\toprule
\textbf{支持} & \textbf{精确定义} & \textbf{见证边界} \\
\midrule
\endfirsthead
\toprule
\textbf{支持} & \textbf{精确定义} & \textbf{见证边界} \\
\midrule
\endhead
\bottomrule
\endfoot
\bottomrule
\endlastfoot
表达 & \(S_\mathrm{exp}(n,e)=C_\mathrm{occ}(e)\) & 表达合法及节点归属 \\
定义展开 & \(S_\mathrm{def}(n,\pi)=C_\mathrm{occ}(\mathrm{last}(\pi))\) & 指定有限展开链，不默认彻底展开递归定义 \\
证明 & \(S_\mathrm{prf}(n,p)=\bigcup_{e\in\mathrm{Expr}_\mathrm{HOL}(p)}C_\mathrm{occ}(e)\) & 有限证明树经回译后的公式、项与实例化参数 \\
结构路线 & \(S_\mathrm{rte}(n,r)=\bigcup_{a\in\mathrm{Anc}_r(n)}C_\mathrm{iface}(a)\) & 第 06 章已选事件契约及其结构见证 \\
\end{longtable}
\endgroup

\(C_\mathrm{iface}(a)\) 取公开输入输出描述及见证中的登记概念。\(\mathrm{Expr}_\mathrm{HOL}(p)\) 是先收集第 02 章认证证明 p 中出现的有限一阶公式、项及实例化参数，再逐个按定义 2.13 的 r 回译所得的 HOL 项集；未经回译的一阶公式不能直接送入 \(C_\mathrm{occ}\)。\(S_\mathrm{prf}\) 记录\textbf{该证明及指定编码}用了什么，不断言所有证明均不可消去它们。实际开放假设 Hyp(p)、路线输入 Req(r) 和概念支持的类型不同，分别用于演绎、先修与概念比较。

\textbf{定义 5.10（统一支持函数）。}


\mcsdisplay{Supp_M(n,i)\in ISet(K_\Sigma),\qquad
ISet(K)=\{Unknown\}\sqcup\{Known(A):A\in\mathcal P(K)\}.}
\noindent 这里 \(\mathrm{ISet}(K)=\mathrm{Info}(\mathcal{P}(K))\)，Info 是第 07 章的一般带未知标签类型。Unknown 表示公开档案缺少完整结构支持见证，不表示某学习者不知道。个人是否理解这些概念由外部模型评价。

\section{5.4 未知、已知空集与部分界}
\textbf{定义 5.11（信息序与安全集合运算）。} \(\mathrm{Unknown}\sqsubseteq\mathrm{Known}(A)\) 对任意 A 成立；不同精确值不可比较。\(\mathrm{Known}(\emptyset)\) 表示已核实该指定用途不涉及登记域中的概念。

用可能值解释安全运算：\(\gamma(\mathrm{Unknown})=\mathcal{P}(K)\)，\(\gamma(\mathrm{Known}(A))=\{A\}\)。计算所有可能输入上的集合运算；若结果唯一返回对应 Known，否则返回 Unknown。例如


\mcsdisplay{Unknown\cap Known(\varnothing)=Known(\varnothing),\qquad
Unknown\cap Known(\{c\})=Unknown\ (c\in K),}

\mcsdisplay{Unknown\cup Known(K)=Known(K).}
\noindent 不能未经语义证明就规定 \(\mathrm{Unknown}\cap X=X\)。

需要保留部分信息时，使用区间 (L,U)，其含义是 \(\{A:L\subseteq A\subseteq U\}\)。完全未知是 \((\emptyset,K)\)，精确值是 (A,A)。并、交的安全界分别为 \((L_{1}\cup L_{2},U_{1}\cup U_{2})\) 与 \((L_{1}\cap L_{2},U_{1}\cap U_{2})\)；这是安全近似，不承诺保留全部输入相关性。

区间属于独立接口 \(\mathrm{Supp}_M^\mathrm{bounds}(n,i)\)，其值域为 \(\mathrm{Int}(K)=\{(L,U):L\subseteq U\subseteq K\}\)，不直接塞入 ISet(K)。遗忘部分信息的投影为 \(\alpha(L,U)=\mathrm{Known}(L)\)（L=U 时），否则为 Unknown。若某条目同时保存区间与精确值，两者必须一致；只有经过完整性认证的精确区间才可提升为 Known。这样 \(\mathrm{Supp}_M\) 保持定义 5.10 的返回类型，区间接口另保留尚未精确的信息。

\section{5.5 方法节点的统一接入}
\textbf{定义 5.12（方法公开支持）。} 方法 m 的共享材料包括任务接口、显式 Use 标注、公共示范与失效实例。对指定任务版本 v，令 \(\mathrm{Mat}_v(m)\) 为这些材料中已登记的有限 HOL 表达列表；材料节点先取其公开表达，不把节点引用直接当作表达项。对不能直接形式扫描的材料，由 \(\mathrm{Use}_v(m)\) 记录其经过公共锚定的概念。定义


\mcsdisplay{Supp_M(m,v)=Known\left(Use_v(m)\cup
\bigcup_{e\in Mat_v(m)}C_{\mathrm{occ}}(e)\right)}
\noindent 仅在该版本材料具有完整锚定见证时成立；否则 \(\mathrm{Supp}_M\) 返回 Unknown，已证实的部分界另存入 \(\mathrm{Supp}_M^\mathrm{bounds}\)。完整性见证针对列明的材料和用途，不能仅凭上述有限并集已算出，就声称已覆盖一个全局方法的所有适用情境。公共示范中的成功描述任务结果，不默认读取个人记录。

“构造反例”全局方法不必有一个涵盖全部用途的有限支持；其“以常值序列反驳趋近不能取等”实例可有支持 \{序列,极限,常值序列,全称命题\}。局部实例与全局方法通过公共实例化描述连接，个人使用效果在外部层。

定义适用于全部节点的概念接口


\mcsdisplay{\mathfrak C_M(n)=
(C_{\mathrm{all}}^{\mathrm{formal}}(n),\ Supp_M(n,-)).}
\noindent 第一分量只作用于已有形式表示，第二分量容纳形式节点、方法和混合材料。\textbf{不为自然语言启发体虚构“全部逻辑等价表达”关系。} 日后加入方法的形式规格，不等于已经形式化原启发体的全部意义。

\section{5.6 充分支持、极小支持和必要性}
\textbf{定义 5.13（相对用途的充分性）。} 选定公共用途 u 和独立见证检查关系 \(\mathrm{Verify}_u(n,A,w)\)，定义


\mcsdisplay{Adeq_u(n,A)\iff\exists w\,Verify_u(n,A,w),\qquad
\mathcal A_u(n)=\{A:Adeq_u(n,A)\}.}
\noindent 用途可以是“以这些假设完成形式证明”或“按声明契约完成结构路线”。经验性的“有助理解”须在外部 E 中定义。根据用途，支持域可以是概念或前提节点，但必须声明。

\textbf{定义 5.14（极小支持族）。}


\mcsdisplay{MinSupp_u(n)=
\{A\in\mathcal A_u(n):
\nexists B\in\mathcal A_u(n)\ (B\subsetneq A)\}.}
\noindent 它是集合族，不是预设唯一集合。包含极小与基数最小是不同优化准则。

\textbf{命题 5.15（存在性与计算条件）。} 若 \(\mathcal{A}_u(n)\) 含一个有限成员 \(A_{0}\)，则存在包含极小成员，无须充分性可判。若候选域有限且充分性可判，则全部极小支持可有限枚举。

\textbf{证明。} 在非空有限族 \(\{A\subseteq A_{0}:A\in\mathcal{A}_u(n)\}\) 中取基数最小者 A*。任何充分真子集仍在该有限族中且基数更小，矛盾。第二结论通过枚举有限幂集并删除含充分真子集者得到。\(\blacksquare\)

非唯一例：充分支持为包含 \{a\} 或 \{b\} 的集合，极小成员是 \{a\} 与 \{b\}。不存在例：自然数集的所有余有限子集，每个成员删去一个元素后仍在族中。若充分族为空，目标相对该用途不可实现，不能对空路线族作全称量化后宣称所有节点均为必要前提。

\textbf{命题 5.16（概念信息不足以决定语义）。} 相同概念出现集不保证两谓词外延存在包含关系。

\textbf{证明。} 在含 0,1 的二元素域上，令 \(P(x)=(x=0)\land(1=1)\)，\(Q(x)=(x=1)\land(0=0)\)。相对于登记的基本概念 \{等号,0,1\}，二者出现集相同，外延却为互不包含的 \{0\} 与 \{1\}。若复合表达也登记为概念，则其直接出现集未必相同，但定理 5.6 的 \(C_\mathrm{all}\) 仍相同；相同全等价概念集依然不能决定包含。\(\blacksquare\)

\section{5.7 文献接口与研究地位}
定义扩张与消去参照本地 EFT VIII §3（印刷 125–128 / PDF 131–134）和 Manin II §8（印刷 66–69 / PDF 80–83），逐项定位见附录 B。\textbf{概念登记、饱和定理、支持信息域与方法接入是本稿提出的定义及条件性推论}，不宣称这些教材已经给出了 MCS 理论。

\clearpage\chapter{06 硬关系、行动契约与结构性路线}
\mcsorigin{原章节：06-硬关系.md}
硬关系描述固定数学背景中的认证推导、指定结构契约下的必要资源，以及有见证的泛化。\textbf{硬不表示“所有人心理上必须先学”}。真实学习所需条件由外部模型评价。本章先定义行动与路线，再定义先修；不引用个性化规划的输出建立自己的前提。

\section{6.1 演绎是带联合前提的关系}
\textbf{定义 6.1（前提接口）。} 固定理论 T、类型环境 \(\Xi\)。令 \(P=(n_{1},\ldots,n_k)\) 为命题节点的有限带标签族，其命题负载为 \(\varphi_{1},\ldots,\varphi_k\)；\(\Gamma_P\) 是对应的开放假设环境。相同数学命题可以有不同来源标签。结论 n 的负载为 \(\psi\)。

\textbf{定义 6.2（认证演绎）。}


\mcsdisplay{Ded_{T,\Xi}(P,n,p)
\iff Check_T(\Xi,\Gamma_P,\psi,p)=1.}
\noindent Check 是第 02 章固定演算的有限证书检查器。对无限公理呈现，证明还须携带可检查的公理成员见证。擦去 p 得到 \(\mathrm{Ded}(P,n)\equiv\exists p\) Ded(P,n,p)，它一般仅可半判定，不承诺可判。

对应高阶特征函数的类型为


\mcsdisplay{Ded:PremFam\to Node\to ProofCode\to o.}
\noindent PremFam 可以在元理论定义为有限带标签节点列表，并通过档案编码接入对象语言。把它画成二元边是特定显示约定，不能改变联合语义。

\textbf{命题 6.3（相对语义可靠性）。} 若 Ded(P,n,p)，且 H 是满足 T 的 Henkin 模型，赋值 v 使全部 \(\varphi_i\) 为真，则 \(\psi\) 在 H,v 中为真。

\textbf{证明。} Check 成功给出第 02 章演算的有限推导。按第 03 章可靠性定理，每个规则保持相应假设下的满足关系。由假设全部成立得到结论成立。\(\blacksquare\)

这里认证的是条件推导，不是无条件真理；背景 T 无模型时不会得到“存在一个真实知识世界”的结论。

\textbf{命题 6.4（弱化与组合）。} 在背景、类型和假设标签相容时，演绎可弱化，也可用前提的证明替换开放假设，合成为一份有限证明。

\textbf{证明。} 弱化保留原证明并增加未使用的允许假设。组合时先 \(\alpha\) 改名局部变元，保证本征变量不被新假设捕获，再以给定证明替换相应开放叶；自然推演的替换/切规则给出合法推导。替换次数有限，故新证书有限。\(\blacksquare\)

因此 P 中某项可能没有被实际使用。另以 Hyp(p) 记录证明真正开放的假设叶；即使某假设出现在 Hyp(p)，它也未必对所有证明必要。\textbf{不能由 Ded(P,n,p) 直接推出 P 每个成员都是 n 的硬先修。}

\section{6.2 行动契约不预设学习路线}
\textbf{定义 6.5（公开行动契约）。} 固定本体 M，\(\Lambda_M\) 是经公开模式生成的契约族。契约


\mcsdisplay{a=(id_a,I_a,O_a,w_a,mode_a)}
\noindent 具有有限输入 \(I_a\subseteq N_M\)、有限输出 \(O_a\subseteq N_M\)、结构见证 \(w_a\) 和模式。输入输出是可引用的知识资源，不是个人掌握布尔值。

本章把这些资源视为可重复引用、使用后不消耗的资源，因而集合接口足够。限用次数、设备占用、时间或相互冲突的行动若影响可执行性，须另给带数量或状态的契约；不能直接套用下面的集合可达性结论。

\begingroup\small\setlength{\tabcolsep}{4pt}
\begin{longtable}{>{\raggedright\arraybackslash}p{\dimexpr 0.1433\linewidth-2\tabcolsep\relax}>{\raggedright\arraybackslash}p{\dimexpr 0.5029\linewidth-2\tabcolsep\relax}>{\raggedright\arraybackslash}p{\dimexpr 0.3537\linewidth-2\tabcolsep\relax}}
\toprule
\textbf{模式} & \textbf{独立的合法性义务} & \textbf{不承诺的内容} \\
\midrule
\endfirsthead
\toprule
\textbf{模式} & \textbf{独立的合法性义务} & \textbf{不承诺的内容} \\
\midrule
\endhead
\bottomrule
\endfoot
\bottomrule
\endlastfoot
认证推导 & 输出命题均有以指定输入和背景为假设的检查通过证明 & 学习者自动理解证明 \\
定义引入 & 新符号、新鲜性、类型及保守定义证书满足规定 & 该定义是唯一教学入口 \\
构造/示例 & 输出满足所给规格的形式证明，或声明的可检查有限见证 & 任意无限模型性质均可判 \\
表征转换 & 输入输出之间有指定层次的翻译/等价证书 & 两种表征认知成本相等 \\
方法/任务呈现 & 公共任务与方法接口匹配，输入输出锚定有效 & 使用方法必然成功或必然学会 \\
\end{longtable}
\endgroup

各模式的合法性是关于有限记录及其见证的独立判断，不以存在一条已经规划的路线为条件。可枚举候选契约，并在检查后加入认证部分；全体数学可行契约不必可判。

公开契约可以\textbf{规定}多项联合输入。对此“必要”意为执行该契约所承担的输入义务，不自动意味着这些条件在数学或心理上全局不可替代。若要求逻辑最小前提，必须调用第 05 章相应用途的极小充分支持定义，另承担其存在与计算义务。

\section{6.3 结构性路线的生成}
\textbf{定义 6.6（有根结构路线）。} 给定公开可用背景 \(B\subseteq N_M\) 和有限目标 \(G\subseteq N_M\)，路线 r 包括：

\begin{enumerate}
\item 有限事件集 \(P_r\)，每个事件 p 标记一个 \(\mathrm{act}(p)\in\Lambda_M\)；
\item 对每个 \(m\in I_\mathrm{act}(p)\)，指派 h(p,m) 为背景来源 B 中的 m，或另一个事件 q 的输出 \(m\in O_\mathrm{act}(q)\)；
\item 由所有事件来源指派产生有向边 \(q\to p\)；要求此图无有向环；
\item 对每个 \(g\in G\)，指定目标来源 \(\zeta(g)\)：背景中的 g 或输出含 g 的事件。
\end{enumerate}
路线见证保存所有指派及契约证书，记其集合为 \(\mathrm{Rte}_M(B,G)\)。允许重复使用同一契约，每次使用有不同事件标识。事件的严格依赖由图的传递闭包给出，自反传递闭包记为 \(\preceq_r\)。

\textbf{命题 6.7（路线偏序与结构可执行性）。} \((P_r,\preceq_r)\) 是有限偏序；任一线性扩张依序执行时，每个事件所需的公开输入均由背景或已完成事件提供。

\textbf{证明。} 自反性、传递性来自闭包；若不同 p,q 相互可达，则原图有环，与条件矛盾，故反对称。有限无环图有入度零顶点；依次删除得到拓扑排序。任何线性扩张都把 h 指派的事件来源排在当前事件之前，背景来源从开始可用，因此按事件数归纳得到全部输入可用。\(\blacksquare\)

这只证明资源接口可满足。把“提供一个证明”当成“学习者已经掌握该证明”是额外认知假设，不由本命题给出。遗忘、疲劳等也不属于这个结构可执行性结论。

\textbf{定义 6.8（目标活跃部分与资源支持）。} 对 \(g\in G\)，若 \(\zeta(g)\) 为事件 p，取 p 及其全部依赖祖先为 \(\mathrm{Act}_r(g)\)，置


\mcsdisplay{Req_r(g)=\bigcup_{q\in Act_r(g)} I_{act(q)}.}
\noindent 若 \(\zeta(g)\) 是背景来源，则 \(\mathrm{Act}_r(g)=\emptyset\)，并以单独背景标签记录 g 已在 B 中；约定 \(\mathrm{Req}_r(g)=\emptyset\)。背景依赖 \(B_r(g)\) 是所有活跃指派中从 B 取得的资源。区分“没有新的学习事件”与“没有背景假设”。

无关旁支不进入 \(\mathrm{Req}_r(g)\)。同一节点的多个事件仍分别保存，其支持可以不同。概念级路线支持由第 05 章把活跃契约映到概念域，不能与这里的节点输入集混同。

\section{6.4 路线内硬先修与路线族共同硬先修}
\textbf{定义 6.9（契约相对的硬先修）。}


\mcsdisplay{HardPre_r(m,g)\iff m\in Req_r(g).}
\noindent 含义是：在该路线已经选定的行动契约与来源指派中，完成目标承担资源 m 的输入义务。若只关心“必须安排学习事件”的先修，另定义


\mcsdisplay{NewPre_r(m,g)\iff HardPre_r(m,g)\land
\text{某个活跃输入的 m 来源是事件而非背景}.}
\noindent \(\mathrm{HardPre}_r\) 不应直接当作节点上的偏序；同一节点复习时可以既是先前产出又是后续输入。偏序位于事件集。

这里的输入义务也不等于“必须在起点额外提供”。例如契约 \(\emptyset\to\{m\}\)、\(\{m\}\to\{g\}\) 给出 \(\mathrm{HardPre}_r(m,g)\)，但 m 可在路线内生成，起点无须提供 m。若研究初始背景的必要性，应对背景来源 \(B_r(g)\) 或允许背景族另行量化，不能用 \(\mathrm{Req}_r(g)\) 的交集回答该问题。

\textbf{定义 6.10（路线族共同必需）。} 对显式声明的\textbf{非空}路线族 \(\emptyset\ne\mathcal{R}\subseteq\mathrm{Rte}_M(B,G)\)，定义


\mcsdisplay{HardPre_{\mathcal R}(m,g)
\iff \forall r\in\mathcal R\ HardPre_r(m,g)
\iff m\in\bigcap_{r\in\mathcal R}Req_r(g).}
\noindent 空族返回“该路线域中没有可行路线”，不作全称命题真空化后生成先修事实。路线族可由全部公共契约、某类允许表征或公理预算限定，其标签必须随关系保存。宣称对全部数学路线必需，远强于对当前登记路线族必需。

\textbf{命题 6.11（族缩小增加共同约束）。} 若 \(\emptyset\ne\mathcal{R}_{1}\subseteq\mathcal{R}_{2}\)，则


\mcsdisplay{\bigcap_{r\in\mathcal R_2}Req_r(g)
\subseteq \bigcap_{r\in\mathcal R_1}Req_r(g).}
\noindent \textbf{证明。} 左边成员属于 \(\mathcal{R}_{2}\) 每条路线的支持，故也属于其子族 \(\mathcal{R}_{1}\) 的每条支持。\(\blacksquare\)

个人状态筛选的 \(\mathcal{R}_{E,s}\) 是派生路线族。其共同必经点可以增加，\textbf{不得反向登记成原 \(\mathcal{R}\) 的普遍硬先修}。这正是本体与外部模型分离后必须保留的下标。

\textbf{命题 6.12（替代先修不能压成联合先修）。} 若 a,b 是互异资源，目标 c 有两条路线的支持分别为 \{a\}、\{b\}，则共同必需集为空，但不存在由空资源执行这两条已声明路线之一的方式。

\textbf{证明。} \(\{a\}\cap\{b\}=\emptyset\)；第一路线的输入义务是 a，第二路线的是 b。没有任一输入时两份契约都不适用。把两者合并成 \{a,b\} 会错误增强要求，把交集空解释成无条件可行则错误减弱要求。\(\blacksquare\)

检查有限路线族的共同必需可枚举；检查无限族中的全称必要性通常不能由一份路线证书完成。只有有限穷尽证书、独立不变量证明或相应决定过程，才支持全族认证。

\section{6.5 硬泛化、特化与不同载体}
\textbf{定义 6.13（同载体硬泛化）。} 设 g,s 是同一背景 T 中、同一对象类型 \(\alpha\) 上的概念，其谓词为 \(P_g,P_s:\alpha\to o\)。定义


\mcsdisplay{HGen_T(g,s,p)\iff
Check_T\bigl(\vdash\forall x^\alpha(P_s(x)\to P_g(x)),p\bigr)=1.}
\noindent 记 \(g\succeq_H\) s 表示“g 比 s 一般”。特化是逆关系：\(\mathrm{HSpec}(s,g,p)\equiv\mathrm{HGen}(g,s,p)\)。其集合值函数形式是 \(\mathrm{Gen}(s)=\{g:g\succeq_H\) s\} 与 \(\mathrm{Spec}(g)=\{s:g\succeq_H\) s\}，不要求单值可逆。

\textbf{命题 6.14（语义包含、预序和商）。} 在满足 T 的模型中，HGen 的证书保证 \(P_s\) 的外延包含于 \(P_g\)。存在证书的关系自反、传递，通常只是预序。若对相互泛化取等价商，则诱导偏序，但本体节点仍保留原身份。

\textbf{证明。} 语义包含由可靠性。自反由 \(P\to P\)；传递由两份全称蕴含的实例化及命题推理合成。若 \(g\succeq s\) 且 \(s\succeq g\)，它们谓词可证等价，记录仍可不同，故节点层未必反对称。在商上，相互可达恰是等价类相等，因此反对称成立。\(\blacksquare\)

\textbf{定义 6.15（跨载体的结构泛化）。} 对不同类型或理论呈现，不强写外延包含。给出明确结构映射 \(F:\mathrm{Obj}(s)\to\mathrm{Obj}(g)\)，并证明它把每个 s-结构送到 g-结构；如忘掉光滑结构而保留底层拓扑流形。若使用范畴，另声明作用于态射和保持恒等/复合的证据。

这种关系携带 F，表示“经指定遗忘/解释得到的一般结构”。\textbf{任意模型论解释不自动等于泛化}：一个理论解释另一个理论，并不保证载体相同、对象类包含，或只删除结构。一般理论解释作为第 08 章独立关系保存。

同载体泛化与跨载体结构泛化可统一为带模式标签的关系和见证族，但不能删掉标签后继承全部同一预序性质。

\textbf{例 6.16（逆关系非函数）。} 在共同的群结构编码类型上，群的特化包括阿贝尔群与有限群。\(\mathbb{Z}\) 是无限阿贝尔群，\(S_{3}\) 是有限非阿贝尔群，所以二者互不包含，Spec(群) 不只有一个元素。有限阿贝尔群给出一个共同特化；本例不声称“不存在共同特化”或“任何意义下均无最小共同特化”。

\section{6.6 关系函数及其证据类型}
\textbf{定义 6.17（特征函数与集合值函数）。} 本章每种关系 R 可视为到 o 的高阶谓词，或擦去见证后的集合值函数。例如


\mcsdisplay{DedSucc(P)=\{n:\exists p\,Ded(P,n,p)\},\qquad
Pre_{\mathcal R}(g)=\bigcap_{r\in\mathcal R}Req_r(g).}
\noindent 它们的数学存在不等于能完整计算输出。概念支持函数提供“这份证据涉及什么”的附加函数，但\textbf{不能仅由概念集反推出关系的真值}。本章的证明检查、契约来源指派和结构蕴含见证，正是所缺的语义条件。

\section{6.7 演绎闭包和本体一致性见证}
\textbf{定义 6.18（有限前提演绎闭包）。} 在固定语句域 \(\mathrm{Sent}_\Sigma\) 上，对 \(X\subseteq\mathrm{Sent}_\Sigma\) 定义


\mcsdisplay{Cn_T(X)=\{\varphi:\exists\Delta\subseteq_{\rm fin}X,\
T;\Delta\vdash\varphi\}.}
\noindent \textbf{定理 6.19（代数闭包）。} \(\mathrm{Cn}_T\) 广延、单调、幂等，而且


\mcsdisplay{Cn_T(X)=\bigcup_{\Delta\subseteq_{\rm fin}X}Cn_T(\Delta).}
\noindent \textbf{证明。} 广延由假设规则，单调由弱化。若 \(\varphi\in\mathrm{Cn}_T(\mathrm{Cn}_T(X))\)，其有限证明仅用有限个中间假设 \(\psi_{1},\ldots,\psi_k\)；每个 \(\psi_i\) 来自 X 的有限子集 \(\Delta_i\)。由命题 6.4 合成得到以有限并 \(\bigcup\Delta_i\subseteq X\) 为假设的 \(\varphi\) 的证明，故 \(\varphi\in\mathrm{Cn}_T(X)\)。反向由广延。有限性等式直接由定义的有限见证成立。\(\blacksquare\)

这给出固定背景下的推演闭包，不是学习者状态闭包。公开行动的资源可达闭包可类似通过有限输入 Horn 规则定义，但它只解释契约资源生产，不能无标签地识别为逻辑后承。

\textbf{命题 6.20（一个非空核心实例）。} 若背景 T 有模型，则存在符合本章接口的非空 MCS 核心实例：取有限个合法节点、若干恒等/已认证推导契约、有限公共模板，概念支持按第 05 章定义，路线按定义 6.6 生成；不加入未经认证的硬关系。

\textbf{证明。} 取 T 中一个合法闭公式及其反身蕴含的证明，登记其 Claim 与 Proof 节点；以一份输入该公式、输出同一公式的契约为例。给定包含该输入的 B，可建立单事件路线，其依赖图为空，所有输入由 B 提供。其余接口取已定义的有限记录或空关系族。没有额外公式要求 T 的模型违反自身公理，故给出一致的结构记录实例；逻辑语义由原 T 模型解释。\(\blacksquare\)

此为相对于 T 可满足性的模型见证，不是在 T 内证明 T 自身一致。完整类型语义域通常无限，有限的是登记网络而非整个 HOL 模型。

\clearpage\chapter{07 通用认知结构、软关系与外部认知模型}
\mcsorigin{原章节：07-软关系与认知状态.md}
本章服从 D　接口契约。本体 \(M\) 保存共享的数学内容、方法、表征、任务、误区模式、行动契约和关系描述；外部模型 \(E\) 解释这些对象对于状态 \(s\) 的学习者意味着什么。认知模型的\textbf{描述}可以成为节点，其运行时个体状态仍在本体之外。

\section{7.1 通用关系与个体评价分离}
\textbf{定义 7.1（通用关系描述）。} \(r\in R_M^0\) 是带类型的关系记录，含源与目标、关系种类、适用任务、公开背景、表征对应、结构见证或候选依据、已知失效条件及来源，不含学习者标识。字段分别标注数学证明、引用结果、经验资料、建模假设、解释材料与未知项，不组成单一可信度全序。

\textbf{定义 7.2（外部评价）。} 接口


\mcsdisplay{\operatorname{Eval}_E(s,r,ctx)}
\noindent 输出该模型声明的评价对象，例如适用性、成本偏序、成功率区间或 Unknown；也可输出集合或分布，不强制为单一实数。\(ctx\) 显式列出任务、目标、表征、预算和证据版本。给定外部判据 \(\pi\)，定义


\mcsdisplay{\operatorname{Soft}^{\pi}_{E,s}(r,ctx)
\Longleftrightarrow
\pi(\operatorname{Eval}_E(s,r,ctx))=\mathsf{Accept}.}
\noindent 通用描述存在不表示此人的评价成立；未知值的处理由 \(\pi\) 明示。派生软关系不进入数学演绎闭包。

\section{7.2 软先修、泛化、应用、类比与对偶}
\textbf{定义 7.3（软先修）。} 公开候选记录 \(m,n\) 与可比较的教学行动及控制条件。外部模型比较先接触 \(m\) 与对照方案对目标 \(n\) 的成本或表现。指定因果模型时可定义


\mcsdisplay{\Delta_E(s)=P_E(Y_n=1\mid do(a_m),s,ctx)
-P_E(Y_n=1\mid do(a_0),s,ctx).}
\noindent 此式不自动识别因果效应；缺乏模型、识别条件或数据时为未知。观察到 \(P(Y_n=1\mid Y_m=1)\) 较大只能支持相关，不能替代上式。软先修不默认传递、反对称或普遍有效，也不产生第 06 章的结构必经资源。

\textbf{定义 7.4（软泛化/特化）。} 通用描述记录实例—模式的抽象或实例化对应、保存结构和丢失特征。特化是同一对应的逆向阅读，不要求成为单值函数；两个方向的教学效果分别评价。有同载体蕴含或合格遗忘见证时另立硬关系，不能用软效果代替证明。

\textbf{定义 7.5（应用）。} 应用描述含源对象/方法接口、目标任务、输入解释、输出解释和条件检查。严格应用可携带代入与证明证书；启发性应用只给候选方案。数学应用合法性留在认证层，使用该应用教学的效果属于外部评价。

\textbf{定义 7.6（类比）。} 类比由带关系结构 \(A,B\)、部分对象对应 \(f\)、拟保持关系族 \(J\)、已认证子族 \(J_{\rm ok}\)、已反驳子族 \(J_{\rm fail}\) 和未决部分组成。各关系的元数、方向、量词和有效域必须列明；未知保持性不计入已认证。仅保持某些关系不能宣称为保持所有结构的同态。类比提出的结论先登记为问题或候选，另行证明。

\textbf{定义 7.7（对偶）。} 描述可引用严格的反变等价、Galois 连接或指定类型的对偶化操作；只有满足相应公理才使用这些名称。严格变换是数学事实，经它迁移学习的效果另由 \(\operatorname{Eval}_E\) 评价。不是所有反向对应都可逆，也不是所有数学对偶都仅是软相似。

\section{7.3 外部状态、观察、转移与多对多锚定}
\textbf{定义 7.8（外部认知模型）。} \(E\) 指定状态空间 \(X_E\)、条目域 \(Item_E\)、结果域 \(Y_E\)、外部行动域 \(A_E\)、观察/转移/评价接口及版本。状态可以是集合、符号结构、偏序状态或潜在向量，\textbf{不规定必须是六元组}。最小关系接口为


\mcsdisplay{\operatorname{Obs}_E\subseteq X_E\times Item_E\times Y_E,\qquad
\operatorname{Trans}_E\subseteq X_E\times A_E\times X_E.}
\noindent 部分函数是特例；采用概率模型时另指定可测结构，并用观察核 \(O_E(dy\mid s,i)\)、转移核 \(K_E(ds'\mid s,a)\) 实现。核的归一性是该模型的公理，不是所有认知模型的公理。允许学习、遗忘、复习、修正和非单调状态变化；不要求学后状态包含全部旧掌握项。

\textbf{定义 7.9（锚定）。}


\mcsdisplay{\operatorname{Anchor}_E\subseteq Item_E\times N_M\times Competence_M.}
\noindent 能力维度可含实例辨认、定义陈述、计算、证明、动机解释、表征转换、迁移和自我诊断。一个条目可锚定多个节点和能力；同一节点—能力对可有多个观测条目。锚定说明观测指向，不自动证明测量效度。个体作答、掌握推断、观测误差均留在 \(E\)。

\textbf{边界 7.10（非全知）。} \(T\vdash\varphi\) 不推出任何人掌握 \(\varphi\)。能够陈述群公理不意味着能够证明 Cayley 定理。数学推演闭包与掌握更新不是同一运算。

\section{7.4 外部相对认知原子性}
\textbf{定义 7.11（可认知与允许分解）。} 固定 \(E,s\)、目标规格 \(g\)、表征族 \(\mathcal R\)、允许分解规则 \(\mathcal D\)，外部模型明确指定：

\begin{enumerate}
\item \(\operatorname{Ind}_E(s,n;g,\mathcal R)\)：存在允许表征，使 \(n\) 的目标接口可在声明背景下单独呈现、评估，并满足模型给定的可认知判据。判据可以是约束或带阈值的预测，须公开；不要求零先修。
\item \(U\mathrel{\mathcal D_g}n\)：\(U\) 是至少两个互异单位组成的有限集合，每个单位是该分解规则规定的真部分，并有重构 \(n\) 目标接口的契约见证。改名、复制 \(n\) 和添加无关单位不算真分解。
\item \(\operatorname{Real}_E(s,U\Rightarrow n;g,\mathcal R)\)：存在由这些单位的独立呈现/评估与重构契约组成的有限行动见证，满足目标的覆盖与效度条件。见证同时指定各单位的局部目标、出现次序及对应外部中间状态；后续单位可使用较早输出。各步的可认知判据在这些局部目标和状态下检验，不能把整体验收目标 g 原样赋给每个部分。
\end{enumerate}
这些是\textbf{模型中的严格谓词}，不声称是可判的心理事实。目标 \(g\) 必须给出能力与覆盖条件，不能把复述名称当成掌握全部意义。

采用非确定或概率转移时，还须固定行动策略和 Real 的成功语义，例如存在成功轨迹、所有允许轨迹均成功，或成功概率达到指定阈值。三者给出的原子性不同；单条成功样本不能认证后两者。有限见证只表示策略/行动的有限描述，不自动提供其全轨迹性质的决定过程。

\textbf{定义 7.12（相对认知原子）。}


\mcsdisplay{\begin{aligned}
\operatorname{Atom}_E(s,n;g,\mathcal R,\mathcal D)
\Longleftrightarrow{}&
\operatorname{Ind}_E(s,n;g,\mathcal R)\\
&\land\neg\exists U\,[U\mathrel{\mathcal D_g}n
\land\operatorname{Real}_E(s,U\Rightarrow n;g,\mathcal R)].
\end{aligned}}
\noindent 这是相对于允许真分解的极小性，不是字数最短，也不要求极小单位唯一。有限候选、可判判据下可穷举；一般原子存在性和可判性需要另证。良基分解关系阻止无限下降，但不能单独保证终端单位满足可认知判据。认证 Atom 需要确认 Ind 并排除所有允许的可实现真分解；确认 Ind 为假或找到一个这样的真分解即可反驳 Atom，其余信息尚不充分时才返回未知。未发现分解不构成极小性证明。

同一流形节点可在一个外部模型中拆为图册、过渡映射和相容性目标，在另一个模型中整体呈现。差异不改变节点身份；自动生成的候选也不自动成为心理原子。

\section{7.5 一般误区模式与个体实例}
\textbf{定义 7.13（误区模式）。} \(p\in Pattern_M\) 保存目标概念、可能的错误规则/表征、触发情境、辨析任务、反例、否定范围和来源。“收敛序列不能取到极限”及常值序列反例是共享模式。数学反驳是结构证据，模式流行程度是另附经验资料。

\textbf{定义 7.14（个体实例）。} \(\operatorname{Inst}_E(s,p)\) 是外部判断，附观测、时间、解释模型和不确定性；不是本体新公理。答对一道题不必消除相关的所有误解。共享图形/比喻属于公共材料，此人当前激活的概念意象及其与定义的协调程度属于外部解释。

\section{7.6 非干扰与外部保守扩张}
\textbf{契约 7.15（只读适配）。} \(D(M,E,s,g,b,\rho)\) 仅读取固定版本 \(M\)，不修改节点身份、数学负载、证书及硬关系。新增公共解释材料是单独版本变更，不是状态更新的隐含副作用。

\textbf{命题 7.16（结构非干扰）。} 固定 \(M\)，替换 \(E,s,g,b,\rho\) 不改变核心记录及已认证硬关系。再固定一个核心解释 A，若两次外部接入均以 A 为约化，则任一核心公式在相同赋值下的解释不变。

\textbf{证明。} 只读契约使核心记录和证明检查的输入不变，故认证关系保持。语义部分的两次核心约化都是 A，且外部状态符号不在核心签名内；对核心项和公式结构归纳即可。固定一个理论档案不意味着该理论只有一个模型，本命题不比较两个不同核心模型中的真值。结论依赖只读契约，不是对任意实现的保证。\(\blacksquare\)

\textbf{命题 7.17（外部扩张的语义保守性）。} 设核心理论 \(T_M\) 的扩张为 \(T_{M,E}\)。固定第 03 章的一种语义，假定新增外部类型/符号不重解释核心符号；\textbf{每个}核心模型 \(A\models T_M\) 都存在扩张 \(A^+\models T_{M,E}\)，且 \(A^+|_{\Sigma_M}=A\)；接口不排除核心原有模型。则对任意核心句子 \(\varphi\)，


\mcsdisplay{T_{M,E}\models\varphi\quad\Longleftrightarrow\quad T_M\models\varphi.}
\noindent \textbf{证明。} 右向左由约化。反之若 \(T_M\not\models\varphi\)，取反模型 \(A\)，按假设扩张为 \(A^+\)；核心解释未变，故 \(A^+\not\models\varphi\)，与扩张蕴含矛盾。\(\blacksquare\)

逐模型可扩张是本命题采用的\textbf{充分条件}，常称模型保守性；它比“没有新增旧语言后果”更强，不是后一性质的一般必要条件。不能仅因某个旧模型不能原样扩张就宣布理论不保守；证明不保守须给出新增的旧语言后果或其他充分反证。

仅把状态放在另一张表中不足以证明保守性。例如外部接口强迫核心节点数有限，就可能排除原有模型。允许单状态、无观测、无有效锚定、评价全未知的哑外部模型，并使接口约束对其成立：当外部层只是元层只读关系记录时，可直接与每个核心模型配对；需要非空条目域时放一个未锚定占位条目。若把外部类型内部化为新的 HOL 基本类型并允许混合箭头，则还须证明相应类型域扩张与抽象闭合，不能仅凭“单状态”省略该义务。禁止哑模型而仍沿用本判据时，必须另证扩张条件；也可另给直接的语义或语法保守证明。以上是语义保守性；据此推出 \(\vdash\) 层面的保守性另需演算可靠性及对应完备性，不能据此声称标准高阶语义具有有效完备演算。

\section{7.7 冷启动、可复现性与外部知识空间实例}
\textbf{定义 7.18（未知值）。}


\mcsdisplay{\mathsf{Info}(A)=\{\mathsf{Unknown}\}\sqcup
\{\mathsf{Known}(a):a\in A\}.}
\noindent 这是外部接口的信息标签，不是对象逻辑新增真值。未作答不等于答错；无误区证据不等于没有误区；未知支持不等于 \(\mathsf{Known}(\varnothing)\)。未知布尔评价既不认证接受，也不认证反驳。冷启动可给出多条公共结构路线并标效果未估计；使用先验须明示来源和版本。

\textbf{命题 7.19（派生可复现性）。} 固定核心、模型及算法版本和全部显式输入，确定性输出相同。随机算法固定种子时样本相同，未固定种子时只要求分布相同。证明是函数/核的外延性；隐含档案读取不满足前提。\(\blacksquare\)

\textbf{定义 7.20（有限知识空间实例）。} 外部条目域为有限集 \(Q\)，状态族 \(\mathcal K\subseteq\mathcal P(Q)\) 含 \(\varnothing,Q\)。并封闭时是知识空间；若满足逐项可达和学习一致性，可使用有限学习空间理论。条目经 Anchor 联系节点与能力，不直接等于所有数学对象。例如


\mcsdisplay{\mathcal K=\{\varnothing,\{a\},\{b\},\{a,b\},\{a,c\},\{b,c\},\{a,b,c\}\}.}
\noindent 它并封闭，每个非空状态可删去一项而仍合法；但 \(\{a,c\}\cap\{b,c\}=\{c\}\notin\mathcal K\)。偏序的全部下闭集对交封闭，故一个条目偏序不能完整表示此 OR 路线实例。规划应先选方案再输出事件偏序，或将状态族按包含排序。并闭、无遗忘及作答模型都是所选 \(E\) 的条件，不是本体公理。

\section{7.8 原作者来源与核验}
\begin{itemize}
\item Doignon–Falmagne，\href{https://arxiv.org/pdf/1511.06757}{作者稿}，§2，PDF 第 3–5 页 Definitions 1、3 给出条目、状态族和学习空间条件；§3，PDF 第 10 页 Definition 9、Theorem 10 区分一般空间与偏序对应的 ordinal space；§9，PDF 第 28 页起区分状态和有噪声作答。这里只用作外部模型实例。
\item Tall–Vinner，\href{https://wrap.warwick.ac.uk/id/eprint/507/1/WRAP\_Tall\_dot1981a-concept-image.pdf}{1981 年作者机构稿}，PDF 第 3 页 “evoked concept image” 段及第 4 页极限段，讨论情境激活与“趋近不能等于”的附加信念。支持意象/定义及模式/个体的区分，不构成本章 Atom 谓词的实证证明。
\item Gentner，\href{https://groups.psych.northwestern.edu/gentner/papers/Gentner83.2b.pdf}{作者机构稿}，印刷第 158 页三条映射规则及后段，提出对象关系、高阶关系与系统性的重要性。本稿据此设计类比接口，个体效果仍待另证。
\end{itemize}
\textbf{验收 7.21。} 两个外部学习者可以获得不同入口与提示，同时共享节点、正式定义和证明；冷启动产生未知；一般误区模式不因某人的错误而改变；相关性不生成硬先修；采用命题 7.17 认证保守性时，必须核查每个核心模型可扩张。

\section{7.9 极小分解存在与唯一性的充分条件}
\textbf{定义 7.22（可行分解偏序）。} 固定前述外部参数，令 F 是达到同一目标的可行分解呈现族。每个呈现保留行动见证、局部目标及中间状态，不能只由单位集合决定。将重命名和已认证的相同分解合并；以 \(U\le V\) 表示 U 是 V 的允许细化。要求这个关系确实是偏序：不能把表征循环误当作严格细化。元素的单位属于本体允许的部件，是否可行由 E 判断。F 可以为空，也可能有无限下降链。

\textbf{命题 7.23（有限存在与有条件唯一）。} 若 F 非空且有限，则每个 \(V\in F\) 下有极小可行分解。若再满足下向有向性：任意 \(U,V\in F\) 有 \(W\in F\) 使 \(W\le U,V\)，则 F 有唯一极小元，且它是最小元。

\textbf{证明。} 从 V 逐次取严格更小者，有限性迫使停止，所得即极小。若 U,V 均极小，共同细化 W 同时小于二者，极小性给 W=U=V。任意元素之下存在极小元，故这个唯一极小元小于所有元素。\(\blacksquare\)

从极小呈现推出 Atom 还需要\textbf{局部替换条件}：对呈现中单位 u 的每次使用，固定该次的局部目标 \(g_u\)、输入状态 \(s_u\)、允许表征和分解规则；每个满足该局部 Real 的真分解都可替换此次使用，并保持后续行动及整体目标可行，且所得呈现在 F 中严格下降。若极小呈现的每个单位又在这些局部参数下满足 Ind，则不存在相应的局部 Real 真分解，否则替换将违反极小性；故各单位在其各自参数下满足 Atom。

不能省略后续可行性的保持：拆解即使能单独教会一个单位，也可能消耗预算或改变状态，使后续行动失败。此时只有“在该整条方案内不能进一步细化”，推不出定义 7.12 的局部原子性。Ind 也不能省略：不可理解的终端字符串同样可能无法再分。

有限性可换为该细化序的良基性并分别证明可行族非空；唯一性仍需共同细化等额外条件。只有两种互不相容的极小分解时，存在性成立而唯一性失败；本稿允许保留二者，不强迫学习者共享唯一心理粒度。

\clearpage\chapter{08 关系运算、生成代数与结构保持}
\mcsorigin{原章节：08-关系代数与结构.md}
本章分开三个任务：泛代数处理生成、闭包与商；模型论处理语言、解释和满足；范畴论处理可组合的结构保持变换。三者均不自动赋予数学变换以教学效果。术语参考 Burris–Sankappanavar，I §5、II §§3、5、6、10；以下 MCS 构造和证明由本稿给出。\href{https://www.math.uwaterloo.ca/\textasciitilde{}snburris/htdocs/ualg.html}{作者原书}

\section{8.1 关系及关系之间的关系}
\subsection{定义 8.1（有类型的关系空间）}
对载体 \(X,Y\)，置 \(\operatorname{Rel}(X,Y)=\mathcal P(X\times Y)\)，高阶类型为 \(X\to Y\to o\)。关系包含也是高阶谓词：


\mcsdisplay{R\sqsubseteq S\iff\forall x:X\,\forall y:Y\;(Rxy\Rightarrow Sxy).}
\noindent \(\forall R:\mathsf{Node}\to\mathsf{Node}\to o\;\operatorname{Trans}(R)\Rightarrow\cdots\) 真正量化关系变量。元层全幂集不等于 Henkin 模型的关系量化域，本章集合构造在 ZFC 元层进行。

硬演绎类型为\(\mathrm{Ded}_T\subseteq\mathsf{PremFam}\times\mathsf{Claim}\times\mathsf{Proof}\)。这里 PremFam 是第 06 章的有限带标签前提族。将联合输入 \(\{p,q\}\) 到 \(r\) 投影为普通二元边只供导航，不能读成“有 \(p\) 或 \(q\) 即足够”。

\subsection{定义 8.2（关系运算）}
对 \(R\subseteq X\times Y,S\subseteq Y\times Z\)，定义


\mcsdisplay{S\circ R=\{(x,z):\exists y\;(xRy\land ySz)\},\quad
R^\smile=\{(y,x):xRy\},\quad 1_X=\{(x,x):x\in X\}.}
\noindent 并、交、补要求相同类型，补注明环境。限制为 \(R|_{A,B}=R\cap(A\times B)\)。自关系上置 \(R^0=1_X,R^{n+1}=R\circ R^n,R^*=\bigcup_{n<\omega}R^n\)。

\subsection{命题 8.3（关系代数）}
复合结合，以恒等关系为单位；逆反转复合次序；复合对任意并分配。每个 \(\operatorname{Rel}(X,X)\) 因而是有单位的关系 quantale。

\textbf{证明。} 两种括号的三重复合均等价于存在 \(y,z\) 使 \(xRy,ySz,zTw\)。单位消去中介，逆交换端点。某一步属于 \(\bigcup_iS_i\) 等价于存在 \(i\) 使之属于 \(S_i\)，交换存在量词得到分配律。\(\blacksquare\)

这不证明类比、软先修等具体关系传递。“类比链可达”必须另立关系，不能仍称直接类比。

\subsection{定义 8.4（两种集合提升）}

\mcsdisplay{R_{\forall\exists}(A,B)\iff\forall a\in A\;\exists b\in B\;aRb,\qquad
R_{\exists\forall}(A,B)\iff\exists b\in B\;\forall a\in A\;aRb.}
\noindent 前者逐个覆盖，后者要求共同承接者。只有 \(a_1Rb_1,a_2Rb_2\) 时，两者不同。逐个取中介见证可证


\mcsdisplay{R_{\forall\exists}(A,B)\land S_{\forall\exists}(B,C)
\Rightarrow(S\circ R)_{\forall\exists}(A,C).}
\noindent 逆命题不成立，因为指定的 \(B\) 未必包含所需中介。空集合保留通常量词语义，空覆盖不是存在性证书。

\section{8.2 证据层}
\subsection{定义 8.5（证据纤维）}
关系端点配见证集 \(\mathsf W_R(x,y)\)，认证关系是此集非空。档案保存实际见证及检查器。复合见证定义为


\mcsdisplay{\mathsf W_{S\circ R}(x,z)=
\coprod_{y\in Y}\mathsf W_R(x,y)\times\mathsf W_S(y,z).}
\noindent 配对见证不自动是另一种证书。推导合成需演算的替换或切规则，解释合成需类型与满足引理，经验相关性复合不产生因果证明。

\subsection{定义 8.6（认证的关系包含）}
一族端点保持的映射\(j_{x,y}:\mathsf W_R(x,y)\to\mathsf W_S(x,y)\)认证 \(R\sqsubseteq S\)。有限数据库暂时未见反例不提供这种变换。

尚未登记见证不表示数学上无见证。数据区分 Known、Unknown 和带反证的 Refuted；空的已知列表若无穷尽性证书只表示列表为空。第 05 章 \(\mathrm{Known}(\emptyset)\) 表示某个认证概念支持确实为空，两者不混用。数学证明、经验结果、建模假设不能通过数值权重互相转成证书。

\section{8.3 生成与闭包}
\subsection{定义 8.7（有限规则生成）}
给定 \(\mathcal Q\subseteq\mathcal P_f(U)\times U\)，令


\mcsdisplay{F(X)=X\cup\{u:\exists(A,u)\in\mathcal Q,\ A\subseteq X\},\qquad
c(X)=\bigcup_{n<\omega}F^n(X).}
\noindent 零前提规则合法。规则形成不得依赖尚未生成的闭包中的否定事实。

\subsection{命题 8.8（代数闭包）}
\(c\) 扩张、单调、幂等，且\(c(X)=\bigcup_{A\subseteq_f X}c(A)\)。

\textbf{证明。} 零次迭代包含 \(X\)。前提包含保持规则适用，故所有迭代单调。一个有限前提集若在 \(c(X)\) 中，则其各元素都出现在有限阶段，取阶段最大值，下一步产生结论。因此 \(c(X)\) 已封闭，再迭代不增加元素。任一有限推导树只用有限个 \(X\)-叶子，证明最后等式的正向包含；反向由单调性。\(\blacksquare\)

该定理分别适用于带上下文语法、固定演算、固定依赖证书和公开行动的结构可达集合。它不把个人掌握集设为演绎闭包，也不使带偏好、预算或遗忘的规划成为闭包。

\subsection{定义 8.9（依赖补全与行动可达）}
固定已经选择的支持/路线证书，以 \(pDn\) 表示它使用 \(p\) 支持 \(n\)，置


\mcsdisplay{c_D^-(X)=\{p:\exists n\in X,\ pD^*n\}.}
\noindent 它是闭包；环不破坏闭包公理。行动可达则以\(\mathcal Q_\Lambda=\{(I_a,o):a\in\Lambda_M,o\in O_a\}\)向前生成。二者分别回答“这份选择依赖什么”和“这些背景能产出什么”。把所有 OR 前提并入只是导航超集，不是共同硬先修。

\textbf{反例 8.10。} 若 \(aDb,bDc\) 而 \(a\not Dc\)，一次补全 \(\{c\}\) 只得到 \(\{b,c\}\)，第二次才出现 \(a\)；一次操作不幂等，必须取可达闭包。

\section{8.4 表达代数与节点档案}
\subsection{定义 8.11（上下文索引的表达代数）}
排序由上下文 \(\Xi\) 和类型 \(\tau\) 索引，载体为 \(\mathrm{Tm}_\Sigma(\Xi;\tau)\)。其运算包括


\mcsdisplay{\mathrm{app}:\mathrm{Tm}(\Xi;\sigma\to\tau)\times
\mathrm{Tm}(\Xi;\sigma)\to\mathrm{Tm}(\Xi;\tau),}

\mcsdisplay{\mathrm{lam}:\mathrm{Tm}(\Xi,x:\sigma;\tau)\to
\mathrm{Tm}(\Xi;\sigma\to\tau),}
\noindent 以及变量、常元、重命名、弱化和无捕获替换。绑定可用 de Bruijn 代码，或先商去 \(\alpha\) 重命名。抽象不应写成忽略上下文的普通单排序操作。

\subsection{命题 8.12（表达的递归泛性质）}
给定能解释变量、常元、应用和抽象，并满足上下文与 \(\alpha\) 改名相容条件的目标表达代数，有唯一保持这些构造子的求值映射。若还要求它保持定义 8.11 中的弱化和无捕获替换，目标的这些运算须满足对应的递归方程：替换在变量上由赋值决定、保持常元与应用，并在抽象下按新变量提升；弱化与重命名的恒等、复合和绑定相容律也须成立。

\textbf{证明。} 叶子由变量和常元解释确定，内节点递归调用子树映像及目标运算。有限树保证终止，\(\alpha\) 相容条件保证代表无关。任一构造子相容映射在叶子一致，结构归纳后处处一致。对源项再次归纳，使用列明的目标递归方程，才得到弱化、替换的保持性。仅给一个任意命名为“替换”的目标运算并不能推出后一结论。\(\blacksquare\)

若再商去 \(\beta\eta\)，目标必须满足对应等式。这个自由对象是表达结构；带来源、角色、动机和证据的整个节点档案并不因此成为自由代数。

\subsection{定义 8.13（同余与关系饱和）}
运算同余是排序内等价关系 \(\theta\)，使每个运算保持逐项等价；它保证商运算良定义。若还要保留关系，需额外满足


\mcsdisplay{x_i\theta y_i\;(1\le i\le k)\Rightarrow
[R(x_1,\ldots,x_k)\leftrightarrow R(y_1,\ldots,y_k)].}
\noindent 固定理论下，可证等价能够成为公式连接词代数的同余，不能笼统断言它从来不是同余。但它通常不保持节点来源、证明成本、动机或路线角色。Lindenbaum 代数因此只作语义投影，保留原节点到等价类的映射。

\textbf{反例 8.14（定理坍缩）。} \(T\vdash\varphi,T\vdash\psi\) 蕴含 \(T\vdash\varphi\leftrightarrow\psi\)。若按此取节点商，所有定理都落在真元，无法保留学习内容身份。

\section{8.5 模型论中的局部与整体}
本地 Prestel–Delzell §§2.2–2.4 已核正文：§2.2 结构态射印刷 65–67 / PDF 75–77；§2.3 子结构印刷 71–74 / PDF 81–84；§2.4 初等扩张与链印刷 75–79 / PDF 85–89，其中生成子结构 p.78/PDF88，初等链定理 p.79/PDF89。这里只直接使用一阶结论；高阶情形必须经指定 Henkin 编码或另行论证。

\subsection{定义 8.15（子结构与初等性）}
固定一阶或多排序编码语言，\(A\subseteq B\) 要求各排序非空，含常元、对函数封闭，关系为限制。若对所有公式及 \(A\) 中参数有 \(A\models\varphi\iff B\models\varphi\)，才写 \(A\preccurlyeq B\)。

\subsection{命题 8.16（证明切片保持）}
若证明 \(p\) 的开放假设、公理实例、推理步骤全部保留且符号解释未变，原检查器仍接受 \(p\)。

\textbf{证明。} 对有限证明树归纳。叶子仍为原假设或公理；内节点的规则与全部子证明保持，故检查结果保持。\(\blacksquare\)

这不保证所有存在见证保留。只含一元谓词 \(P\) 的结构 \(B=\{0,1\},P^B=\{1\}\) 有子结构 \(A=\{0\}\)，但 \(B\models\exists xP(x)\)、\(A\not\models\exists xP(x)\)。筛选甚至未必对函数封闭。

\subsection{命题 8.17（初等链粘合）}
在固定一阶/多排序编码语言中，\(A_0\preccurlyeq A_1\preccurlyeq\cdots\)的并 \(A\) 是每个 \(A_n\) 的初等扩张。

\textbf{证明。} 函数有限元性与兼容解释给出并结构。对公式归纳；原子及布尔情形直接。并结构的存在见证及有限参数同属某个 \(A_m\)，对子公式使用归纳假设，再由 \(A_n\preccurlyeq A_m\) 转回存在断言。反向使用 \(A_n\) 中的见证。\(\blacksquare\)

普通嵌入链不够：有限线序 \(\{0,\ldots,n\}\) 都有最大元，并 \(\mathbb N\) 却没有。标准高阶的全函数域也未必被链并保留。

\section{8.6 语言、理论、模型的 institution 接口}
Goguen–Burstall 用随签名变化相容的满足关系组织逻辑。下面固定一个具体版本，不预设所有翻译保守或所有理论可粘合。\href{https://publish.lfcs.inf.ed.ac.uk/reports/90/ECS-LFCS-90-106/}{原始报告}

\subsection{定义 8.18（具体 institution）}
固定基本类型与逻辑常元，给出四项数据：

\begin{enumerate}
\item \(\mathbf{Sig}\)：对象是有类型非逻辑常元签名；态射 \(\sigma:\Sigma\to\Sigma'\) 是保持类型的符号映射，对逻辑常元恒等。
\item \(\mathrm{Sen}(\Sigma)\)：封闭 \(o\)-项；\(\mathrm{Sen}(\sigma)\) 递归翻译常元、应用和抽象，构成协变函子到 \(\mathbf{Set}\)。
\item \(\mathrm{Mod}_H(\Sigma)\)：对象为指定 Henkin 模型，态射为保持应用、逻辑解释和常元的类型族同构，构成群胚。约化保留类型域，把 \(c\) 解释为原模型中的 \(\sigma(c)\)，同构也约化，构成反变函子 \(\mathbf{Sig}^{op}\to\mathbf{CAT}\)。
\item \(\models_{H,\Sigma}\)：第 03 章满足关系。
\end{enumerate}
约化表达的翻译在原模型已可解释，故抽象闭合保持。类型映射、域限制或商解释必须另扩展接口。标准模型另形成 \(\mathrm{Mod}_{std}\)。取群胚足以定义 institution，并没有声称普通模型同态保持所有公式。

\subsection{命题 8.19（满足条件）}

\mcsdisplay{A'\models_{\Sigma'}\mathrm{Sen}(\sigma)(\varphi)
\iff\mathrm{Mod}(\sigma)(A')\models_\Sigma\varphi.}
\noindent \textbf{证明。} 对项证明更强的赋值相容等式：变量不变，常元是约化定义，应用逐项归纳，抽象逐点归纳并用外延性。封闭 \(o\)-项取真值即得结果。同样归纳证明恒等及复合翻译律。\(\blacksquare\)

\subsection{定义 8.20（理论态射）}
理论对象为 \((\Sigma,T)\)。态射\(\sigma:(\Sigma,T)\to(\Sigma',T')\)要求 \(T'\models_H\sigma(T)\)。满足条件推出 \(T'\)-模型约化为 \(T\)-模型，并推出


\mcsdisplay{T\models_H\varphi\Rightarrow T'\models_H\sigma(\varphi).}
\noindent 可计算的证明输送还要为所用源公理提供目标证明，且翻译每条规则。一般域/商解释另用“表达翻译、模型构造、满足引理”接口，核查非空性、函数全定义、商良定义；任意解释不是硬泛化的自动证据。

\subsection{命题 8.21（模型扩张保证的保守性）}
设 \(\Sigma\subseteq\Sigma'\)、\(T\subseteq T'\)，每个 \(T\) 的 Henkin 模型可扩张为 \(T'\)-模型且约化回自身，则旧句子满足\(T'\models_H\varphi\Rightarrow T\models_H\varphi\)。若演算对指定模型类完备，进一步得到可证性保守。

\textbf{证明。} 任取 \(A\models_HT\)，扩张为 \(A'\models_HT'\)，把 \(A'\models\varphi\) 按满足条件转回 \(A\models\varphi\)。语法结论另用完备性。\(\blacksquare\)

保守性声明须给出相应层次的证明；模型扩张是本节采用的充分判据，不是理论保守性的一般必要条件，也可另给直接的语义或语法证明。只说“增加新排序”不够，新公理仍可能限制旧模型；但排除某些旧模型本身也未必产生新的旧语言后果。

\section{8.7 MCS 态射及其边界}
\subsection{定义 8.22（认证空间态射）}
指定语法/理论翻译后，\(F:M\to M'\) 包括节点、证据、行动映射，保持角色类型、表达负载和认证见证，并满足


\mcsdisplay{I_{F(a)}=F[I_a],\qquad O_{F(a)}=F[O_a].}
\noindent 行动类型也保持。采用严格相等固定一个可用的态射类，其他放宽版本另命名。恒等和逐项复合满足这些条件，因而形成范畴；对象类大小按第 01 章处理。

\subsection{命题 8.23（路线输送）}
对每条 \(r\in\mathrm{Rte}_M(B,G)\)，存在一条\(r'\in\mathrm{Rte}_{M'}(F[B],F[G])\)，其事件沿用原事件标识，行动由 \(F\) 翻译。

\textbf{证明。} 对事件 p 的每个目标输入 \(m'\in I_{F(\mathrm{act}(p))}=F[I_{\mathrm{act}(p)}]\)，选择一个满足 F(m)=m' 的源输入 m，并翻译原来的 h(p,m)。若为背景来源，\(m'\in F[B]\)；若为事件 q 的来源，\(m'\in O_{F(\mathrm{act}(q))}\)。每条所得事件边都是原依赖图的边，故新图是无环子图。对每个 \(g'\in F[G]\) 同样选择一个原目标 g 并翻译 \(\zeta(g)\)。全部选择有限，因此得到合法路线。\(\blacksquare\)

非单射节点映射可能把同一事件的不同输入、或不同目标合并；其原来源可能不同，不能把两份来源直接写成一个函数值。因此此处先证明路线输送的存在性，不声称无需选择便得到唯一的路线映射，也不声称保留每条原依赖边。若要求路线层的函子，须限制为相应接口单射的态射，或另给与恒等、复合相容的来源选择机制。

共同硬先修未必保持，因为目标空间可以新增绕开 \(F(p)\) 的路线。要保持它，另须每条目标路线可提升为源路线，并有足够的节点身份反映“使用 \(p\)”；成本与认知收益更不在上述保持结论中。

\subsection{命题 8.24（兼容标准模型的粘合）}
若两个标准 HOL 模型拥有相同全类型域和逻辑解释，在共享常元上相同，则合并非共享常元解释，得到唯一并签名结构，其约化分别为两侧模型，并满足两侧理论的并。

\textbf{证明。} 解释无冲突，全函数域容纳混合表达的抽象解释，满足条件适用于两侧；每个常元已由一侧确定，故唯一。\(\blacksquare\)

Henkin 模型还要检查混合签名表达的抽象闭合；任意局部模型未必有兼容公共约化。因此没有“一切知识片段都能粘合”的无条件定理，更没有教学效果粘合定理。

\clearpage\chapter{09 话题、领域、学科与元理论的衔接}
\mcsorigin{原章节：09-话题、领域、学科与元理论.md}
层级不是把若干标签放在同一种节点集合里。聚合、数学理论和元理论有不同的形成规则；需要统一表示时使用带类型引用，而不是消除差别。

\section{9.1 关联参数与话题}
\textbf{定义 9.1（公开关联规格）。} 固定一个本体版本与可比较的节点域 \(U\subseteq N_M\)。规格 \(\theta\) 给出公开选择的支持索引、关系种类与用途。精确信息齐备时定义对称亲和函数 \(a_\theta:U\times U\to[0,1]\)，例如在有限已知支持下采用


\mcsdisplay{a_\theta(u,v)=\alpha J(S_u,S_v)+\beta W(u,v)+\gamma Q(u,v),
\quad\alpha,\beta,\gamma\ge0,\quad\alpha+\beta+\gamma=1.}
\noindent J 为集合 Jaccard 比例，空—空时公开约定为 0；W(u,v) 明确定义为两者任一方向存在指定结构关系证据时取 1，否则在确认不存在该登记证据时取 0，因此 W 对称；Q 表示是否参与同一声明任务。信息不完整时另用评价接口返回 Unknown 或包含所有允许值的区间 \([\ell,h]\subseteq[0,1]\)，不把这些信息标签塞进 \(a_\theta\) 的实数值域，也不按零补值。这里的权重是公开组织约定，不是已测得认知强度；不能使用饱和的 \(C_\mathrm{all}\) 作为默认支持。

给定阈值 \(\tau\) 和整数 \(d\ge 1\)，在已知 \(a_\theta(u,v)\ge\tau\)，或安全区间下界 \(\ell\ge\tau\) 的不同节点之间连边。未获得此证据只表示边尚未认证，不能据此断言 \(a_\theta(u,v)<\tau\)。\(\theta\)、\(\tau,d\) 是话题版本的一部分，允许研究其变化，不默认存在唯一自然划分。

\textbf{定义 9.2（话题）。} 话题 t 是 \((\mathrm{id},\theta,S,E_t,\mathrm{bdry})\)，其中 \(S\subseteq U\) 非空，\(E_t\) 是上述阈值图在 S 上的边，要求它连通且图直径不超过 d；单点直径为 0。bdry 保存 S 外的被引用前提。默认只登记有限 S；无限扩展另给大小与直径语义。

这是一种可检验的“强关联节点合并”约定，不宣称聚合直接证明教学连贯。一个节点可属于多个话题；不强制 \(\theta\) 图的所有极大连通分量就是话题，否则短桥可能把巨大领域误并为一块。

\textbf{例 9.3。} “流形”话题可含拓扑流形、图册、\(C^k\) 过渡映射和光滑结构；“张量与张量场”话题可共享流形与截面节点。公共节点重叠不产生重复数学身份。

\section{9.2 领域、学科与覆盖}
\textbf{定义 9.4（多层聚合）。} 取互异排序 Node、Topic、Domain、Discipline。领域 d 登记一组话题及其公共问题族、方法族和跨话题接口；学科登记一组领域及组织准则。形成规则至少要求成员排序正确、引用非悬空、所声明接口有来源。

定义展平映射


\mcsdisplay{Flat(n)=\{n\},\quad
Flat(t)=S_t,\quad
Flat(d)=\bigcup_{t\in Members(d)}Flat(t),}
\noindent 学科同理。每个聚合保留成员对应，不以 Flat 相同判定聚合身份。两个领域可以有相同节点集合却强调不同的问题组织。

\textbf{命题 9.5（展平与覆盖）。} 若聚合 A 的每个成员也是同层聚合 B 的成员，则 \(\mathrm{Flat}(A)\subseteq\mathrm{Flat}(B)\)。多层展开的并不依赖分组括号。

\textbf{证明。} 前者由成员并的包含；后者由集合并的结合律。它们不证明两种分组有相同教学意义。\(\blacksquare\)

覆盖是指定族 \(\{S_i\}\) 满足 \(U=\bigcup S_i\)；分割还要求两两不交。话题通常仅采用覆盖，不采用分割。

\section{9.3 粗粒度关系的逻辑量词}
\textbf{定义 9.6（关系聚合的两种读法）。} 对节点关系 R 和聚合 A,B，分别定义


\mcsdisplay{R_\exists(A,B)\iff\exists x\in Flat(A)\exists y\in Flat(B)\ R(x,y),}

\mcsdisplay{R_{\mathrm{cov}}(A,B)\iff
\forall y\in Flat(B)\exists x\in Flat(A)\ R(x,y).}
\noindent 第二种记为 \(R_\mathrm{cov}\) 以区别第 08 章覆盖源集合的 \(R_\forall\exists\)；它在这里覆盖目标集合。二者不得使用同一个未注明量词的边。前者只说存在联系，后者说覆盖每个目标；都不表示“学完 A 足以学完 B”。联合前提的聚合还需保留完整输入超边和证书。

若允许空聚合，\(\mathrm{Flat}(B)=\emptyset\) 时 \(R_\mathrm{cov}\) 真空为真，\(R_\exists\) 为假。展示“整域关联”时须另外要求目标非空，不能把空覆盖显示为已有一项关系见证。

\textbf{反例 9.7。} A=\{a\}，\(B=\{b_{1},b_{2}\}\)，仅有 \(R(a,b_{1})\)。则 \(R_\exists\) 真、\(R_\mathrm{cov}\) 假。把存在联系解释为整领域先修会错误地给 \(b_{2}\) 添关系。聚合时只能在保留的证据下声明性质，其他结构以边界记录保留。

\section{9.4 理论不是主题集合}
\textbf{定义 9.8（数学理论的双重表示）。} 数学理论对象为 \((\Sigma,T,\)演算,语义选择)，带模型类与推演闭包；其 Theory 节点是对该对象的有限描述或环境引用。用 \(\mathrm{denotes}:\mathrm{TheoryNode}\to\mathrm{TheoryCode}\) 连接，两者类型不同。

一个主题可以同时介绍多个相互不等价理论；一个理论也可跨多个领域使用。话题聚合不自动构造公理并；理论并则须处理签名相容、一致性和模型粘合，见第 08 章。

理论作为学习对象时，可以有动机、主要定理和路线；个人是否掌握理论依然外置。不得用“已阅读理论节点”代表掌握其演绎闭包。

\section{9.5 元理论层级与自描述}
\textbf{定义 9.9（分层引用）。} 设 \(L_{0}\) 为被研究的数学对象语言，\(L_{1}\) 能描述 \(L_{0}\) 的有限语法和解释，后续 \(L_{i+1}\) 可引用 \(L_i\) 的代码。每个层级声明自己的真理/满足讨论范围。MCS 的某个版本可在更高层被编码为研究对象，但不在同层加入满足全部自身句子的无约束真理谓词。

这里的 \(L_i\) 是本节局部“语言层级”记号，不是可构造宇宙 \(L_\alpha\)；为避免混淆也可写 \(\mathrm{MetaLang}_i\)。本稿不以二者的字母相同建立数学对应。

\textbf{命题 9.10（层级引用不改变下层解释）。} 在只添加下层语法代码、保守定义和外部引用，且每个下层模型可扩张的条件下，旧语言语义后果保持。

\textbf{证明。} 这是第 08 章模型扩张保守性命题的实例：取任意旧模型扩张，旧符号解释不变，满足条件把旧句子真值转回。没有扩张条件时不能断言。\(\blacksquare\)

这使节点、知识簇、领域理论以及关于 MCS 自身的研究在统一的有类型引用系统中衔接，同时保持各层的证据义务。

\clearpage\chapter{10 人性化节点模板及外部呈现}
\mcsorigin{原章节：10-节点模板.md}
下列十三种是\textbf{呈现角色}，不强迫它们成为十三个互斥语法构造子。它们与定义 4.4 的十四个构造类型的对应是：概念\(\leftrightarrow\mathrm{Concept}\)、定义\(\leftrightarrow\mathrm{Definition}\)、公理/定理/性质\(\leftrightarrow\)\textbf{Claim（三种角色）}、证明\(\leftrightarrow\mathrm{Proof}\)、例子\(\leftrightarrow\mathrm{Example}\)、反例\(\leftrightarrow\mathrm{Counterexample}\)、问题\(\leftrightarrow\mathrm{Problem}\)、理论\(\leftrightarrow\mathrm{Theory}\)、构造\(\leftrightarrow\mathrm{Construction}\)、全局/局部方法\(\leftrightarrow\)\textbf{Method（两个范围标签）}；Symbol、Term 等形式候选使用共享模板及类型推导附件，\textbf{Representation 与 Miscon\mcsbrk{}ceptio\mcsbrk{}nPatte\mcsbrk{}rn} 使用各自的负载说明与结构字段，无须虚构完整的人类教学故事。

\section{10.1 共享模板契约}
\textbf{定义 10.1。} 每个专用模板继承以下十二组字段。无适用内容标“明确不适用”，未知标 Unknown，空白不代表已查明。

\begin{enumerate}
\item \textbf{身份与角色}：节点引用、版本、角色、背景、相关节点；改变角色不自动合并身份。
\item \textbf{问题动机}：回答什么问题，不引入它有何具体困难；从学科内部、外部应用或结构统一中给出至少两条有内容的动机路径，不能只换说法重复同一需求。历史入口与现代重构分别注明来源，不为凑路径编造历史。
\item \textbf{直觉与表征}：口头图像、图形、代数表示及对应；写出胚子、解决局部困难的中间形式、规范定义的生长链，每个正式定义先配一句直觉。逐项列出直觉不能推出的结论，保留曾经误用它的原因与检验办法。
\item \textbf{正式负载}：理论、签名、类型、变量语境、局部假设、表达与形成见证；非形式启发体单列。
\item \textbf{概念支持}：表达、定义、证明和路线的支持分别引用第 05 章索引；不以 \(C_{\rm all}\) 代替依赖，不把未知填成空集。
\item \textbf{公开入口与输出}：引用 \(\Lambda_M\) 中的 \(I_a,O_a,w_a\) 及行动类型；输出可用不表示学习成功。
\item \textbf{例子与边界}：具体正例、反例与极端例；对定理逐一尝试删除假设，写明结论如何失败。若该条件其实冗余，应修正命题；若暂未确定是否可删，标为待查，不能把所有充分条件假装成必要条件，也不能编造反例。
\item \textbf{一般误区模式}：链接 \(Pattern_M\)、触发条件和诊断题；不填写某人是否具有该误区。
\item \textbf{自检任务}：任务规格、节点—能力锚定、工具、成功条件及评分依据，保留至少一条值得以后回看的问题。只有检查器确实可判时标“可判”，开放证明题不能假装具有通用判定器。
\item \textbf{应用与迁移}：目标任务、条件匹配、可认证结论与失效场景；给出低层重述、高层特例或相似对象联动，类比须标明保持及失效之处。最后用一两句命名核心思想，列出所用方法及其在笔记中的位置；教学效果留作外部评价槽。
\item \textbf{可选结构路线}：背景 \(B\)、目标 \(G\)、契约路线、共同/替代输入与桥接证书；不预设唯一先修链。
\item \textbf{来源与证据状态}：分别登记完整证明、引用结果、经验资料、建模假设、解释与未知项，保留具体出处和核验状态。
\end{enumerate}
以下每一种模板均须填完整的共享部分及专用部分，不能只填专用关键词。未知标签使候选档案可保存，不使其自动达到该角色的认证条件：身份、版本和角色须能确定；宣称为合法定义、已证定理或已认证行动时，其负载、背景及相应合法性证据不能以 Unknown 代替。人工论证、引用认证与机器检查仍按各自状态标注。

\section{10.2 概念模板}
\textbf{定义 10.2（Concept）。} 专用字段包括：

\begin{itemize}
\item \textbf{对象与边界}：对象类型、参数域、分类谓词、量词范围、命名约定。
\item \textbf{等价与判别}：必要条件、充分条件分别登记；已认证等价表述及各自证明；不同模型的解释。
\item \textbf{认知入口}：能说明核心区别的典型例、边界例与非例；从例子到定义的结构对应。
\item \textbf{结构联系}：邻近概念间的蕴含或指定映射，不能据名称推定。
\item \textbf{自检维度}：辨认、构造例子、排除反例、使用定义分别测查。
\end{itemize}
等价表不合并节点身份；“微分流形”应明确当前 \(C^k\) 约定，不让名称替代参数。

\section{10.3 定义模板}
\textbf{定义 10.3（Definition）。} 专用字段包括：

\begin{itemize}
\item \textbf{新符号}：名称、类型、参数及与旧签名的新鲜性条件。
\item \textbf{定义主体}：旧语言公式、变量绑定、无捕获条件、展开与折叠规则。
\item \textbf{合法性}：谓词定义的扩张构造；函数规格的存在唯一性、函数域实现与新抽象闭合条件（第 03 章）；循环引用检查。
\item \textbf{保守性}：结论针对的语言、模型类、证明系统及其适用范围。
\item \textbf{解释与练习}：为何选择此定义、其他约定的互译、一次完全展开、一个非法定义反例。
\end{itemize}
动机不能代替定义合法性证书；定义引入与定理推导使用不同契约类型。

\section{10.4 公理模板}
\textbf{定义 10.4（Axiom）。} 专用字段包括：

\begin{itemize}
\item \textbf{地位}：所属理论、公理还是模式、精确公式及实例化规则。
\item \textbf{结构意图}：希望约束什么；删去它允许哪些额外结构。
\item \textbf{模型证据}：满足模型；违反此公理但满足其余指定条件的模型，缺证时为未知。
\item \textbf{元性质}：相容性、独立性、冗余性与各自证据，不能从列表形式推断独立。
\item \textbf{使用和分支}：实际调用的证明索引；替换该公理产生的理论及解释方向。
\end{itemize}
理论内采用某公理不等于该句在所有数学理论中都为真。

\section{10.5 定理模板}
\textbf{定义 10.5（Theorem）。} 专用字段包括：

\begin{itemize}
\item \textbf{精确命题}：量词顺序、共同背景、局部假设和联合前提。
\item \textbf{证明资源}：完整证明或明确引用；每条证法的实际支持、证据类型和认证状态。
\item \textbf{发现线索}：为什么猜测此结论，哪一特殊情形提示关键步骤。
\item \textbf{变化范围}：加强、削弱、逆命题、等号情形及删除假设的反例。
\item \textbf{应用接口}：输入条件如何匹配，输出能否继续作为其他行动的资源。
\end{itemize}
“概念节点被展示”不是一个可随意塞进演绎前提的命题；其公式负载或定义规则须显式引用。

\section{10.6 证明模板}
\textbf{定义 10.6（Proof）。} 专用字段包括：

\begin{itemize}
\item \textbf{目标与环境}：Claim 引用、理论、开放假设和允许推理规则。
\item \textbf{策略与步骤}：策略概览、关键中间断言、逐步理由、规则/引理实例。
\item \textbf{完整证据}：证明树或完整数学论证；机器检查状态与人工论证状态分开。
\item \textbf{隐含资源核查}：选择、有限性、维数、正则性等条件及使用位置。
\item \textbf{替代与失败}：其他证法的支持差异、失败尝试的确切错误、发现过程的解释材料。
\end{itemize}
本稿中的完整数学证明不自动意味着已编码成第 02 章演算并运行内核检查。

\section{10.7 性质模板}
\textbf{定义 10.7（Property）。} 专用字段包括：

\begin{itemize}
\item \textbf{承担者}：对象类型、性质谓词与参数。
\item \textbf{作用尺度}：局部/整体，内在/依赖表征，以及该区分的定义。
\item \textbf{判据与证据}：必要、充分、等价条件及证明链接。
\item \textbf{保持条件}：在明确列出的子对象、商、乘积、同构或其他变换下是否保持。
\item \textbf{判别任务}：正例、边界例、反例和能够观测的性质差别。
\end{itemize}
名称为“不变量”不意味着对所有变换都保持。

\section{10.8 例子模板}
\textbf{定义 10.8（Example）。} 专用字段包括：

\begin{itemize}
\item \textbf{对象规格}：集合、运算、结构或构造参数，足以重现。
\item \textbf{满足验证}：目标概念每个条件的检查与证明。
\item \textbf{实际计算}：至少一项具体运算、坐标表达或性质验证。
\item \textbf{典型与偶然}：突出的核心特点和不得推广的偶然特点。
\item \textbf{邻近变形}：参数变化、变成反例的边界及替代例子。
\end{itemize}
教学用途可以是有解释依据的候选；“对此人容易”仍属外部评价。

\section{10.9 反例模板}
\textbf{定义 10.9（Counterexample）。} 专用字段包括：

\begin{itemize}
\item \textbf{否定目标}：被反驳的精确句子及量词范围。
\item \textbf{构造与见证}：对象完整规格、满足的前提、失败结论的具体见证。
\item \textbf{失效定位}：哪一推理步骤不能成立，能否去除无关结构。
\item \textbf{未否定范围}：仍可能成立的弱命题和其他对象类型。
\item \textbf{修正与诊断}：修正命题候选、一般误区模式和辨析任务。
\end{itemize}
一个 \(C^1\) 图册不是 \(C^2\) 图册，不等于其底空间不容许任何光滑结构。

\section{10.10 问题模板}
\textbf{定义 10.10（Problem）。} 专用字段包括：

\begin{itemize}
\item \textbf{输入输出}：任务类型、约束、允许工具及解的正确性规格。
\item \textbf{客观状态}：开放/已解/条件解决等状态，来源和核验日期。
\item \textbf{已有进展}：已证结果、部分范围、未知环节。
\item \textbf{探索记录}：尝试、失败原因、可用方法及子问题。
\item \textbf{组合条件}：各子问题的解怎样汇合为原问题的解。
\end{itemize}
客观问题状态是公共档案版本；某人是否解出属于外部实例，个体状态更新不能直接修改前者。

\section{10.11 理论模板}
\textbf{定义 10.11（Theory）。} 专用字段包括：

\begin{itemize}
\item \textbf{语言与逻辑}：签名、类型、演算及语义选择。
\item \textbf{公理呈现}：公理/模式、定义扩张和依赖。
\item \textbf{模型与证据}：模型类、存在性说明、反模型及未知项。
\item \textbf{结构变换}：解释、保守扩张、约化的方向和保持条件。
\item \textbf{组织与学习}：主要定理、目标问题族、局部子理论、引用边界和公开行动路线。
\end{itemize}
有限目录不声称枚举理论的一切后果。

\section{10.12 构造模板}
\textbf{定义 10.12（Construction）。} 专用字段包括：

\begin{itemize}
\item \textbf{输入与输出规格}：对象类型、入口条件及期望性质。
\item \textbf{执行步骤}：辅助对象、选择与每步合法性。
\item \textbf{正确性}：输出存在并满足规格的证明。
\item \textbf{唯一性层次}：相等、同构、自然同构或只保证存在；选择依赖与选择无关性。
\item \textbf{映射相容与实例}：何时形成函子，一个完全计算实例，失败输入和边界。
\end{itemize}
没有自然性验证不称自然构造；“典范”须说相对于什么结构。

\section{10.13 全局方法模板}
\textbf{定义 10.13（GlobalMethod）。} 专用字段包括：

\begin{itemize}
\item \textbf{启发体}：名称、自然语言操作建议及适用任务族。
\item \textbf{严格接口}：输入/候选输出类型、触发线索、可检查前提。
\item \textbf{使用过程}：提示序列、停止信号、失效情境及误用风险。
\item \textbf{跨域实例}：至少两个不同领域的实例，与局部方法的关联。
\item \textbf{支持与评价}：已知概念支持、未知部分、方法呈现行动 \(I/O/w\)、外部效果槽及经验来源。
\end{itemize}
适用于直觉建立与失效、内省、特例切入、反例构造、逆向分析、找不变量、以应用替代练习、可视化等。输出可为猜想或问题；其正确性另证。

\section{10.14 局部方法模板}
\textbf{定义 10.14（LocalMethod）。} 专用字段包括：

\begin{itemize}
\item \textbf{限定任务}：可识别输入形状、数学背景和参数范围。
\item \textbf{步骤与分支}：每步适用条件、合法性依据、候选输出。
\item \textbf{核查与失败}：验证方式、停止信号、成功例和失效例。
\item \textbf{方法联系}：所属全局策略、可替代技巧和选择提示。
\item \textbf{实现边界}：自动化范围、未知支持、外部难度/效果槽。
\end{itemize}
已有完整正确算法时可另建 Construction/Proof 节点；“何时想到使用它”的启发仍是方法。求极限约分保留去心条件，指标升降要求给定非退化配对，积分技巧不保证初等原函数存在。

\section{10.15 外部渲染与模板验收}
\textbf{定义 10.15。} \(\operatorname{Render}_E(s,t,g,ctx)\) 从公共模板实例 \(t\) 选择解释深度、表征、任务和提示，输出带来源的呈现。每个呈现的数学断言须溯源到原负载或具有相应证据的变换；省略证明步骤须显式标为节略，不得省略改变命题的假设或量词。启发图像、候选解释和外部预测分别标注，不能因显示在同一页面而取得数学认证。可显示外部 \(\operatorname{Inst}_E(s,p)\)，但不把个体诊断写回公共误区模式。掌握估计、成本和个体顺序不是模板固定值。

\textbf{验收 10.16。} 十三种专用模板均须填写共享及专用字段，或明确标未知/不适用，并分别检查字段完整性与所声明角色的认证条件。替换两个外部状态，原模板的正式负载、支持、证明和来源不变；输出可选用不同的有证据表征、任务和提示，但须保留到该原模板及变换证据的溯源。不能要求不同呈现逐字相同，也不能以原模板未改动代替对输出断言的核查。字段完整性既不认证数学内容，也不证明教学效果。

\clearpage\chapter{11 局部化、边界与 36 个正式算子}
\mcsorigin{原章节：11-局部化36方案.md}
局部化是对本体的聚焦、投影、聚合或标注，不等于范畴论中“把一族态射变成可逆”的 localization。各算子在 ZFC 中给出语义定义；可计算性另列条件。认知适配算子只接收外部适配器产生的显式参数，不读取隐藏的个人档案。

\section{11.1 统一输出与筛选原语}
\textbf{定义 11.1（带边界视图）。}


\mcsdisplay{Loc_\lambda(M)=(V,\partial V,\tau,Pres).}
\noindent V 的对象有显示类型，可为节点、聚合或事件引用；每条显示关系保留其关系类型、量词读法和证据。\(\tau\) 是溯源对应，普通节点 v 满足 \(\tau(v)=\{v\}\)，聚合显示项映到一个声明的有限节点集合。\(\partial V\) 保存外部引用、未展开的前提与未决判断；Pres 是逐项已证或条件性保持声明，不是任意“保真”标签。

\textbf{定义 11.2（通用切片 Cut）。} 对公开节点集合 \(S\subseteq N_M\)、关系种类 J 和可选的证明/契约子族 W，定义 \(\mathrm{Cut}_M(S,J,W)\)：

\begin{enumerate}
\item 可见对象为 S，保留 J 中端点均在 S 且被选择的关系记录；
\item 对每个可见负载、所选证书或行动，refs(x) 中落在 S 外的引用作为边界；不把外部依赖删除为零；
\item 保留证书原编号及环境版本。若声称证明可在视图内独立检查，还须其所有引用均已提供；否则只声称可追溯到环境中的证书；
\item 未决选择谓词输出“未决候选边界”，不把未知当假；
\item 聚合时另给有限块族及块到原节点的溯源，边标注第 09 章的存在或覆盖读法。
\end{enumerate}
以下记 \(C(S)=\mathrm{Cut}_M(S,J,W)\)，未特别说明时 J,W 取指定参数范围内全部记录。记录数无限时此为数学定义；只对有限登记片段和可判选择条件给算法。

\textbf{定义 11.3（共同记号）。} U 是当前研究的候选域，S(n,i) 是已知概念支持；D 是已选择证书所给输入依赖；\(R_J\) 是指定关系种类擦去见证后的导航关系；Rte、Req、B 与行动 \(I_a,O_a\) 按第 06 章定义。D 的导航可达不冒充共同先修。所有阈值、方向、目标及选定路线都进入 \(\lambda\)。

\section{11.2 逻辑与证据族}
\begingroup\small\setlength{\tabcolsep}{4pt}
\begin{longtable}{>{\raggedright\arraybackslash}p{\dimexpr 0.2208\linewidth-2\tabcolsep\relax}>{\raggedright\arraybackslash}p{\dimexpr 0.3896\linewidth-2\tabcolsep\relax}>{\raggedright\arraybackslash}p{\dimexpr 0.3896\linewidth-2\tabcolsep\relax}}
\toprule
\textbf{编号、算子} & \textbf{输入与数学定义} & \textbf{保持、边界、损失与计算条件} \\
\midrule
\endfirsthead
\toprule
\textbf{编号、算子} & \textbf{输入与数学定义} & \textbf{保持、边界、损失与计算条件} \\
\midrule
\endhead
\bottomrule
\endfoot
\bottomrule
\endlastfoot
\textbf{LC01 签名约化} & 输入类型合法的子签名 \(\Sigma'\)；\(S=\{n\in U:\mathrm{Payload}(n)\) 的全部非逻辑符号及所需类型在 \(\Sigma'\}\)，输出 C(S) 并注明背景约化 & 在固定模型中仅忘掉其他符号的解释、保留相关类型域及赋值时，\(\Sigma'\) 表达的满足关系不变；证据和定义展开若越界仍留边界。不是缩小量词域或自动初等子模型。有限代码可判 \\
\textbf{LC02 理论切片} & 输入公理子集 \(A'\subseteq\mathrm{Ax}(T)\)；选择已有证明 p 且 \(\mathrm{UsedAx}(p)\subseteq A'\) 的节点及 p 所需记录，保留开放假设 \(\Gamma_p\) & 保留原判断 \(A';\Gamma_p\vdash\varphi\) 的证书；只有 \(\Gamma_p=\emptyset\) 才称 A' 的定理。未收录不等于在 A' 中不可证；穷尽全部后果还需有效公理呈现与演算，一般仅可枚举 \\
\textbf{LC03 公理预算} & 输入有限候选公理表及非负结构成本 w；给定已选 A' 且 \(\Sigma_{a\in A'}w(a)\le b\)，再调用 LC02 & 预算选择是外部给定或有限优化，不依靠典范序偷定偏好；删去公理及其后果保留说明。预算值可判时有限检查 \\
\textbf{LC04 证明回溯} & 输入有限证明 DAG p 与结论节点；取结论在 p 中的全部祖先及其表达/公理引用 \(S_p\)，输出 \(C(S_p)\) & 保留该证明；不保留全部可能证明。有限 DAG 可达终止，环境引理可作为带版本边界 \\
\textbf{LC05 依赖补全} & 输入种子 \(S_{0}\) 与已选支持 D；\(S=\{x:\exists y\in S_{0},\mathrm{xD}*\) y\}，输出 C(S) & 是固定 D 的闭包，允许有环；不决定路线可学性。不能把全部 OR 分支当作必修。有限图可算，无界图不保证终止 \\
\textbf{LC06 独立性边界} & 输入候选句 \(\varphi\) 及理论 A；标记有 \(A+\varphi\) 和 \(A+\neg\varphi\) 模型证据的节点，并附两模型证据 & 相对指定语义及模型可信性说明独立；无证据者为 Unknown。一般独立性不可判，不能仅搜索失败就标独立 \\
\end{longtable}
\endgroup

LC02 对一个现成证书作筛选，在它实际调用的有限公理均有可判的 A' 成员检查或有效成员见证时是可判任务；任意给定子集 A' 不自动提供这种算法。“一个节点存在任意符合预算的证明”则是存在量词问题。两者不得写成同一个算法保证。

\section{11.3 概念与结构族}
\begingroup\small\setlength{\tabcolsep}{4pt}
\begin{longtable}{>{\raggedright\arraybackslash}p{\dimexpr 0.2208\linewidth-2\tabcolsep\relax}>{\raggedright\arraybackslash}p{\dimexpr 0.3896\linewidth-2\tabcolsep\relax}>{\raggedright\arraybackslash}p{\dimexpr 0.3896\linewidth-2\tabcolsep\relax}}
\toprule
\textbf{编号、算子} & \textbf{输入与数学定义} & \textbf{保持、边界、损失与计算条件} \\
\midrule
\endfirsthead
\toprule
\textbf{编号、算子} & \textbf{输入与数学定义} & \textbf{保持、边界、损失与计算条件} \\
\midrule
\endhead
\bottomrule
\endfoot
\bottomrule
\endlastfoot
\textbf{LC07 种子邻域} & 输入 \(S_{0}\)、导航关系 J、方向及半径 \(k\in\mathbb{N}\)；\(S=\bigcup_{j\le k}R_J^j[S_{0}]\) & 保留半径内指定记录，外部依赖进边界；\(S_{0}\) 有限且每个相关节点能在有限时间返回完整有限邻接表时计算终止。仅“可枚举且局部有限”不足以识别枚举何时结束 \\
\textbf{LC08 支持交汇} & 输入节点 u,v 及已选已知支持 A=S(u,i),B=S(v,j)；输出 \(C(\{u,v\}\cup(A\cap B))\) & 只反映这两份支持共同概念，不声称全路线共同必需；未知支持使用区间/未决边界 \\
\textbf{LC09 支持差异} & 输入两个已证极小支持 \(A\in\mathrm{MinSupp}_u(n)\)、\(B\in\mathrm{MinSupp}_v(m)\)；输出 \(C(\{n,m\}\cup(A\triangle B))\) 并附原支持引用 & 极小性必须有证据；不假定所有支持均可求极小。一般版可输入普通支持，但须另标“非极小” \\
\textbf{LC10 泛化锥} & 输入 c；\(S=\{c\}\cup\{g:g\succeq_H\) c\}，保留类型匹配的蕴含/映射见证 & 同载体与跨载体模式分开；无限锥未必可穷尽。有限认证关系表上可达可算 \\
\textbf{LC11 特化锥} & 输入 c；\(S=\{c\}\cup\{s:c\succeq_H\) s\} & 与 LC10 取逆关系，不取逆单值函数；丢失锥外对应，结构边界仍保留 \\
\textbf{LC12 对偶对应} & 输入已声明部分对偶操作 F 及有效域 A；\(S=A\cup F[A]\)，保留成对显示及结构见证 & 仅保留 F 已证的性质；F 未定义部分进入边界。反向可恢复必须有逆/等价证据，认知收益不在此保证内 \\
\end{longtable}
\endgroup

\section{11.4 问题与方法族}
\begingroup\small\setlength{\tabcolsep}{4pt}
\begin{longtable}{>{\raggedright\arraybackslash}p{\dimexpr 0.2318\linewidth-2\tabcolsep\relax}>{\raggedright\arraybackslash}p{\dimexpr 0.3841\linewidth-2\tabcolsep\relax}>{\raggedright\arraybackslash}p{\dimexpr 0.3841\linewidth-2\tabcolsep\relax}}
\toprule
\textbf{编号、算子} & \textbf{输入与数学定义} & \textbf{保持、边界、损失与计算条件} \\
\midrule
\endfirsthead
\toprule
\textbf{编号、算子} & \textbf{输入与数学定义} & \textbf{保持、边界、损失与计算条件} \\
\midrule
\endhead
\bottomrule
\endfoot
\bottomrule
\endlastfoot
\textbf{LC13 目标回溯} & 输入目标 G、有限行动候选 \(\Lambda'\)；从 G 向后加入所有满足 \(O_a\cap\)当前目标\(\ne\emptyset\) 的行动及 \(I_a\)，输出带 OR 标签的展开 & 输出的是候选展开，不能把分支并当一条路线。循环保留标记；有限 \(\Lambda'\) 可不动点终止 \\
\textbf{LC14 反例诊断} & 输入错误断言 \(\varphi\) 与反例证书 w；S 取 \(\varphi\)、反例、w 及被否定的假设/结论子式锚点 & 保持明确否定范围；不推断所有类似命题都假。对有限证书切片可计算 \\
\textbf{LC15 策略适用域} & 输入方法 m 与任务域 Q；选择 \(q\in Q\) 且有 \(\mathrm{Match}(\mathrm{In}_m,q,w)\) 和适用条件证据的任务，再取其节点锚点 & 匹配证据不保证学习成功；非形式条件未决则标边界。仅在匹配检查器可判且 Q 有限时有限完成 \\
\textbf{LC16 不变量追踪} & 输入变换族 F、函数 I 与证书 I(Fx)=I(x)；选在指定种子/半径内可达的 x、Fx 和对应 I 节点 & 保留所证变换下不变量，不能扩大到任意变换；计算条件与半径及有效域一致 \\
\textbf{LC17 构造入口} & 输入已有来源指派的有限无环构造 c；对指定输出回溯活跃步骤，选择这些步骤引用的外部输入端口 & 保留“本构造”的入口，不声称所有构造最小入口；零输入构造事件不冒充外部资源，AND/OR 分支按已选指派保存。保留原构造版本供展开 \\
\textbf{LC18 应用迁移} & 输入已认证应用描述 r、目标领域 H、部分输入解释 f；选择与 f 匹配的任务及输出候选 & 未完成条件匹配/正确性证明的迁移保持候选状态；外部有效性评价另附。有限接口匹配可算，任意可用性不保证可判 \\
\end{longtable}
\endgroup

\section{11.5 认知适配族：外部参数产生，核心只读}
\textbf{定义 11.4（适配器）。} Adapt(E,s,g,b) 产生带来源的参数包 \(\theta\)，其中只含显式节点/能力引用、已知或未知标记、可使用表征、诊断模式引用及公开算法版本。\(\theta\) 不写入 M。以下六个算子都是 \(\mathrm{Loc}_\lambda(M)\)，\(\lambda\) 包含该参数包；同一 M 可同时输出不同学习者视图。

\begingroup\small\setlength{\tabcolsep}{4pt}
\begin{longtable}{>{\raggedright\arraybackslash}p{\dimexpr 0.2424\linewidth-2\tabcolsep\relax}>{\raggedright\arraybackslash}p{\dimexpr 0.3788\linewidth-2\tabcolsep\relax}>{\raggedright\arraybackslash}p{\dimexpr 0.3788\linewidth-2\tabcolsep\relax}}
\toprule
\textbf{编号、算子} & \textbf{输入与数学定义} & \textbf{保持、边界、损失与计算条件} \\
\midrule
\endfirsthead
\toprule
\textbf{编号、算子} & \textbf{输入与数学定义} & \textbf{保持、边界、损失与计算条件} \\
\midrule
\endhead
\bottomrule
\endfoot
\bottomrule
\endlastfoot
\textbf{LC19 已知内容折叠} & \(\theta\) 给已认证为适合折叠的有限节点块 \(A_i\)；把每个 \(A_i\) 显示为一块，\(\tau(v_i)=A_i\)，其余节点照常显示 & 保留公共节点身份、跨块边与来源；“已知”标签来源是 E，未知不折叠。需块族处理重叠：默认允许重复显示并共享来源 \\
\textbf{LC20 认知前沿} & \(\theta\) 给当前输入域的互斥分类 K（可用）、F（明确不可用）、\(U_?\)（未知），未分类者归 \(U_?\)；确认前沿为 \(I_a\subseteq K\)，可能前沿为 \(I_a\subseteq K\cup U_?\) 且 \(I_a\cap U_?\ne\emptyset\) & 相冲突的状态标签须先报告并按显式策略处理，不能同时计入确认和不可用。两类不合并成确定掌握；有限行动库及可判分类上可算，E 的估计真实性另验 \\
\textbf{LC21 误区聚焦} & \(\theta\) 给个体关联的公共模式集合 \(P_s\) 及不确定性；S 取各模式的概念、诊断题和反例锚点 & 个体判断只作显示注记；模式本身不被改写。未知模式可作为待诊断而非确诊 \\
\textbf{LC22 表征选择} & \(\theta\) 给允许表征集合 \(R_s\)；对目标节点 n 取 \(\mathrm{Rep}(n)\cap R_s\) 及已认证转换 & 不改变 n 的形式负载；若集合空，报告缺少匹配表征，不能删除目标。有限登记可算 \\
\textbf{LC23 阶段适配} & \(\theta\) 给阶段约束 \(C_\mathrm{stage}(a,g)\)；在公开行动中筛选判定为允许者，保留依赖边界 & 阶段是外部模型解释，不是节点的永久年龄标签；未知可行性单列。有限且判据可判时完成 \\
\textbf{LC24 掌握证据缺口} & \(\theta\) 给目标能力对 \(Q_g\subseteq N\times\mathrm{Competence}\) 及已有观测覆盖 \(\mathrm{Cov}_E(s)\)；缺口 \texttt{D=Q\mcsbrk\_g\textbackslash{}Cov\mcsbrk\_E(s)}，选其节点投影及有公共任务引用的诊断节点；Anchor 覆盖 D 的外部条目另作派生注记 & 外部条目不伪装成 M 节点；缺少证据不等于未掌握，锚定不保证测量效度。有限集合差及匹配可算 \\
\end{longtable}
\endgroup

LC19 是显示折叠，不构造“此人版数学节点”；LC24 只指出覆盖缺口，不伪装为掌握事实。

\section{11.6 多路线比较族}
\begingroup\small\setlength{\tabcolsep}{4pt}
\begin{longtable}{>{\raggedright\arraybackslash}p{\dimexpr 0.2318\linewidth-2\tabcolsep\relax}>{\raggedright\arraybackslash}p{\dimexpr 0.3841\linewidth-2\tabcolsep\relax}>{\raggedright\arraybackslash}p{\dimexpr 0.3841\linewidth-2\tabcolsep\relax}}
\toprule
\textbf{编号、算子} & \textbf{输入与数学定义} & \textbf{保持、边界、损失与计算条件} \\
\midrule
\endfirsthead
\toprule
\textbf{编号、算子} & \textbf{输入与数学定义} & \textbf{保持、边界、损失与计算条件} \\
\midrule
\endhead
\bottomrule
\endfoot
\bottomrule
\endlastfoot
\textbf{LC25 共同核心} & 输入非空指定路线族 \(\mathcal{R}\) 和目标 g；\(S=\{g\}\cup\bigcap_{r\in\mathcal{R}}\mathrm{Req}_r(g)\) & 显示共同资源，不是充分路线；空族另报无路线/未知，不作真空全称推断。有限族可算 \\
\textbf{LC26 路线差异} & 输入 \(r_{1},r_{2}\)；S 为其活跃资源集对称差，另附两份来源指派 & 节点集合差不能代替事件多重性差；同时保留事件对应。有限路线可算 \\
\textbf{LC27 汇合点} & 输入路线及目标输出；选择具有同一节点引用的输出，或有显式桥接证书的输出对 & 同名或逻辑近似不自动汇合；形式汇合不代表外部认知状态相同。有限已知桥库可查 \\
\textbf{LC28 瓶颈} & 输入固定 B,G、有限 \(\Lambda'\)；节点 m 为瓶颈当 G 可达且按下述禁用语义排除 m 后 G 不可达；输出 m 及两次可达验证记录 & 与指定路线域相对；个体筛选得到的瓶颈不得回写全局。有限 Horn 可达可决定；约束规划需用相应有界判定 \\
\textbf{LC29 跨路线桥接} & 输入 \(r_{1},r_{2}\) 及认证的接口翻译 F；选择被 F 对应的节点/行动及翻译证据 & 仅输送已证明的结构关系；不是所有桥接都是 institution 态射，需写明实际类型。成本不随之保持 \\
\textbf{LC30 前提失效后的替代路线} & 输入禁用资源 F、B,G；在不使用 F 的契约路线中重新选取/规划，输出剩余路线及边界 & 不能只删图上节点就保留旧证书。无界搜索失败为未知，有界穷尽才报该界内不可达 \\
\end{longtable}
\endgroup

LC28 默认采用严格资源禁用：对禁用集合 F，把 B 换成 B\textbackslash{}F，排除 \((I_a\cup O_a)\cap F\ne\emptyset\) 的行动，并在 \(G\cap F\ne\emptyset\) 时直接判定不可达。LC30 采用同一约定。因此一个虽未被使用却会产出禁用节点的行动也被排除；若只想禁止消费某资源，应另给允许副产物的规则，不能混用两种瓶颈结论。禁用与普通“隐藏节点”不同；若禁止目标本身，瓶颈结论是平凡的，应明确标注。

\section{11.7 层级与资源族}
\begingroup\small\setlength{\tabcolsep}{4pt}
\begin{longtable}{>{\raggedright\arraybackslash}p{\dimexpr 0.2208\linewidth-2\tabcolsep\relax}>{\raggedright\arraybackslash}p{\dimexpr 0.3896\linewidth-2\tabcolsep\relax}>{\raggedright\arraybackslash}p{\dimexpr 0.3896\linewidth-2\tabcolsep\relax}}
\toprule
\textbf{编号、算子} & \textbf{输入与数学定义} & \textbf{保持、边界、损失与计算条件} \\
\midrule
\endfirsthead
\toprule
\textbf{编号、算子} & \textbf{输入与数学定义} & \textbf{保持、边界、损失与计算条件} \\
\midrule
\endhead
\bottomrule
\endfoot
\bottomrule
\endlastfoot
\textbf{LC31 话题聚合} & 输入第 09 章合法有限话题块族；每块为显示项，\(\tau\) 为成员集，关系按明确 \(\exists\) 或 \(\forall\exists\) 读法提升 & 保存成员与来源，可展开；不能把存在联系解释为整话题必修。重叠允许 \\
\textbf{LC32 跨领域接口} & 输入领域 A,B；选择端点分别在 Flat(A),Flat(B) 的指定关系及其接口前提 & 丢失域内部细节而保留跨域证据；没有跨边不证明两领域无数学联系。有限记录可查 \\
\textbf{LC33 层级投影} & 输入显示层级 l 或有限语法深度界 k；分别投影到该层聚合或选生成深度\(\le k\) 的候选 & 聚合层级与语法深度不是一个量，必须选择模式。有效有限代码能判深度；该层可能仍无限 \\
\textbf{LC34 时间预算} & 输入有限行动库、有限背景/目标与标注域、外部事件成本 \(c_E\)、事件界 h、预算 b；按第 12 章枚举满足总成本\(\le b\) 的结构路线并取其可见资源 & 在第 12 章有限有效条件、成本与预算比较均可判时有限完成；成本未知者单列。仅给 h 或可计算实数成本不足以保证全部比较终止；预算不能删掉依赖后仍宣称合法 \\
\textbf{LC35 证据类别} & 输入允许类别集合 C 与核验条件 \(\pi\)；选择 \(\mathrm{type}(\mathrm{ev})\in C\) 且 \(\pi(\mathrm{ev})\) 成立的关系记录及节点 & 证据类别不是统一可信度排序；经验结果不能按高分变成数学证明。有限记录及可判 \(\pi\) 可算 \\
\textbf{LC36 未知边界} & 输入研究域 U；标出 Supp=Unknown、未完成证明、未定匹配等对象及其公开依赖，输出“已知区/未决区”双层显示 & 未决不是假、空或不存在；只改变标注，不生成新硬关系。公开字段有限时可查，语义上所有未知不能穷尽 \\
\end{longtable}
\endgroup

\section{11.8 组合、视图模式与保持定理}
\textbf{定义 11.5（四类视图变换）。} 限制选子集；聚合给块及溯源；展开沿已保存成员引用恢复指定细节；表征转换使用明确的翻译证书。它们组成带类型的部分操作：后一操作的输入契约必须由前一视图及其环境边界提供。信息不足时返回未定义或待补依赖，不猜测。

\textbf{命题 11.6（纯限制的交换）。} 若 P,Q 是只依赖固定 M 的选择谓词，且关系保留规则相同，则两次纯限制的可见节点集均为 \(\{n:P(n)\land Q(n)\}\)。

\textbf{证明。} 集合交交换。边界须从最终可见对象的原始引用重新计算，才能保证溯源一致；简单丢弃中间边界不满足命题条件。\(\blacksquare\)

\textbf{反例 11.7（一般算子不交换）。} 链 \(a\to b\to c\)，从种子 \{c\} 在完整图中做依赖补全再限制到 \{a,c\} 得 \{a,c\}；先限制为 \{a,c\} 再只在诱导图中补全则得 \{c\}。后者丢失 b 的环境依赖，因此不能作为自包含视图。组合必须说明是在 M 中还是在 V 中求闭包。

\textbf{命题 11.8（适配不改本体）。} 任意外部参数包 \(\theta_{1},\theta_{2}\) 生成的局部化可以不同，但其来源均指向同一 M，且不改变 M 中任一节点或硬关系。

\textbf{证明。} 每个算子的输出仅由选择、聚合、复制引用及标注构成，未定义对 M 的更新操作；其语义是只读函数。逐次组合仍保持只读性。\(\blacksquare\)

上述算子族可按参数、组合和公开见证继续生成新的聚焦方式，其数量不限于 36。只有满足对应定理条件的实例才携带相应 Pres；不把“创造性筛选”升级成未经证明的模型论保持结论。

\section{11.9 聚合与展开的代数接口}
\textbf{定义 11.9（固定块族的展开与完整块选择）。} 对有限公共研究域 U 及块族 \((S_i)_{i\in I}\)，\(S_i\subseteq U\)，定义


\mcsdisplay{\beta:\mathcal P(I)\to\mathcal P(U),\quad
\beta(J)=\bigcup_{i\in J}S_i,
\qquad
\gamma:\mathcal P(U)\to\mathcal P(I),\quad
\gamma(A)=\{i:S_i\subseteq A\}.}
\noindent \(\beta\) 展开所选块，\(\gamma\) 选取内容全部落在可见域中的块。它不同于“只要块与 A 相交就选取”的存在式聚合；不相交或空块按集合定义处理，实际话题块要求非空。

\textbf{命题 11.10（聚合—展开伴随）。} 在包含偏序上有


\mcsdisplay{\beta(J)\subseteq A\iff J\subseteq\gamma(A).}
\noindent 因此 \(\gamma\beta\) 是块选择域上的闭包，\(\beta\gamma\) 是节点域上的内部算子；两者均幂等。\(\beta\gamma(A)\) 一般不等于 A。

\textbf{证明。} 并集落在 A 中当且仅当每个被选块都落在 A 中，即所述双条件。两映射单调，双条件给 \(J\subseteq\gamma\beta(J)\) 及 \(\beta\gamma(A)\subseteq A\)。把第二不等式代入 \(A=\beta(J)\)，再作用 \(\gamma\)，得到 \(\gamma\beta\gamma\beta(J)\subseteq\gamma\beta(J)\)，反向由广延；内部算子的幂等对偶得到。若块仅为 \{a,b\}，则 \(\beta\gamma(\{a\})=\emptyset\)，说明粗粒度不能保留所有细节。\(\blacksquare\)

伴随只描述成员信息的保留与损失。引用边界及证明标签仍依 Cut 重算，认知效用不由此伴随保证。

\textbf{定义 11.11（可组合的视图算子）。} 各种视图契约作为类型，算子为这些类型间带显式参数的部分函数，并携带实际证明的 Pres。若 \(F:V_{1}\rightharpoonup V_{2}\)、\(G:V_{2}\rightharpoonup V_{3}\)，则 \(G\circ F\) 在 F(v) 与 G(F(v)) 均定义时取该值。恒等与此复合满足结合律，构成部分函数范畴中的一个有类型操作体系。Pres 的复合另需目标性质与下一算子的前提匹配，不能只合并标签。

\textbf{命题 11.12（视图复合律）。} 在共同定义域内，\(H\circ(G\circ F)=(H\circ G)\circ F\)；恒等算子不改变结果。纯限制算子还满足交换、幂等及交的规律，聚合与一般补全不必满足交换律。

\textbf{证明。} 两种括号均恰在 F(v)、G(F(v))、H(G(F(v))) 顺次定义时有值，且值相同。恒等直接；纯限制等式化为集合交的相应等式，并按命题 11.6 重算边界。反例 11.7 排除一般交换律。\(\blacksquare\)

\clearpage\chapter{12 从局部空间到学习事件偏序}
\mcsorigin{原章节：12-视图转换与路径规划.md}
本章把第 06 章的结构路线与第 07 章的外部约束组合。结构可行、个人可行、某个预算内可行是三个不同判断。规划结果附条件与见证，不能仅画一张 DAG 就声称学习成功。

\section{12.1 输入、结果与失败状态}
\textbf{定义 12.1（规划输入）。} 固定本体版本 M、带边界视图 V、外部模型 E、状态 s、目标 g、预算 b 和策略 \(\rho\)。g 含目标节点集合 G 及能力规格；E 从显式状态给出已确认可用背景 B、未知入口和外部评价规则。公开候选行动 \(\Lambda'\) 来自 \(\Lambda_M\)，其缺少的前提允许从 \(\partial V\) 引入为明确环境引用。

策略 \(\rho\) 指定允许的结构行动、事件上界 h、外部不确定性的处理方式、排程约束和候选选择规则。方法行动证书只保证呈现接口，不保证状态转移成功。

\textbf{定义 12.2（结果包与偏序投影）。} 成功结果包为


\mcsdisplay{\Pi=(P,\preceq,act,h_{\rm src},\zeta,ev,\chi,\ell,w,\mathrm{meta}).}
\noindent P 是有限事件集；\(\preceq\) 为依赖偏序；act 标记行动；\(h_\mathrm{src}\)、\(\zeta\) 是第 06 章输入与目标来源；\(\mathrm{ev}:P\to N_M\) 是各事件的关注节点，所有实际输入输出另由 act 保留；\(\chi\) 是排程及状态约束；\(\ell\) 是至少一个通过约束的线性扩张；w 包含相应核验见证。meta 记录本体、外部模型、参数和算法版本。事件关注节点可重复。

对固定参数，将视图集记为 \(\mathrm{View}_M\)，将有限偏序的同构类型集合记为 FinPos；所需映射是结果包忘掉附加数据后的主投影：


\mcsdisplay{Plan_{M,E,s,g,b,\rho}:View_M\rightharpoonup FinPos,
\qquad V\longmapsto[(P,\preceq)].}
\noindent 附加结果包须另行保存，否则偏序本身无法说明外部状态约束。候选族函数 Plans 保存指定有限候选域中全部经确认可行的候选及条件性候选的独立标记；Pareto 是随后派生的展示子族，Plan 由 \(\rho\) 选择。确认可行只表示满足所声明的 E 模型判据，不表示现实学习成功已经发生。

\textbf{定义 12.3（分析状态）。} 计算接口另输出 Found、Conditional、Infeas\mcsbrk{}ibleWi\mcsbrk{}thinBo\mcsbrk{}und 或 Unknown。Conditional 带未确认入口/模型假设；Unknown 包括检查不可判、预算外搜索未完成或资料缺失；只有该界内全部候选均已判定失败且无待决项，才输出 Infeas\mcsbrk{}ibleWi\mcsbrk{}thinBo\mcsbrk{}und。不能把后两者统称全局不可规划。

\section{12.2 AND/OR 选择先于排序}
对每个行动，\(I_a\) 是联合输入，即 AND。针对同一目标选择不同合法行动或输入生产者是 OR。一个候选路线必须保存自己的选择，不能先把全部候选前提并成图再排序。

\textbf{例 12.4。} \(\Lambda'\) 有 \(a_{1}:\{u,v\}\to\{g\}\) 与 \(a_{2}:\{w\}\to\{g\}\)。若 B=\{u,v\}，选 \(a_{1}\) 可执行；若 B=\{w\}，选 \(a_{2}\) 可执行；\(B=\emptyset\) 时两者都不能直接执行。将图写成 u,v,w 均先于 g 错误增强要求，将三条边各自解释为充分则错误削弱要求。

个人状态可使某分支当前不适合，但该分支仍保留在公开 \(\Lambda_M\) 中。筛选后的共同必经点带 \((E,s,\rho)\) 标签，不回写为公共硬先修。

\section{12.3 事件图、循环和结构可靠性}
\textbf{定义 12.5（事件展开）。} 为每次行动使用分配新的事件标识。若输入 m 由事件 q 提供给 p，添加 \(q\to p\)；若来自 B，只记录背景引用。拒绝有环的来源指派，取


\mcsdisplay{\preceq=(\to)^*}
\noindent 为自反传递闭包。可另添加策略要求的次序边，但仍须无环。

\textbf{定理 12.6（结构排程可靠性）。} 每个线性扩张都尊重候选中已经选择的公共输入义务。

\textbf{证明。} 第 06 章命题 6.7 已证明：来源事件严格在使用事件之前，背景资源从起点可用。增加无环策略边只缩小线性扩张族，不损坏原依赖。\(\blacksquare\)

这不保证任一扩张满足外部遗忘、容量或时长约束，所以 \(\chi\) 仍需逐排程检验。

\textbf{命题 12.7（循环诊断）。} 对有限行动库，空背景下只有相互等待的行动且没有可触发行动时，不存在非空可执行序列。

\textbf{证明。} 若有序列，其第一个事件全部输入须来自背景，与不可触发矛盾。\(\blacksquare\)

有效的联合入口必须是一份独立合法的新行动或来源见证，并实际生成无环展开。压缩强连通分量只是图分析，不认证学习入口。不可达证书可由有限 Horn 可达闭包缺少目标给出；带额外个人约束时，该闭包可达只是必要结构检查。

\section{12.4 外部可行性与动态状态}
\textbf{定义 12.8（外部排程判据）。} 对线性扩张 \(\ell=(p_{1},\ldots,p_k)\)，外部关系 \(\mathrm{Run}_E(s,\ell,s',z)\) 给出可能终态及观察轨迹；依据第 07 章关系或概率核解释。\(\chi\) 明确采用以下哪种判据：

\begin{itemize}
\item 确定性模型：按转移逐步计算，入口、预算与终态目标全部满足；
\item 有限不确定情景模型：选择存在情景可行，或所有列明情景可行，不能混用量词；
\item 概率模型：在给定模型下成功概率达到声明阈值，需可判的阈值比较或足以分离阈值的误差界证书。下界已达到阈值可认证通过，上界严格小于阈值可认证失败；区间跨越阈值时保留 Unknown。
\end{itemize}
默认缺少外部模型时不赋概率，返回条件性结构路线。无界或连续模型的可行性可能不可判；有界事件数不自动使外部积分或真实心理判断可判。即使概率是可计算实数，对阈值的非严格比较也未必可判，特别是在恰等于阈值且没有等号证书时。

个人掌握允许遗忘。结构资源“一个证明已经存在”不会因某人遗忘而从 M 消失；但该人是否需要复习，由 E 产生新的事件与约束。

\section{12.5 路线比较、等价与转换}
\textbf{定义 12.9（四种比较）。}

\begin{enumerate}
\item 目标等价：同一形式目标，或有指定桥接证书的目标；不是相同终态认知。
\item 事件结构同构：事件双射保持行动标签及偏序；可保留不同来源编号。
\item 资源比较：比较 Req、证明支持和所用公理；包含关系不直接等于认知难度顺序。
\item 外部效用比较：E 输出成本与效果偏序、区间或 Unknown。
\end{enumerate}
若成本向量 c、效果向量 q 均有可比较值，Pareto 支配要求每项成本不大于、效果不小于，并至少一项严格改进。不知道的效果不能用默认数字完成支配。保存未被证实支配者，不宣称它们是真实最优路线。

\textbf{命题 12.10（线性扩张的结构互换）。} 同一有限偏序的任意两个线性扩张可由相邻不可比事件的有限交换互达。

\textbf{证明。} 取目标扩张的首元素 e，在另一扩张中把它向前移动。位于它之前的元素不能小于 e，因为目标扩张把 e 放第一；也不能大于 e，因为当前扩张合法，故均与 e 不可比，可以逐次交换。固定首元素，对剩余偏序归纳。\(\blacksquare\)

若想据此断言认知等价，还须相关外部转移可交换且成本与观察判据保持；一般不成立。结构路线转换与认知效果相同是两个不同定理。

\section{12.6 有限片段的完整搜索}
\textbf{定义 12.11（有限搜索条件）。} 要求：

\begin{itemize}
\item \(\Lambda'\)、候选背景/目标和允许事件关注节点均有可终止返回的完整有限列表；
\item 明确事件上界 \(h\in\mathbb{N}\)；重复行动允许，但次数受 h 限制；
\item 每份来源/关注标注下，允许的策略次序与 \(\chi\) 实例由固定 \(\rho\) 在有限时间给出完整有限列表；固定唯一规则是其特例，不遍历任意约束公式；
\item 契约、来源、标注合法性、预算与所选外部情景检查均可判，包括所用实数阈值的比较；
\item 比较与策略选择在该有限候选集合上可判且终止。
\end{itemize}
证明证书作为已提供输入，或其枚举也有显式长度界；不能隐含调用“寻找任意数学证明”的无界过程。

\textbf{算法 12.12（穷尽而有界的参考算法）。}

\begin{enumerate}
\item 枚举 \(0\le k\le h\) 的行动序列，事件名规范为 1,…,k，并枚举每个事件全部允许的关注节点指派；
\item 对序列中每个输入，枚举背景或此前输出的合法来源指派，拒绝缺输入者；
\item 枚举并保留满足 G 的所有目标来源指派，构造事件来源 DAG；枚举 \(\rho\) 所给策略次序与 \(\chi\) 实例，拒绝与来源边合并后有环的次序，再取偏序闭包；
\item 枚举候选偏序的全部线性扩张，运行 \(\chi\)，分别保存通过的 \(\ell\) 与见证；不同排程可能有不同成本与效果，不能仅保留首个通过者后就声称完成 Pareto 比较。未能判定的结果单列，不能伪装成通过或失败；
\item 仅在同构保持全部外部相关标注、来源和约束时去重；否则保留两项。完整可行集留作 Plans；再比较已知成本/效果，派生 Pareto 展示子族并保留无法判定支配关系的候选；
\item \(\rho\) 明确选一个或展示整个候选族。纯复现用途可按有限代码序选首项，不把该顺序解释为教育最优。
\end{enumerate}
\textbf{定理 12.13（有界终止、可靠与相对完备）。} 定义 12.11 条件下，算法有限终止；认证为 Found 的结果满足所有已声明条件；每个在指定库、标注/约束候选域及事件界内满足条件的路线至少有一个保持这些数据的同构代表进入 Plans。Pareto 展示子族不声称保留所有可行路线。

\textbf{证明。} 若 \(|\Lambda'|=m\)，行动词总数为 \(\Sigma_{k\le h}m^k\)，其中长度零有且仅有空词。每个事件的有限输入只有有限个此前来源，目标来源与关注标注亦有限，策略次序/约束实例来自完整有限列表，每个偏序的线性扩张不超过 k! 个。每次检查终止，因此整个枚举终止。可靠性由逐条拒绝规则与记录的见证。任一候选域内的合法路线有一个尊重其来源与策略次序的拓扑排序，对应枚举中的行动词；其输入来源、目标来源、关注标注及约束实例也全部被枚举，其每个可行线性扩张均被检查，故代表及排程见证进入 Plans。同构去重按上述完整数据进行，不遗漏外部可区分的候选；后续 Pareto 筛选不改变此枚举结论。\(\blacksquare\)

最坏复杂度可以很高；本稿承诺有界参考过程正确性，不承诺工业规模性能或全局最优。

\section{12.7 视图边界、冷启动及非干扰}
\textbf{命题 12.14（环境保持的规划）。} 若 V 只隐藏资源的显示，仍保留契约、证书和环境来源，则在 M 中合法且全部引用可解析的候选可被记录为带边界路线；不能仅因输入不可见而删除其义务。

\textbf{证明。} 输入 \(h_\mathrm{src}\) 保持原环境资源引用，核验在相同 M 版本中完成。显示集合缩小不改变来源可用性。若来源在外部状态中未确认可用，则结果标 Conditional，而不是结构失败或假定已掌握。\(\blacksquare\)

冷启动可展示“经 a”或“经 b”的公共路线及入口条件，建议诊断题；不猜测用户已掌握哪条入口。同一共享空间的两名学习者可以得到不同事件序列、复习安排和提示，M 本身不变。

本章的完成要求是接口、一般证明与有限实例核查。软件产品、真实学习效果实验以及无限域最优规划属于后续工作，不能在交付表中冒称已实现。

\section{12.8 路径表达的代数}
\textbf{定义 12.15（有端口的路线表达）。} 本节接口 B、G 均有限。记 \(r:B\Rightarrow G\) 为带全部输入来源 h、目标端口指派 \(\zeta\) 的有限结构路线，可附结构排程次序 \(K_r\)；\(K_r\) 与来源边的并须无环。\(K_r\) 仅规定事件先后，不含个人状态或概率断言。定义以下有类型运算，事件名在运算前统一改为互异名字。

\begin{enumerate}
\item 单位 \(\varepsilon_B:B\Rightarrow B\) 没有事件，目标均由背景端口提供。
\item 顺序 \(r_{2};r_{1}:B\Rightarrow G\)：若 \(r_{1}:B\Rightarrow C\)、\(r_{2}:C\Rightarrow G\)，定义端口替换 \(\sigma\)，使 \(\sigma(\)背景 \(m\in C)=\zeta_{1}(m)\)，而 \(r_{2}\) 的事件来源保持原事件。对 \(r_{2}\) 的每个输入来源 \(h_{2}(p,m)\) \textbf{及每个最终目标来源 \(\zeta_{2}(g)\)} 同时应用 \(\sigma\)；\(r_{1}\) 的输入来源不变。保留两路的 K，并加入所有第一路事件早于所有第二路事件的结构排程要求。中间端口不匹配则未定义。
\item 并行 \(r_{1}\parallel r_{2}:B\Rightarrow G_{1}\cup G_{2}\)：两路从同一 B 出发，要求本运算的目标端口集合 \(G_{1}\)、\(G_{2}\) 不相交，事件名不相交，各自来源保持，不添加跨路依赖。内部产出可以重叠。目标重叠时必须另行给端口合并操作与来源选择，本运算暂不定义。
\item 选择 \(\mathcal{F}\oplus\mathcal{G}\) 是相同输入/目标接口的路线族并集；它是 OR。单一行动的输入集合仍是 AND，不因选择运算改变。
\end{enumerate}
顺序中的分号按“右边先执行”约定，与函数复合一致。严格排程约束与来源偏序分开：无资源依赖的两步也可因策略被顺序安排。并行保留各路原有 K，意味着结构上允许交错；合成后的外部容量、状态和成功概率须重新由 \(\chi\) 检查，不从各路单独可行直接推出。

\textbf{命题 12.16（路线运算的结构律）。} 上述运算保留结构合法性。模去事件改名且保留端口指派，顺序结合并以 \(\varepsilon\) 为单位；并行结合、交换；路线族选择结合、交换、幂等。在所有接口匹配且尚未按个人成本剪枝时，逐元素顺序组合和并行组合均对族选择分配。

\textbf{证明。} 顺序的所有新来源边及新增次序均从第一路指向第二路，不能形成回边，因此仍无环；每个替换端口仍有相同节点类型的来源。三路组合的任一输入及最终目标沿端口追溯到同一最早来源，代入复合结合；附加排程在两种括号下都要求第一路早于第二路早于第三路。右单位的替换是背景恒等；左单位虽无事件，其目标背景端口经 \(\sigma\) 恰恢复原 \(\zeta\)，故也保持路线，且空事件不添次序。并行是不相交事件并及共同背景引用，三路目标两两不交时，两种括号均有定义；调换顺序仅影响事件命名和输出端口排列。族运算按存在某个分支作元素定义，分配律即存在量词对并集的分配。\(\blacksquare\)

例如 \(r:\{b\}\Rightarrow\{c\}\) 以事件 e 的输出 c 为目标，\(\varepsilon_\{\\mathrm{setminus} {c\}}\) 的目标来源则是背景 c。复合 \(\varepsilon_\{\\mathrm{setminus} {c\}};r\) 时必须把该目标引用替换为事件 e；若只替换输入而遗漏 \(\zeta\)，最终路线会错误引用不属于背景 \{b\} 的 c，连单位律与结构合法性都会失败。

这里没有证明顺序交换，也没有证明认知状态转移交换。若 E 中两行动互相干扰，\(r_{1}\parallel r_{2}\) 的若干线性扩张会被 \(\chi\) 排除；若预算或概率阈值用于中间剪枝，分配后的候选集合也可能改变。比较优化策略时必须指出在哪一步筛选。

\clearpage\chapter{13 四组案例的贯通与交叉检验}
\mcsorigin{原章节：13-贯通案例.md}
案例把同一理论接口应用到四种不同数学情境。每组都给出共享节点、明确数学条件、所用概念支持、公开契约、方法与误区、局部化、事件偏序，以及外部学习者差异。自然语言数学证明与机器证书严格分开：完整桥梁已有正文论证；2026-09-29 新增四个实际数学片段的 ND 子集证书，具体范围见覆盖表，完整桥梁仍未机器认证。

\section{13.1 案例入口}
\begingroup\small\setlength{\tabcolsep}{4pt}
\begin{longtable}{>{\raggedright\arraybackslash}p{\dimexpr 0.3333\linewidth-2\tabcolsep\relax}>{\raggedright\arraybackslash}p{\dimexpr 0.3333\linewidth-2\tabcolsep\relax}>{\raggedright\arraybackslash}p{\dimexpr 0.3333\linewidth-2\tabcolsep\relax}}
\toprule
\textbf{案例} & \textbf{数学桥梁/反例} & \textbf{多路线与外置内容} \\
\midrule
\endfirsthead
\toprule
\textbf{案例} & \textbf{数学桥梁/反例} & \textbf{多路线与外置内容} \\
\midrule
\endhead
\bottomrule
\endfoot
\bottomrule
\endlastfoot
极限与连续 & \(\varepsilon-\delta\) 与序列判别的双向证明；常值序列反驳“趋近不能取到” & 量词入口与序列入口；个人是否持有误区外置；桥梁目标仍有联合输入 \\
\(C^{k}\) 与光滑流形 & 同一图册载体上 \(C\infty\Rightarrow C^{k}\)；最大图册的遗忘映射；\(C^{1}\) 非 \(C^{2}\) 坐标转换 & 拓扑入口与坐标计算入口；显示粒度外置；不把图册不光滑误说成底层空间不可光滑化 \\
张量与张量场 & \(V\otimes V*\cong\mathrm{End}(V)\) 的明确同构；兼容光滑分量与截面；Christoffel 非张量反例 & 矩阵/基入口与多重线性入口；向场的汇合须有丛或等价兼容条件 \\
群概念的多来源 & Cayley 定理、三角形对称、模 5 单位群、\(S_{3}\) 非交换例 & 置换、几何、算术三入口；两名学习者不同筛选但共享定义与证明 \\
\end{longtable}
\endgroup

\section{13.2 共同的数据归属}
每组案例中的同一个节点引用、正式条件、证明论证和反例都在 M 中。群的交换律真假、流形的微分阶数、张量的变换规律，不随学生或展示变化。

每名学习者的确认背景、未决能力、偏好、已激活误区、复习需求及时间成本，都属于 E,s。这里的学习者是\textbf{明确假设的演示实例}，不是基于采样得到的人群规律；文本中“适合某入口”只在给定模型条件下成立。

具体选定的提示、路线、局部图及事件排序属于 D。公开契约只认证资源/呈现接口，不能将输出节点直接解释成个体掌握。若外部模型把一次行动设为可靠掌握转移，须将这条假设写在 E 中。

\section{13.3 跨案例同构与不同点}
四例都可以抽取“两个入口—共同目标”的公共形式，但各例的桥梁不同：

\begin{itemize}
\item 极限：在度量与所列选择背景下，两个定义具有等价证据。
\item 流形：忘却部分正则性是一向结构映射；不可默认为任意相反方向的升级。
\item 张量：有限维假设支持指定线性同构；张量场的局部到整体还要求兼容性。
\item 群：不同来源的实例满足同一公理；Cayley 给忠实置换表示，不把每个群强行等同于某个固定对称群。
\end{itemize}
所以“多路线”的统一层次是公共行动及证据结构，不是宣称所有数学桥梁都同型，更不是所有路线有相同认知效果。第 08 章态射必须检查具体翻译与保持义务。

\section{13.4 局部化组合示例}
群案例可先 LC32 取几何—代数接口，再 LC13 回溯 Group，随后由外部状态产生 LC20 的确认/可能前沿；最终执行第 12 章的 AND/OR 选择与排序。顺序及原空间版本进入 \(\lambda\)。

张量案例先 LC16 追踪坐标变换下对象不变，再 LC14 展示非张量反例，随后 LC29 显示矩阵与抽象路线桥梁。若隐藏联络背景，边界必须保存它，不能把计算反例呈现成无前提结论。

流形案例的 LC10/LC11 必须声明是在“同一图册集合上的性质”还是“不同最大结构的遗忘映射”模式；两种模式不共用一个未注明载体的包含符号。

极限案例的 LC25 对非空选定路线族求共同资源。目标只是陈述一个定义与目标是证明两个标准等价时，路线族不同，不能跨目标偷换共同先修。

\section{13.5 与有限核验的关系}
有限模型脚本 及 报告 对这些模式的小型符号实例检查：联合前提、替代入口、两名合成学习者的群/张量路线与极限提示、状态更新非干扰、冷启动及隐藏依赖。符号行动不是四个数学领域的完整形式库；脚本没有重证极限判别、Cayley 或张量同构。

数学论证由各案例正文给出；逻辑与结构性质由相应章节给出；外部认知效度仍需未来实际观测。三种证据各承担自己的结论。

\clearpage\chapter{案例 01 极限、连续与两种认知入口}
\mcsorigin{原章节：01-极限与连续.md}
\section{1 公共数学背景与节点}
工作背景为实数的度量结构；函数 \(f:D\to\mathbb{R}\)，a 是 D 的聚点。函数极限只讨论 \(x\ne a\) 的去心条件，连续则讨论 \(a\in D\) 时 f(a) 与极限/邻域条件的相容。所有个人状态都在第 5 节外置。

\begingroup\small\setlength{\tabcolsep}{4pt}
\begin{longtable}{>{\raggedright\arraybackslash}p{\dimexpr 0.1930\linewidth-2\tabcolsep\relax}>{\raggedright\arraybackslash}p{\dimexpr 0.4467\linewidth-2\tabcolsep\relax}>{\raggedright\arraybackslash}p{\dimexpr 0.3603\linewidth-2\tabcolsep\relax}}
\toprule
\textbf{节点} & \textbf{类型及正式内容} & \textbf{选定概念支持} \\
\midrule
\endfirsthead
\toprule
\textbf{节点} & \textbf{类型及正式内容} & \textbf{选定概念支持} \\
\midrule
\endhead
\bottomrule
\endfoot
\bottomrule
\endlastfoot
Rdist & 概念：\(d(x,y)=|x-y|\) & 实数、绝对值、距离 \\
SeqConv & 概念：\(x_n\to a\) \(\Leftrightarrow\) \(\forall\varepsilon>0\exists N\forall n\ge N,|x_n-a|<\varepsilon\) & 序列、量词、距离 \\
LimitED & 概念：\(\mathrm{lim}_{x\to a}f(x)=L\) \(\Leftrightarrow\) \(\forall\varepsilon>0\exists\delta>0\forall x\in D,0<|x-a|<\delta\Rightarrow|f(x)-L|<\varepsilon\) & 函数、距离、去心邻域、量词 \\
LimitSeq & 概念：对所有 D\textbackslash{}\{a\} 中趋于 a 的序列，\(f(x_n)\to L\) & 序列、映射、收敛、全称量化 \\
Bridge & 定理：在上述度量域中 \(\mathrm{LimitED}\Leftrightarrow\mathrm{LimitSeq}\) & 两种极限定义及选定证明 \\
Continuous & 概念：\(\forall\varepsilon>0\exists\delta>0\forall x\in D,|x-a|<\delta\Rightarrow|f(x)-f(a)|<\varepsilon\) & 函数、邻域、函数值 \\
BridgeCont & 定理：\(a\in D\) 且为聚点时，连续等价于函数极限为 f(a) & Continuous、LimitED、函数值 \\
ConstantSeq & 例子：\(x_n=a\) 对所有 n；其极限为 a & 序列、极限、恒等 \\
NoNeverEqual & 性质：收敛不能推出各项不等于极限 & ConstantSeq、SeqConv、全称命题 \\
\end{longtable}
\endgroup

量词、函数值、序列表达等基础引用在本案例的环境边界中声明，未假装不存在这些背景。节点身份不同，Bridge 不把两个表达节点合并。\(C_\mathrm{all}\) 在第 05 章条件下均可能饱和，所以表中使用显式表达支持。

\section{2 桥梁定理的完整数学证明}
若 \(\varepsilon-\delta\) 条件成立，任取 \(x_n\in D\\mathrm{setminus} {a}\) 且 \(x_n\to a\)。给 \(\varepsilon>0\)，取对应 \(\delta\)；存在 N，使 \(n\ge N\) 时 \(|x_n-a|<\delta\)。由 \(x_n\ne a\) 得 \(|f(x_n)-L|<\varepsilon\)，故 \(f(x_n)\to L\)。

反之，若 \(\varepsilon-\delta\) 条件不成立，则存在 \(\varepsilon_{0}>0\)，使每个 \(\delta>0\) 都有 \(x\in D\) 满足 \(0<|x-a|<\delta\) 且 \(|f(x)-L|\ge\varepsilon_{0}\)。对 \(\delta=1/n\) 选择一个 \(x_n\)，得到 \(x_n\to a\)，但 \(f(x_n)\) 不趋于 L，矛盾。此选择在本稿 ZFC 元理论中成立；若对象理论削弱选择原则，须另外提供序列选择所需前提。\(\blacksquare\)

该证明给出正式化目标及全部关键假设；尚未编码为第 02 章 ND 证书，Check 状态是“未运行”。不能把“序列入口更直观”当作数学上不再使用 \(\varepsilon-\delta\) 的证明。

连续与极限的桥梁也可直接验证。若 f 在 a 连续，将连续的邻域条件限制到 \(x\ne a\)，就得到以 f(a) 为值的函数极限。反之，极限条件处理所有 \(x\ne a\)；当 x=a 时，\(|f(x)-f(a)|=0\) 自动满足要求，两种情形合起来给出连续。\(\blacksquare\)

聚点条件使这里采用的函数极限概念具有通常的唯一性背景。若 a 是 D 的孤立点，f 在 a 自动连续；不能在未声明约定时仍用去心条件给它强指定一个唯一函数极限。这个边界应写进定义模板。

\section{3 公开行动契约与路线}
下表是待认证形式实现与公共呈现的明确规格。数学结论由上节论证支撑；关于学习成功的断言未被加入。

\begingroup\small\setlength{\tabcolsep}{4pt}
\begin{longtable}{l>{\raggedright\arraybackslash}p{\dimexpr 0.2939\linewidth-2\tabcolsep\relax}>{\raggedright\arraybackslash}p{\dimexpr 0.2365\linewidth-2\tabcolsep\relax}>{\raggedright\arraybackslash}p{\dimexpr 0.3988\linewidth-2\tabcolsep\relax}}
\toprule
\textbf{行动} & \textbf{\(I_a\)} & \textbf{\(O_a\)} & \textbf{模式/见证} \\
\midrule
\endfirsthead
\toprule
\textbf{行动} & \textbf{\(I_a\)} & \textbf{\(O_a\)} & \textbf{模式/见证} \\
\midrule
\endhead
\bottomrule
\endfoot
\bottomrule
\endlastfoot
\(a_D\) & Rdist、量词表达 & LimitED & 定义引入及去心条件说明 \\
\(a_S\) & Rdist、序列表达 & SeqConv、LimitSeq & 定义/表征呈现 \\
\(a_B\) & LimitED、LimitSeq & Bridge & 上节双向证明 \\
\(a_C\) & LimitED、函数值 & Continuous & 连续定义与 f(a) 的目标代入 \\
\(a_\mathrm{BC}\) & LimitED、Continuous、函数值 & BridgeCont & 上述去心/不去心两情形证明 \\
\(a_X\) & SeqConv、ConstantSeq & NoNeverEqual & 常值序列的直接反例 \\
\(a_\mathrm{DS}\) & Bridge、LimitED & LimitSeq & 桥接引用，不等于另一端完整掌握 \\
\(a_\mathrm{SD}\) & Bridge、LimitSeq & LimitED & 同一桥接的反向接口 \\
\end{longtable}
\endgroup

目标若仅是“能用 \(\varepsilon-\delta\) 陈述极限”，\(a_D\) 是一条入口。目标若是“能辨析两种标准”，需获得两种定义并执行 \(a_B\)。目标不同不能假装为同一路线的更优替代。学习者可以先 \(a_D\) 后 \(a_S\)，也可先 \(a_S\) 后 \(a_D\)；完整 Bridge 的共同结构要求仍是两种定义。

\section{4 方法、误区与自检}
全局方法为“直觉的建立与失效”“构造反例”；局部方法为“用常值或最终常值序列检验动态趋近意象”。方法接口输入为一个含全称词的直觉猜测，输出为候选反例及其否定范围，支持由 SeqConv 和 ConstantSeq 锚定。

模式 NeverEqual 的错误规则是\(“x_n\to a\) 蕴含所有 n 都有 \(x_n\ne a”\)。常值序列使 \(|x_n-a|=0<\varepsilon\)，取 N=1 即满足定义，同时每项等于 a，故反驳该规则。它不否定函数极限定义中的去心条件；两者量化对象不同。

自检：\mcsnum{1}给出能取到与不能取到极限的收敛序列；\mcsnum{2}解释函数极限为何不必等于 f(a)；\mcsnum{3}写出 \(a_B\) 反向证明的量词否定。评分依据是定义与上节证明，不根据学习者信心打分。

数学史动机采用 Stillwell 第 9 章 §9.1（印刷 121–122 / PDF 137–138）的面积、切线、运动与级数问题；概念意象区分参照 Tall–Vinner 作者稿。历史来源不充当硬先修。

\section{5 两个外部状态与冷启动}
共享 M 完全相同。\(E_A\) 的示例状态确认“量词陈述、距离”，序列能力未知；\(E_B\) 确认“序列例子、距离”，量词陈述仍需诊断。这是说明接口的假设状态，不是实测结论。

A 的初步呈现选择 \(a_D\)；B 先选择 \(a_S\) 并增加量词澄清任务。两人若要完成 Bridge，最终都需满足相同公开输入。若仅 B 的观测提示 NeverEqual，则 \(\mathrm{Inst}_\mathrm{EB}\) 指向该公共模式并触发 \(a_X\)；公共模式内容不因此改变。

没有个人数据时，两种入口均可展示为 Conditional，不默认用户没有距离知识，也不默认已经掌握量词。题目条目分别锚定“定义陈述”“例子辨认”“证明”能力，不能把一道题答对扩成全部掌握。

\section{6 局部化与事件偏序}
用 LC13 从 Bridge 回溯得到 \(a_D\)、\(a_S\)、\(a_B\) 的候选；LC21 按外部模式选择 \(a_X\)；LC04 展示双向证明切片。隐藏 Rdist 时，\(\partial V\) 仍记录其引用。

设背景提供距离、量词和序列表达，事件为 \(e_D,e_S,e_B\)，依赖 \(e_D<e_B\)、\(e_S<e_B\)。两个线性扩张 \(e_D,e_S,e_B\) 与 \(e_S,e_D,e_B\) 都满足结构要求。外部偏好可以选择顺序，却不能把另一条从本体删除。若增加复习 \(e'_D\)，它与 \(e_D\) 关注同一节点，但有独立事件身份和时序。

验收：两入口合法、Bridge 共同要求不丢失、错误直觉能被精确反驳、个人诊断及复习均不改 M。

\section{7 本轮机器覆盖（2026-09-29）}
常值取值证书已编码并检查 \((\lambda n.a)(n)=a\)，无开放假设。N、R在此证书中仅为排序名，尚未编码实数度量、序列收敛或完整Bridge。自然语言证明、片段认证、公开行动和外部评价分别列在四案例覆盖表，不可把这一取值步骤扩大为收敛桥梁已认证。

\clearpage\chapter{案例 02 \(C^k\) 图册、光滑流形与泛化的载体}
\mcsorigin{原章节：02-Ck与光滑流形.md}
\section{1 约定与公共节点}
固定有限维 n、Hausdorff 且第二可数的拓扑空间 X。一张图 \((U,\varphi)\) 是 \(U\subseteq X\) 开集到 \(\mathbb{R}^n\) 开集的同胚。图册覆盖 X；对图 \((U,\varphi)\)、\((V,\psi)\)，过渡映射 \(\psi\circ\varphi^{-1}\) 的域为 \(\varphi(U\cap V)\)、像为 \(\psi(U\cap V)\)，两者均开。\(C^k\) 条件对所有有序图对成立，因而两个方向都必须 \(C^k\)；仅检验一个方向的 \(C^k\) 同胚不足以推出逆也 \(C^k\)。此处\(“C^k\) 流形”明确 \(k\ge 1\)；不让“微分流形”的名称暗中决定 k。

\begingroup\small\setlength{\tabcolsep}{4pt}
\begin{longtable}{>{\raggedright\arraybackslash}p{\dimexpr 0.2551\linewidth-2\tabcolsep\relax}>{\raggedright\arraybackslash}p{\dimexpr 0.3976\linewidth-2\tabcolsep\relax}>{\raggedright\arraybackslash}p{\dimexpr 0.3474\linewidth-2\tabcolsep\relax}}
\toprule
\textbf{节点} & \textbf{形式内容} & \textbf{选定支持} \\
\midrule
\endfirsthead
\toprule
\textbf{节点} & \textbf{形式内容} & \textbf{选定支持} \\
\midrule
\endhead
\bottomrule
\endfoot
\bottomrule
\endlastfoot
TopMan & 局部欧氏的上述拓扑空间 & 拓扑、同胚、覆盖、\(\mathbb{R}^n\) \\
Chart/Atlas & 图、覆盖图册 & 开集、坐标、覆盖 \\
Transition & \(\psi\circ\varphi^{-1}\) 的重叠域映射 & 复合、逆、定义域 \\
CkAtlas & 所有过渡映射均 \(C^k\) & 图册、多元微分、\(C^k\) \\
SmoothAtlas & 所有过渡映射均 \(C^\infty\) & 图册、各阶微分 \\
\(\mathrm{Compatible}_k\) & 两图册的并为 \(C^k\) 图册 & 图册、兼容 \\
\(\mathrm{Max}_k\) & 给定图册的最大 \(C^k\) 兼容扩张 & 兼容、集合并 \\
\end{longtable}
\endgroup

这些记录保留 k、维数、底空间及图册等参数。光滑结构与 \(C^k\) 结构的极大图册可能不是同一集合。

\section{2 两种精确的硬泛化}
第一种采用共同载体：所有上述拓扑图册编码。谓词 \(P_k(A)\) 表示每个过渡函数 \(C^k\)，\(P_\infty(A)\) 表示每个过渡函数所有阶可微且导数连续。由定义，对每个有限 k 有


\mcsdisplay{P_\infty(A)\Rightarrow P_k(A).}
\noindent 这给出同载体的 \(\mathrm{CkAtlas}\succeq_H\) SmoothAtlas，而非反向。

第二种采用结构对象 \((X,[A]_\infty)\) 与 \((X,[A]_k)\)。将光滑图册视作 \(C^k\) 图册，再取相应最大 \(C^k\) 扩张，得到明确的遗忘映射。若 A、B 表示同一光滑结构，则 \(A\cup B\) 是光滑图册，因而是 \(C^k\) 图册；它们的最大 \(C^k\) 扩张相同，故映射不依赖代表图册的选择。它保留底空间及指定的 \(C^k\) 结构，不宣称光滑极大图册与 \(C^k\) 极大图册字面相等，也不宣称存在典范逆函数。

\section{3 兼容扩张与一个严格边界例}
对一个 \(C^k\) 图册 A，令 \(A_\mathrm{max}\) 为所有与 A 的每张图兼容的图。\(A_\mathrm{max}\) 仍为 \(C^k\) 图册：覆盖由 \(A\subseteq A_\mathrm{max}\) 保证。两张新增图在某重叠点附近都可经过 A 的某张图，过渡映射是两个 \(C^k\) 映射的复合，故局部 \(C^k\)；正则性是局部性质，因而整个过渡域 \(C^k\)。它显然包含每张可再兼容加入的图，所以最大。\(\blacksquare\)

在 \(X=\mathbb{R}\) 上，令 h(x)=x+x|x|。h'(x)=1+2|x| 连续且正，h 严格递增且两端趋于 \(\pm\infty\)，所以为同胚；逆的导数为 \(1/h'(h^{-1}(y))\)，可由差商直接推出并且连续。因此 \{id,h\} 是 \(C^{1}\) 图册。h 在 0 两侧的二阶导数分别为 2 和 \(-2\)，不是 \(C^{2}\)，故这个图册不是 \(C^{2}\) 图册。

反例只证明同一图册的 \(C^{1}\) 条件不推出 \(C^{2}\) 条件；\textbf{不}证明 \(\mathbb{R}\) 不能另取光滑图册。这样避免把表达/结构选择与底空间存在性混为一谈。

上述为完整数学论证，依赖链式法则等分析背景；尚未转成第 02 章 ND 证书。泛化方向本身由所选谓词直接验证，不借用未经展开的光滑化定理。

\section{4 公开契约与多路线}
\begingroup\small\setlength{\tabcolsep}{4pt}
\begin{longtable}{>{\raggedright\arraybackslash}p{\dimexpr 0.1709\linewidth-2\tabcolsep\relax}>{\raggedright\arraybackslash}p{\dimexpr 0.3438\linewidth-2\tabcolsep\relax}>{\raggedright\arraybackslash}p{\dimexpr 0.2515\linewidth-2\tabcolsep\relax}>{\raggedright\arraybackslash}p{\dimexpr 0.2339\linewidth-2\tabcolsep\relax}}
\toprule
\textbf{行动} & \textbf{输入} & \textbf{输出} & \textbf{结构依据} \\
\midrule
\endfirsthead
\toprule
\textbf{行动} & \textbf{输入} & \textbf{输出} & \textbf{结构依据} \\
\midrule
\endhead
\bottomrule
\endfoot
\bottomrule
\endlastfoot
\(m_\mathrm{chart}\) & 拓扑、\(\mathbb{R}^n\)、局部同胚 & Chart、Atlas & 定义引入 \\
\(m_\mathrm{trans}\) & Atlas、复合与逆 & Transition & 重叠域与映射计算 \\
\(m_\mathrm{reg}\) & Transition、多元 \(C^k\) & CkAtlas & 正则性条件 \\
\(m_\mathrm{smooth}\) & Transition、各阶微分 & SmoothAtlas & \(C^\infty\) 条件 \\
\(m_\mathrm{forget}\) & SmoothAtlas、k、所选载体模式 & 同一图册的 \(C^k\) 性质，或相应最大 \(C^k\) 结构 & 两种模式分别使用上节蕴含/遗忘证明，不混合输出类型 \\
\(m_\mathrm{example}\) & 一元微分、图册 & \(C^{1}\) 非 \(C^{2}\) 的图册例 & h 的完整计算 \\
\end{longtable}
\endgroup

此表的定义呈现行动输出公共概念接口；若目标改为认证某个指定图册确为 \(C^k\) 或光滑图册，\(m_\mathrm{reg}\)、\(m_\mathrm{smooth}\) 还必须携带所有相关过渡映射满足正则性的证据。已有\(“C^k”\)概念不等于某个输入图册已满足该性质。

“从曲面参数化进入”与“从抽象拓扑图册进入”是两种公共呈现路线。前者通过球面等例子建立局部坐标动机；后者直接给出图册接口。两者可汇合于同一通用 Transition 定义；具体曲面、图册及坐标对上的过渡映射则保留不同参数与节点身份，不能因使用同一模板就把它们当作同一映射。

全局方法是“局部到整体”“选择表征”“从特例切入”；局部方法是“逐对检查图册过渡正则性”。输入为给定图册，输出为兼容性检查或失败点。启发意义不替代每对过渡映射的证明。

\section{5 外部适配、局部视图和偏序}
\(E_A\) 确认拓扑语言而多元分析未知；\(E_B\) 确认参数曲面和微分计算而抽象拓扑未知。M 保持相同。A 的呈现从 \(m_\mathrm{chart}\) 进入并标记分析入口待补；B 从公共曲面例子进入，诊断同胚/覆盖语言后汇合。个人“看起来像流形”的直觉只触发任务，不认证 TopMan。

LC10 从 SmoothAtlas 展示 \(C^k\) 的泛化锥；LC11 展示特定正则性子类；LC17 只看图册构造入口；LC22 选择具体曲面或抽象图册表征。被隐藏的 Hausdorff、第二可数和过渡域仍在边界中。

事件 \(e_\mathrm{chart}<e_\mathrm{trans}<e_\mathrm{reg}\) 是一条 \(C^k\) 路线；光滑路线为 \texttt{e\mcsbrk\_chart<e\mcsbrk\_trans<e\mcsbrk\_smooth<e\mcsbrk\_forget}。\(e_\mathrm{example}\) 可在图册和一元分析资源可用后插入，与不依赖它的其他呈现事件不可比。预算只能筛掉完整路线，不能删去 \(e_\mathrm{trans}\) 后继续宣称兼容性已证。

自检：\mcsnum{1}写出一对图的过渡域；\mcsnum{2}指出 h 的反例否定哪一命题；\mcsnum{3}解释最大图册与同一底空间的区别。三个条目分别锚定定义、计算、结构辨析能力。

验收：载体和 k 明确，泛化方向正确，\(C^{1}\) 反例不误伤光滑化存在问题，两名学习者不同入口只改变派生呈现。

\section{6 本轮机器覆盖（2026-09-29）}
同图过渡片段展开 \(\lambda u.\varphi(\varphi^{-1}(u))\)，在指定点的右逆等式假设下认证其值为该点。证书保留该开放假设及两个排序，尚未表达真实图的限制域、拓扑、正则性或最大图册。完整泛化桥梁仍未机器认证；精确条件、反例与行动引用见覆盖表。

\clearpage\chapter{案例 03 张量、坐标变换与张量场}
\mcsorigin{原章节：03-张量与张量场.md}
\section{1 公共背景与类型}
固定域 F、有限维向量空间 V，dim V=n，次数 \(r,s\in\mathbb{N}\)（允许 0，零重张量积约定为 F）。对偶 V* 是线性函数 \(V\to F\) 的空间。采用


\mcsdisplay{T^r_s(V)=V^{\otimes r}\otimes(V^*)^{\otimes s}.}
\noindent 它自然作用于 r 个协向量与 s 个向量。约定必须随节点保存，不把上下标习惯当作已知。流形部分另取光滑实流形 X。

\begingroup\small\setlength{\tabcolsep}{4pt}
\begin{longtable}{>{\raggedright\arraybackslash}p{\dimexpr 0.2464\linewidth-2\tabcolsep\relax}>{\raggedright\arraybackslash}p{\dimexpr 0.3992\linewidth-2\tabcolsep\relax}>{\raggedright\arraybackslash}p{\dimexpr 0.3544\linewidth-2\tabcolsep\relax}}
\toprule
\textbf{节点} & \textbf{内容} & \textbf{选定概念支持} \\
\midrule
\endfirsthead
\toprule
\textbf{节点} & \textbf{内容} & \textbf{选定概念支持} \\
\midrule
\endhead
\bottomrule
\endfoot
\bottomrule
\endlastfoot
Dual & V*=Hom(V,F) & 线性映射、域 \\
TensorProduct & 双线性映射的通用对象 & 向量空间、双线性、通用性质 \\
\(\mathrm{Tensor}_\mathrm{rs}\) & 上述张量空间 & 张量积、对偶、次数 \\
MultiLinear & \((V*)^r\times V^s\to F\) 的多线性函数 & 对偶、函数、多线性 \\
BasisLaw & 各指标的换基规律 & 基、逆矩阵、分量 \\
TensorBundle & 各点 \(T^r_s(T_\mathrm{xX})\) 的光滑丛 & 切丛、张量、局部平凡化 \\
TensorField & TensorBundle 的光滑截面 & 丛、截面、光滑 \\
\(\mathrm{NonTensor}\Gamma\) & Christoffel 系数的非张量反例 & 坐标、联络、Leibniz 规则 \\
\end{longtable}
\endgroup

有限维条件不可删除后继续沿用全部自然同构。也不在未给度量时默认 V 与 V* 有自然同构。

\section{2 一个贯通两种表示的完整证明}
考虑 (1,1) 型。定义


\mcsdisplay{\Phi:V\otimes V^*\longrightarrow End(V),\qquad
\Phi(v\otimes\alpha)(w)=\alpha(w)v.}
\noindent \((v,\alpha)\mapsto(w\mapsto\alpha(w)v)\) 对两变量分别线性，因此由张量积通用性质唯一诱导线性 \(\Phi\)。取基 \(e_i\) 与对偶基 \(\varepsilon^j\)，则 \(\Phi(e_i\otimes\varepsilon^j)\) 把 \(e_k\) 映为 \(\delta^j_k\) \(e_i\)，正是矩阵单位 \(E_\mathrm{ij}\)。这些张量构成 \(V\otimes V*\) 的基，矩阵单位构成 End(V) 的基，故 \(\Phi\) 为同构。基只用于证明双射；\(\Phi\) 本身的定义未选择基，因此该同构不依赖所选坐标。\(\blacksquare\)

有限维假设在满射性处实质起作用。无限维时，代数张量仍是有限项和，故 \(\Phi\) 的每个输出都为有限秩算子；无限维 V 的恒等映射为无限秩，不在像中。因此不能把上式不加限制地推广为无限维 End(V) 的表示。

令新基 \(e'_j=\Sigma_i\) \(e_i\) \(A^i_j\)。向量坐标满足 \([v]'=A^{-1}[v]\)；若 L 的矩阵为 [L]，则


\mcsdisplay{[L]'=A^{-1}[L]A.}
\noindent 证明：把 v 的旧坐标 Av' 代入 L，再把输出换成新坐标 \(A^{-1}\mathrm{LAv}'\)，对任意 v' 成立即得矩阵等式。该规律是同一线性映射在两基下的表达，不是产生一个新的 L。对一般 \(T^r_s\)，每个上指标乘 \(A^{-1}\)、每个下指标乘 A；由纯张量逐因子展开并线性延拓得到。

这连接“多线性/通用性质入口”和“线性映射/坐标入口”。数学证明已写全关键步骤，尚未编码运行 HOL 检查器。

\section{3 从张量到张量场}
在光滑流形 X 上，纤维为 \(T^r_s(T_\mathrm{xX})\)。每个坐标图提供切空间基并给出局部张量分量；相邻图的分量按相应 Jacobian 换基规律变换。

\textbf{命题。} 在固定覆盖 X 的图册上，满足每个重叠域张量变换规律的光滑局部分量族，与 X 上的光滑张量场一一对应。若图族不覆盖 X，只能得到其并集上的截面，不能声称已构造全局张量场。

\textbf{证明。} 从截面取局部平凡化分量，平滑性和过渡规律由丛定义得到。反之，对每点选一个包含它的坐标图，用该图分量确定纤维元素；过渡规律保证换图后是同一元素，故良定义。每个图中分量光滑，故得到光滑截面。两种构造逐点互逆。\(\blacksquare\)

前提是采用正确的张量丛过渡函数；任意带指标函数族不满足该前提。第 14 章以 \(y=x^{2}\) 计算 \(\Gamma^y_\mathrm{yy}=-1/(2y)\)，而原坐标 \(\Gamma=0\)，反驳“所有指标数组都是张量”。

\section{4 公开行动与结构路线}
\begingroup\small\setlength{\tabcolsep}{4pt}
\begin{longtable}{>{\raggedright\arraybackslash}p{\dimexpr 0.1468\linewidth-2\tabcolsep\relax}>{\raggedright\arraybackslash}p{\dimexpr 0.3408\linewidth-2\tabcolsep\relax}>{\raggedright\arraybackslash}p{\dimexpr 0.2524\linewidth-2\tabcolsep\relax}>{\raggedright\arraybackslash}p{\dimexpr 0.2600\linewidth-2\tabcolsep\relax}}
\toprule
\textbf{行动} & \textbf{输入} & \textbf{输出} & \textbf{见证} \\
\midrule
\endfirsthead
\toprule
\textbf{行动} & \textbf{输入} & \textbf{输出} & \textbf{见证} \\
\midrule
\endhead
\bottomrule
\endfoot
\bottomrule
\endlastfoot
\(t_\mathrm{dual}\) & 线性映射 & Dual & 对偶定义 \\
\(t_\mathrm{univ}\) & 双线性、Dual & TensorProduct、\(\mathrm{Tensor}_\mathrm{rs}\) & 通用构造及次数声明 \\
\(t_\mathrm{matrix}\) & 基、线性映射 & End 表征、BasisLaw & 换基计算 \\
\(t_\mathrm{bridge}\) & TensorProduct、Dual、End 表征 & \(\Phi\) 同构 & 上节基到基证明 \\
\(t_\mathrm{bundle}\) & 光滑流形、切丛、\(\mathrm{Tensor}_\mathrm{rs}\) & TensorBundle & 过渡函数构造 \\
\(t_\mathrm{field}\) & TensorBundle、截面 & TensorField & 定义与局部分量等价 \\
\(t_\mathrm{guard}\) & 坐标变换、联络 & \(\mathrm{NonTensor}\Gamma\) & 明确反例 \\
\end{longtable}
\endgroup

给定不同公开背景，可以先构造抽象张量再计算分量，也可从线性映射矩阵切入后补通用性质，再汇合于 \(t_\mathrm{bridge}\)。矩阵路线不宣称凭一张数组即可定义全部张量。

方法节点：“寻找不变量”对应换基前后的同一映射；“区分对象与表征”对应张量与分量；局部方法“用变换律排除非张量候选”输出反例或尚未认证的候选。若所有局部变换检查无限，不能声称任意输入都有终止算法。

\section{5 两个外部状态与事件偏序}
\(E_A\) 确认抽象线性代数与双线性，坐标操作熟练度未知；\(E_B\) 确认矩阵计算，通用性质未知。两者共享同一 \(\Phi\)、同一 \(\mathrm{Tensor}_\mathrm{rs}\) 与相同反例。A 先 \(t_\mathrm{univ}\) 再 \(t_\mathrm{matrix}\)；B 先 \(t_\mathrm{matrix}\) 再 \(t_\mathrm{univ}\)。两种入口都必须在 \(t_\mathrm{bridge}\) 前满足其联合输入。

取事件 \(e_\mathrm{dual}\)、\(e_\mathrm{univ}\)、\(e_\mathrm{matrix}\)、\(e_\mathrm{bridge}\)、\(e_\mathrm{bundle}\)、\(e_\mathrm{field}\)。依赖包括 \texttt{e\mcsbrk\_dual<e\mcsbrk\_univ<e\mcsbrk\_bridge}、\texttt{e\mcsbrk\_matrix<e\mcsbrk\_bridge}、\texttt{e\mcsbrk\_univ<e\mcsbrk\_bundle<e\mcsbrk\_field}；若切丛背景未确认，额外加入入口条件而不改数学定义。外部复习可新增 \(e'_\mathrm{matrix}\)，但不创建一个“此人版张量节点”。

LC08 比较两种证明支持的交汇；LC12 展示对偶相关结构；LC16 追踪换基不变量；LC22 选择分量或无坐标表征；LC24 标注通用性质的掌握证据缺口（不判为未掌握）；LC36 另显示公共证明或支持记录中的未决项。

自检条目：\mcsnum{1}算一个 \(2\times 2\) 矩阵的 \(A^{-1}\mathrm{LA}\)；\mcsnum{2}解释为什么 \(\Phi\) 是同构而不仅线性；\mcsnum{3}指出给定 \(\Gamma\) 反例中的非齐次项来源。它们锚定计算、证明、迁移辨析三种能力，一次答对不能覆盖全部。

验收：有限维与指标约定明确，数学对象不依赖坐标或学习者，局部分量汇合有证明，个人适配只影响呈现和排程。

\section{6 本轮机器覆盖（2026-09-29）}
秩一算子可加性证书认证 \(R(w)=\alpha(w)v\) 保持加法，理论仅登记所用标量分配律，开放假设明确为\(\alpha\)可加。该片段是\(\Phi\)构造中的实际代数步骤；标量线性、张量积通用性质、有限维双射及丛/截面仍未编码。覆盖范围与逐条件反例见覆盖表。

\clearpage\chapter{案例 04 群：几何、置换与模乘法的汇合}
\mcsorigin{原章节：04-群概念的多来源.md}
\section{1 同一数学目标与三个入口}
公共群节点 Group 的对象为 \((G,\cdot,e,\mathrm{inv})\)，满足结合律、左右单位及左右逆。使用此带运算签名可以直接表达公理；非空由常元 e 保证。交换律不属于群的定义。

\begingroup\small\setlength{\tabcolsep}{4pt}
\begin{longtable}{>{\raggedright\arraybackslash}p{\dimexpr 0.1653\linewidth-2\tabcolsep\relax}>{\raggedright\arraybackslash}p{\dimexpr 0.4174\linewidth-2\tabcolsep\relax}>{\raggedright\arraybackslash}p{\dimexpr 0.4174\linewidth-2\tabcolsep\relax}}
\toprule
\textbf{入口节点} & \textbf{对象与证明} & \textbf{与抽象群的联系} \\
\midrule
\endfirsthead
\toprule
\textbf{入口节点} & \textbf{对象与证明} & \textbf{与抽象群的联系} \\
\midrule
\endhead
\bottomrule
\endfoot
\bottomrule
\endlastfoot
Perm(X) & X 上双射，运算为复合，单位为恒等，逆为逆映射 & 复合结合且双射封闭，直接满足群公理 \\
TriangleSym & 平面中保持一个给定等边三角形的刚性对称 & 复合/逆仍保持图形，属于变换群 \\
Units5 & \{1,2,3,4\} 的模 5 乘法 & 2 的幂依次为 1,2,4,3，形成循环四元群 \\
Group & 上述公理的抽象结构 & 三个入口分别给出例子/抽象路线 \\
Cayley & 每个群嵌入其底集的置换群 & 严格汇合桥梁，不只是相似度 \\
\end{longtable}
\endgroup

Units5 中以 \(2^i\cdot 2^j=2^\{i+j\) mod4\} 验证闭合与结合，单位为 1，逆由指数取负给出。三角形对称可通过其对三个顶点的作用作置换表示；有三个旋转和三个反射，作用忠实。局部几何图像提供动机，不替代群公理检查。

\section{2 Cayley 桥梁的完整证明}
对群 G 的每个 g，定义 \(L_g:G\to G\)，\(L_g(x)=\mathrm{gx}\)。\(L_{g^{-1}}\) 是其逆，故 \(L_g\) 是双射。并且对所有 x，


\mcsdisplay{(L_g\circ L_h)(x)=g(hx)=(gh)x=L_{gh}(x).}
\noindent 故 \(g\mapsto L_g\) 是群同态。若 \(L_g=L_h\)，在单位元 e 处取值得 g=h，故同态单射。像是 Perm(G) 的一个子群，因此每个群同构于某个置换群的子群。\(\blacksquare\)

结论不是“每个群都是有限对称群”，也不是“每个群与某个完整的 \(S_n\) 同构”。当 G 无限时，其底集置换群也相应无限。

非交换反例取 \(S_{3}\)：按右侧先作用约定，(12)(23) 与 (23)(12) 对 1 的像分别为 2 与 3，故不相等。Units5 的交换性因此不能推广到所有群。

\section{3 概念支持与公开行动}
Group 的选定定义支持为 \{集合,二元运算,结合律,单位,逆\}；Perm 的支持为 \{映射,复合,双射,恒等,逆映射\}；Cayley 的证明支持是群公理、左乘映射及单射。三个支持族的交和差提供路线比较，\(C_\mathrm{all}\) 不承担这个功能。

\begingroup\small\setlength{\tabcolsep}{4pt}
\begin{longtable}{>{\raggedright\arraybackslash}p{\dimexpr 0.1682\linewidth-2\tabcolsep\relax}>{\raggedright\arraybackslash}p{\dimexpr 0.3054\linewidth-2\tabcolsep\relax}>{\raggedright\arraybackslash}p{\dimexpr 0.2616\linewidth-2\tabcolsep\relax}>{\raggedright\arraybackslash}p{\dimexpr 0.2648\linewidth-2\tabcolsep\relax}}
\toprule
\textbf{行动} & \textbf{输入 I} & \textbf{输出 O} & \textbf{见证类型} \\
\midrule
\endfirsthead
\toprule
\textbf{行动} & \textbf{输入 I} & \textbf{输出 O} & \textbf{见证类型} \\
\midrule
\endhead
\bottomrule
\endfoot
\bottomrule
\endlastfoot
\(g_\mathrm{perm}\) & 映射、复合、双射 & Perm 例子 & 公理逐项检查 \\
\(g_\mathrm{geom}\) & 刚性变换、三角形 & TriangleSym 例子 & 对称闭合及顶点作用 \\
\(g_\mathrm{mod}\) & 同余、模乘法 & Units5 例子 & 幂表及群律 \\
\(g_{\mathrm{abs}_p}\) & Perm 例子、抽象规则 & Group 呈现 & 从例子提炼同一公理模板 \\
\(g_{\mathrm{abs}_g}\) & TriangleSym 例子、抽象规则 & Group 呈现 & 相同公理模板 \\
\(g_{\mathrm{abs}_m}\) & Units5 例子、抽象规则 & Group 呈现 & 明确不保留偶然交换性 \\
\(g_\mathrm{cayley}\) & Group、映射、双射 & Cayley & 上节证明 \\
\(g_\mathrm{noncomm}\) & Perm 三元素例子 & NoncommExample & 两个复合的显式计算 \\
\end{longtable}
\endgroup

三个 \(g_\mathrm{abs}\) 行动认证的是一个共同定义的呈现与例子对应，不认证“看过一个例子就学会所有群”。要证明 Cayley 仍需其声明输入。

公共方法节点包括“从特例切入”“抽象共同运算”“寻找不变量”和“构造反例”。以 Units5 的偶然交换性为失败边界，以 \(S_{3}\) 为反例；抽象时既记录保留条件，也记录丢弃条件。

\section{4 外部模型、局部化与路线}
\(E_A\) 确认几何对称及图形操作，符号复合待诊断；\(E_B\) 确认同余计算，图形操作未知。A 获得 \(g_\mathrm{geom}\to g_{\mathrm{abs}_g}\)，B 获得 \(g_\mathrm{mod}\to g_{\mathrm{abs}_m}\)。若 \(E_C\) 确认函数复合，则可选 \(g_\mathrm{perm}\to g_{\mathrm{abs}_p}\)。公共 M 的三条路线始终同时存在。

从同一目标 Group 回溯，LC13 保留三种 OR 入口；LC25 的共同核心主要是抽象规则与目标接口，不因此声称零背景可学；LC26 显示几何、模算术和映射的专用支持；LC27 标明三路输出是同一 Group 引用。Cayley 另构成严格桥接，并不使三个入口的认知状态相等。

对固定某一路线，事件 \texttt{e\mcsbrk\_entry<e\mcsbrk\_abs<e\mcsbrk\_cayley}；\(e_\mathrm{cayley}\) 的映射/双射输入若不在背景，就由额外事件先提供。NoncommExample 可以在相应置换资源可用时执行，用来修正 \(E_B\) 的外部交换性误推断。该误区不是 Group 节点的属性。

历史来源是 Stillwell 第 19 章 §19.1（印刷 291–292 / PDF 307–308），讨论几何、数论和方程论的多来源。此处路线为现代重构，不宣称复现历史顺序，也不把历史先后设为硬先修。

\section{5 模板与验收}
Group 概念模板填写定义、典型例、非例、变体、动机、支持、自检和路线；Cayley 使用定理与证明模板；\(S_{3}\) 使用反例模板；“从特例切入”使用全局方法模板。开放问题式任务可问“给定有限运算表是否为群”，检查有限表的群公理可判；这不意味着任意无限结构群性质都可判。

自检：\mcsnum{1}在 Units5 中求 3 的逆；\mcsnum{2}写出 \(L_g\) 为双射的逆；\mcsnum{3}指出 Cayley 单射证明为何使用 e；\mcsnum{4}解释为什么模乘法例子的交换性不是群公理。

所有个人确认、误区、成本和作答留在 E。数学证明在本案例给出，完整Cayley的机器 Check 尚未完成；结构模型核查仅检验声明 I/O 路线，不替代上述数学论证或教学效果验证。

\section{6 本轮机器覆盖（2026-09-29）}
Cayley单射步骤证书已从 \(L_g=L_h\) 和右单位律推导 \(g=h\)，包含有类型\(\lambda\)项、提升、\(\beta\)公理实例、等式替换及对象等式回接。这里右单位已足够，不把结合律、逆元或完整Cayley登记成新公理。左乘双射、同态和像子群仍为后续义务；覆盖表给出片段证明、条件反例及公开行动接口。

\clearpage\chapter{14 反例与边界审计}
\mcsorigin{原章节：14-反例集.md}
本章是理论的回归规范。每个反例写明所否定的推断；不以一个实例否定条件更强的正确命题。第 15 章记录哪些属于一般证明、有限验证或引用事实。

\section{14.1 全等价概念集饱和}
\textbf{反例 14.1。} 在有闭表示 \(t_c\) 的语言中，e 与 \((\lambda z.e)t_c\) \(\beta\) 等价，后者却直接出现任意选定概念 c。第 05 章定理 5.6 因而使 \(C_\mathrm{all}(e)\) 等于全部登记概念，而不是某个有区分力的“真实依赖集合”。

\textbf{否定范围：} 否定“搜遍所有等价表达后，就精确得到不可缺少的概念”。不否定该函数的数学定义，也不否定指定表达的出现集。修复为保留 \(C_\mathrm{all}\) 的退化结论并使用有用途与见证的支持族。

\textbf{反例 14.2（类型错误的闭包）。} \(C_\mathrm{all}\) 的输入是节点，输出是概念集合；将输出再作为节点输入构成类型错误。即使另定义集合提升，也必须重新检验广延、单调和幂等，不能由“取了一个并集”直接得出。

\section{14.2 逻辑等价不宜作为节点身份}
\textbf{反例 14.3。} 若 \(T\vdash\varphi\) 且 \(T\vdash\psi\)，则命题逻辑给出 \(T\vdash\varphi\leftrightarrow\psi\)。于是对可证明等价取商，会把同一 T 的\textbf{所有已证命题}放入真命题等价类，而不只是把某定理的不同写法合并。

两个证明项的相等属于另一类型，不能从命题等价推出。定义节点身份时保留表达、证据、角色与来源，再另设内容等价关系。语义商可以用于特定分析视图，不得充当全部知识记录的唯一身份。

\section{14.3 概念集合与逻辑内容不是同一信息}
\textbf{反例 14.4。} 二元素域上的 \(P(x)=(x=0)\land(1=1)\) 与 \(Q(x)=(x=1)\land(0=0)\) 使用相同基本登记概念 \{等号,0,1\}，外延却互不包含。全等价概念集更会相同。

\textbf{否定范围：} 否定“概念集相同/包含就能推出概念等价、泛化或演绎”。修复为附加蕴含或结构映射证据，见 §06。

\section{14.4 筛选不是初等子模型}
\textbf{反例 14.5。} 语言仅含一元谓词 P。结构 A 的域为 \{a,b\}，\(P^A=\{b\}\)。取诱导子结构 B，其域为 \{a\}，\(P^B=\emptyset\)。没有函数闭合问题，但


\mcsdisplay{A\models\exists x\,P(x),\qquad B\not\models\exists x\,P(x).}
\noindent 所以即使合法子结构也未必是初等子结构；一般显示筛选更不能自动获得这种性质。若仅是带外部依赖的视图，它仍可引用环境中的 b，但不能把“环境满足”写成“筛选域自身满足”。

\section{14.5 AND、OR 与空路线族}
\textbf{反例 14.6（联合不能拆开）。} 在经典命题逻辑中 \(\{p,p\to q\}\vdash q\)；仅由 \{p\} 或 \(\{p\to q\}\) 均不能推出 q。前者的反模型取 p 真、q 假；后者取 p 假、q 假。把联合超边拆成两个独立充分边改变了语义。

\textbf{反例 14.7（替代不能合并）。} 公开契约为 \(\{a\}\to\{c\}\) 与 \(\{b\}\to\{c\}\)。目标 c 可以由 a 或 b 分别支持，不要求同时具备两者；两条路线的共同先修为空，但不存在空输入的契约。故并集当必修集和交集当充分集都错误。

\textbf{反例 14.8（空族真空化）。} 若 \(\mathcal{R}=\emptyset\)，公式 \(\forall r\in\mathcal{R}\) \(\mathrm{HardPre}_r(m,c)\) 对每个 m 都真。这不表示所有数学知识都是不可达目标 c 的先修。第 06 章将共同先修限制在非空路线族，空族另报告其不可行或尚未发现状态。

“未发现路线”与“证明所有路线均不存在”亦不同；无界搜索超时只能报告未知。

\section{14.6 固定节点偏序不能一般地表达替代路线}
\textbf{反例 14.9。} 取知识状态族


\mcsdisplay{\mathcal K=\{\varnothing,\{a\},\{b\},\{a,b\},
\{a,c\},\{b,c\},\{a,b,c\}\}.}
\noindent 它表达 c 可经 a 或 b 到达。任何偏序的全部下闭集对交封闭，但


\mcsdisplay{\{a,c\}\cap\{b,c\}=\{c\}\notin\mathcal K.}
\noindent 所以没有一个节点偏序恰以 \(\mathcal{K}\) 为全部下闭集。它不否定每条已经选定路线存在事件偏序，也不否定状态集自身按包含关系构成偏序。知识空间模型属于外部认知模型的一个选择，不能作为本体唯一公理。

\section{14.7 循环、复习与有限搜索}
\textbf{反例 14.10（循环压缩不产生入口）。} 只有契约 \(\{a\}\to\{b\}\) 和 \(\{b\}\to\{a\}\)，背景为空。若存在非空可执行事件序列，其第一个事件就需要背景中并不存在的 a 或 b，矛盾。把两个节点压成一块不会提供缺少的输入。

若另有合法契约 \(\emptyset\to\{a\}\)，则可以选择该事件再执行 \(\{a\}\to\{b\}\)；正确修复是生成一份新的无环事件见证，而不是原循环自动变成偏序。加入不相关入口 \(\emptyset\to\{d\}\) 不改变不可行性。

\textbf{反例 14.11（重复节点不等于重复事件）。} 背景含 a，事件 \(e_{1}\) 为学习 a 的一个应用，\(e_{2}\) 为复习 a，且 \(e_{1}<e_{2}\)。映射 \(\mathrm{ev}(e_{1})=\mathrm{ev}(e_{2})=a\) 完全合法。若在节点层保留该前后关系则出现 a<a；因此学习顺序放在事件上，不能据此改变节点身份。

\textbf{反例 14.12（有限节点不保证有限搜索）。} 只有一个节点 a、一种可重复行动 \(a\to a\)，也存在长度为 1,2,3,… 的无限多个事件序列。有限节点及有限行动并不足以使枚举全部计划终止。须另限事件数或规划时域，并使检查条件可判。

\section{14.8 证据使用不等于全局必要}
\textbf{反例 14.13。} 有证明 p 以允许假设集 \{A,B\} 得到 A，而证明实际仅用 A。若由 Ded(\{A,B\},A,p) 把 B 标成硬必需，则错误。即使一个证明使用 B，另一证明也可能绕过 B。第 06 章分别保存允许假设、实际叶假设、指定路线契约要求和路线族共同要求。

\textbf{反例 14.14（私有筛选增加的必经点）。} 原路线族为经 a 或经 b 达到 c；某外部模型筛掉 b 路线后，a 成为个性化子族的共同前提。但 a 仍不是原族共同前提。对路线子族的全称判断不能回写为原本体关系。

\section{14.9 未知既不是空，也不是否定}
\textbf{反例 14.15。} 某方法 m 的支持尚未锚定时为 Unknown，不是 \(\mathrm{Known}(\emptyset)\)。将它当空支持会生成“无概念输入即可使用”的错误判断。

对一个非空集合 A，\(\mathrm{Unknown}\cap\mathrm{Known}(A)\) 可能是 \(\mathrm{Known}(\emptyset)\)，也可能是 Known(A)，不能固定等于 Known(A)。安全信息语义及区间界见 §05。

同理，外部状态没有关于节点 n 的观察，不推出“未掌握 n”，也不推出“已经掌握 n”。冷启动只能采用显式的未知政策与公开路线入口。

\section{14.10 泛化的逆不是函数；一般解释也不必是泛化}
\textbf{反例 14.16。} 群有阿贝尔群、有限群等不同特化。\(\mathbb{Z}\) 与 \(S_{3}\) 分别见证两个子类互不包含。特化是集合值关系。有限阿贝尔群是它们的共同特化，故不能用本例声称“没有共同特化”。

\textbf{边界 14.17。} 把一种数学结构编码进另一种理论可以是合法解释，但编码往往改变载体及对象表示，不因此得到同载体子类包含。严格结构泛化须保留指定的遗忘/解释模式及其证书。

同样，\(C^k\) 流形与光滑流形比较必须说明采用共同的图册编码、非极大图册规格，或从光滑结构到 \(C^k\) 结构的遗忘映射。不能把不同意义的极大图册当成同一集合直接比较。

\section{14.11 认知状态不能反射为数学真理}
\textbf{反例 14.18。} 外部记录 \(\mathrm{Believe}(s,\varphi)\) 和 \(\mathrm{Believe}(s,\neg\varphi)\) 描述一个可能矛盾的信念状态。若加上无条件反射 \(\mathrm{Believe}(s,\psi)\Rightarrow\psi\)，就能在对象理论推出 \(\varphi\) 与 \(\neg\varphi\)，再由经典爆炸推出任意结论。

正确接口把信念作为外部数据；只有独立检查通过的证明才能认证数学命题。个体答对、信心高或推荐效果好均不是这种证明。

\textbf{反例 14.19（新符号不自动保证保守扩张）。} 核心语言有谓词 P。扩展引入 Known，并加上


\mcsdisplay{\exists x\,Known(x),\qquad
\forall x(Known(x)\to P(x)).}
\noindent 扩展推出核心句 \(\exists\mathrm{xP}(x)\)，而取 P 为空的核心模型即可验证原核心不能推出它（前提是原理论允许该模型且演算可靠）。虽然 Known 是新符号，桥接公理仍限制了核心模型。证明\textbf{每个合法核心模型都能扩张}是模型保守性的充分判据；结合相应语义的可靠性与完备性，可由此推出旧语言的证明论保守性。它不是一般证明论保守性的必要条件，也不能不加区分地移用于不完备的标准高阶语义。另一条有效路线是直接证明旧语言推导可消去新增设施，不能只说新增部分是外部模块。

\section{14.12 数学保守性与认知效果不能等同}
\textbf{边界 14.20。} 新增一个可消去的定义缩写，可以保持旧语言全部可证明后果，却改变表达长度和可读性。反过来，数学上等价的两个表征也不必对某人同样容易。此区分见 Manin II §8（印刷 66 / PDF 80）。MCS 保存共同形式内容与表征关系，个体成本与收益留在 E 中。

观察“先学 a 的人更容易学会 b”也可能反映背景差异；不是对学习 a 的因果效应证明。因果声明需另给干预定义、假设或研究证据。

\section{14.13 有指标的量不必是张量}
\textbf{反例 14.21。} 实直线正半轴使用坐标 x，取平坦联络 \(\Gamma^x_\mathrm{xx}=0\)；改用 \(y=x^{2}\)。由 \(\nabla\) 的 Leibniz 规则，


\mcsdisplay{\Gamma^y_{yy}=\frac{dy}{dx}\frac{d^2x}{dy^2}
=2\sqrt y\left(-\frac1{4y^{3/2}}\right)
=-\frac1{2y}\ne0.}
\noindent 若 \(\Gamma\) 是一个张量，其在 x 坐标下为零就应在所有坐标下为零，因此 Christoffel 系数不是张量分量。此例只否定“有指标且有变换公式就一定是张量”，不否定联络本身有严格几何定义。张量与张量场的正例见 cases/03。

\section{14.14 生成、模型与验证边界}
\textbf{反例 14.22（可生成不等于真）。} 有效语法可以生成 \(\varphi\) 和 \(\neg\varphi\) 的合法候选。类型正确只说明表达有意义；是否可证由证书检查确定。自动生成不是全体数学真理判定。

\textbf{边界 14.23。} 有限阶、有限元构造规则作用于可数种子所产生的有限语法在 \(\omega\) 步并中封闭，不需借用可构造宇宙 L 的超穷层级、凝结或可定义良序。可定义也不蕴含可计算；L 中的定理不能因为“都是分层生成”就移植到 MCS。

\textbf{边界 14.24。} 对有限行动库的枚举核验不是 HOL 全语义有效性的决定程序。有限网络实例也不是完整类型语义域有限的证明。第 15 章把结构实例验证、手写一般证明与引用元定理分开报告。

\textbf{边界 14.25（有限候选不使阈值判断自动可判）。} 即使只有一条路线，且其成功概率 p 是可计算实数，也不自动有决定 \(p\ge\theta\) 的算法。近似区间可以持续夹住 \(\theta\)，不能把尚未分离阈值当作失败；完整搜索只在每个所需比较均可判或已有有效分离/等号证书时推出有界可行或不可行结论。一个已通过的候选足以证明有解，但若所有候选均无通过证据且仍有待决项，只能报告 Unknown。

这不是纯粹的精度实现问题。对程序 e，若它在第 \(t\ge 1\) 步停机，令 \(p_e=1/2-2^{-t-1}\)；若不停机，令 \(p_e=1/2\)。模拟前 n 步，若尚未停机便输出 1/2，其误差小于 \(2^{-n}\)，故 \(p_e\) 可统一计算。但 \(p_e\ge 1/2\) 当且仅当 e 不停机；若存在对这些输入总会终止的精确阈值判断，就能决定停机问题。因此第 12 章必须把“比较可判”独立列为条件。

\clearpage\chapter{15 验证、完成状态与证明义务}
\mcsorigin{原章节：15-验证与完成标准.md}
本章把“已有何种证据”与“还不能声称什么”并列记录。定义写出不等于已有实现；一般数学证明不等于完成机器检查；合成学习者的运行不等于认知模型实证有效。

\section{15.1 证据状态}
\begingroup\small\setlength{\tabcolsep}{4pt}
\begin{longtable}{>{\raggedright\arraybackslash}p{\dimexpr 0.1935\linewidth-2\tabcolsep\relax}>{\raggedright\arraybackslash}p{\dimexpr 0.3584\linewidth-2\tabcolsep\relax}>{\raggedright\arraybackslash}p{\dimexpr 0.4480\linewidth-2\tabcolsep\relax}}
\toprule
\textbf{标签} & \textbf{含义} & \textbf{可支持的结论} \\
\midrule
\endfirsthead
\toprule
\textbf{标签} & \textbf{含义} & \textbf{可支持的结论} \\
\midrule
\endhead
\bottomrule
\endfoot
\bottomrule
\endlastfoot
DEF & 类型、定义或明确建模约定已给出 & 可检视语义、前提及接口，不自动证明模型存在 \\
PROOF & 正文给出数学证明及条件 & 相对于所列元理论、引用定理及假设的结论 \\
REF & 直接采用经典来源结论 & 仅限附录 B 定位与说明的适用范围 \\
FINITE & 脚本对明确有限域完成核验 & 该有限实例/枚举域内成立，不替代一般证明 \\
ILLUSTRATION & 显式规定的案例或外部参数 & 展示接口可用，不推断群体规律 \\
NOT-CLAIMED & 本次没有完成或不作该承诺 & 完整案例认证、检查器机器元证明、实际教学收益、无限最优规划等 \\
\end{longtable}
\endgroup

这不是一个可信度数值排序。例如经验测量可能与数学定义一样重要，但不能用它代替某条证明的公理叶。

\section{15.2 形式基础检查}
\begingroup\small\setlength{\tabcolsep}{4pt}
\begin{longtable}{>{\raggedright\arraybackslash}p{\dimexpr 0.2327\linewidth-2\tabcolsep\relax}>{\raggedright\arraybackslash}p{\dimexpr 0.3836\linewidth-2\tabcolsep\relax}>{\raggedright\arraybackslash}p{\dimexpr 0.3836\linewidth-2\tabcolsep\relax}}
\toprule
\textbf{检查} & \textbf{证据} & \textbf{条件与范围} \\
\midrule
\endfirsthead
\toprule
\textbf{检查} & \textbf{证据} & \textbf{条件与范围} \\
\midrule
\endhead
\bottomrule
\endfoot
\bottomrule
\endlastfoot
元理论与计数边界 & 定义 1.9、引用事实 1.10、反例 1.11、边界声明 1.12、命题 1.13，REF/PROOF & 定位见附录 B 的 R02-b（自指与真理不可定义）、R02-c（L 的生成背景）、R02-d（可计算描述与程序视角）；L 只作限界比较，不作算法承诺 \\
类型、常元、变元、连接词、各型量词与高阶关系量化 & 第 02 章定义 2.1–2.7，DEF & 基本类型有效，箭头型闭合；不是只对关系名称编码后仍只量化个体 \\
替换、类型与语法判定 & 命题 2.8–2.9，PROOF & 有限有类型语法；\(\alpha\) 正规编码 \\
推理规则完整指定 & 定义 2.10–2.13，DEF & 采用 \(\lambda\) 提升、多排序 ND、经典反证、等式及本征变量条件 \\
给定证书检查终止 & 命题 2.14，PROOF；子集实现另记 FINITE & 实际开放假设显式有限；理论公理引用带可判见证，嵌入任意大假设族另需成员见证；已实现有限理论上的 ND 白名单子集，完整 HOL 检查器尚未实现 \\
Henkin 解释与编码对应 & 定义 3.6，DEF；命题 3.5、3.7、3.8，PROOF & 两值、外延及所有赋值下抽象闭合；\(\ell_a\) 是函数符号，其值可被对应排序量化 \\
可靠性/Henkin 完备性 & 定理 3.9，PROOF+REF & 一阶完备性是引用元定理；给出实际模型/翻译匹配 \\
标准语义边界 & 推论 3.10–3.11，REF+PROOF & 未宣称全函数标准有效性可有效穷尽 \\
相对一致性前提 & 约定 1.1，DEF & 分别写明推理所用元理论、被讨论的理论及相对一致性假设；不能把任意模型存在假设统一替换成 Con(ZFC) \\
定义扩张保守 & 命题 3.13、定义 3.14 及其论证，PROOF & 旧项定义，或函数在旧域中的实现；逐点唯一不能独自替代实现条件 \\
非空结构实例 & 例 3.17、命题 6.20，PROOF & 基础核的标准模型；相对于领域 T 可满足性的登记实例，不是 T 自证一致 \\
有效生成边界 & 命题 1.7，PROOF+REF；命题 1.8、定理 4.7、命题 4.9，PROOF & 有限可描述候选与有限见证；不枚举全部语义对象；生成最小性依赖规则的单调性；计数事实定位见 R07-b \\
\end{longtable}
\endgroup

\section{15.3 概念、关系、层级、视图与路径}
\begingroup\small\setlength{\tabcolsep}{4pt}
\begin{longtable}{>{\raggedright\arraybackslash}p{\dimexpr 0.2273\linewidth-2\tabcolsep\relax}>{\raggedright\arraybackslash}p{\dimexpr 0.3863\linewidth-2\tabcolsep\relax}>{\raggedright\arraybackslash}p{\dimexpr 0.3863\linewidth-2\tabcolsep\relax}}
\toprule
\textbf{检查} & \textbf{证据} & \textbf{防止的误读} \\
\midrule
\endfirsthead
\toprule
\textbf{检查} & \textbf{证据} & \textbf{防止的误读} \\
\midrule
\endhead
\bottomrule
\endfoot
\bottomrule
\endlastfoot
全等价概念集饱和 & 定理 5.6、推论 5.7，PROOF；V08 有限见证 & 相同概念并集不等于相同知识或依赖 \\
支持族与未知语义 & 定义 5.9–5.12，DEF；V07 & Unknown 不等于已知空集；证明代码先回译为 HOL 表达再扫描；部分界用独立接口保存 \\
极小支持存在与非唯一 & 命题 5.15、反例，PROOF & 区分包含极小、基数最小及无极小情形 \\
联合前提与演绎组合 & 定义 6.2、命题 6.3–6.4，PROOF & 不把多个前提拆成各自充分的二元边 \\
先修无循环定义 & \(6.5\to 6.6\to 6.8\to 6.9-6.10\to\)第 12 章 & 公共契约及路线先于先修，先修先于个性化规划 \\
路线族必要性 & 命题 6.11–6.12，PROOF；V02 & 非空族、OR 替代、个体筛选的作用域 \\
硬泛化与逆向特化 & 定义 6.13–6.15，PROOF/DEF & 同载体蕴含与跨载体映射分模式；逆关系不是单值逆函数 \\
闭包与生成代数 & 定理 6.19、命题 8.8、8.12，PROOF；V01 & 形式闭包不等于掌握增长；绑定代数保留上下文 \\
关系及其关系 & 定义 8.1–8.6、命题 8.3，PROOF & 任意复合导航关系不自动仍为直接类比或因果关系 \\
模型/范畴接口 & 命题 8.19、8.21、8.23–8.24，PROOF；V14 & 非单射输送合并输入时须重选来源，所得依赖可能只是原图子图；满足、保守和粘合另有条件 \\
话题/领域/学科/理论/元理论 & 第 09 章，DEF+PROOF & 可重叠聚合，不混为同一类数学对象 \\
13 种角色模板 & 第 10 章，DEF & 共享题目、动机和失效范围；个人填写与诊断外置 \\
36 种局部化逐项定义 & LC01–LC36，DEF；V11 数量/链接核验 & 每项有参数、选集/变换、边界及计算条件 \\
视图组合与保持 & 命题 11.6、11.8 及反例 11.7，PROOF & 纯限制交换不推广至任意聚合/闭包/适配组合 \\
聚合—展开与视图代数 & 命题 11.10、11.12，PROOF；V12 & 固定块族的伴随；有类型部分函数的复合；不许丢失边界 \\
事件偏序及多线性扩张 & 命题 6.7、定理 12.6、命题 12.10，PROOF；V03/V04 & 节点可重复；结构等价不等于认知效果相同 \\
有限算法终止与相对完备 & 定义 12.11、算法 12.12、定理 12.13，PROOF；有限实例 V01/V05 & 全部候选维度有限且能完整枚举、事件有界、约束可判；仅概率可计算不足以保证阈值比较可判；V01/V05 不是完整规划器认证 \\
Pareto 与未知 & 定义 12.9，DEF；V10 & 未知效果不以默认数值支配；保留不可比较候选 \\
路径表达的代数 & 定义 12.15、命题 12.16，PROOF；V13 & 顺序复合同时替换输入与目标端口；并行目标不相交；外部剪枝另算 \\
\end{longtable}
\endgroup

\section{15.4 外置版的专门验收}
\begingroup\small\setlength{\tabcolsep}{4pt}
\begin{longtable}{>{\raggedright\arraybackslash}p{\dimexpr 0.1792\linewidth-2\tabcolsep\relax}>{\raggedright\arraybackslash}p{\dimexpr 0.4104\linewidth-2\tabcolsep\relax}>{\raggedright\arraybackslash}p{\dimexpr 0.4104\linewidth-2\tabcolsep\relax}}
\toprule
\textbf{要求} & \textbf{形式依据} & \textbf{案例/运行检查} \\
\midrule
\endfirsthead
\toprule
\textbf{要求} & \textbf{形式依据} & \textbf{案例/运行检查} \\
\midrule
\endhead
\bottomrule
\endfoot
\bottomrule
\endlastfoot
本体非干扰 & 第 04 章身份/负载无 s；命题 7.16；命题 11.8 & V09 公共结构哈希及两人派生结果不同 \\
状态更新非干扰 & 7.8、7.15 的 Trans 与只读接口 & V09 模拟学习、遗忘、答错、纠正、重校准后核心不变 \\
保守接入 & 命题 7.17、8.21；逐模型可扩张的充分判据 & 第 14 章相信反映原则的反例；V07 只作命题层见证 \\
显式输入 & 定义 7.2、7.8、11.4、12.1；命题 7.19 & V10 逐项检查核心、模型、状态、目标、预算、策略与随机种子对复现键的作用 \\
随机模型 & 7.8 的概率核和 12.8 的概率判据 & DEF；本次不为真实效果拟合概率，不声称经验校准 \\
冷启动 & Unknown 与 Conditional 的不同层次 & V09 保留未知掌握，展示三种群入口条件 \\
多对多观测锚定 & 定义 7.9 & 四案例题目区分计算、定义、证明及迁移能力 \\
相对认知原子性 & 定义 7.11–7.12、7.22；命题 7.23 & 极小存在/唯一条件分别给出；真实原子性为模型判断 \\
多学习者案例 & 四篇案例中的显式外部状态 & 群入口、极限误区、张量路线差异；公共数学证明相同 \\
方法节点外置评价 & 定义 4.12、5.12、7.2–7.7 & 通用方法接口共享；个人成功率和适用性不写回硬关系 \\
接口与效果的识别边界 & 反例 16.1，PROOF & V15 为两个确定模型的有限见证；相同已有记录容许未试行动的相反预测 \\
\end{longtable}
\endgroup

非干扰定理的实质前提是类型及数据依赖的隔离；任意另写一个会修改 M 的程序当然不满足它。有限夹具通过不是所有未来实现的安全证明。

\section{15.5 原理论回归与负例}
第 14 章保留 25 个具体反例或边界实例，覆盖：概念并集饱和、全等价取商坍缩、概念包含不足、普通子结构不初等、联合输入、替代路线、空族、知识空间非单一节点偏序、无入口循环、复习事件、有限节点无限搜索、冗余前提、个人筛选误回写、未知支持、泛化逆关系、一般解释、相信反映、新符号不保守、数学/认知保守分离、非张量数组、语法/真理、L/可计算性、有限测试/一般证明及可计算概率/可判阈值。

可执行部分见 核验报告 与 脚本。报告只在全部断言通过后生成；出现失败须修改理论实例或程序并重新核验，不允许只改报告文字。

\section{15.6 交付范围}
交付包括完整模块正文、四篇贯通案例、双页码文献证据表、定义/定理与记号索引、13 种模板、36 个局部化方案、有限模型核验输入与脚本。九份本地资料的相关正文均有对应证据条目；没有将查过目录扩写成读完全文，也没有把在线摘要标成已获得原文全文。

本次未交付通用 HOL 证明助理、实际学习者测评系统或经验上已经最优的路线。正文中的 Check 是严格指定的数学关系；现有 Python 实现只覆盖其列明子集，四案例各有局部数学证书，完整桥梁尚未编码完毕。站点现有结构规划功能也不因此变成数学证明助理。这些边界限制“已经计算验证全部数学及认知结论”的说法。

\section{15.7 “完整框架”六类对象的定义与验收入口}
第 01 章 1.1 所说的“完整框架”是本稿约定的章节覆盖：节点、关系、视图、认知接口、路径和宏观层级各有定义处与验收处。这不等于工程实现已经完备，也不表示数学真理、认知机制或未来文献已被穷尽。逐类对应如下。

\begingroup\small\setlength{\tabcolsep}{4pt}
\begin{longtable}{>{\raggedright\arraybackslash}p{\dimexpr 0.0905\linewidth-2\tabcolsep\relax}>{\raggedright\arraybackslash}p{\dimexpr 0.3286\linewidth-2\tabcolsep\relax}>{\raggedright\arraybackslash}p{\dimexpr 0.3286\linewidth-2\tabcolsep\relax}>{\raggedright\arraybackslash}p{\dimexpr 0.2523\linewidth-2\tabcolsep\relax}}
\toprule
\textbf{对象类} & \textbf{定义处} & \textbf{验收处} & \textbf{证据状态} \\
\midrule
\endfirsthead
\toprule
\textbf{对象类} & \textbf{定义处} & \textbf{验收处} & \textbf{证据状态} \\
\midrule
\endhead
\bottomrule
\endfoot
\bottomrule
\endlastfoot
节点 & 第 04 章定义 4.1–4.12 & 15.2“有效生成边界”；V01、V11 & DEF+PROOF+FINITE \\
关系 & 第 05–08 章（概念支持、硬关系、软关系、关系代数） & 15.3 支持族与未知语义、联合前提与演绎组合、硬泛化与逆向特化、闭包与生成代数、关系及其关系各行 & DEF+PROOF+FINITE \\
视图 & 第 11 章定义 11.1–11.12 与 LC01–LC36 & 15.3 视图组合与保持、聚合—展开与视图代数；V11、V12 & DEF+PROOF+FINITE \\
认知接口 & 第 07 章定义 7.1–7.23 & 15.4 外置版专门验收各要求；V09、V10 & DEF+PROOF+FINITE \\
路径 & 第 12 章定义 12.11、算法 12.12、定理 12.13 & 15.3“有限算法终止与相对完备”；V03、V04、V05 & PROOF+FINITE \\
宏观层级 & 第 09 章定义 9.1–9.10 & 15.3“话题/领域/学科/理论/元理论”一行 & DEF+PROOF \\
\end{longtable}
\endgroup

表中“证据状态”沿用 15.1 的标签，不构成新的可信度排序；某类对象在具体条目上只有 DEF 时，不得据本表推断它已有一般证明。

对整体贡献还须作三项独立判断：接口是否自洽、是否比简化表示增加可验证的能力、具体模型是否通过独立观测或干预检验。前两者与最后一项不能互相替代。第 16.7 节给出相同接口与观测记录允许相反行动预测的反例；本轮完成了明确片段上的表示比较，并未确立一般方法优势或真实学习效果。

\section{15.8 2026-09-29 最小实现的验收记录}
证据标签与检查状态分列：claims.json保存稳定ID、原文位置、假设、依赖、核对层次、证书/输出哈希、有限域及剩余义务；清单检查拒绝悬空、缺项、过期版本和无证书机器声明。检查状态不是标签的高低排序。

\begingroup\small\setlength{\tabcolsep}{4pt}
\begin{longtable}{>{\raggedright\arraybackslash}p{\dimexpr 0.1667\linewidth-2\tabcolsep\relax}>{\raggedright\arraybackslash}p{\dimexpr 0.4167\linewidth-2\tabcolsep\relax}>{\raggedright\arraybackslash}p{\dimexpr 0.4167\linewidth-2\tabcolsep\relax}}
\toprule
\textbf{实际产物} & \textbf{检查入口与记录} & \textbf{精确能力与边界} \\
\midrule
\endfirsthead
\toprule
\textbf{实际产物} & \textbf{检查入口与记录} & \textbf{精确能力与边界} \\
\midrule
\endhead
\bottomrule
\endfoot
\bottomrule
\endlastfoot
最小 ND 内核 & 设计、报告 & 类型、绑定、提升、公理重建及规则白名单；负例与受控缺陷被检出；无一般机器元证明 \\
四案例/公共行动 & 覆盖表、行动报告 & 群单射步骤、常值取值、同图过渡片段、秩一可加性；理论/假设/类型/版本/结论匹配；完整桥梁未认证 \\
独立有限语义 & 参考实现 & 35项/175次比较、42案例赋值、11背景实例；不是一般模型完备性证明 \\
合成研究流程 & 研究方案、报告 & 固定输入、两种隔离划分、每表示2295候选；增强行动图与MCS在此片段相同；无真实效果 \\
修订与重放 & 交付入口、进度 & 保存本轮前基线、命令/退出码/日志和哈希；TeX/PDF未同步，仍为旧导出 \\
\end{longtable}
\endgroup

“文中有证明”“某证书被某版检查器接受”“检查器已形式验证”是三个独立字段。当前只有前两者在各自列明范围内有记录；第三者未完成。

\clearpage\chapter{16 边界与进一步研究}
\mcsorigin{原章节：16-开放问题.md}
核心定义、接口与本稿证明范围已在正文明确。以下是扩大理论及验证强度的后续问题，不用它们替代当前章节必要的定义，也不把尚未验证的结论标作已完成。

\section{16.1 机器形式化}
第 02 章指定有限语法、认证翻译和 ND 规则；第 03 章证明模型对应。下一步可实现或接入可信证明内核，编码类型替换、抽象提升、源变量条件及实际数学案例。需独立检查认证翻译的实现和内核，而不能用文档编号测试代替证明认证。

2026-09-29 已新增独立 ND 子集检查器、四案例局部证书和独立有限语义/负例验证，原结构回归另行保留。下一步的缺口是完整规则与数学库、实现保真的机器元证明以及完整案例桥梁。未经翻译证明仍不能把现有 HOL 系统或 Lean 的证书直接混用；Python 模块不具有 LCF 抽象定理类型的封装保证。

\section{16.2 概念支持的可用性}
全等价概念并集的饱和说明它不能独自承担先修发现。值得研究的是受限变换族 \(\Theta\)、任务相关支持及极小支持的计算：哪些限制同时保留数学解释力、可计算性和稳定性？增加符号别名如何改变 \(\kappa\)，如何记录版本迁移？

极小支持不一定唯一；无限充分族可能没有极小成员。后续算法必须声明自己计算的是有界近似、一个包含极小者还是基数最优者，不让 UI 隐去差别。

\section{16.3 认知模型的识别与比较}
Atom、Eval、Obs、Trans 是明确的外部模型接口，不是已被实验确认的心理定律。应研究观察锚定的测量效度、误区判断的可识别性、不同模型是否能给出相同可观测预测，以及干预效应能否从设计中识别。

外部模型可用知识空间、概率状态或其他表示。比较模型时，应固定 M 及任务、保留未知和误差界。即使预测等价，也不推出内部心理状态相同。真实数据采集、个体隐私管理和教学实验不属于本次理论专稿已经完成的工作。

当前合成流程与真实研究方案已给出可运行的模型、参与者/时间隔离入口、留出任务预测与失败判据。这完成研究准备，不填补真实测量效度、因果收益或一般可识别性的缺口。

\section{16.4 高效规划与无限情形}
本稿给出有限行动、事件界及可判约束下的完整搜索证明。扩展方向包括因子化 AND/OR 搜索、约束求解、有限记忆近似、启发式 Pareto 前沿与可验证剪枝。

允许无限复习、连续时间、不可计算概率或无限证明搜索后，原终止定理不适用。即使概率可计算，与阈值的比较也须另外可判；没有分离阈值的误差界时应保留 Unknown。研究应分别给出新的紧致性、良基性、误差或收敛条件，不能只用“节点有限”重复保证终止。

\section{16.5 理论翻译与结构粘合}
第 08 章只证明特定 institution 的满足条件、指定态射的路线输送和兼容标准模型粘合。更一般的 Henkin 模型粘合、商解释、不同逻辑间的翻译，需要检查混合语言抽象闭合、旧模型可扩张及证据输送。

结构同构、理论保守性、可计算性、证明成本和认知效用是不同保持目标。未来研究可以找出同时保持其中若干项的态射类，但不存在由“范畴化”三个字自动得到的全部保持定理。

\section{16.6 层级与内容扩展}
话题允许重叠且由明确结构关联参数构造，领域/学科由有型展开相连。如何选取有长期研究价值的参数、如何在不同数学传统间比较分类，需要公共修订和具体案例，不能被任意一次聚类输出定义为唯一自然分类。

本稿的四组案例可扩展到测度、代数几何、概率、逻辑与基础等，但扩展必须提供实际数学内容及桥梁证据。自动生成覆盖的是有限可描述候选；其研究重要性、教学适切性及所有语义真理均不由枚举本身保证。

\section{16.7 对本稿解释力的自我批判}
\textbf{为什么需要再退一步。} 从数学内部看，我们需要分清一条关系是定义出来的、由证明得到的，还是仅由记录估计的；否则“已经写成公式”容易被当成“已经说明原因”。从学习应用看，我们希望同一知识允许不同入口，却还需要回答怎样判断某条入口确实适合某个人。两条动机相遇，才要求把共享知识、个人模型与派生路线分别记录。

可以把本稿的设计生长链读成：知识条目及连线 \(\to\) 带联合输入、替代路线和来源的知识结构 \(\to\) 外置个人状态的 M/E/D 接口 \(\to\) 对特定模型作独立检验。前三步主要解决表示与检查问题，第四步才承担教学预测的检验责任。这是本稿的逻辑构造动机，不冒称经过了同样次序的历史或心理实验。

\textbf{接口覆盖不等于解释完成。} 把学习者状态外置解决了公共知识与个体记录混用的问题，却未决定状态应含哪些变量、观察如何测量理解、行动怎样改变能力。因此，本稿首先是知识组织和模型接入规范；其中的数学定理说明指定结构具有哪些性质，不足以把整套规范确认为学习机制理论。非干扰在固定公共结构和只读依赖下成立，不能单独说明某个软件遵守隔离，更不能说明隔离提高学习收益。

\textbf{反例 16.1（相同记录不能唯一决定未试行动的效果，PROOF）。} 证明思路是让两个模型在已试行动上完全一致，只在未试行动上分叉；若二者都符合接口，接口就不能替我们选定其中一个。固定同一个公共本体 M，其中有两个呈现行动 a、b，二者公共输入均为空、输出均为目标节点 g；此公共输出只表示提供了 g 的内容，不表示学习成功。令两个外部模型均有状态集 \(X=\{s_{0},s_{1}\}\)、初态 \(s_{0}\)，观察函数 \(\mathrm{Obs}(s_{i})=i\)，把观察值 1 约定为本例中的成功指标。对每个 \(s\in X\)，定义确定转移


\mcsdisplay{Trans_{E_1}(s,a)=Trans_{E_2}(s,a)=s_1,\qquad
Trans_{E_1}(s,b)=s_1,\quad Trans_{E_2}(s,b)=s_0.}
\noindent 二者都符合本稿的接口与本体非干扰约定。对任意仅由 a 组成的已收集行动序列，它们给出完全相同的观察记录；但对随后实施 b 的结果，一个给出 1，另一个给出 0。因此，从该记录及本稿接口不能确定 b 的效果，也不能据此确立两个行动等效。该例是确定模型，允许概率核不会自动消除这种不识别。\(\blacksquare\)

\textbf{需要增加什么。} 要使具体认知模型可被数据反驳，须事先限制可接受的模型族、观测含义与参数更新方式，并给出对未参与拟合的任务或行动的预测；不能看过所有结果后再自由改写 Eval 或 Trans。预测准确也不单独证明干预有效。若要比较路线的教学收益，应另行给出能识别行动效应的实验或明确识别假设，固定相同目标、初始条件、时间预算和测量标准，报告不确定性及失败情形。

\textbf{需要比较什么。} 四组案例展示了表达能力，但尚未证明必须使用本稿全部结构。后续可固定相同内容与外部模型，比较普通依赖图、带联合输入和替代方案的行动图、本稿完整表示；分别测量哪些合法路线被遗漏、哪些无证据路线被误接受、核验与维护成本如何变化。只有数据或构造反例支持的差异才能作为贡献。现有 AND/OR 反例说明单一节点偏序的一项限制，不等于证明所有图模型都不足。

因此，章节覆盖、一般数学证明、有限回归、机器内核认证和认知实证应分别报告完成状态。尚未完成的一般模型识别、更大范围方法比较与真实干预检验，是本稿解释力的当前缺口；不能仅把它们列作未来扩展后，就把当前工作称为已完成的认知理论。

\textbf{本轮有限比较的更新。} 比较记录固定四行动、三事件界、预算3及27种信息状态，分别核验2295候选。带类型行动图、增强二部图与MCS片段一致，因而没有据此得到MCS独有优势。普通依赖图的不足只针对报告中明确的编码惯例；更大表示比较、真实学习干预及识别仍未完成。

\textbf{保留错误路径。} 本轮复核推翻了三种看似顺畅的想法：“路线只要接好下一步的输入就能组合”，遗漏了最终目标也可能引用中间端口（12.15）；“每一步的概率都能近似计算，有界规划就能决定”，遗漏了阈值相等的判定问题（14.25）；“两个学习者共用理论档案，所以全部核心真值相同”，遗漏了档案可以有不同模型，真值保持还须固定核心解释（7.16、E.2）。错误的共同来源，是把某一层已有的性质未经检查地搬到另一层。

\textbf{思想与方法。} 将这个检验习惯称为“保持义务”：每次变换都明确保留什么、需要什么条件、由什么见证负责。低层做法是核对输入、输出、类型和量词；高层可与第 08 章的满足保持和模型约化联动，但不能只换上范畴术语就省掉这些核对。16.1 使用“同记录、异预测”构造反模型，12.15 用空事件单位检验复合，14.25 用停机归约检验可判性；相应有限见证见 V13、V15，一般不可判结论不由有限测试证明。

\textbf{回看问题。} 当一项新规则加入 M 或 E 时，它究竟排除了哪个原先允许的反例？它增加的是可检查的约束，还是只增加了一个名称？若没有排除任何错误预测，本稿暂不据此声称获得新的解释力。

\clearpage\chapter{E 核心公理的分型汇总与模型义务}
\mcsorigin{原章节：E-核心公理与验证条件.md}
本附录把各章分散的接口汇总为一套有类型的理论模式，便于后续实现或证明助理编码。它不另引入一个与第 02 章竞争的逻辑。核心仍是 ZFC 中明确指定的记录结构；其可内部化部分用所选 HOL 表达，语法代码及检查器的意义按指定解释固定。

\section{E.1 签名和构造层}
互异基本类型包括 Node、Kind、Code、TheoryCode、ProofCode、PremFam、Action、RouteCode、ViewCode、Evt、Template、Aggregate，以及角色标号域 Role。它们是引用/代码域。某类合法记录为空时，仍可使用非空代码域加 Valid 谓词，避免与第 03 章各排序非空的约定冲突。

核心非逻辑常元至少有：


\mcsdisplay{\begin{aligned}
&kind:Node\to Kind,\quad payload:Node\to Code,\quad theory:Node\to TheoryCode,\\
&WF:Node\to o,\quad Registered,External,Concept:Node\to o,\\
&ValidA:Action\to o,\quad In,Out:Action\to Node\to o,\\
&Check:TheoryCode\to PremFam\to Node\to ProofCode\to o,\\
&FiniteSetCode:Code\to o,\quad Ref:Code\to Node\to o.
\end{aligned}}
\noindent Code 和 Node 不互相默认强制转换。PremFam 解释为有限带标签命题节点列表的代码，事件/路线/视图也先检查代码，再通过有型投影得到有效数据。谓词本身是箭头型项，能够被抽象和量化；命名该谓词的 Node 是另一对象。

TypeOK、Check、WF 等符号在\textbf{合法结构类}中不是任意赋值：它们分别解释为第 02、04 章指定的有限语法关系或由该关系构成的接口。只有一个名叫 Check 的任意关系，不足以认证数学证明。将这些递归关系完全内部化需另给数码表示与表示正确性证明，不能默认有限 Python 检查已经完成这一步。

\section{E.2 核心公理模式}
\textbf{定义 E.1（核结构条件的逻辑形式）。} 下列模式与第 04 章定义 4.2 联用，未展开的谓词按对应章节的定义解释。

\begingroup\small\setlength{\tabcolsep}{4pt}
\begin{longtable}{>{\raggedright\arraybackslash}p{\dimexpr 0.5984\linewidth-2\tabcolsep\relax}>{\raggedright\arraybackslash}p{\dimexpr 0.4016\linewidth-2\tabcolsep\relax}}
\toprule
\textbf{公理/定义模式} & \textbf{含义及检查位置} \\
\midrule
\endfirsthead
\toprule
\textbf{公理/定义模式} & \textbf{含义及检查位置} \\
\midrule
\endhead
\bottomrule
\endfoot
\bottomrule
\endlastfoot
\(\forall n.\) \(\mathrm{Registered}(n)\Rightarrow\mathrm{WF}(n)\) & 每个登记节点具有合法类型和负载 \\
\(\forall n,m.\) \(\mathrm{Registered}(n)\land\mathrm{Ref}(\mathrm{payload}(n),m)\Rightarrow\mathrm{Registered}(m)\lor\mathrm{External}(m)\) & 负载引用有登记或显式环境来源 \\
\(\forall a,m.\) \(\mathrm{ValidA}(a)\land(\mathrm{In}(a,m)\lor\mathrm{Out}(a,m))\Rightarrow\mathrm{Registered}(m)\lor\mathrm{External}(m)\) & 行动端点无悬空引用 \\
\(\forall a.\) \(\mathrm{ValidA}(a)\Rightarrow\mathrm{FiniteIn}(a)\land\mathrm{FiniteOut}(a)\land\mathrm{WitnessOK}(a)\) & 有限联合接口和模式相符的结构见证 \\
\(\forall P,n,p.\) \(\mathrm{Ded}(P,n,p)\Leftrightarrow\mathrm{Check}(\mathrm{theory}(n),P,n,p)\) & 硬演绎的定义，背景及类型相容检查包含在 Check \\
\(\forall n,i.\) \(\mathrm{ValidSuppIndex}(n,i)\Rightarrow\mathrm{Supp}(n,i)\in\mathrm{ISet}(K_\Sigma)\) & 支持的值域与未知标签相容 \\
\(\forall t.\) \(\mathrm{ValidTemplate}(t)\Rightarrow\mathrm{PublicSchema}(t)\) & 模板字段仅描述共享内容；个体填写是外部实例 \\
\(\forall a.\) \(\mathrm{ValidAggregate}(a)\Rightarrow\mathrm{TypedMembers}(a)\land\mathrm{Traceable}(a)\) & 聚合具有层级类型及可展开引用 \\
\end{longtable}
\endgroup

身份功能性由 payload、kind 的函数类型保证；版本迁移另用版本参数化引用。WF 的构造定义保证不同记录类型形成规则不冲突。FiniteIn、WitnessOK、TypedMembers 等不是自然语言标签：它们分别展开为第 06 章有限集合及对应证书、第 09 章成员类型判断。对于启发性方法，WitnessOK 只认证公共接口，不认证启发体等价于算法。

Supp 的上述值域只容纳精确已知集合和 Unknown。部分信息另存于第 05 章的 \(\mathrm{Supp}^\mathrm{bounds}\)，值为 \(L\subseteq U\subseteq K_\Sigma\)；仅 L=U 时能投影为 Known(L)。不把一个真部分区间直接塞入 \(\mathrm{ISet}(K_\Sigma)\)。

不能给所有节点内容加入统一的无条件“都真”公理：Claim 节点允许未决甚至已反驳命题；相应证据状态单独记录。不同领域理论亦以 theory 标签隔离，不能把两个不相容背景的结论悄悄并成一个理论。

\section{E.3 路线、先修与关系量化}
对有效路线代码 r，由解码得到有限事件谓词 \(E_r\)、行动标记 \(\mathrm{act}_r\)、来源边 \(S_r\) 及背景 \(B_r\)。其有效性包括：


\mcsdisplay{\forall p,m\;
[E_r(p)\land In(act_r(p),m)]
\Rightarrow
[BSource_r(p,m)\lor
\exists q(E_r(q)\land EventSource_r(p,m,q)\land Out(act_r(q),m))].}
\noindent 来源指派是函数性选择，不能只保存“某个生产者可能存在”。\(S_r\) 的每条边来自该来源选择，且无环；源为背景时须另有 \(m\in B_r\)。目标也有指定来源。于是第 06 章得到 \(\preceq_r\) 的偏序性质。

对指定非空路线族 \(\mathcal{R}\) 与目标 g：


\mcsdisplay{HardPre_{\mathcal R}(m,g)
\iff
\forall r\,[r\in\mathcal R\Rightarrow m\in Req_r(g)].}
\noindent \(\mathcal{R}\in\mathcal{P}(\mathrm{Rte}_M(B,G))\) 在元层解释；在 Henkin 模型内部量化某个路线集型变量时，只量化该模型实际域内的集合表示，不能把它自动扩大为元层全幂集。相同提醒适用于全概念支持族与自动生成的最小闭包。

令 \(\mathrm{Rel}=\mathrm{Node}\to\mathrm{Node}\to o\)。关系间包含、逆及复合用真正高阶表达：


\mcsdisplay{Sub=\lambda R^{Rel}S^{Rel}.
\forall x^{Node}y^{Node}(Rxy\Rightarrow Sxy),}

\mcsdisplay{Inv=\lambda R^{Rel}x^{Node}y^{Node}.Ryx,\qquad
Comp=\lambda S^{Rel}R^{Rel}x^{Node}z^{Node}.
\exists y^{Node}(Rxy\land Syz).}
\noindent 它们分别具有 \(\mathrm{Rel}\to\mathrm{Rel}\to o\)、\(\mathrm{Rel}\to\mathrm{Rel}\)、\(\mathrm{Rel}\to\mathrm{Rel}\to\mathrm{Rel}\) 的类型。可量化 \(F:\mathrm{Rel}\to\mathrm{Rel}\) 来陈述保单调：


\mcsdisplay{\forall R,S:Rel.\ Sub(R,S)\Rightarrow Sub(FR,FS).}
\noindent 因此本体研究的不只是端点之间的边，也包括关系运算、关系包含与证据变换。对于具体硬/软关系，只有各自的定义或见证支持相应公理，不能因上述通用运算存在就宣布每种关系传递。

\section{E.4 外部接口的相对模型类}
E 的状态类型 \(X_E\) 不要求属于上述核心签名。最小接入在元层构成有兼容引用的对 (M,E)，并保留投影


\mcsdisplay{\pi:\operatorname{ExtMCS}\to\operatorname{CoreMCS},
\qquad \pi(M,E)=M.}
\noindent 固定 M 后的所有外部模型、状态及派生结果处在 \(\pi\) 的同一纤维中。学习转移、遗忘和个体诊断只能改变纤维内的数据；公共修订 \(M\to M'\) 是另一种显式操作。把 E 的模型描述编码为核心节点，并不会把其全部运行状态加入核心载体。

\textbf{命题 E.2（固定核心解释下的不可区分性）。} 两个接入对象投影为同一个档案 M 时，核心记录相同。再固定 M 的一个核心解释 A，并要求两个语义接入均以 A 为约化，则任何只使用核心签名的项或公式在相同核心赋值下解释相同。

\textbf{证明。} 档案部分由投影及只读契约直接得到。语义部分的变量赋值取相同核心对象；A 固定核心常元、应用及量词域的解释，对项和公式逐层归纳即得。若内部化扩张改变旧箭头量词域，则以同一 A 为约化的前提失败。相同理论档案可有不同模型，不能只凭档案相同就断言任意未决句子在那些模型中真值相同。\(\blacksquare\)

该非干扰性质本身不能推出\textbf{理论保守性}。一个充分判据是：核心理论的每个模型都有兼容扩张，并在指定语义中保持全部旧解释；第 07 章给出该条件下的证明。这种逐模型可扩张性比句子层的保守性更强，不能未经证明把它称为所有保守扩张的必要条件。

\section{E.5 派生结构及合法模型的四种检查}
\begingroup\small\setlength{\tabcolsep}{4pt}
\begin{longtable}{>{\raggedright\arraybackslash}p{\dimexpr 0.1573\linewidth-2\tabcolsep\relax}>{\raggedright\arraybackslash}p{\dimexpr 0.4213\linewidth-2\tabcolsep\relax}>{\raggedright\arraybackslash}p{\dimexpr 0.4213\linewidth-2\tabcolsep\relax}}
\toprule
\textbf{结构} & \textbf{数学定义层} & \textbf{核验义务} \\
\midrule
\endfirsthead
\toprule
\textbf{结构} & \textbf{数学定义层} & \textbf{核验义务} \\
\midrule
\endhead
\bottomrule
\endfoot
\bottomrule
\endlastfoot
节点/本体 & 有类型记录、形成规则及上述核心条件 & 良构、引用、证据含义、版本及类型一致 \\
关系及关系间运算 & 认证纤维、\(\lambda\) 特征函数、复合与包含 & 见证保持、量词域、方向及背景相容 \\
局部/显示结构 & Loc、Cut、聚合伴随与部分函数复合 & 隐藏依赖、溯源、逐项 Pres，不误称初等子模型 \\
学习计划 & 公共路线、外部约束、来源 DAG 和事件偏序 & AND/OR 选择、排程可行性、预算、不确定性、界内搜索 \\
\end{longtable}
\endgroup

本稿的“统一”是共享有类型引用、证据与保持义务，不是把所有东西都塞进一个无差别二元图。通用网络是这些结构的关系呈现；具体软件若要自称此理论的模型，必须按声明的片段履行相应义务。

\clearpage\chapter{A 定义、定理与反例索引}
\mcsorigin{原章节：A-定义与定理索引.md}
由 \texttt{validation/verify\mcsbrk\_finite\mcsbrk\_models.py} 按正文编号生成。索引列出定义及声明位置，不把“被索引”视为已经证明或机器认证；证据状态见第 15 章。

\begingroup\small\setlength{\tabcolsep}{4pt}
\begin{longtable}{l>{\raggedright\arraybackslash}p{\dimexpr 0.0836\linewidth-2\tabcolsep\relax}>{\raggedright\arraybackslash}p{\dimexpr 0.3618\linewidth-2\tabcolsep\relax}>{\raggedright\arraybackslash}p{\dimexpr 0.4705\linewidth-2\tabcolsep\relax}}
\toprule
\textbf{编号} & \textbf{种类} & \textbf{标题/内容} & \textbf{文档} \\
\midrule
\endfirsthead
\toprule
\textbf{编号} & \textbf{种类} & \textbf{标题/内容} & \textbf{文档} \\
\midrule
\endhead
\bottomrule
\endfoot
\bottomrule
\endlastfoot
1.1 & 约定 & 元理论，DEF & 01-元理论与可构造宇宙，第 11 行 \\
1.2 & 定义 & 两种完备性与判定性，DEF & 01-元理论与可构造宇宙，第 15 行 \\
1.3 & 定义 & 四类角色，DEF & 01-元理论与可构造宇宙，第 31 行 \\
1.4 & 定义 & 大小约定，DEF & 01-元理论与可构造宇宙，第 46 行 \\
1.5 & 命题 & 外置的基本非干扰，PROOF & 01-元理论与可构造宇宙，第 50 行 \\
1.6 & 约定 & 有效签名，DEF & 01-元理论与可构造宇宙，第 71 行 \\
1.7 & 命题 & 有限描述的计数边界，PROOF+REF & 01-元理论与可构造宇宙，第 75 行 \\
1.8 & 命题 & 有限前提生成在 \(\omega\) 达到闭包，PROOF & 01-元理论与可构造宇宙，第 89 行 \\
1.9 & 定义 & 可构造层级，DEF：引用背景 & 01-元理论与可构造宇宙，第 116 行 \\
1.10 & 引用事实 & 内模型背景，REF & 01-元理论与可构造宇宙，第 124 行 \\
1.11 & 反例 & 可定义不等于可计算，REF & 01-元理论与可构造宇宙，第 126 行 \\
1.12 & 边界声明 & 本稿的采用范围，NOT-CLAIMED & 01-元理论与可构造宇宙，第 128 行 \\
1.13 & 命题 & 真理与证明，PROOF+REF & 01-元理论与可构造宇宙，第 150 行 \\
2.1 & 定义 & 简单类型、签名与字母表 & 02-对象语言：简单类型论，第 9 行 \\
2.2 & 定义 & 三种判断与一项元层检查 & 02-对象语言：简单类型论，第 19 行 \\
2.3 & 定义 & 项形成 & 02-对象语言：简单类型论，第 23 行 \\
2.4 & 定义 & 原始逻辑常元 & 02-对象语言：简单类型论，第 34 行 \\
2.5 & 例 & 关系和关系之间的量化 & 02-对象语言：简单类型论，第 48 行 \\
2.6 & 定义 & 自由变量与 \(\alpha\) 等价 & 02-对象语言：简单类型论，第 64 行 \\
2.7 & 定义 & 无捕获替换 & 02-对象语言：简单类型论，第 74 行 \\
2.8 & 命题 & 替换保持类型 & 02-对象语言：简单类型论，第 83 行 \\
2.9 & 命题 & 语法有效性 & 02-对象语言：简单类型论，第 89 行 \\
2.10 & 定义 & 认证翻译 & 02-对象语言：简单类型论，第 102 行 \\
2.11 & 定义 & Boolean、外延与抽象公理 & 02-对象语言：简单类型论，第 128 行 \\
2.12 & 定义 & 多排序经典自然推演 ND & 02-对象语言：简单类型论，第 172 行 \\
2.13 & 定义 & 本稿的 HOL 可证性 & 02-对象语言：简单类型论，第 188 行 \\
2.14 & 命题 & 认证检查终止 & 02-对象语言：简单类型论，第 201 行 \\
2.15 & 命题 & 常用规则的认证 & 02-对象语言：简单类型论，第 228 行 \\
2.16 & 定义 & 显式可选公理 & 02-对象语言：简单类型论，第 241 行 \\
2.17 & 例 & \(\beta\) 与谓词节点 & 02-对象语言：简单类型论，第 256 行 \\
3.1 & 定义 & 函数框架 & 03-语义：Henkin与标准，第 7 行 \\
3.2 & 定义 & 抽象闭合 & 03-语义：Henkin与标准，第 15 行 \\
3.3 & 定义 & 逻辑解释与满足 & 03-语义：Henkin与标准，第 24 行 \\
3.4 & 定义 & 两类模型 & 03-语义：Henkin与标准，第 28 行 \\
3.5 & 命题 & 赋值局部性及语义替换 & 03-语义：Henkin与标准，第 36 行 \\
3.6 & 定义 & 认证语言 & 03-语义：Henkin与标准，第 54 行 \\
3.7 & 命题 & Henkin 到认证模型 & 03-语义：Henkin与标准，第 56 行 \\
3.8 & 命题 & 认证模型的外延化 & 03-语义：Henkin与标准，第 66 行 \\
3.9 & 定理 & 指定演算的可靠性与 Henkin 完备性 & 03-语义：Henkin与标准，第 80 行 \\
3.10 & 推论 & 语义的单向关系 & 03-语义：Henkin与标准，第 96 行 \\
3.11 & 推论 & Henkin 紧致性 & 03-语义：Henkin与标准，第 104 行 \\
3.12 & 定义 & 项定义 & 03-语义：Henkin与标准，第 112 行 \\
3.13 & 命题 & 项定义保守 & 03-语义：Henkin与标准，第 117 行 \\
3.14 & 定义 & 规格定义的额外义务 & 03-语义：Henkin与标准，第 123 行 \\
3.15 & 边界 & 边界 3.15（无独立标题；见正文） & 03-语义：Henkin与标准，第 141 行 \\
3.16 & 例 & 边界 3.15 的具体反模型 & 03-语义：Henkin与标准，第 143 行 \\
3.17 & 例 & 基础核的标准模型；存在性见证 & 03-语义：Henkin与标准，第 160 行 \\
4.1 & 定义 & 本体版本 & 04-节点生成与身份，第 7 行 \\
4.2 & 定义 & 合法本体条件 & 04-节点生成与身份，第 17 行 \\
4.3 & 定义 & 不可变节点记录 & 04-节点生成与身份，第 36 行 \\
4.4 & 定义 & 节点构造与形成条件 & 04-节点生成与身份，第 44 行 \\
4.5 & 命题 & 良构与真理分离 & 04-节点生成与身份，第 65 行 \\
4.6 & 定义 & 候选宇宙 & 04-节点生成与身份，第 73 行 \\
4.7 & 定理 & 最小生成 & 04-节点生成与身份，第 82 行 \\
4.8 & 定义 & 公平候选生成 & 04-节点生成与身份，第 88 行 \\
4.9 & 命题 & 枚举覆盖与边界 & 04-节点生成与身份，第 90 行 \\
4.10 & 定义 & 两种相等 & 04-节点生成与身份，第 98 行 \\
4.11 & 命题 & 学习者替换保持身份 & 04-节点生成与身份，第 104 行 \\
4.12 & 定义 & 方法接口 & 04-节点生成与身份，第 112 行 \\
5.1 & 定义 & 概念登记 & 05-概念函数与支持族，第 7 行 \\
5.2 & 定义 & 语法概念出现 & 05-概念函数与支持族，第 11 行 \\
5.3 & 命题 & 有限性与相对可计算性 & 05-概念函数与支持族，第 20 行 \\
5.4 & 定义 & 相对等价 & 05-概念函数与支持族，第 28 行 \\
5.5 & 定义 & 全等价形式概念并集 & 05-概念函数与支持族，第 37 行 \\
5.6 & 定理 & 全等价概念集饱和 & 05-概念函数与支持族，第 47 行 \\
5.7 & 推论 & 不能恢复依赖 & 05-概念函数与支持族，第 65 行 \\
5.8 & 定义 & 允许变换系统 & 05-概念函数与支持族，第 71 行 \\
5.9 & 定义 & 四类结构支持 & 05-概念函数与支持族，第 81 行 \\
5.10 & 定义 & 统一支持函数 & 05-概念函数与支持族，第 92 行 \\
5.11 & 定义 & 信息序与安全集合运算 & 05-概念函数与支持族，第 103 行 \\
5.12 & 定义 & 方法公开支持 & 05-概念函数与支持族，第 123 行 \\
5.13 & 定义 & 相对用途的充分性 & 05-概念函数与支持族，第 145 行 \\
5.14 & 定义 & 极小支持族 & 05-概念函数与支持族，第 154 行 \\
5.15 & 命题 & 存在性与计算条件 & 05-概念函数与支持族，第 164 行 \\
5.16 & 命题 & 概念信息不足以决定语义 & 05-概念函数与支持族，第 170 行 \\
6.1 & 定义 & 前提接口 & 06-硬关系，第 7 行 \\
6.2 & 定义 & 认证演绎 & 06-硬关系，第 9 行 \\
6.3 & 命题 & 相对语义可靠性 & 06-硬关系，第 26 行 \\
6.4 & 命题 & 弱化与组合 & 06-硬关系，第 32 行 \\
6.5 & 定义 & 公开行动契约 & 06-硬关系，第 40 行 \\
6.6 & 定义 & 有根结构路线 & 06-硬关系，第 64 行 \\
6.7 & 命题 & 路线偏序与结构可执行性 & 06-硬关系，第 73 行 \\
6.8 & 定义 & 目标活跃部分与资源支持 & 06-硬关系，第 79 行 \\
6.9 & 定义 & 契约相对的硬先修 & 06-硬关系，第 91 行 \\
6.10 & 定义 & 路线族共同必需 & 06-硬关系，第 108 行 \\
6.11 & 命题 & 族缩小增加共同约束 & 06-硬关系，第 118 行 \\
6.12 & 命题 & 替代先修不能压成联合先修 & 06-硬关系，第 129 行 \\
6.13 & 定义 & 同载体硬泛化 & 06-硬关系，第 137 行 \\
6.14 & 命题 & 语义包含、预序和商 & 06-硬关系，第 146 行 \\
6.15 & 定义 & 跨载体的结构泛化 & 06-硬关系，第 150 行 \\
6.16 & 例 & 逆关系非函数 & 06-硬关系，第 156 行 \\
6.17 & 定义 & 特征函数与集合值函数 & 06-硬关系，第 160 行 \\
6.18 & 定义 & 有限前提演绎闭包 & 06-硬关系，第 171 行 \\
6.19 & 定理 & 代数闭包 & 06-硬关系，第 178 行 \\
6.20 & 命题 & 一个非空核心实例 & 06-硬关系，第 188 行 \\
7.1 & 定义 & 通用关系描述 & 07-软关系与认知状态，第 7 行 \\
7.2 & 定义 & 外部评价 & 07-软关系与认知状态，第 9 行 \\
7.3 & 定义 & 软先修 & 07-软关系与认知状态，第 23 行 \\
7.4 & 定义 & 软泛化/特化 & 07-软关系与认知状态，第 30 行 \\
7.5 & 定义 & 应用 & 07-软关系与认知状态，第 32 行 \\
7.6 & 定义 & 类比 & 07-软关系与认知状态，第 34 行 \\
7.7 & 定义 & 对偶 & 07-软关系与认知状态，第 36 行 \\
7.8 & 定义 & 外部认知模型 & 07-软关系与认知状态，第 40 行 \\
7.9 & 定义 & 锚定 & 07-软关系与认知状态，第 47 行 \\
7.10 & 边界 & 非全知 & 07-软关系与认知状态，第 53 行 \\
7.11 & 定义 & 可认知与允许分解 & 07-软关系与认知状态，第 57 行 \\
7.12 & 定义 & 相对认知原子 & 07-软关系与认知状态，第 67 行 \\
7.13 & 定义 & 误区模式 & 07-软关系与认知状态，第 83 行 \\
7.14 & 定义 & 个体实例 & 07-软关系与认知状态，第 85 行 \\
7.15 & 契约 & 只读适配 & 07-软关系与认知状态，第 89 行 \\
7.16 & 命题 & 结构非干扰 & 07-软关系与认知状态，第 91 行 \\
7.17 & 命题 & 外部扩张的语义保守性 & 07-软关系与认知状态，第 95 行 \\
7.18 & 定义 & 未知值 & 07-软关系与认知状态，第 108 行 \\
7.19 & 命题 & 派生可复现性 & 07-软关系与认知状态，第 115 行 \\
7.20 & 定义 & 有限知识空间实例 & 07-软关系与认知状态，第 117 行 \\
7.21 & 验收 & 验收 7.21（无独立标题；见正文） & 07-软关系与认知状态，第 129 行 \\
7.22 & 定义 & 可行分解偏序 & 07-软关系与认知状态，第 132 行 \\
7.23 & 命题 & 有限存在与有条件唯一 & 07-软关系与认知状态，第 134 行 \\
8.1 & 定义 & 有类型的关系空间 & 08-关系代数与结构，第 7 行 \\
8.2 & 定义 & 关系运算 & 08-关系代数与结构，第 21 行 \\
8.3 & 命题 & 关系代数 & 08-关系代数与结构，第 32 行 \\
8.4 & 定义 & 两种集合提升 & 08-关系代数与结构，第 40 行 \\
8.5 & 定义 & 证据纤维 & 08-关系代数与结构，第 58 行 \\
8.6 & 定义 & 认证的关系包含 & 08-关系代数与结构，第 69 行 \\
8.7 & 定义 & 有限规则生成 & 08-关系代数与结构，第 79 行 \\
8.8 & 命题 & 代数闭包 & 08-关系代数与结构，第 90 行 \\
8.9 & 定义 & 依赖补全与行动可达 & 08-关系代数与结构，第 99 行 \\
8.10 & 反例 & 反例 8.10（无独立标题；见正文） & 08-关系代数与结构，第 111 行 \\
8.11 & 定义 & 上下文索引的表达代数 & 08-关系代数与结构，第 115 行 \\
8.12 & 命题 & 表达的递归泛性质 & 08-关系代数与结构，第 130 行 \\
8.13 & 定义 & 同余与关系饱和 & 08-关系代数与结构，第 138 行 \\
8.14 & 反例 & 定理坍缩 & 08-关系代数与结构，第 149 行 \\
8.15 & 定义 & 子结构与初等性 & 08-关系代数与结构，第 155 行 \\
8.16 & 命题 & 证明切片保持 & 08-关系代数与结构，第 159 行 \\
8.17 & 命题 & 初等链粘合 & 08-关系代数与结构，第 167 行 \\
8.18 & 定义 & 具体 institution & 08-关系代数与结构，第 181 行 \\
8.19 & 命题 & 满足条件 & 08-关系代数与结构，第 192 行 \\
8.20 & 定义 & 理论态射 & 08-关系代数与结构，第 201 行 \\
8.21 & 命题 & 模型扩张保证的保守性 & 08-关系代数与结构，第 213 行 \\
8.22 & 定义 & 认证空间态射 & 08-关系代数与结构，第 225 行 \\
8.23 & 命题 & 路线输送 & 08-关系代数与结构，第 235 行 \\
8.24 & 命题 & 兼容标准模型的粘合 & 08-关系代数与结构，第 246 行 \\
9.1 & 定义 & 公开关联规格 & 09-话题、领域、学科与元理论，第 7 行 \\
9.2 & 定义 & 话题 & 09-话题、领域、学科与元理论，第 18 行 \\
9.3 & 例 & 例 9.3（无独立标题；见正文） & 09-话题、领域、学科与元理论，第 22 行 \\
9.4 & 定义 & 多层聚合 & 09-话题、领域、学科与元理论，第 26 行 \\
9.5 & 命题 & 展平与覆盖 & 09-话题、领域、学科与元理论，第 38 行 \\
9.6 & 定义 & 关系聚合的两种读法 & 09-话题、领域、学科与元理论，第 46 行 \\
9.7 & 反例 & 反例 9.7（无独立标题；见正文） & 09-话题、领域、学科与元理论，第 60 行 \\
9.8 & 定义 & 数学理论的双重表示 & 09-话题、领域、学科与元理论，第 64 行 \\
9.9 & 定义 & 分层引用 & 09-话题、领域、学科与元理论，第 72 行 \\
9.10 & 命题 & 层级引用不改变下层解释 & 09-话题、领域、学科与元理论，第 76 行 \\
10.1 & 定义 & 定义 10.1（无独立标题；见正文） & 10-节点模板，第 7 行 \\
10.2 & 定义 & Concept & 10-节点模板，第 26 行 \\
10.3 & 定义 & Definition & 10-节点模板，第 38 行 \\
10.4 & 定义 & Axiom & 10-节点模板，第 50 行 \\
10.5 & 定义 & Theorem & 10-节点模板，第 62 行 \\
10.6 & 定义 & Proof & 10-节点模板，第 74 行 \\
10.7 & 定义 & Property & 10-节点模板，第 86 行 \\
10.8 & 定义 & Example & 10-节点模板，第 98 行 \\
10.9 & 定义 & Counterexample & 10-节点模板，第 110 行 \\
10.10 & 定义 & Problem & 10-节点模板，第 122 行 \\
10.11 & 定义 & Theory & 10-节点模板，第 134 行 \\
10.12 & 定义 & Construction & 10-节点模板，第 146 行 \\
10.13 & 定义 & GlobalMethod & 10-节点模板，第 158 行 \\
10.14 & 定义 & LocalMethod & 10-节点模板，第 170 行 \\
10.15 & 定义 & 定义 10.15（无独立标题；见正文） & 10-节点模板，第 182 行 \\
10.16 & 验收 & 验收 10.16（无独立标题；见正文） & 10-节点模板，第 184 行 \\
11.1 & 定义 & 带边界视图 & 11-局部化36方案，第 7 行 \\
11.2 & 定义 & 通用切片 Cut & 11-局部化36方案，第 15 行 \\
11.3 & 定义 & 共同记号 & 11-局部化36方案，第 25 行 \\
11.4 & 定义 & 适配器 & 11-局部化36方案，第 64 行 \\
11.5 & 定义 & 四类视图变换 & 11-局部化36方案，第 103 行 \\
11.6 & 命题 & 纯限制的交换 & 11-局部化36方案，第 105 行 \\
11.7 & 反例 & 一般算子不交换 & 11-局部化36方案，第 109 行 \\
11.8 & 命题 & 适配不改本体 & 11-局部化36方案，第 111 行 \\
11.9 & 定义 & 固定块族的展开与完整块选择 & 11-局部化36方案，第 119 行 \\
11.10 & 命题 & 聚合—展开伴随 & 11-局部化36方案，第 131 行 \\
11.11 & 定义 & 可组合的视图算子 & 11-局部化36方案，第 143 行 \\
11.12 & 命题 & 视图复合律 & 11-局部化36方案，第 145 行 \\
12.1 & 定义 & 规划输入 & 12-视图转换与路径规划，第 7 行 \\
12.2 & 定义 & 结果包与偏序投影 & 12-视图转换与路径规划，第 11 行 \\
12.3 & 定义 & 分析状态 & 12-视图转换与路径规划，第 28 行 \\
12.4 & 例 & 例 12.4（无独立标题；见正文） & 12-视图转换与路径规划，第 34 行 \\
12.5 & 定义 & 事件展开 & 12-视图转换与路径规划，第 40 行 \\
12.6 & 定理 & 结构排程可靠性 & 12-视图转换与路径规划，第 48 行 \\
12.7 & 命题 & 循环诊断 & 12-视图转换与路径规划，第 54 行 \\
12.8 & 定义 & 外部排程判据 & 12-视图转换与路径规划，第 62 行 \\
12.9 & 定义 & 四种比较 & 12-视图转换与路径规划，第 74 行 \\
12.10 & 命题 & 线性扩张的结构互换 & 12-视图转换与路径规划，第 83 行 \\
12.11 & 定义 & 有限搜索条件 & 12-视图转换与路径规划，第 91 行 \\
12.12 & 算法 & 穷尽而有界的参考算法 & 12-视图转换与路径规划，第 101 行 \\
12.13 & 定理 & 有界终止、可靠与相对完备 & 12-视图转换与路径规划，第 110 行 \\
12.14 & 命题 & 环境保持的规划 & 12-视图转换与路径规划，第 118 行 \\
12.15 & 定义 & 有端口的路线表达 & 12-视图转换与路径规划，第 128 行 \\
12.16 & 命题 & 路线运算的结构律 & 12-视图转换与路径规划，第 137 行 \\
14.1 & 反例 & 反例 14.1（无独立标题；见正文） & 14-反例集，第 7 行 \\
14.2 & 反例 & 类型错误的闭包 & 14-反例集，第 11 行 \\
14.3 & 反例 & 反例 14.3（无独立标题；见正文） & 14-反例集，第 15 行 \\
14.4 & 反例 & 反例 14.4（无独立标题；见正文） & 14-反例集，第 21 行 \\
14.5 & 反例 & 反例 14.5（无独立标题；见正文） & 14-反例集，第 27 行 \\
14.6 & 反例 & 联合不能拆开 & 14-反例集，第 37 行 \\
14.7 & 反例 & 替代不能合并 & 14-反例集，第 39 行 \\
14.8 & 反例 & 空族真空化 & 14-反例集，第 41 行 \\
14.9 & 反例 & 反例 14.9（无独立标题；见正文） & 14-反例集，第 47 行 \\
14.10 & 反例 & 循环压缩不产生入口 & 14-反例集，第 64 行 \\
14.11 & 反例 & 重复节点不等于重复事件 & 14-反例集，第 68 行 \\
14.12 & 反例 & 有限节点不保证有限搜索 & 14-反例集，第 70 行 \\
14.13 & 反例 & 反例 14.13（无独立标题；见正文） & 14-反例集，第 74 行 \\
14.14 & 反例 & 私有筛选增加的必经点 & 14-反例集，第 76 行 \\
14.15 & 反例 & 反例 14.15（无独立标题；见正文） & 14-反例集，第 80 行 \\
14.16 & 反例 & 反例 14.16（无独立标题；见正文） & 14-反例集，第 88 行 \\
14.17 & 边界 & 边界 14.17（无独立标题；见正文） & 14-反例集，第 90 行 \\
14.18 & 反例 & 反例 14.18（无独立标题；见正文） & 14-反例集，第 96 行 \\
14.19 & 反例 & 新符号不自动保证保守扩张 & 14-反例集，第 100 行 \\
14.20 & 边界 & 边界 14.20（无独立标题；见正文） & 14-反例集，第 111 行 \\
14.21 & 反例 & 反例 14.21（无独立标题；见正文） & 14-反例集，第 117 行 \\
14.22 & 反例 & 可生成不等于真 & 14-反例集，第 129 行 \\
14.23 & 边界 & 边界 14.23（无独立标题；见正文） & 14-反例集，第 131 行 \\
14.24 & 边界 & 边界 14.24（无独立标题；见正文） & 14-反例集，第 133 行 \\
14.25 & 边界 & 有限候选不使阈值判断自动可判 & 14-反例集，第 135 行 \\
16.1 & 反例 & 相同记录不能唯一决定未试行动的效果，PROOF & 16-开放问题，第 51 行 \\
E.1 & 定义 & 核结构条件的逻辑形式 & E-核心公理与验证条件，第 27 行 \\
E.2 & 命题 & 固定核心解释下的不可区分性 & E-核心公理与验证条件，第 101 行 \\
\end{longtable}
\endgroup

36 个局部化方案用 LC01–LC36 独立编号，见 第 11 章。四组案例的节点/行动标识是案例内局部标识，不与章节定理编号合并。

\clearpage\chapter{B 文献证据表与引用边界}
\mcsorigin{原章节：B-文献证据表.md}
核验对象为 E:\textbackslash{}MCS\textbackslash{}reference 中的九份资料。下列 PDF 页一律从 1 开始计数；印刷页取页面上实际页码，\textbf{不使用全书统一偏移公式}。表中“正文”表示读过对应段落；“原典定位”与“本稿证明”分开。扫描 OCR 只作检索辅助，符号疑点以图像复核。正文节选检索记录在 evidence 目录，不是认证证明。

\section{B.1 本地资料登记}
\begingroup\small\setlength{\tabcolsep}{4pt}
\begin{longtable}{l>{\raggedright\arraybackslash}p{\dimexpr 0.3894\linewidth-2\tabcolsep\relax}>{\raggedright\arraybackslash}p{\dimexpr 0.3894\linewidth-2\tabcolsep\relax}r}
\toprule
\textbf{编号} & \textbf{作者、简称} & \textbf{本地文件名} & \textbf{PDF 总页数} \\
\midrule
\endfirsthead
\toprule
\textbf{编号} & \textbf{作者、简称} & \textbf{本地文件名} & \textbf{PDF 总页数} \\
\midrule
\endhead
\bottomrule
\endfoot
\bottomrule
\endlastfoot
R01 & Ebbinghaus–Flum–Thomas，Mathematical Logic & \texttt{1994\mcsbrk\_Book\mcsbrk\_Mathem\mcsbrk{}atical\mcsbrk{}Logic.pdf} & 290 \\
R02 & Manin，A Course in Mathematical Logic for Mathematicians & GTM 053 - ISBN978-1-4419-0615-1 - Yu. I. Manin - A Course in Mathematical Logic for Mathematicians..pdf & 389 \\
R03 & Prestel–Delzell，Mathematical Logic and Model Theory & UTX - Mathematical Logic and Model Theory.pdf & 198 \\
R04 & 汪芳庭，《数理逻辑》 & 数理逻辑 (汪芳庭) (z-library.sk, 1lib.sk, z-lib.sk).pdf & 214 \\
R05 & 余俊伟、赵晓玉等，《数理逻辑》 & 数理逻辑 (Junwei Yu (余俊伟), Xiaoyu Zhao (赵晓玉) etc.) (z-library.sk, 1lib.sk, z-lib.sk).pdf & 382 \\
R06 & 郝兆宽、杨睿之、杨跃，《数理逻辑：证明及其限度》第二版 & 数理逻辑：证明及其限度（第二版） (郝兆宽, 杨睿之, 杨跃) (z-library.sk, 1lib.sk, z-lib.sk).pdf & 275 \\
R07 & 冯琦，《数理逻辑导引》 & 数理逻辑导引 (冯琦) (z-library.sk, 1lib.sk, z-lib.sk).pdf & 534 \\
R08 & 李未，《数理逻辑基本原理与形式演算》 & 数理逻辑基本原理与形式演算 (李未著) (z-library.sk, 1lib.sk, z-lib.sk).pdf & 270 \\
R09 & Stillwell，《数学及其历史》 & 数学及其历史 (（美）斯狄瓦著) (Z-Library).pdf & 474 \\
\end{longtable}
\endgroup

对应原文件均在上述 reference 目录。本表只声明下列相关正文的核对范围，不声称逐字审校了九本全书。

\section{B.2 逻辑基础、定义扩张与有效性的正文证据}
\begingroup\small\setlength{\tabcolsep}{4pt}
\begin{longtable}{l>{\raggedright\arraybackslash}p{\dimexpr 0.2349\linewidth-2\tabcolsep\relax}>{\raggedright\arraybackslash}p{\dimexpr 0.0955\linewidth-2\tabcolsep\relax}>{\raggedright\arraybackslash}p{\dimexpr 0.1026\linewidth-2\tabcolsep\relax}>{\raggedright\arraybackslash}p{\dimexpr 0.2349\linewidth-2\tabcolsep\relax}>{\raggedright\arraybackslash}p{\dimexpr 0.2349\linewidth-2\tabcolsep\relax}}
\toprule
\textbf{证据号} & \textbf{来源及定位} & \textbf{印刷页} & \textbf{PDF 页} & \textbf{正文所支持的内容} & \textbf{在本稿的采用与限制} \\
\midrule
\endfirsthead
\toprule
\textbf{证据号} & \textbf{来源及定位} & \textbf{印刷页} & \textbf{PDF 页} & \textbf{正文所支持的内容} & \textbf{在本稿的采用与限制} \\
\midrule
\endhead
\bottomrule
\endfoot
\bottomrule
\endlastfoot
R01-a & EFT，VIII §3，Definition 3.1、Theorem 3.2 & 125–128 & 131–134 & 定义扩张的保守性、定义符号消去；函数定义须存在唯一 & 第 03 章旧项定义和扩张条件。原书为一阶；HOL 函数域实现、新抽象闭合另证，不能直接照搬 \\
R01-b & EFT，IX §1，二阶逻辑 & 138–142 & 142–146 & 谓词/关系变量量化与二阶标准解释的表达能力 & 第 02–03 章高阶语言与语义差别；\(\lambda\) 类型结构由本稿另定 \\
R01-c & EFT，X §5，Trahtenbrot 定理与二阶不完全性 & 170–172 & 173–175 & 有限可满足性、有限有效性不可枚举及二阶有效式不可枚举 & 第 01、03、14 章：有限结构逐个检查不等于有效性可穷尽 \\
R02-a & Manin，II §8 Language Extensions，8.2–8.3 & 66–69 & 80–83 & 引入可定义名称并给出消去翻译；旧语言可证性保持 & 第 03、08 章保守性；第 10 章命名与表征可有认知价值，但不能由逻辑保守推知效果相同 \\
R02-b & Manin，II §9，§11 & 69；74–77 & 83；88–91 & SELF、自指与真理不可定义的展开 & 第 09 章对象/元语言分层；没有在同层加入无条件真理谓词 \\
R02-c & Manin，IV §1–3：可构造宇宙、可定义性与绝对性、L 作为集合论模型 & 151–160 & 164–173 & 可构造层级与八个集合运算；可定义性与绝对性；L 满足集合论公理的逐条核验（§3，印刷 159–160 / PDF 172–173） & 第 01 章定义 1.9 与引用事实 1.10 的定位；只作边界比较，不把 L 的序数索引当作有限节点算法，不追加 V=L。IV §4（GCH 在 L 中成立，印刷 161–166 / PDF 174–179）本稿明确不使用 \\
R02-d & Manin，IX §§1–3，计算理论中的 worlds/constructive worlds & 285–287；289–293 & 293–295；297–301 & 可计算描述、编号与程序视角 & 区分计算理论的 constructive world 与集合论 L；节点用有限代码及公平枚举 \\
R03-a & Prestel–Delzell，§1.2 & 8–13 & 18–23 & 对象语言/元语言、字母表、项和公式的递归形成 & 第 01–02 章层次和结构归纳；字母表、类型与签名见定义 2.1，项与公式的递归形成见定义 2.3–2.4；“字符串存在”不等于“命题成立” \\
R03-b & Prestel–Delzell，§1.5 & 36–46 & 46–56 & 满足、可靠性与一阶完备性的模型构造 & 第 03 章明确调用多排序一阶完备性；本稿给出与所选 HOL 编码的双向对应，不冒称逐行重证经典完备性 \\
R03-c & Prestel–Delzell，§1.5 末的 theory/deduction 区分 & 47 & 57 & 固定模型的理论与给定公理的后承区别 & 第 04、06 章背景理论、主张、认证证据分离 \\
R04-a & 汪芳庭，§2.1.1 & 61–63 & 69–71 & 从常元、变元、函数、谓词建项；项集作为生成代数；有限形成层次 & 第 02、04、08 章语法生成。书中这里是无绑定的一阶项；本稿 \(\lambda\) 抽象必须带上下文排序 \\
R04-b & 汪芳庭，§2.1.2 & 64 & 72 & 从原子公式、连接词和量词生成公式 & 第 02 章归纳形成。这里的逻辑“原子”不等于心理认知原子 \\
R05-a & 余俊伟等，§3.2.1，定义 3.2.1 及评注 & 71–72 & 82–83 & 逻辑/非逻辑符号、元数、尚未赋义的语法符号 & 第 02 章类型签名。高阶不靠“把谓词名字存成个体”实现，实际给出关系型量词 \\
R05-b & 余俊伟等，§3.3.1，定义 3.3.1 & 88–89 & 99–100 & 结构为论域加解释，赋值再解释自由变元 & 第 03 章结构/赋值分开；多排序和函数域条件由本稿明确化 \\
R06-a & 郝等，§3.1.1 & 63 & 82 & 一阶语言的基本符号和形成起点 & 第 02 章符号层的交叉核对，不据此声称已有完整 HOL 内核 \\
R06-b & 郝等，§6.1，定理 6.1.1、引理 6.1.2 & 105–107 & 124–126 & 逐规则可靠性、语义替换引理 & 第 03 章可靠性证明组织：公理成立和规则保真分别检查 \\
R06-c & 郝等，§6.2，定理 6.2.1 开始 & 107–109 & 126–128 & 可数一阶语言的 Henkin 完备性构造起点 & 用于核对经典定理的假设；不能把“一阶 Henkin 证明法”与“高阶一般模型”两种含义混为同一陈述 \\
R07-a & 冯琦，§13.3.1，定义 13.39、定理 13.18 & 461–462 & 478–479 & 谓词定义、函数存在唯一、消去翻译和保守扩充 & 对第 03 章再次交叉核查。旧稿把此段落错误归入第 4 章，本版修正 \\
R07-b & 冯琦，§3.1「预备知识：可数与不可数」（定义 3.1、性质条 (7)(8)、定理 3.1（康托）、定理 3.2） & 93–95 & 110–112 & 可数集定义；\(\mathbb{N}\) 的有限次幂与有限序列全体可数（项集与表达式集随之可数）；实数集不可数（康托定理）；可数个可数集之并可数 & 第 01 章命题 1.7 的标准计数事实定位。只支持计数性结论，不支持节点生成、可判性或算法结论；本稿对有限自然数序列采用“长度加元素和”的有限权重层显式枚举，不靠耗尽无限长度层或可数选择来获得该枚举 \\
R08-a & 李未，§3.2，定义 3.10 & 52–55 & 65–68 & 推理树、有限证明树、可证序贯，区分未闭合推理树与证明 & 第 02、06 章有限证书，不能把未完成推演当作证明 \\
R08-b & 李未，§3.3–3.4 相关段 & 56–58 & 69–71 & 推理系统可靠性及协调性接口 & 本稿采用明确的 ND 认证，不声称与该书 G 系统具有逐条相同规则 \\
\end{longtable}
\endgroup

\textbf{实际校正。} Manin 的 I、II、III、IV 章题与页码不能用旧稿大范围描述代替；EFT VIII 的 125 页位于 PDF 131，IX 的 138 页位于 PDF 142；Prestel §2.3 从印刷 71/PDF 81 开始，§2.4 从印刷 75/PDF 85 开始，只有逐页核对的段落才写进上表；余俊伟此处差值为 11；冯琦此处差值为 17。只保留逐段核验的映射。

R04、R06、R08 的 OCR 可能把量词、负号、序贯符号漏掉；本稿公式采用重新写出的定义与证明，不把 OCR 乱码当成公理。R08 PDF66 的印刷 53 页和 R06 PDF126 的印刷 107 页另保存图像锚点于 evidence/page-checks。

\section{B.3 结构、模型与局部—整体的正文证据}
\begingroup\small\setlength{\tabcolsep}{4pt}
\begin{longtable}{l>{\raggedright\arraybackslash}p{\dimexpr 0.3309\linewidth-2\tabcolsep\relax}>{\raggedright\arraybackslash}p{\dimexpr 0.1204\linewidth-2\tabcolsep\relax}>{\raggedright\arraybackslash}p{\dimexpr 0.1204\linewidth-2\tabcolsep\relax}>{\raggedright\arraybackslash}p{\dimexpr 0.3309\linewidth-2\tabcolsep\relax}}
\toprule
\textbf{证据号} & \textbf{来源及定位} & \textbf{印刷页} & \textbf{PDF 页} & \textbf{采用内容与边界} \\
\midrule
\endfirsthead
\toprule
\textbf{证据号} & \textbf{来源及定位} & \textbf{印刷页} & \textbf{PDF 页} & \textbf{采用内容与边界} \\
\midrule
\endhead
\bottomrule
\endfoot
\bottomrule
\endlastfoot
R03-d & Prestel–Delzell，§2.2 & 65–67 & 75–77 & 项结构与模型构造；表达代数不自动等于带来源的节点代数 \\
R03-e & 同书，§2.3 & 71–74 & 81–84 & 子结构、嵌入、初等性必须分开；任意节点子图不是初等子模型 \\
R03-f & 同书，§2.4 & 75–79 & 85–89 & 生成子结构 p.78/PDF88、初等链 p.79/PDF89；第 08 章给出闭包与链并证明，注明固定一阶编码语言 \\
R03-g & 同书，Appendix B & 177–183 & 184–190 & 标准二阶语义与有效完备演算的障碍；不否定 Henkin 完备性 \\
R05-c & 余俊伟等，§4.2.1、§4.2.2 & 169–172 & 180–183 & 结构同态、同构与子结构。该书“同态”的关系保持条件须按原定义读，不能与别处仅单向保持的约定混用 \\
\end{longtable}
\endgroup

模型论只负责所声明的满足/保持问题。学习路线从一种表征翻译到另一种表征后是否更容易，是 E 的任务；本表任何模型论结果都不证明个体认知效用保持。

\section{B.4 数学历史作为动机来源}
\begingroup\small\setlength{\tabcolsep}{4pt}
\begin{longtable}{l>{\raggedright\arraybackslash}p{\dimexpr 0.2253\linewidth-2\tabcolsep\relax}>{\raggedright\arraybackslash}p{\dimexpr 0.1134\linewidth-2\tabcolsep\relax}>{\raggedright\arraybackslash}p{\dimexpr 0.1134\linewidth-2\tabcolsep\relax}>{\raggedright\arraybackslash}p{\dimexpr 0.2253\linewidth-2\tabcolsep\relax}>{\raggedright\arraybackslash}p{\dimexpr 0.2253\linewidth-2\tabcolsep\relax}}
\toprule
\textbf{证据号} & \textbf{Stillwell 定位} & \textbf{印刷页} & \textbf{PDF 页} & \textbf{本稿采用} & \textbf{边界} \\
\midrule
\endfirsthead
\toprule
\textbf{证据号} & \textbf{Stillwell 定位} & \textbf{印刷页} & \textbf{PDF 页} & \textbf{本稿采用} & \textbf{边界} \\
\midrule
\endhead
\bottomrule
\endfoot
\bottomrule
\endlastfoot
R09-a & 序言、§1.1 & vii–viii；1–2 & 8–9；17–18 & 历史与不同数学支线的相互联系，可作为模板的动机入口 & 历史出现次序不规定人人的先修次序 \\
R09-b & §9.1 & 121–122 & 137–138 & 面积、切线、运动和级数等微积分问题来源 & 案例一采用多种问题动机；严格极限判别由案例自行证明 \\
R09-c & §14.1–14.2 附近；§17.1 & 207；255 & 223；271 & 代数结构与几何对象的历史连接、微分几何动机 & 不把这些页视为现代光滑流形的完整定义或证明 \\
R09-d & §19.1 及后续群论起源讨论 & 291–295 & 307–311 & 方程、置换与几何对称等群论来源 & 案例四的 Cayley 定理另给证明；不将现代模算术教学路线误称该页逐字历史结论 \\
R09-e & §22.1–22.2 & 347–348 & 363–364 & 拓扑与曲面研究动机 & 案例二独立约定 Hausdorff、第二可数和微分阶数；不借旧版历史书裁定今天的未解问题 \\
\end{longtable}
\endgroup

此扫描本有其出版年代。任何带时代性的“尚未解决”不直接纳入本体当前数学事实；本稿没有据此声称几何化猜想仍未解决。

\section{B.5 补充的原作者文献}
\begingroup\small\setlength{\tabcolsep}{4pt}
\begin{longtable}{l>{\raggedright\arraybackslash}p{\dimexpr 0.3097\linewidth-2\tabcolsep\relax}>{\raggedright\arraybackslash}p{\dimexpr 0.3097\linewidth-2\tabcolsep\relax}>{\raggedright\arraybackslash}p{\dimexpr 0.3097\linewidth-2\tabcolsep\relax}}
\toprule
\textbf{编号} & \textbf{原作者来源与可访问链接} & \textbf{已核范围/采用方式} & \textbf{本稿接口与适用限制} \\
\midrule
\endfirsthead
\toprule
\textbf{编号} & \textbf{原作者来源与可访问链接} & \textbf{已核范围/采用方式} & \textbf{本稿接口与适用限制} \\
\midrule
\endhead
\bottomrule
\endfoot
\bottomrule
\endlastfoot
S01 & Henkin, \emph{Completeness in the theory of types}, JSL 15(2), 1950, 81–91，\href{https://www.cambridge.org/core/journals/journal-of-symbolic-logic/article/abs/completeness-in-the-theory-of-types1/551D2896D3C6F197030D3D8EEA911F1E}{出版方条目}、\href{https://aaronxyliu.github.io/download/PL/Henkin.pdf}{原论文扫描副本}，DOI 10.2307/2266967 & 2026-09-27 增核扫描正文：印刷84页/PDF5页的一般模型定义、印刷85页/PDF6页的 Theorem 1（一致闭公式集有一般模型）；未声称逐页审校全文 & 一般模型思想与原定理定位。原文演算包含本稿基础核未默认加入的选择设施，不能直接照搬；本稿完备性仍依赖自己的编码及一阶完备性匹配 \\
S02 & Burris–Sankappanavar, \emph{A Course in Universal Algebra}，1981/作者 2012 重排版，\href{https://www.math.uwaterloo.ca/\textasciitilde{}snburris/htdocs/ualg.html}{作者书页} & 原作者版本与章目；I §5 闭包、II §§3、5、6、10 的子代数/同余/商/项代数组织 & 第 08 章使用标准泛代数框架并给出所需有限闭包及递归泛性质证明；不把该书当认知效果证据 \\
S03 & Goguen–Burstall, \emph{Institutions: Abstract Model Theory for Specification and Programming}, JACM 39(1), 1992, 95–146；\href{https://publish.lfcs.inf.ed.ac.uk/reports/90/ECS-LFCS-90-106/}{原始报告}；\href{https://cseweb.ucsd.edu/\textasciitilde{}goguen/projs/inst.html}{Goguen 作者说明} & 作者说明中签名、句子、模型、满足条件及原始论文定位；报告站点访问间歇失败 & 第 08 章给出具体有类型签名 institution 和自身满足引理。作者说明网页个别公式方向的排版不作为类型检查依据 \\
S04 & Doignon–Falmagne, \emph{Knowledge Spaces and Learning Spaces}, 2015，\href{https://arxiv.org/abs/1511.06757}{作者稿}、\href{https://arxiv.org/pdf/1511.06757}{PDF} & §2，PDF3–5 定义 1、3；§3，PDF10 定义 9/定理 10；§8 投影；§9，PDF28 起带噪声反应 & 外部状态族、条目与节点锚定、OR 路线。此文不是“1999 年书的 arXiv 版”，也不使所有学习过程并闭或无遗忘 \\
S05 & Tall–Vinner, \emph{Concept image and concept definition in mathematics with particular reference to limits and continuity}, 1981，\href{https://wrap.warwick.ac.uk/id/eprint/507/1/WRAP\_Tall\_dot1981a-concept-image.pdf}{作者机构稿} & PDF3 的激活意象与 PDF4 的极限附加信念 & 共享误区模式与个体激活状态分开。不能由此推出本稿 Atom 判据已具有测量效度 \\
S06 & Gentner, \emph{Structure-mapping: A theoretical framework for analogy}, Cognitive Science 7, 1983, 155–170，\href{https://groups.psych.northwestern.edu/gentner/papers/Gentner83.2b.pdf}{作者机构稿} & 印刷158 页的映射规则与关系系统性段落 & 类比保留部分关系及失效位置。不能把某种类比的结构合法性当作对所有人的有效性 \\
S07 & Harrison, \emph{Towards Self-verification of HOL Light}，\href{https://www.cl.cam.ac.uk/\textasciitilde{}jrh13/papers/holhol.pdf}{作者稿} & 有限内核、等式与抽象的证明工程对照 & 仅为后续机器形式化参考。本稿 ND 编码不是 HOL Light 实现，也未由该内核检查 \\
S08 & Jan von Plato, \emph{Order in open intervals of computable reals}, MSCS 9(1), 1999, 103–108，\href{https://www.cambridge.org/core/journals/mathematical-structures-in-computer-science/article/abs/order-in-open-intervals-of-computable-reals/1E54DC6F43183479E82F78BE81A828BE}{出版方摘要}，DOI 10.1017/S09601\mcsbrk{}295980\mcsbrk{}02564 & 2026-09-27 核对作者、书目和摘要中实数次序不可判的背景说明；PDF 请求重定向至摘要页，未核全文 & 第 11–12 章区分可计算概率与可判阈值比较；具体规划边界由本稿另行说明，不把摘要当作完整证明 \\
\end{longtable}
\endgroup

补充书与论文不是本稿新增节点、局部化或外置接口的现成定理来源；这些定义和证明明确归本稿。原作者站点若未返回全文，只报告实际核对的摘要、说明或书目，不把访问失败写成已读全文。

\section{B.6 从文献到本稿的依赖链}
\begin{enumerate}
\item R03-a、R04-a、R05-a、R08-a 支持语法/证明区分；第 02 章实际指定语言及有限检查规则。
\item R03-b、R06-b/c 是一阶可靠性/完备性基准；第 03 章用类型编码和函数外延化建立自己的 Henkin 对应。
\item R01-a、R02-a、R07-a 检查定义扩张条件；第 03 章补足高阶域实现义务，第 07–08 章用逐模型可扩张证明外部保守性。
\item R03-e/f、R05-c、S02、S03 组织结构保持；第 08、11–12 章分别证明哪些运算保留哪些性质。
\item R09、S04–S06 提供历史与认知建模素材；个体状态、概率、时间和收益仍外置，其真实效果没有被演绎证明取代。
\end{enumerate}
\textbf{审计规则。} 引用只承担此表所写的有限作用。“定义选择”不由引用变成自然定律，“有限回归通过”不由引用变成一般定理，“作者谈过某主题”不构成具体页码上的完整证明。

\section{B.7 2026-09-29 执行中的增核范围}
本轮复核发现，S01已有具体原典定位，S07也已有作者稿登记；建议中的检索未命中不能改写为仓库缺少来源。本轮重新读取S01扫描的印刷83–85页/PDF4–6页局部可抽取正文，核对选择设施、一般模型定义和Theorem 1；重新读取S07作者稿摘要及§2开头局部文本。未逐页审阅这两篇全文，未核对HOL Light版本化源码，不据此认证本稿Python内核。

来源、实际核对层次及没有取得的原文文件哈希另列在sources-audit.json。其他本地文献沿用本表原登记，不假称本轮重读全部原页。正文3.9仍以自己的翻译/模型对应连接一阶元定理，未直接照搬含选择设施的原演算。

\clearpage\chapter{C 记号、类型与方向}
\mcsorigin{原章节：C-记号索引.md}
\section{C.1 基本判断}
\begingroup\small\setlength{\tabcolsep}{4pt}
\begin{longtable}{>{\raggedright\arraybackslash}p{\dimexpr 0.1817\linewidth-2\tabcolsep\relax}>{\raggedright\arraybackslash}p{\dimexpr 0.7220\linewidth-2\tabcolsep\relax}>{\raggedright\arraybackslash}p{\dimexpr 0.0963\linewidth-2\tabcolsep\relax}}
\toprule
\textbf{记号} & \textbf{类型/含义} & \textbf{位置} \\
\midrule
\endfirsthead
\toprule
\textbf{记号} & \textbf{类型/含义} & \textbf{位置} \\
\midrule
\endhead
\bottomrule
\endfoot
\bottomrule
\endlastfoot
ZFC & 元理论工作背景；不假设 V=L；相对一致性结论分别声明元理论、相关理论及实际假设，不把一致性与传递模型存在混同 & 01 \\
\(\Sigma\)、\(\mathcal{T}\) & 有类型符号签名、句子集理论（公理集或其演绎闭包），不指某一结构的真理论 & 01–03 \\
字体约定 & 高层对象用花体：\(\mathcal{L}_\Sigma\) 对象语言、\(\mathcal{K}\) 模型类、\(\mathcal{T}\) 理论、\(\mathcal{M}\) 本体、\(\mathcal{E}\) 外部模型。结构层对象用哥特体：\(\mathfrak A\) 模型（满足对象语言的结构，第 03 章）、\(\mathfrak C_M\) 概念函数（第 05 章）；外部认知模型 \(\mathcal{E}\) 不是这一意义上的模型，而是状态空间与接口的规格。对象层的具体项（节点、状态、参数）与集合论层级 L 保持普通字体。\(\mathcal{D}\) 已用于分解规则，派生结果暂记普通 D。花体与普通体表示不同对象：\(\mathcal{K}\) 是模型类，\(K_\Sigma\) 是登记概念域 & 01 起 \\
\(\Xi\) & 类型环境，形如 \(x:\sigma\) & 02 \\
\(\Gamma\)、\(\Gamma_{0}\)、\(\Gamma_P\) & 公式假设族；与类型环境不同 & 02、05、06 \\
\(\Xi\vdash\Sigma\) \(t:\tau\) & 良类型项判断 & 02 \\
\(\Gamma\vdash\mathcal{T}\) \(\varphi\) & 固定认证演算的推导；\(\Xi\) 需要时另列 & 02 \\
t\#、\(\hat{\varphi}\)、\(F_\Sigma\)、\(H_\Sigma\) & \(\lambda\) 提升的多排序项/公式翻译、认证语言及背景公理 & 02 \\
\(A,v\vDash\varphi\) & 模型与赋值下的满足 & 03 \\
\(\vDash H\)、\(\vDash\mathrm{std}\) & Henkin 与全函数标准语义；两者不混用 & 03 \\
\(\mathrm{Check}_\mathcal{T}\) & 对显式有限实际假设检查证书；理论公理及任意大假设族的成员资格须有可判依据或见证 & 02、06 \\
\end{longtable}
\endgroup

类型 \(\sigma\to\tau\) 是函数类型，箭头右结合；谓词 \(\alpha\to o\) 也可作为变量取值并被量化。关系类型 \(\mathrm{Node}\to\mathrm{Node}\to o\) 不同于记录该关系的节点代码。

\section{C.2 本体、生成与概念}
\begingroup\small\setlength{\tabcolsep}{4pt}
\begin{longtable}{>{\raggedright\arraybackslash}p{\dimexpr 0.1957\linewidth-2\tabcolsep\relax}>{\raggedright\arraybackslash}p{\dimexpr 0.4327\linewidth-2\tabcolsep\relax}>{\raggedright\arraybackslash}p{\dimexpr 0.3715\linewidth-2\tabcolsep\relax}}
\toprule
\textbf{记号} & \textbf{含义} & \textbf{注意} \\
\midrule
\endfirsthead
\toprule
\textbf{记号} & \textbf{含义} & \textbf{注意} \\
\midrule
\endhead
\bottomrule
\endfoot
\bottomrule
\endlastfoot
\(\mathcal{M}\)、\(N_\mathcal{M}\)、\(\partial\mathcal{M}\) & 本体版本、节点引用域、环境版本边界 & 不含具体学习者 \\
n=(id,k,d) & 不可变节点记录 & 逻辑等价不决定身份 \\
\(U_\Sigma\)、F、\(U_n\) & 生成候选全集、有限元规则、有限生成阶段 & 下标 n 是自然数，不是任意可计算序数 \\
Payload、Rep、Cert & 正式负载、公开表征、证据档案 & 自然语言证明与已检查代码分标签 \\
\(K_\Sigma\)、\(\kappa\) & 登记概念域、有限代码到概念锚定 & 每个形式概念有闭表示 \(t_c\) \\
\(C_\mathrm{occ}\) & 直接语法出现函数 & 扫描归约前代码 \\
\(C_\mathrm{all}\) & 全等价形式概念并集 & 条件下等于 \(K_\Sigma\)；不是集合自映射闭包 \\
\(\Theta\)、\(S_\mathrm{exp}\)、\(S_\mathrm{def}\)、\(S_\mathrm{prf}\)、\(S_\mathrm{rte}\) & 指定变换规则及四种带见证支持 & 各记录自己的来源与用途 \\
Info(A) & Unknown 或 Known(a)，\(a\in A\) & 一般带标签值域 \\
\(\mathrm{ISet}(K)=\mathrm{Info}(\mathcal{P}(K))\) & 概念集合的未知/精确已知值域 & \(\mathrm{Known}(\emptyset)\) 不等于 Unknown \\
\(\mathrm{Supp}^\mathrm{bounds}\)、\(\alpha\) & 部分支持界接口及其到 ISet(K) 的保守投影 & 仅在 L=U 时投影为 Known(L)，否则为 Unknown \\
(L,U) & 可能集合的上下界 \(L\subseteq A\subseteq U\) & 不是把未知变为假 \\
\(\mathrm{MinSupp}_u(n)\) & 指定用途的包含极小充分支持族 & 一般非唯一，可为空 \\
\(\mathrm{Cn}_\mathcal{T}\) & 固定理论下有限前提演绎闭包 & 不是掌握状态更新 \\
\end{longtable}
\endgroup

\(L_\alpha\)、L 只在第 01 章作为集合论背景比较出现。符号 L 在局部公式里也可能指下界，均由上下文明示，不能据同字母推断两者有关；L 保持普通字体，不写成花体 \(\mathcal{L}\)，以免与对象语言 \(\mathcal{L}_\Sigma\) 混淆。

\section{C.3 关系与路线}
\begingroup\small\setlength{\tabcolsep}{4pt}
\begin{longtable}{>{\raggedright\arraybackslash}p{\dimexpr 0.3309\linewidth-2\tabcolsep\relax}>{\raggedright\arraybackslash}p{\dimexpr 0.6691\linewidth-2\tabcolsep\relax}}
\toprule
\textbf{记号} & \textbf{类型/方向} \\
\midrule
\endfirsthead
\toprule
\textbf{记号} & \textbf{类型/方向} \\
\midrule
\endhead
\bottomrule
\endfoot
\bottomrule
\endlastfoot
Ded(P,n,p) & 有限带标签前提族 P、结论节点 n、证明代码 p 的三元认证关系 \\
\(\Lambda_\mathcal{M}\)、\(a=(\mathrm{id},I_a,O_a,w,\mathrm{mode})\) & 公共行动库及有限联合输入/输出契约 \\
\(\mathrm{Rte}_\mathcal{M}(B,G)\) & 给定背景 B 与目标 G 的结构路线族；定义先于个性化规划 \\
\(P_r\)、act、\(h_\mathrm{src}\)、\(\zeta\) & 事件集、行动标记、输入来源、目标生产者 \\
\(\mathrm{Req}_r(g)\) & g 活跃事件所需的输入资源并集；与概念支持不同 \\
\(\mathrm{HardPre}_r\)、\(\mathrm{HardPre}_\mathcal{R}\) & 单路线义务、非空指定路线族的共同义务 \\
\(\mathrm{NewPre}_r\) & 由事件提供而非背景直接给定的活跃输入 \\
\(g\succeq_H\) s & g 比 s 一般；同载体时由 \(P_s\Rightarrow P_g\) 认证 \\
HSpec、Spec(g) & 泛化的逆关系、特化的集合值函数；不是逆单值函数 \\
\(R^{0}_\mathcal{M}\)、\(\mathrm{Eval}_\mathcal{E}\)、\(\mathrm{Soft}^\pi_\mathcal{E},s\) & 公共软关系描述、外部评价、由明确判据激活的结果 \\
\(S\circ R\) & 先 R 后 S；中介节点存在量化 \\
\(R^\mathrm{smile}\)、R* & 逆关系、自反传递闭包 \\
\(F:\mathcal{M}\to\mathcal{M}'\) & 带语法及证据输送的认证空间态射 \\
\end{longtable}
\endgroup

\section{C.4 外部认知、局部化与规划}
\begingroup\small\setlength{\tabcolsep}{4pt}
\begin{longtable}{>{\raggedright\arraybackslash}p{\dimexpr 0.3051\linewidth-2\tabcolsep\relax}>{\raggedright\arraybackslash}p{\dimexpr 0.6949\linewidth-2\tabcolsep\relax}}
\toprule
\textbf{记号} & \textbf{含义与归属} \\
\midrule
\endfirsthead
\toprule
\textbf{记号} & \textbf{含义与归属} \\
\midrule
\endhead
\bottomrule
\endfoot
\bottomrule
\endlastfoot
\(\mathcal{E}\)、\(X_\mathcal{E}\)、s & 外部模型、其状态空间及状态 \\
\(\mathrm{Obs}_\mathcal{E}\)、\(\mathrm{Trans}_\mathcal{E}\)；\(O_\mathcal{E}\)、\(K_\mathcal{E}\) & 观察/转移关系；可选概率核 \\
\texttt{Anchor\mcsbrk\_\ensuremath{\mathcal{E}}\ensuremath{\subseteq}Item\mcsbrk\_\ensuremath{\mathcal{E}}\ensuremath{\times}N\mcsbrk\_\ensuremath{\mathcal{M}}\ensuremath{\times}Competence\mcsbrk\_\ensuremath{\mathcal{M}}} & 条目、节点、能力的多对多观测锚定 \\
\(\mathrm{Atom}_\mathcal{E}(s,n;g,\mathcal{R},\mathcal{D})\) & 相对于目标、表征和分解规则的认知原子性；此处 \(\mathcal{R}\) 是表征族 \\
\(\mathrm{Inst}_\mathcal{E}(s,p)\) & 该状态与公共误区模式的关联判断 \\
\(D(\mathcal{M},\mathcal{E},s,g,b,\rho)\) & 明确输入决定的派生结果（含适配视图、评价、路线与呈现）；只读 \(\mathcal{M}\)。自变量是输入，\(\mathrm{Loc}_\lambda(\mathcal{M})\) 与 Plan(V) 是它的分量而非自变量 \\
Flat & 话题/领域/学科对象到节点域的展开 \\
\(\mathrm{Loc}_\lambda(\mathcal{M})=(V,\partial V,\tau,\mathrm{Pres})\) & 可见结构、隐藏义务、溯源、逐项保持声明 \\
\(\mathrm{Adapt}(\mathcal{E},s,g,b)\) & 从外部输入产生局部化参数 \\
\(\mathrm{Plan}_M,\mathcal{E},s,g,b,\rho:\mathrm{View}_M\) \(\rightharpoonup\) FinPos & 部分规划函数；视图映到有限事件偏序的同构类型，行动/来源/外部约束另存结果包 \\
\(\Pi\)、\(\chi\)、\(\ell\)、ev & 完整规划包、排程约束、通过的线性扩张、事件关注节点 \\
h & 事件数量/规划时域上界；与来源函数 \(h_\mathrm{src}\) 不同 \\
\end{longtable}
\endgroup

符号 \(\mathcal{R}\) 在硬先修章节表示路线族，在 Atom 的显示式中按用户原式表示表征族；每次都带类型声明。可实施编码中应分为 RouteFamily 与 Repres\mcsbrk{}entati\mcsbrk{}onFami\mcsbrk{}ly 两个类型，不能共用未分型变量。

证明长度、所用公理数量等结构量可以在 \(\mathcal{M}\) 中；个人感知难度、预测时长和预期收益在 \(\mathcal{E}\) 中。原始数学代码、知识节点、图形显示点与学习事件是四种不同对象。

\clearpage\chapter{有限模型与文档核验报告}
\mcsorigin{原章节：REPORT.md}
执行命令：python \texttt{E:\textbackslash{}MCS\textbackslash{}mcs-foundations\textbackslash{}validation\textbackslash{}verify\mcsbrk\_finite\mcsbrk\_models.py}

使用 Python 标准库。以下结果验证有限结构实例与文档完整性；不替代正文一般证明、HOL 内核认证或认知实证。

\begingroup\small\setlength{\tabcolsep}{4pt}
\begin{longtable}{>{\raggedright\arraybackslash}p{\dimexpr 0.3713\linewidth-2\tabcolsep\relax}>{\raggedright\arraybackslash}p{\dimexpr 0.1069\linewidth-2\tabcolsep\relax}r>{\raggedright\arraybackslash}p{\dimexpr 0.3713\linewidth-2\tabcolsep\relax}}
\toprule
\textbf{检查} & \textbf{状态} & \textbf{有限实例数} & \textbf{结果} \\
\midrule
\endfirsthead
\toprule
\textbf{检查} & \textbf{状态} & \textbf{有限实例数} & \textbf{结果} \\
\midrule
\endhead
\bottomrule
\endfoot
\bottomrule
\endlastfoot
\texttt{V01\mcsbrk\_closure\mcsbrk\_exhaustive} & passed & 512 & 64 个规则系统 \(\times\) 8 个背景：与全部封闭超集的交这一独立语义计算一致；满足单调、幂等及有界可达性。 \\
\texttt{V02\mcsbrk\_and\mcsbrk\_or\mcsbrk\_prerequisites} & passed & 8 & 联合输入保持 AND，替代入口保持 OR；空路线族不产生先修；个体筛选不改变公共路线族。 \\
\texttt{V03\mcsbrk\_event\mcsbrk\_posets} & passed & 65 & 四事件的 64 个按序有向无环图，核验 315 个线性扩张及相邻不可比交换连通性；另核对来源与策略分别无环但合并成环的二事件反例。 \\
\texttt{V04\mcsbrk\_boundary\mcsbrk\_and\mcsbrk\_sources} & passed & 6 & 三事件界内全部来源指派尊重每个线性扩张；隐藏 a、b 后两项义务均留在边界。 \\
\texttt{V05\mcsbrk\_cycle\mcsbrk\_repeat\mcsbrk\_budget} & passed & 6 & 无入口循环不可执行；独立零输入入口可解除循环；重复访问同一节点保留不同事件；零预算拒绝单位成本行动。 \\
\texttt{V06\mcsbrk\_or\mcsbrk\_knowledge\mcsbrk\_space} & passed & 49 & 七状态外部实例并闭且可达，但不交闭，因此不是任何条目偏序的全部下闭集族。 \\
\texttt{V07\mcsbrk\_unknown\mcsbrk\_and\mcsbrk\_no\mcsbrk\_reflection} & passed & 4 & 未知支持不等于已知空集，集合运算遍历可能值；命题可为假而相信它的判断为真。 \\
\texttt{V08\mcsbrk\_padding\mcsbrk\_witness} & passed & 1 & 闭常元出现在未使用实参中，\(\beta\) 结果仍为 e；一般饱和结论由正文证明。 \\
\texttt{V09\mcsbrk\_external\mcsbrk\_noninterference} & passed & 9 & 两名合成学习者得到不同群/张量路线和极限提示；学习、遗忘、答错、纠正及重校准不改变本体哈希；冷启动保留未知。 \\
\texttt{V10\mcsbrk\_explicit\mcsbrk\_inputs\mcsbrk\_and\mcsbrk\_pareto} & passed & 11 & 逐项改变本体、模型、状态、目标、预算、策略及种子均改变复现键；未知成本保留为不可比较候选，不用默认数值支配。 \\
\texttt{V11\mcsbrk\_document\mcsbrk\_integrity} & passed & 26 & 26 篇正文文档：本地链接、编号唯一性、LC01–LC36、13 类模板及四篇案例均通过检查。 \\
\texttt{V12\mcsbrk\_aggregation\mcsbrk\_adjunction} & passed & 17496 & 三元素域上的全部 128 个非空块子族选择，穷尽检查展开/完整块选择的伴随及闭包、内部算子的幂等性。 \\
\texttt{V13\mcsbrk\_route\mcsbrk\_port\mcsbrk\_substitution} & passed & 10 & 顺序复合同时替换事件输入与最终目标端口；核验左右单位、三路结合、背景透传及无资源依赖时的顺序约束，并拒绝悬空中间目标的旧反例。 \\
\texttt{V14\mcsbrk\_noninjective\mcsbrk\_route\mcsbrk\_transport} & passed & 5 & 两个不同输入合并为同一节点时须选一个合法生产者；输送后依赖为原图子图，不能宣称保留全部原依赖。 \\
\texttt{V15\mcsbrk\_observational\mcsbrk\_nonide\mcsbrk{}ntific\mcsbrk{}ation} & passed & 6 & 两个确定模型在长度0至4的纯a记录上一致，在未试行动b上预测相反；见正文反例16.1的一般论证。 \\
\end{longtable}
\endgroup

公共模型 SHA-256：ece7b3\mcsbrk{}f6266a\mcsbrk{}7ed648\mcsbrk{}3a2140\mcsbrk{}7322cd\mcsbrk{}9c0cd4\mcsbrk{}812fd2\mcsbrk{}ab08c1\mcsbrk{}751e03\mcsbrk{}c26a4f\mcsbrk{}6609

索引条目：224；局部化方案：36；贯通案例：4。

详细结果及两名合成学习者的派生路线见 report.json。本报告仅在所有断言通过后写出；失败时进程非零退出。


\end{document}
